% ============================================================
% 电动力学自学教材 —— 专业审查优化版
% 基于出版编辑与物理学者的双视角审查修订
% ============================================================
\documentclass[11pt,a4paper,openany,oneside]{book}

\usepackage[UTF8]{ctex}
\ctexset{punct=quanjiao}
\linespread{0.98}
\setlength{\parskip}{2pt plus 1pt}
\usepackage{geometry}
\geometry{left=2.2cm,right=2.8cm,top=2.2cm,bottom=2.2cm}
\setlength{\headheight}{15pt}
\usepackage{amsmath,amssymb,amsthm}
\usepackage{mathtools}
\usepackage{bm}
\usepackage{esint}
\usepackage{siunitx}
\usepackage{tikz,pgfplots}
\pgfplotsset{compat=1.18}
\usetikzlibrary{arrows.meta,decorations.pathmorphing,patterns,calc,positioning,shapes,fadings}
\usepackage{xcolor}
\usepackage{graphicx}
\usepackage{enumitem}
\usepackage{fancyhdr}
\usepackage{titlesec}
\usepackage{tcolorbox}
\tcbuselibrary{theorems,skins,breakable}
\usepackage{booktabs}
\usepackage{array}
\usepackage{caption}
\usepackage{subcaption}
\usepackage{hyperref}
\hypersetup{
  pdftitle={电动力学 v2.15（2026-09）},
  pdfauthor={Han \& Kimi \& glm},
  pdfsubject={电动力学自学教材},
  pdfkeywords={电动力学, 电磁场, 麦克斯韦方程组, 相对论, 自学教材}
}
\usepackage{cleveref}
\usepackage{braket}
\usepackage{float}
\usepackage{multirow}
\usepackage{marginnote}
\usepackage{epigraph}
\usepackage{makeidx}
\usepackage[printocg=always]{ocgx2}  % 可折叠内容（线上PDF交互）
                                      % printocg=always：为图层写入 /Print<</PrintState/ON>>，
                                      % 使折叠内容在**打印/导出为PDF时始终输出**。
                                      % 否则（缺 /Usage 字典）图层按当前可见状态打印，
                                      % 默认关闭的答案层与推导层会在纸上变成空白。
\usepackage{pdfcomment}       % PDF tooltip注释

% ---------- marginnote 设置 ----------
\setlength{\marginparwidth}{1.2cm}
\setlength{\marginparsep}{0.4cm}

% ---------- 页眉页脚 ----------
\pagestyle{fancy}
\fancyhf{}
\fancyhead[L]{\nouppercase{\leftmark}}
\fancyhead[R]{\thepage}
\renewcommand{\headrulewidth}{0.4pt}

% ---------- 标题格式 ----------
\titleformat{\chapter}[display]
  {\normalfont\huge\bfseries}{\chaptertitlename\ \thechapter}{20pt}{\Huge}
\titleformat{\section}{\normalfont\Large\bfseries}{\thesection}{1em}{}
\titleformat{\subsection}{\normalfont\large\bfseries}{\thesubsection}{1em}{}

% ---------- 颜色定义 ----------
\definecolor{thmblue}{RGB}{41, 98, 255}
\definecolor{defgreen}{RGB}{46, 125, 50}
\definecolor{cororange}{RGB}{230, 81, 0}
\definecolor{exampurple}{RGB}{106, 27, 154}
\definecolor{insightred}{RGB}{183, 28, 28}
\definecolor{warnyellow}{RGB}{255, 160, 0}
\definecolor{historybrown}{RGB}{121, 85, 72}
\definecolor{reviewblue}{RGB}{25, 118, 210}
\definecolor{tipgreen}{RGB}{56, 142, 60}

% ---------- 定理环境（带编号）----------
\newtcbtheorem[number within=chapter]{theorem}{定理}{
  colback=thmblue!5,colframe=thmblue,fonttitle=\bfseries,breakable,
  left=4pt,right=4pt,top=4pt,bottom=4pt}{thm}
\newtcbtheorem[number within=chapter]{definition}{定义}{
  colback=defgreen!5,colframe=defgreen,fonttitle=\bfseries,breakable,
  left=4pt,right=4pt,top=4pt,bottom=4pt}{def}
\newtcbtheorem[number within=chapter]{corollary}{推论}{
  colback=cororange!5,colframe=cororange,fonttitle=\bfseries,breakable,
  left=4pt,right=4pt,top=4pt,bottom=4pt}{cor}
\newtcbtheorem[number within=chapter]{example}{例题}{
  colback=exampurple!5,colframe=exampurple,fonttitle=\bfseries,breakable,
  left=4pt,right=4pt,top=4pt,bottom=4pt}{ex}
\newtcbtheorem[number within=chapter]{insight}{物理洞见}{
  colback=insightred!5,colframe=insightred!80,fonttitle=\bfseries,breakable,
  left=4pt,right=4pt,top=4pt,bottom=4pt}{ins}
\newtcbtheorem[number within=chapter]{warning}{常见误区}{
  colback=warnyellow!8,colframe=warnyellow,fonttitle=\bfseries,breakable,
  left=4pt,right=4pt,top=4pt,bottom=4pt}{warn}
\newtcbtheorem[number within=chapter]{history}{物理史话}{
  colback=historybrown!5,colframe=historybrown,fonttitle=\bfseries,breakable,
  left=4pt,right=4pt,top=4pt,bottom=4pt}{hist}
\newtcbtheorem[number within=chapter]{review}{本节要点}{
  colback=reviewblue!5,colframe=reviewblue,fonttitle=\bfseries,breakable,
  left=4pt,right=4pt,top=4pt,bottom=4pt}{rev}
\newtcbtheorem[number within=chapter]{tip}{学习提示}{
  colback=tipgreen!5,colframe=tipgreen,fonttitle=\bfseries,breakable,
  left=4pt,right=4pt,top=4pt,bottom=4pt}{tip}

% ---------- 随堂自测小题环境 ----------
% ---------- 随堂自测小题环境 ----------
\newtcolorbox{miniquiz}[1][]{
  colback=reviewblue!3,colframe=reviewblue,fonttitle=\bfseries,
  title={随堂自测},breakable,
  left=4pt,right=4pt,top=4pt,bottom=4pt,#1}

% ---------- 本章要点回顾环境 ----------
\newtcolorbox{chapterreview}[1][]{
  colback=tipgreen!5,colframe=tipgreen,fonttitle=\bfseries,
  title={本章要点回顾},breakable,
  left=4pt,right=4pt,top=4pt,bottom=4pt,#1}

\newcommand{\diff}{\mathrm{d}}

% ---------- 彩色矢量符号（线上PDF专属） ----------
\newcommand{\vecE}{\textcolor{red}{\vect{E}}}
\newcommand{\vecB}{\textcolor{blue}{\vect{B}}}
\newcommand{\vecK}{\textcolor{green!60!black}{\vect{k}}}
\newcommand{\vecJ}{\textcolor{cororange}{\vect{j}}}
\newcommand{\vecA}{\textcolor{blue!70!black}{\vect{A}}}

% ---------- 可折叠内容（ocgx2） ----------
\newcommand{\toggleButton}[2]{%
  \actionsocg{#2}{}{#2}{\fcolorbox{gray}{gray!10}{\small #1}}%
}
\newenvironment{foldable}[1]{%
  \begin{ocg}{#1}{#1}{off}%
}{%
  \end{ocg}%
}

% ---------- 难度标签（彩色） ----------
\newcommand{\diffTag}[1]{%
  \textbf{[#1]}%
}

% ---------- PDF tooltip 注释 ----------
\newcommand{\physTooltip}[2]{%
  \pdfcomment[icon=Note,subject={物理注释},color={yellow!50}]{#2}{#1}%
}

\newcommand{\ii}{\mathrm{i}}
\newcommand{\ee}{\mathrm{e}}
\newcommand{\eps}{\varepsilon}
\newcommand{\vac}{\varepsilon_0}
\newcommand{\muz}{\mu_0}
\newcommand{\vect}[1]{\bm{#1}}
\newcommand{\unitvec}[1]{\hat{\vect{#1}}}
\newcommand{\grad}{\nabla}
\newcommand{\curl}{\nabla\times}
\newcommand{\divg}{\nabla\cdot}
\newcommand{\lap}{\nabla^2}
\newcommand{\difficulty}[1]{\marginnote{\textcolor{red}{\textbf{[#1]}}}}

% ---------- 核心方程方框高亮（快速翻阅复习用） ----------
% 用法：\begin{impEqbox}{标题} ... \end{impEqbox}
\newtcolorbox{impEqbox}[1][]{
  colback=exampurple!4,colframe=exampurple!70,fonttitle=\bfseries,
  title={#1},breakable,
  boxrule=1.2pt,arc=2mm,
  left=6pt,right=6pt,top=5pt,bottom=5pt}


\begin{document}

% ============ 封面 ============
\thispagestyle{empty}
\begin{tikzpicture}[remember picture, overlay]

% ---------- 封面配色 ----------
\definecolor{covTop}{RGB}{26,50,108}
\definecolor{covMid}{RGB}{10,20,52}
\definecolor{covBot}{RGB}{2,3,10}
\definecolor{covGlow}{RGB}{80,205,255}

% ---------- 背景：三段竖向渐变（对比更强） ----------
\shade[top color=covTop, bottom color=covMid]
  (current page.north west) rectangle ([yshift=-0.40\paperheight]current page.north east);
\shade[top color=covMid, bottom color=covBot]
  ([yshift=-0.40\paperheight]current page.north west) rectangle (current page.south east);

% ---------- 中央冷光晕（分层同心圆模拟柔光） ----------
\foreach \r/\o in {12/0.009, 10/0.012, 8.5/0.015, 7/0.018, 5.8/0.022, 4.6/0.027, 3.5/0.033, 2.5/0.039, 1.6/0.045, 0.9/0.054} {
  \fill[covGlow, opacity=\o] ($(current page.center)+(0,3.0)$) circle (\r);
}

% ---------- 背景网格 ----------
\foreach \x in {-9,-8,...,9} {
  \draw[white, opacity=0.030, line width=0.25pt]
    ($(current page.center)+(\x,-13.5)$) -- ($(current page.center)+(\x,13.5)$);
}
\foreach \y in {-13,-12,...,13} {
  \draw[white, opacity=0.030, line width=0.25pt]
    ($(current page.center)+(-9.5,\y)$) -- ($(current page.center)+(9.5,\y)$);
}

% ---------- 星点（固定坐标） ----------
\foreach \p in {(-8.4,8.6),(-7.2,6.4),(-8.9,4.2),(-6.3,9.4),(-5.1,7.1),(-3.4,8.9),(-1.8,6.6),(-0.6,9.8),(0.9,7.8),(2.3,9.1),(3.6,6.3),(5.2,8.4),(6.7,5.6),(8.1,7.9),(-8.2,-3.4),(-6.8,-5.2),(-5.4,-2.8),(-3.9,-4.6),(-2.2,-6.4),(-0.8,-3.8),(1.4,-5.6),(2.9,-3.3),(4.4,-6.0),(6.1,-3.1),(7.6,-4.9),(8.8,-2.4)} {
  \fill[white, opacity=0.5] \p circle (0.55pt);
}

% ---------- 主视觉：三维透视电磁波 ----------
\begin{scope}[shift={($(current page.center)+(0,3.0)$)}]

  % 底纹一：偶极辐射角分布 sin^2(theta)
  \begin{scope}[scale=3.3, opacity=0.15]
    \draw[covGlow, line width=0.9pt, samples=220, domain=0:360, smooth]
      plot ({sin(\x)*sin(\x)*sin(\x)}, {sin(\x)*sin(\x)*cos(\x)});
    \draw[covGlow, line width=0.35pt, samples=220, domain=0:360, smooth]
      plot ({1.26*sin(\x)*sin(\x)*sin(\x)}, {1.26*sin(\x)*sin(\x)*cos(\x)});
  \end{scope}

  % 底纹二：极坐标网格（辐射图风格）
  \begin{scope}[opacity=0.085]
    \foreach \rr in {1.5,2.4,3.3,4.2,5.1} { \draw[covGlow, line width=0.3pt] (0,0) circle (\rr); }
    \foreach \ang in {0,30,...,330} { \draw[covGlow, line width=0.3pt] (0,0) -- (\ang:5.1); }
  \end{scope}

  % 波前弧（自源点向外扩散）
  \foreach \rr in {4.4,5.3,6.2,7.1} {
    \draw[covGlow, opacity=0.10, line width=0.5pt] (-8.6,0) ++(22:\rr) arc (22:-22:\rr);
  }

  % 传播轴
  \draw[green!70!white, opacity=0.85, line width=1.2pt, -{Latex[length=5.5mm,width=3.2mm]}]
    (-8.6,0) -- (8.6,0);
  \node[green!80!white, font=\normalsize, anchor=west] at (8.65,0) {$\vecK$};

  % 偶极源（+ / - 一对电荷）
  \fill[red!70!white] (-8.6,0.32) circle (0.13);
  \node[black, font=\tiny] at (-8.6,0.32) {$+$};
  \fill[blue!60!white] (-8.6,-0.32) circle (0.13);
  \node[white, font=\tiny] at (-8.6,-0.32) {$-$};

  % ---- E 波（红，竖直振荡，带辉光） ----
  \foreach \lw/\op in {12pt/0.05, 6.5pt/0.09, 3.4pt/0.16, 1.8pt/0.96} {
    \draw[red!70!white, opacity=\op, line width=\lw, samples=150, domain=-7.7:7.7, smooth]
      plot (\x, {2.45*sin(52*\x)});
  }
  \foreach \x in {-7.2,-6.6,...,7.2} {
    \draw[red!60!white, opacity=0.32, line width=0.55pt] (\x,0) -- (\x,{2.45*sin(52*\x)});
  }
  \node[red!80!white, font=\large, anchor=south] at (0.95,2.35) {$\vecE$};

  % ---- B 波（蓝，斜向进纸平面，带辉光） ----
  \foreach \lw/\op in {12pt/0.05, 6.5pt/0.09, 3.4pt/0.16, 1.8pt/0.96} {
    \draw[blue!60!white, opacity=\op, line width=\lw, samples=150, domain=-7.7:7.7, smooth]
      plot ({\x+1.30*sin(52*\x)}, {-0.68*sin(52*\x)});
  }
  \foreach \x in {-7.2,-6.6,...,7.2} {
    \draw[blue!55!white, opacity=0.30, line width=0.55pt]
      (\x,0) -- ({\x+1.30*sin(52*\x)},{-0.68*sin(52*\x)});
  }
  \node[blue!70!white, font=\large, anchor=north] at (-4.3,-1.30) {$\vecB$};

\end{scope}

% ---------- 四角暗角（艺术收边） ----------
\fill[black, opacity=0.55, path fading=south]
  ($(current page.north west)$) rectangle ($(current page.north east)+(0,-3.2cm)$);
\fill[black, opacity=0.55, path fading=north]
  ($(current page.south west)$) rectangle ($(current page.south east)+(0,3.2cm)$);
\fill[black, opacity=0.45, path fading=east]
  ($(current page.north west)$) rectangle ($(current page.north west)+(3.0cm,-29.7cm)$);
\fill[black, opacity=0.45, path fading=west]
  ($(current page.north east)$) rectangle ($(current page.north east)+(-3.0cm,-29.7cm)$);

% ---------- 四角径向弧线（辐射纹理） ----------
\begin{scope}[opacity=0.075]
  \foreach \rr in {3.6,5.0,6.4,7.8} {
    \draw[covGlow, line width=0.5pt] ($(current page.south west)+(0,0)$) ++(8:\rr) arc (8:82:\rr);
    \draw[covGlow, line width=0.5pt] ($(current page.north east)+(0,0)$) ++(188:\rr) arc (188:262:\rr);
  }
\end{scope}

% ---------- 书名（柔光 + 主标题） ----------
\foreach \rr/\oo in {3.6/0.018, 2.8/0.022, 2.1/0.026, 1.4/0.030} {
  \fill[white, opacity=\oo] ([yshift=-2.3cm]current page.north) ellipse ({\rr*3.6} and {\rr*0.9});
}
\node[anchor=north, align=center] at ([yshift=-2.4cm]current page.north) {
  {\fontsize{50}{58}\selectfont\bfseries\textcolor{white}{电动力学}}\\[0.32cm]
  {\fontsize{12.5}{15}\selectfont\textcolor{white!78}{从矢量分析到相对论协变的完整电磁理论}}
};

% ---------- 标题下的渐变装饰线 ----------
\draw[white, opacity=0.10, line width=1.6pt] ([yshift=-5.7cm]current page.north) ++(-3.6,0) -- ++(1.9,0);
\draw[covGlow, opacity=0.50, line width=1.6pt] ([yshift=-5.7cm]current page.north) ++(-1.7,0) -- ++(1.9,0);
\draw[white, opacity=0.10, line width=1.6pt] ([yshift=-5.7cm]current page.north) ++(0.2,0) -- ++(1.9,0);

% ---------- 麦克斯韦方程组（全书灵魂） ----------
\node[anchor=south, align=center] at ([yshift=9.6cm]current page.south) {
  {\footnotesize\textcolor{white!66}{$\divg\vecE=\dfrac{\rho}{\eps_0}$\quad$\divg\vecB=0$\quad$\curl\vecE=-\dfrac{\partial\vecB}{\partial t}$\quad$\curl\vecB=\muz\vecJ+\muz\eps_0\dfrac{\partial\vecE}{\partial t}$}}
};
\draw[white, opacity=0.13, line width=0.5pt] ([yshift=8.9cm]current page.south) ++(-6.4,0) -- ++(12.8,0);

% ---------- 作者与版本 ----------
\node[anchor=south, align=center] at ([yshift=2.7cm]current page.south) {
  {\LARGE\bfseries\textcolor{white}{Han \& Kimi \& glm}}\\[0.3cm]
  {\footnotesize\textcolor{white!58}{AI 协作生成 · 人工审校}}\\[0.45cm]
  {\small\bfseries\textcolor{white!80}{v2.15}}\\[0.15cm]
  {\scriptsize\textcolor{white!48}{2026 年 9 月}}
};

\end{tikzpicture}
\clearpage
% ============ 目录 ============
\tableofcontents
\newpage

% ---------- 核心目标快速导航 ----------
\section*{各章核心目标}
\begin{itemize}
\item \hyperlink{ch1-goal}{第1章 数学预备知识：梯度、散度、旋度与场论基础}
\item \hyperlink{ch2-goal}{第2章 静电学基础：高斯定理、环路定理与泊松方程}
\item \hyperlink{ch3-goal}{第3章 电介质与静电边界条件：极化、电位移与边界条件}
\item \hyperlink{ch4-goal}{第4章 静电学边值问题：分离变量、镜像法与格林函数}
\item \hyperlink{ch5-goal}{第5章 静磁学：毕奥-萨伐尔定律、安培定理与矢势}
\item \hyperlink{ch6-goal}{第6章 磁介质：磁化强度、磁场强度与磁学边界条件}
\item \hyperlink{ch7-goal}{第7章 麦克斯韦方程组：位移电流、法拉第定律与电磁波}
\item \hyperlink{ch8-goal}{第8章 电磁波的传播：平面波、反射折射与波导}
\item \hyperlink{ch9-goal}{第9章 电磁波的辐射：推迟势、电偶极辐射与天线}
\item \hyperlink{ch10-goal}{第10章 相对论电动力学：四维协变形式与场张量}
\end{itemize}
\newpage


% ============ 符号说明 ============
\chapter*{符号说明}
\addcontentsline{toc}{chapter}{符号说明}

\begin{center}
\small
\begin{tabular}{@{}llp{0.38\textwidth}@{}}
\toprule
\textbf{符号} & \textbf{含义} & \textbf{说明} \\
\midrule
\hypertarget{symA}{$\vecA, \bm{A}$} & 一般矢量 / 磁矢势（上下文区分） & 粗体或带箭头；磁矢势时 $\vecB = \curl\vecA$，单位 \si{T.m} = \si{Wb/m} \\
\hypertarget{symScalarA}{$A$} & 标量 & 斜体 \\
\hypertarget{symnhat}{$\hat{\vect{n}}, \unitvec{n}$} & 单位矢量 & 模长为1 \\
\hypertarget{symr}{$\vect{r}$} & 场点位置矢量 & 观察点的坐标 \\
\hypertarget{symrprime}{$\vect{r}'$} & 源点位置矢量 & 电荷/电流所在位置 \\
\hypertarget{symR}{$R = |\vect{r}-\vect{r}'|$} & 源点到场点距离 & 推迟势中用 \\
\hypertarget{symgrad}{$\grad, \nabla$} & 梯度算子 & 第1章\S1.3 \\
\hypertarget{symdivg}{$\divg$} & 散度 & $\nabla\cdot$ \\
\hypertarget{symcurl}{$\curl$} & 旋度 & $\nabla\times$ \\
\hypertarget{symlap}{$\lap$} & 拉普拉斯算子 & $\nabla^2$ \\
\hypertarget{symdiff}{$\diff$} & 微分符号 & 正体$\mathrm{d}$ \\
\hypertarget{symii}{$\ii$} & 虚数单位 & $\mathrm{i}=\sqrt{-1}$ \\
\hypertarget{symee}{$\ee$} & 自然对数底 & $\mathrm{e}$ \\
\hypertarget{symeps0}{$\eps_0$} & 真空介电常数 & $\approx 8.854\times 10^{-12}\,\mathrm{F/m}$ \\
\hypertarget{symmuz0}{$\mu_0$} & 真空磁导率 & $4\pi\times 10^{-7}\,\mathrm{H/m}$ \\
\hypertarget{symc}{$c$} & 光速 & $1/\sqrt{\mu_0\eps_0}\approx 3\times 10^8\,\mathrm{m/s}$ \\
\midrule
\multicolumn{3}{l}{\textit{—— 电磁场量（第 3、6 章引入）——}} \\
\hypertarget{symE}{$\vecE$} & 电场强度 & 基本场，单位 \si{V/m}，红色标注 \\
\hypertarget{symD}{$\vect{D}$} & 电位移矢量 & $\eps_0\vecE + \vect{P}$，单位 \si{C/m^2} \\
\hypertarget{symP}{$\vect{P}$} & 极化强度 & 偶极矩密度，单位 \si{C/m^2} \\
\hypertarget{symB}{$\vecB$} & 磁感应强度 & 基本场，单位 \si{T}，蓝色标注 \\
\hypertarget{symH}{$\vect{H}$} & 磁场强度 & $\vecB/\mu_0 - \vect{M}$，单位 \si{A/m} \\
\hypertarget{symM}{$\vect{M}$} & 磁化强度 & 偶极矩密度，单位 \si{A/m} \\
\midrule
\multicolumn{3}{l}{\textit{—— 势与能量（第 2、5、7 章引入）——}} \\
\hypertarget{symphi}{$\phi$} & 标量电势 & $\vecE = -\grad\phi - \partial\vecA/\partial t$，单位 \si{V} \\
\hypertarget{symu}{$u$} & 电磁场能量密度 & $u = \frac{1}{2}(\eps_0 E^2 + B^2/\mu_0)$，单位 \si{J/m^3} \\
\hypertarget{symS}{$\vect{S}$} & 坡印廷矢量（能流密度） & $\vecE \times \vect{H}$，单位 \si{W/m^2} \\
\hypertarget{symchi}{$\chi_e,\ \chi_m$} & 电极化率 / 磁化率 & 无量纲，$\eps_r = 1+\chi_e$，$\mu_r = 1+\chi_m$ \\
\bottomrule
\end{tabular}
\end{center}

\vspace{1em}
\noindent\textbf{约定}：
\begin{itemize}[nosep]
\item 三维空间矢量用粗体（$\vecE$），四维时空矢量用上标（$A^\mu$）
\item 源点坐标带撇（$\vect{r}'$），场点坐标不带撇（$\vect{r}$）
\item 偏导数记号：$\partial_x = \partial/\partial x$，$\partial_t = \partial/\partial t$
\item 重复指标默认求和（爱因斯坦求和约定，第10章启用）
\end{itemize}

\newpage
\chapter*{全书高频公式速查表}
\addcontentsline{toc}{chapter}{全书高频公式速查表}
% ============================================================

按章节分块，方便快速查阅与复习。所有公式均采用SI单位制。

\begin{tip}{使用说明}{quickref-tip}
本表置于全书最前，供复习时\textbf{快速翻阅}使用；各公式的推导与物理意义在对应章节。外文教材（Griffiths、Jackson）部分公式采用高斯单位制，转换方法见附录"单位制转换手册"。
\end{tip}

\section{第1章：矢量分析}

\begin{center}
\begin{tabular}{@{}ll@{}}
\toprule
\textbf{公式} & \textbf{说明} \\
\midrule
$\grad\phi = \unitvec{x}\partial\phi/\partial x + \cdots$ & 直角坐标梯度 \\
$\divg\vecA = \partial A_x/\partial x + \cdots$ & 直角坐标散度 \\
$\curl\vecA = \begin{vmatrix}\unitvec{x}&\unitvec{y}&\unitvec{z}\\\partial_x&\partial_y&\partial_z\\A_x&A_y&A_z\end{vmatrix}$ & 直角坐标旋度 \\
$\lap\phi = \divg\grad\phi$ & 拉普拉斯算子 \\
$\nabla\times(\nabla\times\vecA) = \grad(\divg\vecA) - \lap\vecA$ & 矢量三重积展开（BAC-CAB）应用于 $\nabla$ 算子 \\
$\int_V \divg\vecA\,dV = \oint_S \vecA\cdot d\vect{S}$ & 高斯散度定理 \\
$\int_S (\curl\vecA)\cdot d\vect{S} = \oint_C \vecA\cdot d\vect{l}$ & 斯托克斯定理 \\
$\delta^{(3)}(\vect{r}) = \frac{1}{(2\pi)^3}\int e^{i\vecK\cdot\vect{r}}d^3k$ & $\delta$函数傅里叶表示 \\
\bottomrule
\end{tabular}
\end{center}

\section{第2--4章：静电学}

\begin{center}
\begin{tabular}{@{}ll@{}}
\toprule
\textbf{公式} & \textbf{说明} \\
\midrule
$\vecE=\frac{1}{4\pi\eps_0}\frac{q\unitvec{r}}{r^2}$ & 点电荷电场（库仑） \\
$\divg\vecE=\rho/\eps_0$ & 高斯定律微分形式 \\
$\oint\vecE\cdot d\vect{S}=Q/\eps_0$ & 高斯定律积分形式 \\
$\vecE=-\grad\phi$ & 电场与电势关系 \\
$\lap\phi=-\rho/\eps_0$ & 泊松方程 \\
$\lap\phi=0$ & 拉普拉斯方程（无源区） \\
$\phi_P=\frac{1}{4\pi\eps_0}\frac{\vect{p}\cdot\unitvec{r}}{r^2}$ & 偶极子电势 \\
$\phi(\vect{r})=\frac{1}{4\pi\eps_0}\left[\frac{Q}{r}+\frac{\vect{p}\cdot\unitvec{r}}{r^2}+\cdots\right]$ & 多极展开 \\
$\vect{D}=\eps_0\vecE+\vect{P}=\eps\vecE$ & 各向同性介质 \\
$D_{1\perp}-D_{2\perp}=\sigma_f$ & 法向边界条件 \\
$E_{1\parallel}=E_{2\parallel}$ & 切向边界条件 \\
\bottomrule
\end{tabular}
\end{center}

\section{第5--6章：静磁学}

\begin{center}
\begin{tabular}{@{}ll@{}}
\toprule
\textbf{公式} & \textbf{说明} \\
\midrule
$d\vecB=\frac{\muz}{4\pi}\frac{I d\vect{l}\times\unitvec{r}}{r^2}$ & 毕奥-萨伐尔定律 \\
$\oint\vecB\cdot d\vect{l}=\muz I_{\text{enc}}$ & 安培环路定理 \\
$\vecB=\curl\vecA$ & 磁矢势定义 \\
$\vecA'=\vecA+\grad\chi$ & 规范变换 \\
$\nabla\cdot\vecA=0$ & 库仑规范 \\
$\vecB=\muz(\vect{H}+\vect{M})=\muz\mu_r\vect{H}$ & 线性介质 \\
$\vecJ_M=\nabla\times\vect{M}$ & 磁化电流 \\
$B_{1\perp}=B_{2\perp}$ & 法向边界条件 \\
$H_{1\parallel}-H_{2\parallel}=(\vect{K}_f\times\unitvec{n})$ & 切向边界条件（详见下方注） \\
\bottomrule
\end{tabular}
\end{center}

\noindent{\small\textbf{注}：磁学边界条件中 $\unitvec{n}$ 由区域 2 指向区域 1，
$\vect{K}_f$ 为自由面电流密度（A/m），与电学边界条件 $D_{1\perp}-D_{2\perp}=\sigma_f$ 配对使用。}

\section{第7章：麦克斯韦方程组}

\begin{center}
\begin{tabular}{@{}ll@{}}
\toprule
\textbf{公式} & \textbf{说明} \\
\midrule
$\divg\vecE=\rho/\eps_0$ & 高斯定律（电场有源） \\
$\divg\vecB=0$ & 磁场无源（无磁单极） \\
$\curl\vecE=-\partial\vecB/\partial t$ & 法拉第电磁感应定律 \\
$\curl\vecB=\muz\vecJ+\muz\eps_0\partial\vecE/\partial t$ & 安培-麦克斯韦定律 \\
$\vecE=-\grad\phi-\partial\vecA/\partial t$ & 时变电磁势 \\
$\lap\phi-\muz\eps_0\partial^2\phi/\partial t^2=-\rho/\eps_0$（即 $\Box\phi=-\rho/\eps_0$） & 标势达朗贝尔方程 \\
$\lap\vecA-\muz\eps_0\partial^2\vecA/\partial t^2=-\muz\vecJ$（即 $\Box\vecA=-\muz\vecJ$） & 矢势达朗贝尔方程 \\
$\vect{S}=\vecE\times\vect{H}$ & 坡印廷矢量 \\
$-\oint\vect{S}\cdot d\vecA=\int_V \vecJ\cdot\vecE\,dV+\frac{d}{dt}\int_V w\,dV$ & 坡印廷定理 \\
\bottomrule
\end{tabular}
\end{center}

\section{第8章：电磁波}

\begin{center}
\begin{tabular}{@{}ll@{}}
\toprule
\textbf{公式} & \textbf{说明} \\
\midrule
$\vecE=\vecE_0 e^{i(\vecK\cdot\vect{r}-\omega t)}$ & 平面波一般形式 \\
$k=\omega\sqrt{\muz\eps}=n\omega/c$ & 介质中波数 \\
$\vecB=\frac{1}{\omega}\vecK\times\vecE$ & 平面波 E-B 矢量关系 \\
$|\vecB|=\frac{k}{\omega}|\vecE|=\frac{1}{v}|\vecE|$ & 平面波 E-B 模长关系（$v$ 为相速） \\
$R=\left|\frac{n_1-n_2}{n_1+n_2}\right|^2$ & 正入射反射率 \\
$T=\frac{4n_1 n_2}{(n_1+n_2)^2}$ & 正入射透射率 \\
$\delta=\sqrt{\frac{2}{\mu\sigma\omega}}$ & 趋肤深度（详见下方注） \\
$\omega_c=v\pi\sqrt{(m/a)^2+(n/b)^2}$ & 矩形波导截止角频率（详见下方注） \\
$\theta_c=\arcsin(n_2/n_1)$ & 全反射临界角 \\
\bottomrule
\end{tabular}
\end{center}

\noindent{\small\textbf{本表适用条件注}：除注明外均为 SI 单位制。\textbf{(1)} 涉及介质的公式中 $c$ 指真空光速，
介质中须以相速 $v=1/\sqrt{\mu\eps}$（或 $v=c/n$）替代；平面波 E-B 关系中 $|\vecB|=|\vecE|/v$，
仅真空时才有 $|\vecB|=|\vecE|/c$。\textbf{(2)} 趋肤深度 $\delta=\sqrt{2/(\mu\sigma\omega)}$ 取材料磁导率
$\mu=\muz\mu_r$（非磁性材料 $\mu=\muz$），前提为良导体近似 $\sigma\gg\omega\eps$。
\textbf{(3)} 波导截止角频率中 $v$ 为介质中光速，真空填充时 $v=c$；TE$_{10}$ 模（$m=1,n=0$）化为 $f_c=v/2a$。}

\section{第9章：辐射}

\begin{center}
\begin{tabular}{@{}ll@{}}
\toprule
\textbf{公式} & \textbf{说明} \\
\midrule
$\phi(\vect{r},t)=\frac{1}{4\pi\eps_0}\int\frac{\rho(\vect{r}',t_r)}{|\vect{r}-\vect{r}'|}d^3r'$ & 推迟标势 \\
$\vecA(\vect{r},t)=\frac{\muz}{4\pi}\int\frac{\vecJ(\vect{r}',t_r)}{|\vect{r}-\vect{r}'|}d^3r'$ & 推迟矢势 \\
$t_r=t-\frac{|\vect{r}-\vect{r}'|}{c}$ & 推迟时间 \\
$\vecE=\frac{\muz\omega^2 p_0}{4\pi r}\sin\theta\,\unitvec{\theta}\,e^{i(kr-\omega t)}$ & 电偶极辐射远场 \\
$P=\frac{\muz p_0^2\omega^4}{12\pi c}$ & 电偶极辐射功率 \\
$P=\frac{\muz m_0^2\omega^4}{12\pi c^3}$ & 磁偶极辐射功率 \\
$P=\frac{q^2 a^2}{6\pi\eps_0 c^3}$ & 拉莫尔公式 \\
\bottomrule
\end{tabular}
\end{center}

\section{第10章：相对论电动力学}

\begin{center}
\begin{tabular}{@{}ll@{}}
\toprule
\textbf{公式} & \textbf{说明} \\
\midrule
$x^\mu=(ct,x,y,z)$ & 四维逆变坐标 \\
$\eta_{\mu\nu}=\text{diag}(1,-1,-1,-1)$ & 闵氏度规（本书采用） \\
$x_\mu=\eta_{\mu\nu}x^\nu$ & 协变坐标 \\
$A^\mu=(\phi/c,\vecA)$ & 四维电磁势 \\
$j^\mu=(c\rho,\vecJ)$ & 四维电流 \\
$F^{\mu\nu}=\partial^\mu A^\nu-\partial^\nu A^\mu$ & 电磁场张量 \\
$\partial_\mu F^{\mu\nu}=\muz j^\nu$ & 协变麦克斯韦方程（含源） \\
$\partial_\lambda F_{\mu\nu}+\partial_\mu F_{\nu\lambda}+\partial_\nu F_{\lambda\mu}=0$ & 协变麦克斯韦方程（齐次） \\
\bottomrule
\end{tabular}
\end{center}

\newpage

% ============================================================




% ============ 前言 ============
\chapter*{前言}
\addcontentsline{toc}{chapter}{前言}

电动力学（Electrodynamics）是研究电磁场的基本规律及其与物质相互作用的物理学分支。它是经典物理学的支柱之一，其理论体系不仅奠定了现代电工学、电子学、光学和通信技术的基础，更是狭义相对论诞生的摇篮，并为量子电动力学（QED）的发展提供了必要的理论框架。

\subsection*{本书的组织结构}

本书按照经典的教学顺序组织内容，从最简单、最基础的部分出发，逐步构建出完整的电磁理论大厦：

\begin{enumerate}
\item \textbf{第1章：数学预备知识}。电磁场是矢量场，描述它们的语言是矢量分析与场论。本章从最基本的矢量运算出发，逐步建立起梯度、散度、旋度的完整图像，并介绍高斯定理和斯托克斯定理。

\item \textbf{第2章：静电学基础}。从库仑定律出发，引入电场和电势的概念，建立高斯定理和环路定理，最终得到泊松方程。这是理解整个电磁理论的起点。

\item \textbf{第3章：电介质与静电边界条件}。补充极化强度、电位移矢量、介电常数的物理来源，讲解边界条件——这是边值问题求解的基石。

\item \textbf{第4章：静电学边值问题}。学习求解泊松/拉普拉斯方程的系统方法：镜像法、分离变量法和多极展开。

\item \textbf{第5章：静磁学}。研究恒定电流产生的磁场。毕奥-萨伐尔定律、安培环路定理和磁矢势构成了静磁学的核心框架。

\item \textbf{第6章：磁介质}。补充磁化强度、磁场强度、磁导率的物理来源，讲解磁学边界条件。

\item \textbf{第7章：麦克斯韦方程组}。将静电学和静磁学的方程推广到随时间变化的一般情形。法拉第电磁感应定律和位移电流的引入，标志着经典电磁理论的最终完成。

\item \textbf{第8章：电磁波的传播}。从麦克斯韦方程组导出电磁波方程，研究平面波在真空、介质和导体中的传播特性，以及反射、折射和波导等现象。

\item \textbf{第9章：电磁波的辐射}。研究加速运动的电荷如何产生电磁辐射。推迟势、电偶极辐射和天线理论构成了本章的核心内容。

\item \textbf{第10章：相对论电动力学}。将电动力学纳入狭义相对论框架，揭示电场和磁场是同一物理实体的不同表现——电磁场张量的统一描述。
\end{enumerate}

\subsection*{阅读建议}

\textbf{零基础读者的最小可行动路径}：第1章只看 \S1.1+\S1.3（矢量与梯度/散度/旋度）$\to$ 第2章（高斯定理、唯一性定理）$\to$ 第3章（电介质与边界条件）$\to$ 第4章（边值问题）$\to$ 第5章（静磁学）$\to$ 第7章（麦克斯韦方程组）$\to$ 第8章（电磁波）$\to$ 按需回看其余（第6章磁介质、第9章辐射、第10章相对论可最后补）。第1章开头的"快速通道"盒给出了同样路径与三档划分。

\begin{itemize}
\item \textbf{数学预备}：第1章是后续所有内容的基石，务必扎实掌握。特别要理解梯度、散度、旋度的几何意义。
\item \textbf{循序渐进}：静电学和静磁学是理解麦克斯韦方程组的必要阶梯，不可跳过。
\item \textbf{动手计算}：电动力学是"算"出来的，不是"看"出来的。每道例题都建议先独立尝试，再看解答。
\item \textbf{反复回顾}：各章之间存在深刻的内在联系，学完后续内容后回头再看，常会有新的领悟。
\item \textbf{注意误区}：书中专门设置了"常见误区"板块，这些都是学生最容易犯错的地方，务必重视。
\end{itemize}

\subsection*{难度标识说明}

\begin{itemize}
\item \textcolor{green}{\textbf{[基础]}}：必须掌握的核心内容
\item \textcolor{orange}{\textbf{[提高]}}：深化理解，建议完成
\item \textcolor{red}{\textbf{[挑战]}}：拓展视野，学有余力者尝试
\end{itemize}

\vspace{1em}
\noindent\rule{\textwidth}{0.4pt}
\vspace{1em}

让我们开始这段美妙的电动力学之旅。


% ============================================================
% 第1章：数学预备知识（优化版）
% ============================================================
\chapter{数学预备知识}
\epigraph{自然界最不可理解的，就是自然是可以被理解的。}{--- 阿尔伯特·爱因斯坦}

\hypertarget{ch1-goal}{}\begin{insight}{本章核心目标}{ch1-goal-v2}
掌握描述电磁场所需的数学语言——矢量分析与场论。理解梯度、散度、旋度的几何物理意义，以及它们如何通过积分定理相互联系。
\tcblower
\textbf{学习路径建议：}

\begin{itemize}[leftmargin=*]
\item[\textcolor{red}{$\blacksquare$}] \textbf{必学（必须掌握）}：\S1.1 标量/矢量/张量基础概念；\S1.3 梯度/散度/旋度的定义与计算；\S1.5 高斯定理与斯托克斯定理（积分定理是整个电动力学的基石！）
\item[\textcolor{orange}{$\blacksquare$}] \textbf{重要（建议掌握）}：\S1.7 正交曲线坐标系（球/柱坐标下的微分算子会在第2--4章频繁使用）；\S1.6 $\delta$函数的定义与基本性质
\item[\textcolor{green}{$\blacksquare$}] \textbf{可延后（第一次读可跳过）}：\S1.2.4 并矢与张量（第4/10章才需要）；\S1.6.2--\S1.6.3 $\delta$函数的复合性质与傅里叶表示（用到时再回看）
\end{itemize}

\textbf{本章学完你应该能够：}
\begin{itemize}[label=$\square$]
\item 不查表写出梯度、散度、旋度在直角/柱/球坐标下的表达式
\item 验证高斯定理和斯托克斯定理对简单矢量场成立
\item 计算点电荷密度分布的$\delta$函数表示
\item（可选）理解并矢$\vecA\vecB$的物理意义
\end{itemize}

\textbf{自学提示}：如果你被数学"卡住"了，\textbf{可以先跳过并矢和$\delta$函数的进阶内容}，直接进入第2章。后续章节用到时会提醒你"复习第1.X节"，届时带着具体需求回来补即可。
\end{insight}

\begin{tip}{快速通道：最小可行动阅读路径}{ch1-fast-track}
只想要"最快看到电磁波"？按下表走，其余内容用到时再回看：

\textbf{最小路径}：\S1.1 标量/矢量基础 $\to$ \S1.3 梯度/散度/旋度 $\to$ \S1.5 高斯/斯托克斯定理 $\to$ 直接进\textbf{第2章}。

\textbf{三档划分}：
\begin{itemize}[nosep]
\item \textcolor{green}{\textbf{必看}}：\S1.1、\S1.3、\S1.5（后续每一章都在用）
\item \textcolor{orange}{\textbf{用到再看}}：\S1.7 正交曲线坐标系（第2--4章计算需要时回看）、\S1.6.1--\S1.6.3 $\delta$ 函数基本性质（第2章点电荷要用）
\item \textcolor{red}{\textbf{可跳}}：\S1.2.4 并矢（第4/10章用到）、\S1.6.2--\S1.6.3 $\delta$ 函数复合性质与傅里叶表示（第8/9章用到）
\end{itemize}

被卡住就跳——本书会在用到跳过内容的地方标注"回顾 §X.Y"，届时带着具体需求回来，学习效率最高。
\end{tip}

电动力学处理的核心对象是电场 $\vecE(\vect{r},t)$ 和磁场 $\vecB(\vect{r},t)$，它们是随空间和时间变化的\textbf{矢量场}。为了精确描述这些场及其演化规律，我们需要一套系统的数学工具——矢量分析与场论。

\begin{history}{矢量分析的发展史}{vector-history}
矢量分析的现代形式主要归功于19世纪的两组独立工作。英国数学家\textbf{吉布斯}（J. Willard Gibbs）和物理学家\textbf{亥维赛}（Oliver Heaviside）在1880年代将哈密顿的四元数理论简化为我们今天使用的矢量分析。与此同时，德国数学家\textbf{格拉斯曼}（Hermann Grassmann）发展了自己的\textbf{扩张理论}（Ausdehnungslehre），虽然当时未被广泛接受，但其思想深刻影响了现代微分几何和张量分析。麦克斯韦本人在1873年的《电磁学通论》中尚未使用现代矢量记号，而是采用分量形式——这说明我们今天拥有的数学工具是几代数学家"打磨"的结果。
\end{history}

\section{标量、矢量与张量}

\begin{definition}{标量场}{scalar-v2}
只有大小、没有方向的物理量称为\textbf{标量（Scalar）}。温度 $T$、电势 $\phi$、电荷密度 $\rho$ 都是标量。在三维欧几里得空间中，标量场是一个从 $\mathbb{R}^3$ 到 $\mathbb{R}$ 的映射：
\[
\phi = \phi(x, y, z)
\]
标量在坐标旋转下\textbf{不变}——这是它的定义特征。
\end{definition}

\begin{definition}{矢量场}{vector-v2}
既有大小又有方向，并满足特定合成法则的物理量称为\textbf{矢量（Vector）}。电场 $\vecE$、磁场 $\vecB$、力 $\vect{F}$ 都是矢量。在直角坐标系中：
\[
\vecA = A_x \,\unitvec{x} + A_y \,\unitvec{y} + A_z \,\unitvec{z}
\]
其中 $\unitvec{x}, \unitvec{y}, \unitvec{z}$ 是沿 $x, y, z$ 轴正方向的单位矢量，满足右手定则：$\unitvec{x} \times \unitvec{y} = \unitvec{z}$。
\end{definition}

\begin{insight}{为什么电磁场必须是矢量场？}{why-vector-v2}
从物理角度看，空间中每一点的电磁场不仅有大小，还有明确的方向。如果只用一个标量来描述，根本无法表达这种方向性。更深刻地说，电磁场作为矢量场这一事实，与空间的三维各向同性密切相关：物理定律在旋转下不变，而矢量正是在旋转变换下按特定规则变换的量。在更高维的空间中，电磁场可能是张量场——这实际上启发了弦理论和卡鲁扎-克莱因理论的发展。
\end{insight}

\begin{definition}{张量}{tensor-v2}
\textbf{零阶张量}（标量）只有1个数，在坐标旋转下不变，如温度 $T$、电势 $\phi$。
\textbf{一阶张量}（矢量）有3个分量，如电场 $\vecE$、磁场 $\vecB$。
\textbf{二阶张量}有9个分量，可以初步理解为一个 $3\times 3$ 的矩阵。

电动力学中常见的二阶张量包括：
\begin{itemize}
\item \textbf{麦克斯韦应力张量} $T_{ij}$：描述电磁场对物体施加的应力（单位面积上的力）。例如，$T_{xx}$ 表示 $x$ 方向上垂直于 $x$ 轴的单位面积所受的力。
\item \textbf{各向异性介电张量} $\varepsilon_{ij}$：在晶体中，电位移 $\vect{D}$ 和电场 $\vecE$ 方向不同，关系为 $D_i = \sum_j \varepsilon_{ij} E_j$。对于各向同性介质，$\varepsilon_{ij} = \varepsilon \delta_{ij}$ 退化为一个标量。
\end{itemize}

二阶张量在坐标变换下按照特定规则变换：
\[
T'_{ij} = \sum_{k,l} R_{ik} R_{jl} T_{kl}
\]
其中 $R_{ij}$ 是旋转矩阵。

在第10章相对论电动力学中，电磁场将用\textbf{四维二阶反对称张量} $F^{\mu\nu}$ 统一描述。初学者遇到张量时，不妨先将其理解为"带指标的数"，随着学习的深入再逐步理解其几何意义。
\end{definition}

\section{矢量代数}

\subsection{点积（标量积）}

\begin{definition}{点积}{dot-product-v2}
两个矢量 $\vecA$ 和 $\vecB$ 的\textbf{点积}定义为：
\[
\vecA \cdot \vecB = |\vecA| \, |\vecB| \cos\theta = A_x B_x + A_y B_y + A_z B_z
\]
其中 $\theta$ 是两矢量之间的夹角。点积的结果是标量，满足交换律：$\vecA \cdot \vecB = \vecB \cdot \vecA$。
\end{definition}

\begin{insight}{点积的物理意义}{dot-meaning-v2}
从几何上看，$\vecA \cdot \vecB$ 等于 $\vecA$ 在 $\vecB$ 方向上的投影乘以 $\vecB$ 的模长。这在物理中无处不在：
\begin{itemize}
\item 功的定义 $W = \vect{F} \cdot \vect{d}$——只有力在运动方向上的分量才做功
\item 电场通量 $\Phi_E = \vecE \cdot d\vect{S}$——只有垂直于曲面的电场分量才贡献通量
\item 功率 $P = \vect{F} \cdot \vect{v}$——只有力在速度方向上的分量才产生功率
\item 坡印廷矢量 $\vect{S} = \vecE \times \vecB/\mu_0$ 的大小也与 $E$ 和 $B$ 的"配合"有关
\end{itemize}
\end{insight}

\begin{example}{点积的计算与应用}{dot-example-v2}
\difficulty{基础}
已知 $\vecA = 3\unitvec{x} + 4\unitvec{y}$，$\vecB = 2\unitvec{x} - \unitvec{y}$。求：(a) $\vecA \cdot \vecB$；(b) 两矢量夹角；(c) $\vecA$ 在 $\vecB$ 方向上的投影。
\tcblower
\textbf{解答：}

(a) $\vecA \cdot \vecB = 3 \times 2 + 4 \times (-1) = 6 - 4 = 2$

(b) $|\vecA| = \sqrt{3^2+4^2} = 5$，$|\vecB| = \sqrt{2^2+(-1)^2} = \sqrt{5}$

$\cos\theta = \dfrac{\vecA \cdot \vecB}{|\vecA|\,|\vecB|} = \dfrac{2}{5\sqrt{5}} \approx 0.179$

$\theta \approx \arccos(0.179) \approx 79.7°$

(c) 投影 $= \vecA \cdot \unitvec{B} = \dfrac{\vecA \cdot \vecB}{|\vecB|} = \dfrac{2}{\sqrt{5}} \approx 0.894$
\end{example}

\subsection{叉积（矢量积）}

\begin{definition}{叉积}{cross-product-v2}
两个矢量的\textbf{叉积}定义为：
\[
\vecA \times \vecB = |\vecA| \, |\vecB| \sin\theta \, \unitvec{n}
\]
其中 $\unitvec{n}$ 是垂直于 $\vecA$ 和 $\vecB$ 所在平面的单位矢量，方向由右手定则确定。

在分量形式下，可用行列式表示：
\[
\vecA \times \vecB = \begin{vmatrix}
\unitvec{x} & \unitvec{y} & \unitvec{z} \\
A_x & A_y & A_z \\
B_x & B_y & B_z
\end{vmatrix}
\]
展开后：
\[
\vecA \times \vecB = (A_y B_z - A_z B_y)\unitvec{x} + (A_z B_x - A_x B_z)\unitvec{y} + (A_x B_y - A_y B_x)\unitvec{z}
\]
\end{definition}

\begin{insight}{叉积与空间定向}{cross-orientation-v2}
叉积的方向依赖于坐标系的定向（右手系或左手系），这看似是个约定问题，但背后有深刻的物理意义。电磁学中的洛伦兹力 $\vect{F} = q\vect{v} \times \vecB$、安培力、法拉第电磁感应定律等都涉及叉积。物理定律本身必须是手征不变的——如果实验结果会因为用左手系还是右手系而改变，那就违背了空间各向同性的基本原理。

实际上，当同时考虑电磁力和弱相互作用时，自然界确实展现出微弱的手征不对称性（宇称不守恒，1957年李政道和杨振宁提出，吴健雄实验验证），但这超出了经典电动力学的范畴。
\end{insight}

\begin{warning}{叉积的常见错误}{cross-mistake}
\begin{enumerate}
\item 忘记叉积是\textbf{反交换}的：$\vecA \times \vecB = -\vecB \times \vecA$。很多学生习惯性地认为叉积也交换，导致符号错误。
\item 平行矢量的叉积为零：$\vecA \times \vecA = \vect{0}$。这在计算 $\nabla \times \nabla\phi$ 时很重要。
\item 叉积的模 $|\vecA \times \vecB| = |\vecA|\,|\vecB|\sin\theta$ 是以 $\vecA, \vecB$ 为邻边的平行四边形的面积。
\end{enumerate}
\end{warning}

\begin{example}{叉积的计算}{cross-example-v2}
\difficulty{基础}
求 $\vecA = \unitvec{x} + 2\unitvec{y} + 3\unitvec{z}$ 和 $\vecB = 2\unitvec{x} - \unitvec{y} + \unitvec{z}$ 的叉积，并验证结果与 $\vecA$ 和 $\vecB$ 都垂直。
\tcblower
\textbf{解答：}

\[
\vecA \times \vecB = \begin{vmatrix}
\unitvec{x} & \unitvec{y} & \unitvec{z} \\
1 & 2 & 3 \\
2 & -1 & 1
\end{vmatrix}
= (2\cdot1 - 3\cdot(-1))\unitvec{x} - (1\cdot1 - 3\cdot2)\unitvec{y} + (1\cdot(-1) - 2\cdot2)\unitvec{z}
\]
\[
= 5\unitvec{x} + 5\unitvec{y} - 5\unitvec{z}
\]

验证垂直性：
\begin{align*}
(\vecA \times \vecB) \cdot \vecA &= 5\cdot1 + 5\cdot2 + (-5)\cdot3 = 5 + 10 - 15 = 0 \\
(\vecA \times \vecB) \cdot \vecB &= 5\cdot2 + 5\cdot(-1) + (-5)\cdot1 = 10 - 5 - 5 = 0
\end{align*}
验证通过。
\end{example}

\subsection{混合积与BAC-CAB法则}

\textbf{标量三重积}：
\[
\vecA \cdot (\vecB \times \vect{C}) = \begin{vmatrix}
A_x & A_y & A_z \\
B_x & B_y & B_z \\
C_x & C_y & C_z
\end{vmatrix}
\]
几何意义：以 $\vecA, \vecB, \vect{C}$ 为棱的平行六面体的有向体积。满足循环置换不变性。

\begin{theorem}{BAC-CAB法则}{bac-cab-v2}
\[
{\vecA}\times({\vecB}\times{\vect{C}}) = \vecB(\vecA \cdot \vect{C}) - \vect{C}(\vecA \cdot \vecB)
\]
\end{theorem}

\begin{tip}{记忆技巧}{bac-cab-tip}
BAC-CAB法则的名字就是记忆口诀：外面是A，中间是B和C。结果的第一项是$\vecB$乘$\vecA\cdot\vect{C}$（即"BAC"），第二项是$\vect{C}$乘$\vecA\cdot\vecB$（即"CAB"）。注意第二项有负号！
\end{tip}

\begin{insight}{BAC-CAB在电动力学中的应用}{bac-cab-phy-v2}
BAC-CAB法则在推导电磁波方程时至关重要。例如：
\[
\curl(\curl\vecE) = \grad(\divg\vecE) - \lap\vecE
\]
这里 $\nabla$ 算子扮演双重角色——既对后面的场做微分，又与前面的场形成"点积"——正是BAC-CAB法则的类比。没有这一恒等式，就无法从麦克斯韦方程组导出波动方程。
\end{insight}

\subsection{并矢基础}

\begin{tip}{学习预警：本节可以跳过}{dyadic-skip-tip}
\textbf{第一次阅读可以跳过本节}，不影响后续章节的连贯性。并矢（dyadic）是张量的记号化表示，在本书中只在两处真正用到：
\begin{itemize}[nosep]
\item \textbf{第4章}多极展开中的电四极矩 $Q_{ij}$（本节"并矢的物理意义"中列出的应用之一）
\item \textbf{第10章}相对论电动力学中的四维张量 $F^{\mu\nu}$
\end{itemize}
建议的学习方式：先学完第 1--5 章的矢量内容，等第 4 章遇到 $Q_{ij}$ 或第 10 章遇到场张量时，\textbf{再回头读本节}，此时带着具体需求学习，效率最高。读的时候只需抓住一个核心：并矢 $\vecA\vecB$ 就是"带两个方向的数"，左方向来自 $\vecA$、右方向来自 $\vecB$，其余全是记号细节。
\end{tip}

\begin{tip}{极简通俗版：并矢是什么}{dyadic-plain-tip}
在接触任何指标记号之前，先建立直观图像：

\textbf{一个数}（标量）没有方向——$5$ 就是 $5$。

\textbf{一个矢量}带\textbf{一个}方向——$\vect{v}$ 说"朝某方向、多大"。

\textbf{一个并矢}带\textbf{两个}方向——$\vecA\vecB$ 说"左方向来自 $\vecA$、右方向来自 $\vecB$"。它回答的问题是："\textbf{沿} $\vecB$ 方向\textbf{看过去}，什么量沿 $\vecA$ 方向？"

最常用的物理例子：\textbf{麦克斯韦应力张量} $T_{ij} = \eps_0(E_iE_j - \frac{1}{2}\delta_{ij}E^2) + \frac{1}{\muz}(B_iB_j - \frac{1}{2}\delta_{ij}B^2)$ 就是一个并矢结构（完整式含磁场项，见第7章 §7.4.3）——它回答"沿 $\unitvec{n}$ 方向穿过一个面，电磁场沿哪个方向、传递多大的动量"。左指标 $i$ 是"传递方向"，右指标 $j$ 是"穿面方向"。两个方向缺一不可，这正是为什么需要两个指标的量（张量/并矢），而不是矢量。

\textbf{一句话}：矢量 = 1 个方向的数；并矢 = 2 个方向的数；点乘缩掉一个方向（$\vecA\vecB\cdot\vect{C} = \vecA(\vecB\cdot\vect{C})$，剩 1 个方向）。下面的正式定义只是把这套图像写成 $(\vecA\vecB)_{ij} = A_iB_j$ 的分量形式。
\end{tip}

\begin{definition}{并矢}{dyadic-v2}
两个矢量 $\vecA$ 和 $\vecB$ 的\textbf{并矢}记为 $\vecA\vecB$，是一个二阶张量，其分量形式为：
\[
(\vecA\vecB)_{ij} = A_i B_j
\]
\end{definition}

\begin{theorem}{并矢与矢量的点乘}{dyadic-dot-v2}
并矢与矢量的点乘满足：
\begin{align*}
\vecA\vecB \cdot \vect{C} &= \vecA(\vecB \cdot \vect{C}) \\
\vect{C} \cdot \vecA\vecB &= (\vect{C} \cdot \vecA)\vecB
\end{align*}
注意左点乘和右点乘的结果不同，前者沿 $\vecA$ 方向，后者沿 $\vecB$ 方向。
\end{theorem}

\begin{definition}{单位并矢}{unit-dyadic-v2}
\textbf{单位并矢}定义为：
\[
\bm{I} = \unitvec{x}\unitvec{x} + \unitvec{y}\unitvec{y} + \unitvec{z}\unitvec{z}
\]
它满足：
\[
\bm{I} \cdot \vecA = \vecA \cdot \bm{I} = \vecA
\]
对任意矢量 $\vecA$ 成立。单位并矢的分量形式为 $\delta_{ij}$（Kronecker delta）。
\end{definition}

\begin{insight}{并矢的物理意义}{dyadic-phy-v2}
并矢在电动力学中有重要应用：
\begin{itemize}
\item \textbf{麦克斯韦应力张量}可以表示为电磁场的并矢组合（详见第7章 §7.4.3）
\item \textbf{四极矩}：第4章多极展开中的电四极矩 $Q_{ij}$ 本质上是一个并矢结构，描述电荷分布偏离球对称的程度
\item 在各向异性介质中，介电张量 $\varepsilon_{ij}$ 也可以用并矢基展开
\end{itemize}
\end{insight}

\begin{definition}{并矢的缩并（迹）}{dyadic-trace-v2}
两个矢量 $\vecA$ 和 $\vecB$ 的并矢 $\bm{A}\bm{B}$（记法上同于 $\vecA\vecB$）的\textbf{迹}（trace）定义为其对角分量之和：
\[
\text{tr}(\bm{A}\bm{B}) = \sum_i (\bm{A}\bm{B})_{ii} = \sum_i A_i B_i = \bm{A}\cdot\bm{B}
\]
更一般地，两个并矢之间的\textbf{双点积}（double dot product，记为 ``$:$''）定义为对应分量相乘后求和：
\[
\bm{A}\bm{B} : \bm{I} = \sum_{i,j} A_i B_j \delta_{ij} = \sum_i A_i B_i = \bm{A}\cdot\bm{B}
\]
其中 $\bm{I}$ 为单位并矢。双点积的结果是一个标量，它将并矢的"张量特性"收缩为单一的数值。
\end{definition}

\begin{insight}{并矢缩并的物理意义}{dyadic-trace-phy-v2}
并矢的缩并（迹）在电动力学中频繁出现：
\begin{itemize}
\item \textbf{麦克斯韦应力张量} $T_{ij}$ 是一个二阶张量（并矢结构），通过适当的缩并可得到电磁场对物体施加的应力分量（详见第7章 §7.4.3）。
\item \textbf{电四极矩} $Q_{ij}$ 本质上是一个对称并矢，其迹的形式出现在多极展开中，描述电荷分布偏离球对称的程度。
\item 在连续介质力学中，应力张量的迹（即压力的三倍）与体积膨胀直接相关；在电磁场中，$T_{ij}$ 的迹与电磁场的能量密度密切相关。
\end{itemize}
\end{insight}

\begin{example}{计算并矢的迹}{dyadic-trace-example-v2}
\difficulty{基础}
计算 $\text{tr}(\unitvec{x}\unitvec{y})$，并验证它与单位并矢的双点积结果。
\tcblower
\textbf{解答：}

由并矢分量定义，$(\unitvec{x}\unitvec{y})_{ij} = (\unitvec{x})_i (\unitvec{y})_j$。其中 $(\unitvec{x})_i = \delta_{xi}$，$(\unitvec{y})_j = \delta_{yj}$，因此：
\[
(\unitvec{x}\unitvec{y})_{ij} = \delta_{xi} \delta_{yj}
\]
迹为对角分量之和：
\[
\text{tr}(\unitvec{x}\unitvec{y}) = \sum_{i=1}^{3} (\unitvec{x}\unitvec{y})_{ii} = \sum_{i=1}^{3} \delta_{xi} \delta_{yi} = 0
\]
因为 $x \neq y$，所以 $\delta_{xi}\delta_{yi} = 0$ 对所有 $i$ 成立。

另一方面，双点积：
\[
\unitvec{x}\unitvec{y} : \bm{I} = \unitvec{x}\cdot\unitvec{y} = 0
\]
两者结果一致。若改为 $\text{tr}(\unitvec{x}\unitvec{x})$，则：
\[
\text{tr}(\unitvec{x}\unitvec{x}) = \sum_i \delta_{xi}\delta_{xi} = 1
\]
这与 $\unitvec{x}\cdot\unitvec{x} = 1$ 一致。
\end{example}

\section{梯度、散度与旋度}

这是矢量分析中最核心的三个微分算子。从物理上看，它们分别回答了以下三个问题：
\begin{itemize}
\item \textbf{梯度}：标量场增长最快的方向和速率是什么？
\item \textbf{散度}：矢量场在某点是"源"还是"汇"？有多强？
\item \textbf{旋度}：矢量场在某点周围的"涡旋"有多强？转轴朝哪个方向？
\end{itemize}

\subsection{梯度（Gradient）}

\begin{definition}{梯度}{gradient-v2}
考虑标量场 $\phi(\vect{r}) = \phi(x, y, z)$。它在某一点沿不同方向的变化率可能不同。\textbf{梯度} $\grad\phi$ 是一个矢量，其方向指向 $\phi$\textbf{增长最快}的方向，其大小等于该方向上的方向导数：
\[
\grad\phi = \frac{\partial\phi}{\partial x}\unitvec{x} + \frac{\partial\phi}{\partial y}\unitvec{y} + \frac{\partial\phi}{\partial z}\unitvec{z}
\]
\textbf{Nabla算子}定义为：
\[
\grad = \unitvec{x}\frac{\partial}{\partial x} + \unitvec{y}\frac{\partial}{\partial y} + \unitvec{z}\frac{\partial}{\partial z}
\]
\end{definition}

\begin{insight}{梯度的几何直觉}{grad-intuition-v2}
想象你站在一座山上，高度函数为 $h(x,y)$。梯度 $\grad h$ 指向最陡的"上坡"方向，其大小等于坡度。如果你沿着等高线走（$h$ 不变的方向），梯度与该方向垂直——$\grad h \perp$（等高线）。

在电动力学中，电场 $\vecE = -\grad\phi$，这意味着正电荷会沿着电势\textbf{下降最快}的方向运动。电势的等高线就是等势面，电场总是垂直于等势面。
\end{insight}

% ---------- 图：梯度与等势面 ----------
\begin{center}
\begin{tikzpicture}[scale=1.0]
% 等高线（椭圆）
\foreach \a in {1,1.5,2,2.5,3} {
  \draw[thick,blue] (0,0) ellipse (\a*0.8 and \a*0.5);
}
\node[blue,font=\small,anchor=south] at (-1.8,1.7) {等势面 $\phi=C$};

% 梯度方向（垂直于等高线，指向增长方向）
\draw[->,thick,red] (0,0) -- (1.8,1.08) node[right,font=\small] {$\grad\phi$};
\draw[->,thick,red] (0,0) -- (-1.8,-1.08) node[left,font=\small] {$-\grad\phi$};

% 标注
\node[align=center,font=\small] at (0,-2.5) {$\grad\phi$ 垂直于等势面\\指向$\phi$增长最快的方向};
\end{tikzpicture}
\end{center}


标量场 $\phi$ 沿单位矢量 $\unitvec{n}$ 方向的\textbf{方向导数}为：
\[
\frac{\partial\phi}{\partial n} = \unitvec{n} \cdot \grad\phi
\]

\begin{example}{计算梯度与方向导数}{grad-example-v2}
\difficulty{基础}
求标量场 $\phi(x,y,z) = x^2 y + z^3$ 在点 $(1, 2, -1)$ 处的梯度，以及沿方向 $\unitvec{n} = \frac{1}{\sqrt{3}}(\unitvec{x} + \unitvec{y} + \unitvec{z})$ 的方向导数。
\tcblower
\textbf{解答：}

首先计算各偏导数：
\[
\frac{\partial\phi}{\partial x} = 2xy, \quad \frac{\partial\phi}{\partial y} = x^2, \quad \frac{\partial\phi}{\partial z} = 3z^2
\]

因此梯度为：
\[
\grad\phi = 2xy \,\unitvec{x} + x^2 \,\unitvec{y} + 3z^2 \,\unitvec{z}
\]

在点 $(1, 2, -1)$ 处：
\[
\grad\phi\big|_{(1,2,-1)} = 4\unitvec{x} + \unitvec{y} + 3\unitvec{z}
\]

方向导数为：
\[
\frac{\partial\phi}{\partial n} = \unitvec{n} \cdot \grad\phi = \frac{1}{\sqrt{3}}(4 + 1 + 3) = \frac{8}{\sqrt{3}} \approx 4.62
\]
\end{example}

\subsection{散度（Divergence）}

\begin{definition}{散度}{divergence-v2}
散度作用于矢量场，结果为标量场：
\[
\divg\vecA = \frac{\partial A_x}{\partial x} + \frac{\partial A_y}{\partial y} + \frac{\partial A_z}{\partial z}
\]
\end{definition}

\begin{insight}{散度的物理图像}{div-intuition-v2}
散度描述矢量场在某一点的"发散程度"。想象一个微小的闭合曲面包围某一点：
\begin{itemize}
\item 若 $\divg\vecA > 0$：从曲面内流出的通量大于流入的通量——该点是\textbf{源}（如正电荷产生的发散电场）
\item 若 $\divg\vecA < 0$：流入大于流出——该点是\textbf{汇}（如负电荷）
\item 若处处 $\divg\vecA = 0$：矢量场是\textbf{无源场}（\textbf{螺线管场}）——场线没有起点和终点。注意：没有起点和终点\textbf{并不等于}场线一定闭合成环，例如 $\vecA=\unitvec{z}$ 的场线是无限长直线，而 $\vecA=-y\unitvec{x}+x\unitvec{y}$ 的场线才是闭合圆
\end{itemize}

静电学中的高斯定理告诉我们 $\divg\vecE = \rho/\eps_0$：电荷密度 $\rho$ 正是电场 $\vecE$ 的"源"！这个方程的微分形式比积分形式更深刻地揭示了电场与电荷之间的\textbf{局部}联系。
\end{insight}

% ---------- 图：散度的源与汇 ----------
\begin{center}
\begin{tikzpicture}[scale=0.85]
% 源 (div > 0)
\begin{scope}[shift={(-5,0)}]
\draw[thick] (0,0) circle (1.5);
\foreach \a in {0,45,90,135,180,225,270,315} {
  \draw[->,thick,red] (\a:0.3) -- (\a:1.3);
}
\fill[red] (0,0) circle (3pt);
\node at (0,-2.1) {\textbf{(a) 源}：$\divg\vecA>0$};
\node[align=center,font=\small] at (0,-2.9) {场线从点发出\\净流出};
\end{scope}

% 汇 (div < 0)
\begin{scope}[shift={(0,0)}]
\draw[thick] (0,0) circle (1.5);
\foreach \a in {0,45,90,135,180,225,270,315} {
  \draw[->,thick,blue] (\a:1.3) -- (\a:0.3);
}
\fill[blue] (0,0) circle (3pt);
\node at (0,-2.1) {\textbf{(b) 汇}：$\divg\vecA<0$};
\node[align=center,font=\small] at (0,-2.9) {场线向点会聚\\净流入};
\end{scope}

% 无源 (div = 0)
\begin{scope}[shift={(5,0)}]
\draw[thick] (0,0) circle (1.5);
\draw[->,thick,green!60!black] (-1.3,0.5) -- (1.3,0.5);
\draw[->,thick,green!60!black] (-1.3,-0.5) -- (1.3,-0.5);
\node at (0,-2.1) {\textbf{(c) 无源}：$\divg\vecA=0$};
\node[align=center,font=\small] at (0,-2.9) {场线平行穿过\\净通量为零};
\end{scope}
\end{tikzpicture}
\end{center}


\subsection{旋度（Curl）}

\begin{definition}{旋度}{curl-v2}
旋度作用于矢量场，结果为矢量场：
\[
\curl\vecA = \begin{vmatrix}
\unitvec{x} & \unitvec{y} & \unitvec{z} \\
\partial/\partial x & \partial/\partial y & \partial/\partial z \\
A_x & A_y & A_z
\end{vmatrix}
\]
展开得：
\[
\curl\vecA = \left(\frac{\partial A_z}{\partial y} - \frac{\partial A_y}{\partial z}\right)\unitvec{x} + \left(\frac{\partial A_x}{\partial z} - \frac{\partial A_z}{\partial x}\right)\unitvec{y} + \left(\frac{\partial A_y}{\partial x} - \frac{\partial A_x}{\partial y}\right)\unitvec{z}
\]
\end{definition}

\begin{insight}{旋度的物理图像}{curl-intuition-v2}
旋度描述矢量场在某一点的"涡旋程度"。想象将一个微小的叶轮放入场中：
\begin{itemize}
\item 若 $\curl\vecA \neq 0$：叶轮会旋转，其转轴方向就是 $\curl\vecA$ 的方向，转速正比于 $|\curl\vecA|$
\item 若在某区域处处 $\curl\vecA = 0$：矢量场是\textbf{无旋场}或\textbf{保守场}——叶轮不会旋转
\end{itemize}

在流体力学中，河流中的漩涡正是 $\curl\vect{v} \neq 0$ 的区域。在电磁学中：
\begin{itemize}
\item 静磁学：$\curl\vecB = \muz\vecJ$——电流是磁场的"涡旋源"
\item 法拉第定律：$\curl\vecE = -\partial\vecB/\partial t$——变化的磁场会在周围空间中激发涡旋电场
\end{itemize}
\end{insight}

% ---------- 图：旋度与微小叶轮 ----------
\begin{center}
\begin{tikzpicture}[scale=0.9]
% 有旋场 (curl != 0)
\begin{scope}[shift={(-3.5,0)}]
\foreach \a in {0,60,120,180,240,300} {
  \draw[->,thick,red] ({\a-15}:0.8) arc (\a-15:\a+15:0.8);
}
\draw[thick] (0,0) circle (0.6);
\draw[->,thick,blue] (0,0) -- (0,1.2) node[above] {$\curl\vecA$};
\node[font=\small] at (0,-1.0) {叶轮};
\node at (0,-1.7) {\textbf{(a) 有旋场}};
\node[align=center,font=\small] at (0,-2.4) {$\curl\vecA\neq\vect{0}$\\叶轮旋转};
\end{scope}

% 无旋场 (curl = 0)
\begin{scope}[shift={(3.5,0)}]
\foreach \y in {-1,-0.5,0,0.5,1} {
  \draw[->,thick,green!60!black] (-1.5,\y) -- (1.5,\y);
}
\draw[thick] (0,0) circle (0.6);
\node[font=\small,fill=white,inner sep=1pt] at (0,0) {叶轮};
\node at (0,-1.7) {\textbf{(b) 无旋场}};
\node[align=center,font=\small] at (0,-2.4) {$\curl\vecA=\vect{0}$\\叶轮不转};
\end{scope}
\end{tikzpicture}
\end{center}


\begin{example}{计算散度与旋度}{div-curl-example-v2}
\difficulty{基础}
求矢量场 $\vecA = x^2 \,\unitvec{x} + yz \,\unitvec{y} + xy \,\unitvec{z}$ 的散度和旋度。
\tcblower
\textbf{解答：}

\textbf{散度}：
\[
\divg\vecA = \frac{\partial(x^2)}{\partial x} + \frac{\partial(yz)}{\partial y} + \frac{\partial(xy)}{\partial z} = 2x + z + 0 = 2x + z
\]

\textbf{旋度}：
\[
\curl\vecA = \begin{vmatrix}
\unitvec{x} & \unitvec{y} & \unitvec{z} \\
\partial/\partial x & \partial/\partial y & \partial/\partial z \\
x^2 & yz & xy
\end{vmatrix}
\]

计算各分量：
\begin{align*}
(\curl\vecA)_x &= \frac{\partial(xy)}{\partial y} - \frac{\partial(yz)}{\partial z} = x - y \\
(\curl\vecA)_y &= \frac{\partial(x^2)}{\partial z} - \frac{\partial(xy)}{\partial x} = 0 - y = -y \\
(\curl\vecA)_z &= \frac{\partial(yz)}{\partial x} - \frac{\partial(x^2)}{\partial y} = 0 - 0 = 0
\end{align*}

因此：
\[
\curl\vecA = (x-y)\unitvec{x} - y\unitvec{y}
\]
\end{example}

\begin{example}{验证恒等式}{identity-verify-v2}
\difficulty{提高}
对例题中的 $\vecA$，验证 $\divg(\curl\vecA) = 0$。
\tcblower
\textbf{解答：}

$\curl\vecA = (x-y)\unitvec{x} - y\unitvec{y}$

$\divg(\curl\vecA) = \dfrac{\partial(x-y)}{\partial x} + \dfrac{\partial(-y)}{\partial y} + \dfrac{\partial(0)}{\partial z} = 1 - 1 + 0 = 0$

验证通过。这是一个普遍成立的恒等式——任何矢量场的旋度场都是无源的。
\end{example}

\section{矢量微分算子的重要恒等式}

以下恒等式在电动力学推导中频繁出现，务必熟记并理解其物理含义。

\begin{theorem}{梯度的旋度为零}{grad-curl-zero-v2}
\[
\curl(\grad\phi) = \vect{0}
\]
对任意标量场 $\phi$ 成立。几何上，梯度场没有"涡旋"——它总是从低值指向高值，不可能形成闭合回路。
\end{theorem}

\begin{theorem}{散度的旋度为零}{curl-div-zero-v2}
\[
\divg(\curl\vecA) = 0
\]
对任意矢量场 $\vecA$ 成立。几何上，旋度场是天然的无源场——它的场线没有起点和终点（旋度场本身处处无散，不可能在某处"发出"或"终止"）。
\end{theorem}

\begin{theorem}{梯度的散度（拉普拉斯算子）}{laplacian-v2}
\[
\divg(\grad\phi) = \lap\phi = \frac{\partial^2\phi}{\partial x^2} + \frac{\partial^2\phi}{\partial y^2} + \frac{\partial^2\phi}{\partial z^2}
\]
$\lap$ 称为\textbf{拉普拉斯算子（Laplacian）}。
\end{theorem}

\begin{theorem}{旋度的旋度}{curl-curl-v2}
\[
\curl(\curl\vecA) = \grad(\divg\vecA) - \lap\vecA
\]
这是BAC-CAB法则的"算子版本"。对矢量场 $\lap\vecA$ 作用时，理解为对各分量分别作用：
\[
\lap\vecA = (\lap A_x)\unitvec{x} + (\lap A_y)\unitvec{y} + (\lap A_z)\unitvec{z}
\]
\end{theorem}

\begin{theorem}{时间和空间微分的可交换性}{time-space-commute-v2}
对任意标量场 $f(\vect{r},t)$ 和矢量场 $\vecA(\vect{r},t)$，时间偏导与空间微分算子可交换：
\[
\nabla\times(\partial_t \vecA) = \partial_t(\nabla\times\vecA), \qquad
\nabla(\partial_t f) = \partial_t(\nabla f)
\]
其中 $\partial_t \equiv \partial/\partial t$。这些恒等式成立的前提是场量 $f$ 和 $\vecA$ 对时空变量具有足够的光滑性（通常物理场均满足），使得混合偏导数的求导顺序可交换。
\end{theorem}

\begin{insight}{时间-空间微分可交换性的物理意义}{time-space-commute-phy-v2}
以上定理表面上是数学恒等式，但它在电动力学中具有核心重要性：
\begin{itemize}
\item \textbf{第7章推导麦克斯韦方程组时频繁使用}：例如从法拉第定律的微分形式推导波动方程时，需要交换 $\nabla\times$ 与 $\partial_t$ 的顺序。
\item 它保证了电磁场随时间的演化与空间分布的耦合是自洽的——场在空间某点的"涡旋"（旋度）随时间的变化率，等于该点场量时间变化率的"涡旋"。
\item 在相对论电动力学（第10章）中，这一性质升华为四维协变导数的对称性，是电磁场张量 $F^{\mu\nu}$ 满足比安基恒等式的数学基础。
\end{itemize}
初学者可能在后续推导中对这些交换步骤感到"理所当然"，但应意识到它的成立依赖于场量的光滑性，在存在奇点或突变边界（如理想导体表面）时需要特别小心。
\end{insight}

\begin{insight}{拉普拉斯算子的物理意义}{laplacian-meaning-v2}
拉普拉斯算子 $\lap\phi$ 有深刻的物理含义。在一维情况下，$\lap\phi = \diff^2\phi/\diff x^2$ 衡量函数曲线的"弯曲程度"。在三维中，它可以写成：
\[
\lap\phi = \lim_{V\to 0} \frac{1}{V} \oint_S \grad\phi \cdot \diff\vect{S}
\]
即 $\phi$ 的梯度通过包围该点的无穷小球面的通量密度。如果某点 $\lap\phi > 0$，说明该点的值低于周围平均值（如热源处的温度分布）；如果 $\lap\phi < 0$，则高于周围平均值（如"山峰"）。

在静电学中，泊松方程 $\lap\phi = -\rho/\eps_0$ 告诉我们：正电荷所在处 $\lap\phi<0$，由上面的判据可知该点电势\textbf{高于}周围平均值（$\phi$ 呈"山峰"形）；负电荷处 $\lap\phi>0$，电势低于周围平均值（呈"山谷"形）。这正是点电荷电势 $\phi = q/4\pi\eps_0 r$ 的图像：$q>0$ 时越靠近电荷电势越高。
\end{insight}

\begin{example}{验证旋度的旋度恒等式}{curl-curl-verify-v2}
\difficulty{提高}
计算 $\nabla\times(\nabla\times(r^2\unitvec{z}))$，并验证它等于 $\nabla(\nabla\cdot(r^2\unitvec{z})) - \nabla^2(r^2\unitvec{z})$。
\tcblower
\textbf{解答：}

设 $\vecA = r^2\unitvec{z} = (x^2+y^2)\unitvec{z}$。

\textbf{方法一：直接计算 $\nabla\times(\nabla\times\vecA)$}

先计算 $\nabla\times\vecA$：
\[
\nabla\times\vecA = \begin{vmatrix}
\unitvec{x} & \unitvec{y} & \unitvec{z} \\
\partial/\partial x & \partial/\partial y & \partial/\partial z \\
0 & 0 & x^2+y^2
\end{vmatrix}
= 2y\unitvec{x} - 2x\unitvec{y}
\]

再计算 $\nabla\times(\nabla\times\vecA)$：
\[
\nabla\times(2y\unitvec{x} - 2x\unitvec{y}) = \begin{vmatrix}
\unitvec{x} & \unitvec{y} & \unitvec{z} \\
\partial/\partial x & \partial/\partial y & \partial/\partial z \\
2y & -2x & 0
\end{vmatrix}
= -4\unitvec{z}
\]

\textbf{方法二：使用恒等式 $\nabla(\nabla\cdot\vecA) - \nabla^2\vecA$}

$\nabla\cdot\vecA = \partial(x^2+y^2)/\partial z = 0$，所以 $\nabla(\nabla\cdot\vecA) = \vect{0}$。

$\nabla^2\vecA = \nabla^2(r^2)\unitvec{z}$（因为 $\unitvec{z}$ 是常矢量）：
\[
\nabla^2(r^2) = \frac{\partial^2(x^2+y^2)}{\partial x^2} + \frac{\partial^2(x^2+y^2)}{\partial y^2} + \frac{\partial^2(x^2+y^2)}{\partial z^2} = 2 + 2 + 0 = 4
\]

因此：
\[
\nabla(\nabla\cdot\vecA) - \nabla^2\vecA = 0 - 4\unitvec{z} = -4\unitvec{z}
\]

两种方法结果一致，验证了旋度的旋度恒等式。
\end{example}

\begin{example}{验证散度的乘积法则}{div-product-verify-v2}
\difficulty{基础}
对 $f(r) = r^2$，$\vecA = \unitvec{z}$，验证乘积法则：
\[
\nabla\cdot(f\vecA) = \vecA\cdot(\nabla f) + f(\nabla\cdot\vecA)
\]
\tcblower
\textbf{解答：}

\textbf{左边}：$f\vecA = r^2\unitvec{z} = (x^2+y^2)\unitvec{z}$
\[
\nabla\cdot(f\vecA) = \frac{\partial(0)}{\partial x} + \frac{\partial(0)}{\partial y} + \frac{\partial(x^2+y^2)}{\partial z} = 0
\]

\textbf{右边}：
\begin{align*}
\nabla f &= \nabla(r^2) = 2x\unitvec{x} + 2y\unitvec{y} \\
\vecA\cdot(\nabla f) &= \unitvec{z}\cdot(2x\unitvec{x} + 2y\unitvec{y}) = 0 \\
\nabla\cdot\vecA &= \nabla\cdot\unitvec{z} = 0
\end{align*}

因此：
\[
\vecA\cdot(\nabla f) + f(\nabla\cdot\vecA) = 0 + r^2\times 0 = 0
\]

左边 = 右边 = 0，验证通过。
\end{example}
\section{矢量积分定理}

梯度、散度、旋度是\textbf{局部}的微分算子。积分定理建立了这些局部量与\textbf{全局}行为之间的联系。

\subsection{高斯散度定理（Gauss's Divergence Theorem）}

\begin{theorem}{高斯散度定理}{gauss-theorem-v2}
设 $V$ 是空间中的一个闭合区域，其边界为闭合曲面 $S$。对于任意光滑矢量场 $\vecA$：
\[
\oint_S \vecA \cdot \diff\vect{S} = \int_V (\divg\vecA) \, \diff V
\]
其中 $\diff\vect{S} = \unitvec{n} \, \diff S$，$\unitvec{n}$ 是曲面 $S$ 的\textbf{外法向}单位矢量。
\end{theorem}

\begin{insight}{高斯定理的物理含义}{gauss-insight-v2}
高斯定理建立了"面积分"与"体积分"之间的联系。从物理上看，它说：矢量场穿过闭合曲面的总通量，等于其散度在曲面所围体积内的积分。

想象流体流动：\textbf{流出}一个闭合区域的总流量，等于区域内所有"源"的强度之和。如果区域内没有源（$\divg\vecA = 0$），那么流出多少就流入多少，净通量为零。

在静电学中，高斯定理 $\displaystyle\oint_S \vecE \cdot \diff\vect{S} = Q_{\text{enc}}/\eps_0$ 告诉我们：穿过任意闭合曲面的电通量只与曲面内的净电荷有关——这是电磁学中最强大的计算工具之一。
\end{insight}

\subsection{斯托克斯定理（Stokes' Theorem）}

\begin{theorem}{斯托克斯定理}{stokes-theorem-v2}
设 $S$ 是一个以闭合曲线 $C$ 为边界的曲面。对于任意光滑矢量场 $\vecA$：
\[
\oint_C \vecA \cdot \diff\vect{l} = \int_S (\curl\vecA) \cdot \diff\vect{S}
\]
其中 $\diff\vect{l}$ 沿曲线 $C$ 的切向，$\diff\vect{S}$ 的方向与 $C$ 的绕行方向满足\textbf{右手定则}。
\end{theorem}

\begin{insight}{斯托克斯定理的物理含义}{stokes-insight-v2}
斯托克斯定理建立了"线积分"与"面积分"之间的联系。它说：矢量场沿闭合曲线的\textbf{环流}，等于其\textbf{旋度}穿过以该曲线为边界的任意曲面的\textbf{通量}。

想象将一张纸放入水中漩涡：沿纸边界的环流，等于穿过纸面的涡量总和。如果在\textbf{单连通}区域内流体处处无旋（$\curl\vecA = 0$），那么沿区域内任何闭合曲线的环流都为零——这正是保守场的特征。（区域有"洞"时该结论不成立，见本节末的常见误区。）

在电磁学中，法拉第定律的积分形式 $\displaystyle\oint_C \vecE \cdot \diff\vect{l} = -\dfrac{\diff\Phi_B}{\diff t}$ 就是斯托克斯定理的直接应用：变化的磁通量会在周围空间中激发涡旋电场。
\end{insight}

\begin{example}{利用高斯定理}{gauss-example-v2}
\difficulty{基础}
证明半径为 $R$ 的球面上，位置矢量 $\vect{r}$ 的通量为 $4\pi R^3$。
\tcblower
\textbf{解答：}

由高斯定理：
\[
\oint_S \vect{r} \cdot \diff\vect{S} = \int_V (\divg\vect{r}) \, \diff V
\]

首先计算 $\divg\vect{r}$：
\[
\vect{r} = x\unitvec{x} + y\unitvec{y} + z\unitvec{z} \quad\Rightarrow\quad \divg\vect{r} = \frac{\partial x}{\partial x} + \frac{\partial y}{\partial y} + \frac{\partial z}{\partial z} = 3
\]

因此：
\[
\oint_S \vect{r} \cdot \diff\vect{S} = \int_V 3 \, \diff V = 3 \times \frac{4}{3}\pi R^3 = 4\pi R^3
\]

\textbf{直接验证}：在球面上，$\vect{r}$ 与 $\diff\vect{S}$ 同向，$\vect{r} \cdot \diff\vect{S} = R \, \diff S$，积分得 $R \times 4\pi R^2 = 4\pi R^3$。两种方法结果一致，验证了高斯定理。
\end{example}

\begin{example}{利用斯托克斯定理}{stokes-example-v2}
\difficulty{提高}
设 $\vecA = -y\unitvec{x} + x\unitvec{y}$。计算沿 $xy$ 平面上半径为 $R$ 的圆的环流，分别用直接线积分和斯托克斯定理验证。
\tcblower
\textbf{解答：}

\textbf{方法一：直接线积分}

参数化圆周：$x = R\cos\theta$，$y = R\sin\theta$，$\theta \in [0, 2\pi]$。

$\diff\vect{l} = (-R\sin\theta\,\diff\theta)\unitvec{x} + (R\cos\theta\,\diff\theta)\unitvec{y}$

$\vecA = -R\sin\theta\unitvec{x} + R\cos\theta\unitvec{y}$

$\vecA \cdot \diff\vect{l} = R^2\sin^2\theta\,\diff\theta + R^2\cos^2\theta\,\diff\theta = R^2\,\diff\theta$

$\displaystyle\oint_C \vecA \cdot \diff\vect{l} = \int_0^{2\pi} R^2\,\diff\theta = 2\pi R^2$

\textbf{方法二：斯托克斯定理}

$\curl\vecA = (0 - 0)\unitvec{x} + (0 - 0)\unitvec{y} + (1 - (-1))\unitvec{z} = 2\unitvec{z}$

$\displaystyle\int_S (\curl\vecA) \cdot \diff\vect{S} = 2 \times \pi R^2 = 2\pi R^2$

两种方法结果一致。
\end{example}

\section{Dirac \texorpdfstring{$\delta$}{δ} 函数}

\subsection{定义与基本性质}

\begin{definition}{Dirac $\delta$ 函数}{delta-function-v2}
Dirac $\delta$ 函数（严格说是广义函数或分布）是描述点源（如点电荷）的数学工具。它的定义性质为：
\begin{enumerate}
\item $\delta(x-a) = 0$，当 $x \neq a$
\item $\displaystyle\int_{-\infty}^{\infty} \delta(x-a) \, \diff x = 1$
\end{enumerate}
最核心的性质是\textbf{筛选性质}：
\[
\int_{-\infty}^{\infty} f(x) \, \delta(x-a) \, \diff x = f(a)
\]
\end{definition}

\begin{insight}{$\delta$ 函数的物理图像}{delta-intuition-v2}
$\delta$ 函数可以看作一系列普通函数在极限下的产物。例如，高斯函数族：
\[
\delta_\eps(x) = \frac{1}{\sqrt{2\pi}\eps} \ee^{-x^2/(2\eps^2)}
\]
当 $\eps \to 0$ 时，函数变得越来越窄、越来越高，但始终保持总面积为1。在极限下，它"收缩"成一个在 $x=0$ 处无穷大、在其他地方为零、积分为1的"尖峰"。

在物理上，点电荷的电荷密度、点质量的密度分布、冲量力的作用等都天然地用 $\delta$ 函数描述。它是从连续分布过渡到理想化点源的桥梁。
\end{insight}

% ---------- 图：δ函数的高斯逼近序列 ----------
\begin{center}
\begin{tikzpicture}[scale=0.9]
% ε=0.5
\begin{scope}[shift={(-4.5,0)}]
\draw[->] (-1.5,0) -- (1.5,0) node[right] {$x$};
\draw[->] (0,0) -- (0,2.5);
\draw[thick,blue,domain=-1.5:1.5,samples=50] plot(\x,{2*exp(-\x*\x/0.5)});
\node[font=\small] at (0,-0.5) {$\eps=0.5$};
\node[font=\small] at (0,-1) {宽而矮};
\end{scope}

% ε=0.2
\begin{scope}[shift={(-1.5,0)}]
\draw[->] (-1.5,0) -- (1.5,0) node[right] {$x$};
\draw[->] (0,0) -- (0,2.5);
\draw[thick,red,domain=-1.5:1.5,samples=50] plot(\x,{2*exp(-\x*\x/0.08)});
\node[font=\small] at (0,-0.5) {$\eps=0.2$};
\node[font=\small] at (0,-1) {变窄变高};
\end{scope}

% ε→0
\begin{scope}[shift={(1.5,0)}]
\draw[->] (-1.5,0) -- (1.5,0) node[right] {$x$};
\draw[->] (0,0) -- (0,2.5);
\draw[thick,red] (0,0) -- (0,2.5);
\fill[red] (0,2.5) circle (2pt);
\node[font=\small] at (0,-0.5) {$\eps\to 0$};
\node[font=\small] at (0,-1) {无限尖峰};
\end{scope}

% 面积守恒说明
\begin{scope}[shift={(4.5,0)}]
\node[align=center,font=\small] at (0,1.5) {面积恒为1};
\draw[<->,thick] (-0.5,0.8) -- (0.5,0.8);
\node[align=center,font=\small] at (0,0) {$\int_{-\infty}^{\infty}\delta_\eps(x)\diff x=1$};
\node[align=center,font=\small] at (0,-1) {对所有$\eps$成立};
\end{scope}

\node[align=center,font=\small] at (0,-2.5) {$\delta(x)=\lim_{\eps\to 0}\delta_\eps(x)$\\宽度$\to 0$，高度$\to\infty$，面积$=1$};
\end{tikzpicture}
\end{center}


三维 $\delta$ 函数定义为：$\delta^{(3)}(\vect{r} - \vect{r}') = \delta(x-x')\delta(y-y')\delta(z-z')$

\textbf{三维筛选性质}：$\displaystyle\int f(\vect{r}') \, \delta^{(3)}(\vect{r} - \vect{r}') \, \diff^3r' = f(\vect{r})$

\subsection{\texorpdfstring{$\delta$}{δ}函数的缩放性质与复合函数}

\begin{theorem}{$\delta$函数的缩放公式}{delta-scaling}
对于非零常数 $a$：
\[
\delta(ax) = \frac{1}{|a|}\delta(x)
\]
更一般地，若 $f(x)$ 只有\textbf{单根} $x_i$（即 $f(x_i)=0$ 且 $f'(x_i)\neq 0$），则：
\[
\delta(f(x)) = \sum_i \frac{\delta(x-x_i)}{|f'(x_i)|}
\]
\end{theorem}

\begin{example}{$\delta$函数缩放的应用}{delta-scaling-example}
\difficulty{基础}
求 $\delta(3x-2)$ 的表达式。
\tcblower
\textbf{解答：}令 $f(x)=3x-2$，则 $f'(x)=3$，零点 $x_0=2/3$。
由缩放公式：
\[
\delta(3x-2) = \frac{1}{3}\delta\left(x-\frac{2}{3}\right)
\]
\end{example}

\subsection{\texorpdfstring{$\delta$}{δ} 函数与点电荷}

位于 $\vect{r}'$ 处的点电荷 $q$ 的电荷密度可写为：
\[
\rho(\vect{r}) = q \, \delta^{(3)}(\vect{r} - \vect{r}')
\]

验证总电荷：$Q = \int \rho(\vect{r}) \, \diff^3r = q \int \delta^{(3)}(\vect{r} - \vect{r}') \, \diff^3r = q$。

\subsubsection{一维 $\delta$ 函数}

线电荷密度可以用一维 $\delta$ 函数表示。位于 $x = x_0$ 处的线电荷，其密度为：
\[
\lambda(x) = \lambda_0 \, \delta(x - x_0)
\]
验证总电荷：
\[
Q = \int_{-\infty}^{\infty} \lambda(x) \, \diff x = \lambda_0 \int_{-\infty}^{\infty} \delta(x - x_0) \, \diff x = \lambda_0
\]

\subsubsection{二维 $\delta$ 函数}

面电荷密度可以用二维 $\delta$ 函数表示：
\[
\sigma(x, y) = \sigma_0 \, \delta^{(2)}(\vect{r} - \vect{r}_0)
\]
其中 $\delta^{(2)}(\vect{r} - \vect{r}_0) = \delta(x - x_0)\delta(y - y_0)$。

\subsubsection{$\delta$ 函数的傅里叶积分表示}

\begin{theorem}{$\delta$ 函数的傅里叶积分表示}{delta-fourier-v2}
\[
\delta(x) = \frac{1}{2\pi} \int_{-\infty}^{\infty} e^{ikx} \, \diff k
\]
\end{theorem}

\begin{tip}{重要提示}{fourier-hint-v2}
此表示在第9章推迟势、第8章平面波展开中经常使用。它体现了 $\delta$ 函数"所有频率分量等权叠加"的特性。
\end{tip}

\subsubsection{$\delta$ 函数的缩放性质}

\begin{theorem}{$\delta$ 函数的缩放性质}{delta-scaling-v2}
设 $a$ 为非零实常数（$a \neq 0$），则：
\[
\delta(ax) = \frac{1}{|a|}\delta(x)
\]
\end{theorem}

\begin{proof}
考虑积分 $\displaystyle\int_{-\infty}^{\infty} f(x)\,\delta(ax)\,\diff x$，其中 $f(x)$ 为任意检验函数。

作变量替换 $u = ax$，则 $\diff u = a\,\diff x$，即 $\diff x = \diff u/a$。当 $a > 0$ 时，积分限不变：
\[
\int_{-\infty}^{\infty} f(x)\,\delta(ax)\,\diff x = \int_{-\infty}^{\infty} f\!\left(\frac{u}{a}\right)\delta(u)\,\frac{\diff u}{a} = \frac{1}{a} f(0)
\]
当 $a < 0$ 时，积分限反转，多出负号：
\[
\int_{-\infty}^{\infty} f(x)\,\delta(ax)\,\diff x = \int_{\infty}^{-\infty} f\!\left(\frac{u}{a}\right)\delta(u)\,\frac{\diff u}{a} = \frac{1}{|a|} f(0)
\]
综合两种情况，有：
\[
\int_{-\infty}^{\infty} f(x)\,\delta(ax)\,\diff x = \frac{1}{|a|} f(0) = \int_{-\infty}^{\infty} f(x)\,\frac{1}{|a|}\delta(x)\,\diff x
\]
由检验函数的任意性，即得 $\delta(ax) = |a|^{-1}\delta(x)$。
\end{proof}

\begin{insight}{缩放性质的物理意义}{delta-scaling-phy-v2}
在物理上，$\delta(ax)$ 表示源点位置从 $x=0$ 被"拉伸"或"压缩"了 $a$ 倍。为了保持总"强度"（积分值）为1，$\delta$ 函数的峰值高度必须反比于拉伸因子 $|a|$ 进行调节。例如，在坐标变换 $x' = 2x$ 下，点电荷密度 $\delta(x)$ 变为 $\delta(2x) = \frac{1}{2}\delta(x)$——看似矛盾的表述实际上是因为新坐标系下的积分测度也发生了变化。
\end{insight}

\subsubsection{复合 $\delta$ 函数}

\begin{theorem}{复合 $\delta$ 函数}{delta-composite-v2}
设函数 $f(x)$ 在点 $x_i$ 处有单根（即 $f(x_i)=0$ 且 $f'(x_i) \neq 0$），则：
\[
\delta(f(x)) = \sum_i \frac{\delta(x-x_i)}{|f'(x_i)|}
\]
其中求和遍历 $f(x)$ 的所有单根。
\end{theorem}

\begin{proof}
对任意检验函数 $g(x)$，考虑积分：
\[
I = \int_{-\infty}^{\infty} g(x)\,\delta(f(x))\,\diff x
\]
在每个单根 $x_i$ 的邻域内，$f(x)$ 近似线性：$f(x) \approx f'(x_i)(x-x_i)$。令 $u = f(x)$，则当 $x$ 在 $x_i$ 附近时，$\diff u = f'(x_i)\,\diff x$，即 $\diff x = \diff u/f'(x_i)$。积分化为：
\[
I = \sum_i \int_{f(x_i-\eps)}^{f(x_i+\eps)} g(x(u))\,\delta(u)\,\frac{\diff u}{|f'(x_i)|} = \sum_i \frac{g(x_i)}{|f'(x_i)|}
\]
其中 $|f'(x_i)|$ 保证了 $f'(x_i)<0$ 时积分限反转的符号处理。另一方面：
\[
\int_{-\infty}^{\infty} g(x) \sum_i \frac{\delta(x-x_i)}{|f'(x_i)|}\,\diff x = \sum_i \frac{g(x_i)}{|f'(x_i)|}
\]
两者相等，由检验函数 $g$ 的任意性即得定理成立。
\end{proof}

\begin{example}{复合 $\delta$ 函数的典型应用}{delta-composite-example-v2}
\difficulty{基础}
设 $a \neq 0$，证明：
\[
\delta(x^2 - a^2) = \frac{1}{2|a|}\bigl[\delta(x-a) + \delta(x+a)\bigr]
\]
\tcblower
\textbf{解答：}

令 $f(x) = x^2 - a^2$，其零点为 $x_1 = a$ 和 $x_2 = -a$（$a \neq 0$ 保证两根不相重）。求导数：
\[
f'(x) = 2x \quad\Rightarrow\quad f'(a) = 2a, \quad f'(-a) = -2a
\]
由复合 $\delta$ 函数定理：
\[
\delta(x^2-a^2) = \frac{\delta(x-a)}{|2a|} + \frac{\delta(x+a)}{|-2a|} = \frac{1}{2|a|}\bigl[\delta(x-a) + \delta(x+a)\bigr]
\]
这正是球坐标系中处理 $r^2$ 因子时常见的形式。例如，若需要分离径向与角向变量，此公式可将 $r=a$ 球面上的源分布分解为两个对称平面上的等效源。
\end{example}

\subsection{\texorpdfstring{$\delta$}{δ} 函数的重要表示}

\begin{theorem}{格林函数的泊松方程}{green-poisson-v2}
\[
\lap \frac{1}{|\vect{r} - \vect{r}'|} = -4\pi \, \delta^{(3)}(\vect{r} - \vect{r}')
\]
\end{theorem}

\begin{insight}{为什么这个公式如此重要}{green-important-v2}
这个公式在静电学中具有核心重要性。考虑泊松方程：
\[
\lap\phi = -\frac{\rho}{\eps_0}
\]
如果 $\rho$ 是位于原点的点电荷 $q$，则解为 $\phi = \dfrac{q}{4\pi\eps_0 r}$。将此解代入泊松方程左边，得到 $\lap\dfrac{1}{r} = -4\pi\delta^{(3)}(\vect{r})$——这正是上面的公式在原点的特例。

更一般地，这个公式告诉我们：$1/r$ 函数是拉普拉斯算子在分布意义下的"逆"（差一个常数因子）。这直接导出了电势的积分表达式：
\[
\phi(\vect{r}) = \frac{1}{4\pi\eps_0} \int \frac{\rho(\vect{r}')}{|\vect{r} - \vect{r}'|} \, \diff^3r'
\]
即把电荷分布分解为无数点电荷的叠加，每个点电荷贡献 $1/r$ 势。

\textbf{格林函数的构造（提高）}：上面的 $1/|\vect{r}-\vect{r}'|$ 是\textbf{自由空间}格林函数。带边界条件时，格林函数要"加上齐次解"来满足边界：$G = \dfrac{1}{|\vect{r}-\vect{r}'|} + F(\vect{r},\vect{r}')$，其中 $F$ 满足拉普拉斯方程、恰好使边界条件成立。镜像法构造的正是这个 $F$——例如接地导体平面：$G = \dfrac{1}{|\vect{r}-\vect{r}'|} - \dfrac{1}{|\vect{r}-\vect{r}'_{\text{镜像}}|}$（镜像点在边界另一侧对称位置）。格林函数的完整构造方法是提高内容，初学时掌握"镜像法 = 最简单的格林函数构造"这一对应即可。
\end{insight}

\toggleButton{挑战题}{ocg挑战1}
\begin{ocg}{挑战}{ocg挑战1}{off}
\begin{example}{验证格林函数公式}{green-verify-v2}
\difficulty{挑战}
验证 $\lap(1/r) = -4\pi\delta^{(3)}(\vect{r})$（在原点处）。
\tcblower
\textbf{解答：}

\textbf{情况一：$r \neq 0$}

在球坐标系中，对于仅依赖于 $r$ 的函数，拉普拉斯算子为：
\[
\lap f(r) = \frac{1}{r^2}\frac{\partial}{\partial r}\left(r^2 \frac{\partial f}{\partial r}\right)
\]

令 $f = 1/r$，则 $\partial f/\partial r = -1/r^2$，所以：
\[
r^2 \frac{\partial f}{\partial r} = -1
\]
对 $r$ 求导得 $0$。因此当 $r \neq 0$ 时，$\lap(1/r) = 0$。这与 $\delta^{(3)}(\vect{r})$ 在 $r \neq 0$ 处为零一致。

\textbf{情况二：在 $r = 0$ 附近积分}

取半径为 $\eps$ 的小球 $V_\eps$，计算体积分：
\[
\int_{V_\eps} \lap\frac{1}{r} \, \diff V = \int_{V_\eps} \divg\left(\grad\frac{1}{r}\right) \, \diff V
\]
由高斯定理转化为面积分：
\[
= \oint_{S_\eps} \grad\frac{1}{r} \cdot \diff\vect{S}
\]

在球面上，$\grad(1/r) = -\unitvec{r}/r^2 = -\unitvec{r}/\eps^2$，$\diff\vect{S} = \unitvec{r} \, \eps^2 \diff\Omega$。

\[
= \oint_{S_\eps} \left(-\frac{\unitvec{r}}{\eps^2}\right) \cdot (\unitvec{r} \, \eps^2 \diff\Omega) = -\oint \diff\Omega = -4\pi
\]

因此：
\[
\int \lap\frac{1}{r} \, \diff V = -4\pi
\]
而 $\int (-4\pi\delta^{(3)}(\vect{r})) \, \diff V = -4\pi$。两者积分一致，且两者在 $r \neq 0$ 处都为零，故在分布意义下相等。
\end{example}
\end{ocg}

\section{正交曲线坐标系}

除直角坐标系外，电动力学中常用的坐标系还有\textbf{球坐标系}和\textbf{柱坐标系}。当问题具有球对称性或轴对称性时，使用这些坐标系可以极大地简化计算。


% ---------- 图：三种常用正交曲线坐标系 ----------
\begin{center}
\begin{tikzpicture}[scale=1.0]
% 直角坐标系
\begin{scope}[shift={(-5,0)}]
\draw[->,thick] (0,0,0) -- (2,0,0) node[right] {$x$};
\draw[->,thick] (0,0,0) -- (0,2,0) node[above] {$y$};
\draw[->,thick] (0,0,0) -- (0,0,2) node[below left] {$z$};
\node at (0,-1.2) {\textbf{(a) 直角坐标系}};
\node[align=center,font=\small] at (0,-2.0) {$\vecA=A_x\unitvec{x}+A_y\unitvec{y}+A_z\unitvec{z}$\\单位矢量为常矢量};
\end{scope}

% 柱坐标系
\begin{scope}[shift={(0,0)}]
\draw[->,thick] (0,0) -- (2,0) node[right] {$x$};
\draw[->,thick] (0,0) -- (0,2) node[above] {$z$};
\draw[dashed] (0,0) -- (1.2,0.8);
\fill (1.2,0.8) circle (1.5pt);
\draw[->,thick,red] (1.2,0.8) -- (1.8,1.2) node[right] {$\unitvec{r}$};
\draw[->,thick,blue] (1.2,0.8) -- (0.6,1.4) node[above] {$\unitvec{\phi}$};
\draw[->,thick] (1.2,0.8) -- (1.2,1.6) node[right] {$\unitvec{z}$};
\node at (1,-1.2) {\textbf{(b) 柱坐标系}};
\node[align=center,font=\small] at (1,-2.0) {$(\rho,\phi,z)$\\$\unitvec{\rho},\unitvec{\phi}$随位置变化};
\end{scope}

% 球坐标系
\begin{scope}[shift={(5,0)}]
\draw[dashed] (0,0) -- (1.5,1);
\fill (1.5,1) circle (1.5pt);
\draw[->,thick,red] (1.5,1) -- (2.2,1.47) node[right] {$\unitvec{r}$};
\draw[->,thick,blue] (1.5,1) -- (1.9,0.4) node[right] {$\unitvec{\theta}$};
\draw[->,thick,green!60!black] (1.5,1) -- (0.8,1.3) node[above] {$\unitvec{\phi}$};
\draw (0.5,0) arc (0:35:0.5) node[midway,right,font=\small] {$\theta$};
\node at (0.8,-1.2) {\textbf{(c) 球坐标系}};
\node[align=center,font=\small] at (0.8,-2.0) {$(r,\theta,\phi)$\\适用于球对称问题};
\end{scope}
\end{tikzpicture}
\end{center}

\subsection{球坐标系 \texorpdfstring{$(r, \theta, \phi)$}{(r, θ, φ)}}

\begin{definition}{球坐标与直角坐标的关系}{spherical-coords-v2}
\begin{align*}
x &= r\sin\theta\cos\phi \\
y &= r\sin\theta\sin\phi \\
z &= r\cos\theta
\end{align*}
其中 $r \geq 0$，$0 \leq \theta \leq \pi$，$0 \leq \phi < 2\pi$。

单位矢量：$\unitvec{r}$ 沿径向向外，$\unitvec{\theta}$ 沿 $\theta$ 增加方向（由北极指向赤道平面的方向），$\unitvec{\phi}$ 沿 $\phi$ 增加方向（沿纬线向东）。它们满足右手定则：$\unitvec{r} \times \unitvec{\theta} = \unitvec{\phi}$。
\end{definition}

\begin{theorem}{球坐标系中的微分算子}{spherical-operators-v2}
\textbf{梯度}：
\[
\grad\psi = \frac{\partial\psi}{\partial r}\unitvec{r} + \frac{1}{r}\frac{\partial\psi}{\partial\theta}\unitvec{\theta} + \frac{1}{r\sin\theta}\frac{\partial\psi}{\partial\phi}\unitvec{\phi}
\]

\textbf{散度}：
\[
\divg\vecA = \frac{1}{r^2}\frac{\partial(r^2 A_r)}{\partial r} + \frac{1}{r\sin\theta}\frac{\partial(\sin\theta \, A_\theta)}{\partial\theta} + \frac{1}{r\sin\theta}\frac{\partial A_\phi}{\partial\phi}
\]

\textbf{拉普拉斯算子}：
\[
\lap\psi = \frac{1}{r^2}\frac{\partial}{\partial r}\left(r^2\frac{\partial\psi}{\partial r}\right) + \frac{1}{r^2\sin\theta}\frac{\partial}{\partial\theta}\left(\sin\theta\frac{\partial\psi}{\partial\theta}\right) + \frac{1}{r^2\sin^2\theta}\frac{\partial^2\psi}{\partial\phi^2}
\]

\textbf{旋度}（\textcolor{red}{高查阅频率，务必记住}）：
\begin{align*}
(\curl\vecA)_r &= \frac{1}{r\sin\theta}\left[\frac{\partial}{\partial\theta}(\sin\theta A_\phi) - \frac{\partial A_\theta}{\partial\phi}\right] \\
(\curl\vecA)_\theta &= \frac{1}{r}\left[\frac{1}{\sin\theta}\frac{\partial A_r}{\partial\phi} - \frac{\partial(rA_\phi)}{\partial r}\right] \\
(\curl\vecA)_\phi &= \frac{1}{r}\left[\frac{\partial(rA_\theta)}{\partial r} - \frac{\partial A_r}{\partial\theta}\right]
\end{align*}
\end{theorem}

\begin{insight}{球坐标拉普拉斯的结构}{spherical-lap-insight-v2}
球坐标中的拉普拉斯算子可以写成：
\[
\lap\psi = \underbrace{\frac{1}{r^2}\frac{\partial}{\partial r}\left(r^2\frac{\partial\psi}{\partial r}\right)}_{\text{径向部分}} + \underbrace{\frac{1}{r^2}\left[\frac{1}{\sin\theta}\frac{\partial}{\partial\theta}\left(\sin\theta\frac{\partial\psi}{\partial\theta}\right) + \frac{1}{\sin^2\theta}\frac{\partial^2\psi}{\partial\phi^2}\right]}_{\text{角向部分}/r^2}
\]
角向部分方括号内的微分算子正是角动量平方算子的 $-\hbar^2$ 倍，即 $-\hat{L}^2/\hbar^2$；因此拉普拉斯算子的完整角向项可写为 $-\hat{L}^2/(\hbar^2 r^2)$。这种分离是球对称问题可以求解的关键——径向和角向变量可以完全分离。
\end{insight}

\begin{theorem}{三种坐标系中的体积元与面积元}{volume-elements}
\begin{center}
\small
\setlength{\tabcolsep}{2pt}
\begin{tabular}{@{}cccc@{}}
\toprule
\textbf{坐标系} & \textbf{线元} & \textbf{面积元} & \textbf{体积元} \\
\midrule
直角 & $\diff l = \sqrt{\diff x^2+\diff y^2+\diff z^2}$ & $\diff S = \diff x\diff y$ 等 & $\diff V = \diff x\diff y\diff z$ \\
柱 & $\diff l = \sqrt{\diff r^2+r^2\diff\phi^2+\diff z^2}$ & $\diff S = r\diff r\diff\phi$ 等 & $\diff V = r\diff r\diff\phi\diff z$ \\
球 & $\diff l = \sqrt{\diff r^2+r^2\diff\theta^2+r^2\sin^2\theta\diff\phi^2}$ & $\diff S = r^2\sin\theta\diff\theta\diff\phi$ 等 & $\diff V = r^2\sin\theta\diff r\diff\theta\diff\phi$ \\
\bottomrule
\end{tabular}
\end{center}
\end{theorem}

\begin{insight}{为什么球坐标体积元有 $r^2\sin\theta$？}{volume-element-insight}
从几何直观理解：\\
- $r$ 方向：线元就是 $\diff r$\\
- $\theta$ 方向：弧线长为 $r\diff\theta$\\
- $\phi$ 方向：在纬度圈上，半径为 $r\sin\theta$，弧线长为 $r\sin\theta\diff\phi$\\
三者相乘即得体积元 $r^2\sin\theta\diff r\diff\theta\diff\phi$。面积元同理。
\end{insight}


\begin{example}{球坐标系中的梯度}{spherical-grad-example-v2}
\difficulty{提高}
在球坐标系中，求标量场 $\psi(r, \theta, \phi) = r^2 \sin\theta \cos\phi$ 的梯度。
\tcblower
\textbf{解答：}

计算各偏导数：
\begin{align*}
\frac{\partial\psi}{\partial r} &= 2r\sin\theta\cos\phi \\
\frac{\partial\psi}{\partial\theta} &= r^2\cos\theta\cos\phi \\
\frac{\partial\psi}{\partial\phi} &= -r^2\sin\theta\sin\phi
\end{align*}

代入梯度公式：
\[
\grad\psi = 2r\sin\theta\cos\phi \,\unitvec{r} + r\cos\theta\cos\phi \,\unitvec{\theta} - r\sin\phi \,\unitvec{\phi}
\]
\end{example}

\subsection{柱坐标系 \texorpdfstring{$(\rho, \phi, z)$}{(ρ, φ, z)}}

\begin{theorem}{柱坐标系中的微分算子}{cylindrical-operators-v2}
与直角坐标的关系：
\[
x = \rho\cos\phi, \quad y = \rho\sin\phi, \quad z = z
\]
其中 $\rho \geq 0$，$0 \leq \phi < 2\pi$。

\textbf{梯度}：
\[
\grad\psi = \frac{\partial\psi}{\partial\rho}\unitvec{\rho} + \frac{1}{\rho}\frac{\partial\psi}{\partial\phi}\unitvec{\phi} + \frac{\partial\psi}{\partial z}\unitvec{z}
\]

\textbf{散度}：
\[
\divg\vecA = \frac{1}{\rho}\frac{\partial(\rho A_\rho)}{\partial\rho} + \frac{1}{\rho}\frac{\partial A_\phi}{\partial\phi} + \frac{\partial A_z}{\partial z}
\]

\textbf{拉普拉斯算子}：
\[
\lap\psi = \frac{1}{\rho}\frac{\partial}{\partial\rho}\left(\rho\frac{\partial\psi}{\partial\rho}\right) + \frac{1}{\rho^2}\frac{\partial^2\psi}{\partial\phi^2} + \frac{\partial^2\psi}{\partial z^2}
\]

\textbf{旋度}（\textcolor{red}{高查阅频率，务必记住}）：
\begin{align*}
(\curl\vecA)_\rho &= \frac{1}{\rho}\frac{\partial A_z}{\partial\phi} - \frac{\partial A_\phi}{\partial z} \\
(\curl\vecA)_\phi &= \frac{\partial A_\rho}{\partial z} - \frac{\partial A_z}{\partial\rho} \\
(\curl\vecA)_z &= \frac{1}{\rho}\left[\frac{\partial}{\partial\rho}(\rho A_\phi) - \frac{\partial A_\rho}{\partial\phi}\right]
\end{align*}
\end{theorem}

\subsection{单位矢量的偏导数}

在球坐标系和柱坐标系中，单位矢量本身也是坐标的函数，对坐标求导\textbf{不为零}。这是许多初学者容易忽视的关键点。

\subsubsection{球坐标系中的单位矢量偏导数}

\begin{tip}{初学者关键差异：曲线坐标系的单位矢量随位置变化}{curvilinear-unit-compare}
\textbf{直角坐标系}中 $\unitvec{x}, \unitvec{y}, \unitvec{z}$ 是\textbf{常矢量}——无论站在空间哪一点，$\unitvec{x}$ 永远指向同一个方向。所以 $\partial\unitvec{x}/\partial x = \vect{0}$，求导永远只作用于场分量。

\textbf{曲线坐标系}（柱、球）完全不同：单位矢量\textbf{随点的位置转动}。站在赤道的 $\unitvec{\theta}$ 和站在北极的 $\unitvec{\theta}$ 指向完全不同的方向。一个直观例子：
\begin{center}
\small
\begin{tabular}{@{}lll@{}}
\toprule
 & 单位矢量 & 对位置求导 \\
\midrule
直角坐标 & $\unitvec{x}, \unitvec{y}, \unitvec{z}$（固定方向） & $\partial\unitvec{x}_i/\partial x_j = 0$（全部为零） \\
柱坐标 & $\unitvec{\rho}, \unitvec{\phi}$ 随 $\phi$ 转动 & $\partial\unitvec{\rho}/\partial\phi = \unitvec{\phi} \neq \vect{0}$ \\
球坐标 & $\unitvec{r}, \unitvec{\theta}, \unitvec{\phi}$ 随 $(\theta,\phi)$ 转动 & $\partial\unitvec{r}/\partial\theta = \unitvec{\theta} \neq \vect{0}$ \\
\bottomrule
\end{tabular}
\end{center}

\textbf{这就是为什么}球坐标的 $\lap\phi$ 比直角坐标多出 $\frac{2}{r}\partial_r$、$\frac{1}{r^2}\partial^2_\theta$ 等额外项——对单位矢量求导产生的"几何项"。$\grad$、$\curl$ 在曲线坐标系中的公式里也全是这些几何项。初学时不必死记偏导数表，但要记住结论：\textbf{曲线坐标系求导 = 对分量求导 + 对单位矢量求导}，后者就是公式复杂化的全部来源。
\end{tip}

\begin{theorem}{球坐标单位矢量的偏导数}{spherical-unit-derivatives-v2}
\begin{align*}
\frac{\partial\unitvec{r}}{\partial\theta} &= \unitvec{\theta}, &
\frac{\partial\unitvec{r}}{\partial\phi} &= \sin\theta \, \unitvec{\phi} \\
\frac{\partial\unitvec{\theta}}{\partial\theta} &= -\unitvec{r}, &
\frac{\partial\unitvec{\theta}}{\partial\phi} &= \cos\theta \, \unitvec{\phi} \\
\frac{\partial\unitvec{\phi}}{\partial\phi} &= -\sin\theta \, \unitvec{r} - \cos\theta \, \unitvec{\theta}
\end{align*}
\end{theorem}

% ---------- 图：球坐标单位矢量的方向变化 ----------
\begin{center}
\begin{tikzpicture}[scale=1.0]
% 坐标原点
\fill (0,0) circle (2pt);
\node[below left,font=\small] at (0,0) {$O$};

% 点P
\fill (2,1.5) circle (2pt);
\node[right,font=\small] at (2,1.5) {$P$};

% r方向（从原点指向P）
\draw[->,thick,red] (0,0) -- (2,1.5) node[midway,above left,font=\small] {$\unitvec{r}$};

% θ方向（垂直于r，在竖直平面内）
\draw[->,thick,blue] (2,1.5) -- (2.8,0.9) node[right,font=\small] {$\unitvec{\theta}$};

% φ方向（垂直于纸面向外）
\draw[->,thick,green!60!black] (2,1.5) -- (2,2.3) node[right,font=\small] {$\unitvec{\phi}$};

% 角度标注
\draw (0.5,0) arc (0:35:0.5);
\node[font=\small] at (0.8,0.2) {$\theta$};

% 标注
\node[align=center,font=\small] at (1,-1) {$\unitvec{r}$：径向向外\\$\unitvec{\theta}$：$\theta$增加方向（向南）\\$\unitvec{\phi}$：$\phi$增加方向（向东）};

\node[align=center,font=\small] at (1,-2.5) {关键：$\unitvec{r},\unitvec{\theta},\unitvec{\phi}$都随$(\theta,\phi)$变化\\$\partial\unitvec{r}/\partial\theta=\unitvec{\theta}$，$\partial\unitvec{r}/\partial\phi=\sin\theta\,\unitvec{\phi}$};
\end{tikzpicture}
\end{center}


\subsubsection{柱坐标系中的单位矢量偏导数}

\begin{theorem}{柱坐标单位矢量的偏导数}{cylindrical-unit-derivatives-v2}
\begin{align*}
\frac{\partial\unitvec{\rho}}{\partial\phi} &= \unitvec{\phi}, &
\frac{\partial\unitvec{\phi}}{\partial\phi} &= -\unitvec{\rho}
\end{align*}
对 $\rho$ 和 $z$ 的偏导数为零。
\end{theorem}

\begin{example}{验证单位矢量的偏导数}{unit-derivative-verify-v2}
\difficulty{提高}
验证 $\dfrac{\partial}{\partial\theta}(r\unitvec{r}) = r\unitvec{\theta}$。
\tcblower
\textbf{解答：}

使用乘积法则：
\[
\frac{\partial}{\partial\theta}(r\unitvec{r}) = \frac{\partial r}{\partial\theta}\unitvec{r} + r\frac{\partial\unitvec{r}}{\partial\theta}
\]

由于 $r$ 与 $\theta$ 无关，$\partial r/\partial\theta = 0$。由定理：
\[
\frac{\partial\unitvec{r}}{\partial\theta} = \unitvec{\theta}
\]

因此：
\[
\frac{\partial}{\partial\theta}(r\unitvec{r}) = r\unitvec{\theta}
\]
\end{example}

\begin{review}{常用坐标系体积元与面积元}{coordinate-volume-element}
在电动力学中，正确写出体积元和面积元是计算积分的基础。以下汇总三种常用坐标系的结果：

\begin{center}
\small
\setlength{\tabcolsep}{2pt}
\begin{tabular}{@{}llll@{}}
\toprule
\textbf{坐标系} & \textbf{体积元} & \textbf{面积元} & \textbf{说明} \\
\midrule
直角 $(x,y,z)$ & $\diff x\,\diff y\,\diff z$ & $dS_x=\diff y\,\diff z$ 等 & $h_x=h_y=h_z=1$ \\
\addlinespace
柱 $(\rho,\phi,z)$ & $\rho\,\diff\rho\,\diff\phi\,\diff z$ & $dS_\rho=\rho\,\diff\phi\,\diff z$ 等 & $h_\rho=1,\;h_\phi=\rho,\;h_z=1$ \\
\addlinespace
球 $(r,\theta,\phi)$ & $r^2\sin\theta\,\diff r\,\diff\theta\,\diff\phi$ & $dS_r=r^2\sin\theta\,\diff\theta\,\diff\phi$ 等 & $h_r=1,\;h_\theta=r,\;h_\phi=r\sin\theta$ \\
\bottomrule
\end{tabular}
\end{center}

一般正交曲线坐标系 $(q_1,q_2,q_3)$ 中，若尺度因子为 $h_1,h_2,h_3$，则体积元和面积元为：
\[
\diff V = h_1 h_2 h_3 \,\diff q_1\,\diff q_2\,\diff q_3, \qquad
\diff S_i = \frac{h_1 h_2 h_3}{h_i}\,\diff q_j\,\diff q_k \;(i\neq j\neq k)
\]
直角坐标系中 $h_x=h_y=h_z=1$，柱坐标系中 $h_\rho=1,\;h_\phi=\rho,\;h_z=1$，球坐标系中 $h_r=1,\;h_\theta=r,\;h_\phi=r\sin\theta$。
\end{review}

\begin{warning}{常见错误}{unit-derivative-mistake}
初学者常犯的错误是假设单位矢量是"常矢量"，从而认为它们的偏导数为零。\textbf{只有在直角坐标系中}，$\unitvec{x}, \unitvec{y}, \unitvec{z}$ 不随位置变化。在球坐标和柱坐标中，单位矢量的方向依赖于空间位置，必须考虑其导数。这在计算拉普拉斯算子和旋度时尤为重要。
\end{warning}

\begin{review}{第1章要点回顾}{ch1-review}
\begin{center}
\small
\setlength{\tabcolsep}{2pt}
\begin{tabular}{@{}llll@{}}
\toprule
\textbf{算子} & \textbf{作用于} & \textbf{结果} & \textbf{符号} \\
\midrule
梯度 & 标量场 & 矢量场 & $\grad\phi$ \\
散度 & 矢量场 & 标量场 & $\divg\vecA$ \\
旋度 & 矢量场 & 矢量场 & $\curl\vecA$ \\
拉普拉斯 & 标量/矢量场 & 同类型场 & $\lap$ \\
\bottomrule
\end{tabular}
\end{center}

\textbf{核心定理}：
\begin{itemize}
\item 高斯定理：$\displaystyle\oint_S \vecA \cdot \diff\vect{S} = \int_V \divg\vecA \, \diff V$
\item 斯托克斯定理：$\displaystyle\oint_C \vecA \cdot \diff\vect{l} = \int_S \curl\vecA \cdot \diff\vect{S}$
\end{itemize}
\end{review}

\begin{tip}{本章核心框架}{ch1-summary}
\textbf{一句话总结}：矢量分析是描述电磁场的数学语言。\\
\textbf{四大核心}：梯度（标量→矢量）、散度（矢量→标量，源/汇判断）、旋度（矢量→矢量，旋转程度）、拉普拉斯（二阶导数，波动/扩散方程核心）。\\
\textbf{三大工具}：高斯定理（体积分 $\leftrightarrow$ 面积分）、斯托克斯定理（面积分 $\leftrightarrow$ 线积分）、$\delta$函数（点源描述）。\\
\textbf{衔接}：第2章将用$\grad$描述电势与电场关系，用$\lap$建立泊松方程。
\end{tip}

% ---------- 第1章随堂自测 ----------
\section*{第1章随堂自测}
\addcontentsline{toc}{section}{第1章随堂自测}

\begin{miniquiz}
\textbf{一、选择题（每题只有一个正确答案）}
\begin{enumerate}[label=\arabic*.]
\item 标量场 $\psi(x,y,z)$ 的梯度是一个：
  \begin{enumerate}[label=(\Alph*),nosep]
  \item 标量
  \item 矢量
  \item 张量
  \item 取决于坐标系
  \end{enumerate}
\begin{ocg}{答案}{ocg答案1}{off}
  \textit{答案：B。梯度将标量场映射为矢量场。}
\end{ocg}

\item 矢量场 $\vecA$ 的旋度的散度 $\divg(\curl\vecA)$ 恒等于：
  \begin{enumerate}[label=(\Alph*),nosep]
  \item $|\vecA|$
  \item $\lap\vecA$
  \item $0$
  \item $\curl(\divg\vecA)$
  \end{enumerate}
\begin{ocg}{答案}{ocg答案1}{off}
  \textit{答案：C。任何光滑矢量场的旋度都是无散的。}
\end{ocg}

\item 在球坐标系中，$\unitvec{r} \times \unitvec{\theta}$ 等于：
  \begin{enumerate}[label=(\Alph*),nosep]
  \item $\unitvec{r}$
  \item $\unitvec{\theta}$
  \item $\unitvec{\phi}$
  \item $-\unitvec{\phi}$
  \end{enumerate}
\begin{ocg}{答案}{ocg答案1}{off}
  \textit{答案：C。球坐标单位矢量满足右手定则。}
\end{ocg}

\item 狄拉克$\delta$函数 $\delta(x)$ 的本质是：
  \begin{enumerate}[label=(\Alph*),nosep]
  \item 普通函数
  \item 广义函数（分布）
  \item 仅在 $x=0$ 处非零
  \item 积分等于0
  \end{enumerate}
\begin{ocg}{答案}{ocg答案1}{off}
  \textit{答案：B。$\delta$函数是广义函数，不是普通函数。}
\end{ocg}
\end{enumerate}

\vspace{0.5em}
\textbf{二、判断题（判断正误并简要说明理由）}
\begin{enumerate}[label=\arabic*.,resume]
\item 任意标量场的梯度的旋度恒为零：$\curl(\grad\psi) \equiv \vect{0}$。
\begin{ocg}{答案}{ocg答案1}{off}
  \textit{答案：正确。保守场必无旋。}
\end{ocg}

\item 高斯定理只适用于静电场，不适用于静磁场。
\begin{ocg}{答案}{ocg答案1}{off}
  \textit{答案：错误。高斯定理是数学定理，适用于任何光滑矢量场。}
\end{ocg}

\item 柱坐标系中的单位矢量 $\unitvec{r}, \unitvec{\phi}$ 是常矢量。
\begin{ocg}{答案}{ocg答案1}{off}
  \textit{答案：错误。曲线坐标系的单位矢量随位置变化。}
\end{ocg}
\end{enumerate}
\end{miniquiz}
\toggleButton{显示答案}{ocg答案1}

\begin{chapterreview}
\begin{itemize}[nosep]
\item \textbf{核心概念}：标量场、矢量场、张量；梯度、散度、旋度；拉普拉斯算子；并矢；$\delta$函数。
\item \textbf{核心定理}：高斯定理、斯托克斯定理；$\curl(\grad\psi)\equiv\vect{0}$ 和 $\divg(\curl\vecA)\equiv 0$。
\item \textbf{坐标系}：熟练掌握直角、柱、球三种坐标系中的微分算子。
\item \textbf{自检验}：能否不查表写出球坐标拉普拉斯？能否证明 $\curl(\grad\psi)=0$？
\item \textbf{衔接预告}：第2章将把这些数学工具应用于静电场——标量势$\phi$和电场$\vecE$。
\end{itemize}
\end{chapterreview}

\begin{tip}{本章自检清单}{ch1-checklist}
学完本章后对照下列条目逐条自测，全部通过即可进入下一章；有未通过项时，回到对应小节复习：
\begin{itemize}[label=$\square$, nosep]
\item 不查表写出梯度、散度、旋度在直角/柱/球坐标下的表达式
\item 用高斯定理和斯托克斯定理验证简单矢量场
\item 理解 $\delta$ 函数的筛选性质并完成含 $\delta$ 的积分
\item 知道并矢可以跳过、何时回头补
\end{itemize}
\end{tip}

% ---------- 第1章综合练习题 ----------
\section*{第1章综合练习题}
\addcontentsline{toc}{section}{第1章综合练习题}
\begin{enumerate}[label=\textbf{\arabic*.}]
\item \textbf{基础} 计算 $\grad(r^2)$、$\divg(\vect{r})$ 和 $\curl(\vect{r})$，其中 $\vect{r}=x\unitvec{x}+y\unitvec{y}+z\unitvec{z}$。
\item \textbf{基础} 用高斯散度定理验证：对于半径为 $R$ 的球体，$\oint_S \vect{r}\cdot d\vect{S}=4\pi R^3$。
\item \textbf{提高} 证明并矢 $\bm{A}\bm{B}$ 与矢量 $\bm{C}$ 的点乘满足：$(\bm{A}\bm{B})\cdot\bm{C}=\bm{A}(\bm{B}\cdot\bm{C})$。
\item \textbf{提高} 在柱坐标系中，设 $\psi(\rho,\phi,z)=\rho^2\cos\phi$，计算 $\lap\psi$。
\item \textbf{挑战} 利用傅里叶积分表示证明：$\int_{-\infty}^{\infty} \delta(x-a)\delta(x-b)\,dx = \delta(a-b)$。
\end{enumerate}
\newpage
\section*{习题解答}
\addcontentsline{toc}{section}{习题解答}

\textbf{题1解析：}

[分析] 本题考查直角坐标系中梯度、散度、旋度的直接计算，是矢量分析的基本功。

[解答] 设 $\vect{r} = x\unitvec{x} + y\unitvec{y} + z\unitvec{z}$，则 $r^2 = x^2 + y^2 + z^2$。

梯度：$\grad(r^2) = \unitvec{x}\frac{\partial r^2}{\partial x} + \unitvec{y}\frac{\partial r^2}{\partial y} + \unitvec{z}\frac{\partial r^2}{\partial z} = 2x\unitvec{x} + 2y\unitvec{y} + 2z\unitvec{z} = 2\vect{r}$。

散度：$\divg\vect{r} = \frac{\partial x}{\partial x} + \frac{\partial y}{\partial y} + \frac{\partial z}{\partial z} = 1 + 1 + 1 = 3$。

旋度：$\curl\vect{r} = \begin{vmatrix}\unitvec{x} & \unitvec{y} & \unitvec{z} \\ \partial_x & \partial_y & \partial_z \\ x & y & z\end{vmatrix} = \unitvec{x}(0-0) - \unitvec{y}(0-0) + \unitvec{z}(0-0) = \vect{0}$。

[量纲校验] $[\grad(r^2)] = [L^{-1}][L^2] = [L]$，$[2\vect{r}] = [L]$，一致；$[\divg\vect{r}]$ 无量纲，结果为3无量纲，一致；$[\curl\vect{r}] = [L^{-1}][L] = 1$（无量纲），$\vect{0}$ 无矛盾。\qedsymbol

\textbf{题2解析：}

[分析] 本题用高斯散度定理将面积分转化为体积分，验证 $\vect{r}$ 的散度与球体积的关系。

[解答] 由高斯散度定理：$\displaystyle\oint_S \vect{r}\cdot\diff\vect{S} = \int_V \divg\vect{r}\,\diff V$。

由题1知 $\divg\vect{r} = 3$，故：
\[
\int_V 3\,\diff V = 3 \cdot \frac{4}{3}\pi R^3 = 4\pi R^3
\]

直接验证：在球面上 $\diff\vect{S} = R^2\sin\theta\,\diff\theta\,\diff\phi\,\unitvec{r}$，于是
\[
\vect{r}\cdot\diff\vect{S} = R^3\sin\theta\,\diff\theta\,\diff\phi
\]
积分得 $\displaystyle\int_0^{2\pi}\diff\phi\int_0^\pi R^3\sin\theta\,\diff\theta = 2\pi R^3 \cdot 2 = 4\pi R^3$。\qedsymbol

[量纲校验] $[\oint\vect{r}\cdot\diff\vect{S}] = [L][L^2] = [L^3]$，$[4\pi R^3] = [L^3]$，一致。\qedsymbol

\textbf{题3解析：}

[分析] 本题考查并矢（dyad）与矢量的点乘运算，关键在于理解爱因斯坦求和约定下指标缩并的顺序。

[解答] 并矢 $\bm{A}\bm{B}$ 的矩阵元为 $(AB)_{ij} = A_i B_j$。与矢量 $\bm{C}$ 点乘（右缩并）：
\[
[(\bm{A}\bm{B})\cdot\bm{C}]_i = \sum_j (AB)_{ij} C_j = \sum_j A_i B_j C_j = A_i (\bm{B}\cdot\bm{C})
\]
即 $(\bm{A}\bm{B})\cdot\bm{C} = \bm{A}(\bm{B}\cdot\bm{C})$，其中 $\bm{B}\cdot\bm{C}$ 是标量，因此结果是矢量 $\bm{A}$ 的标量倍数。\qedsymbol

[量纲校验] 设 $[A]=[B]=[C]=[L]$，则左端 $[(AB)C] = [L]^2[L] = [L]^3$，右端 $[A(B\cdot C)] = [L][L]^2 = [L]^3$，一致。\qedsymbol

\textbf{题4解析：}

[分析] 本题要求在柱坐标系中计算拉普拉斯算子，必须使用柱坐标下的正确表达式。

[解答] 柱坐标下拉普拉斯算子：$\lap\psi = \dfrac{1}{\rho}\dfrac{\partial}{\partial\rho}\left(\rho\dfrac{\partial\psi}{\partial\rho}\right) + \dfrac{1}{\rho^2}\dfrac{\partial^2\psi}{\partial\phi^2} + \dfrac{\partial^2\psi}{\partial z^2}$。

给定 $\psi = \rho^2\cos\phi$，计算各项：
\begin{align*}
\frac{\partial\psi}{\partial\rho} &= 2\rho\cos\phi, \quad \frac{\partial}{\partial\rho}\left(\rho\cdot 2\rho\cos\phi\right) = \frac{\partial}{\partial\rho}(2\rho^2\cos\phi) = 4\rho\cos\phi \\
\frac{1}{\rho}\frac{\partial}{\partial\rho}\left(\rho\frac{\partial\psi}{\partial\rho}\right) &= \frac{1}{\rho}\cdot 4\rho\cos\phi = 4\cos\phi \\
\frac{\partial^2\psi}{\partial\phi^2} &= \frac{\partial}{\partial\phi}(-\rho^2\sin\phi) = -\rho^2\cos\phi, \quad \frac{1}{\rho^2}\frac{\partial^2\psi}{\partial\phi^2} = -\cos\phi \\
\frac{\partial^2\psi}{\partial z^2} &= 0
\end{align*}

因此：$\lap\psi = 4\cos\phi - \cos\phi + 0 = 3\cos\phi$。

[量纲校验] $[\lap\psi] = [L^{-2}][L^2] = 1$（无量纲），结果 $3\cos\phi$ 无量纲，一致。\qedsymbol

\textbf{题5解析：}

[分析] 本题利用 $\delta$ 函数的傅里叶表示，通过交换积分顺序完成证明。严格处理广义函数需要谨慎对待积分顺序。

[解答] $\delta$ 函数的傅里叶表示：$\delta(x) = \dfrac{1}{2\pi}\int_{-\infty}^{\infty} e^{ikx}\,\diff k$。

于是：
\begin{align*}
\int_{-\infty}^{\infty} \delta(x-a)\delta(x-b)\,\diff x &= \int_{-\infty}^{\infty} \left[\frac{1}{2\pi}\int e^{ik(x-a)}\,\diff k\right] \left[\frac{1}{2\pi}\int e^{ik'(x-b)}\,\diff k'\right] \diff x \\
&= \frac{1}{(2\pi)^2}\int\int\diff k\,\diff k'\, e^{-ika-ik'b} \int_{-\infty}^{\infty} e^{i(k+k')x}\,\diff x \\
&= \frac{1}{(2\pi)^2}\int\int\diff k\,\diff k'\, e^{-ika-ik'b} \cdot 2\pi\delta(k+k') \\
&= \frac{1}{2\pi}\int\diff k\, e^{-ika+ikb} = \frac{1}{2\pi}\int\diff k\, e^{ik(b-a)} = \delta(a-b)
\end{align*}

[量纲校验] $[\delta(x)] = [L^{-1}]$，$[\diff x] = [L]$，左端量纲 $[L^{-1}]$；右端 $\delta(a-b)$ 量纲 $[L^{-1}]$，一致。\qedsymbol



\begin{warning}{本章常见误区总结}{ch1-warnings}
\begin{enumerate}
\item $\grad\phi$ 指向 $\phi$\textbf{增长}最快的方向，不是下降方向。电场 $\vecE = -\grad\phi$ 才指向下降方向。
\item $\divg\vecA = 0$ 不意味着 $\vecA = \vect{0}$，只意味着场线没有起点和终点（\textbf{不}意味着场线一定闭合，如 $\vecA=\unitvec{z}$）。
\item $\curl\vecA = \vect{0}$ 不意味着 $\vecA = \vect{0}$，只意味着在\textbf{单连通}区域内 $\vecA$ 可以写成某个标量场的梯度。
\item \textbf{$\nabla$ 算子不是普通矢量}！$\nabla\cdot\vecA$（散度，标量）$\neq$ $\vecA\cdot\nabla$（算子，作用于其右侧的场）。$\nabla\times\vecA$（旋度，矢量）$\neq$ $-\vecA\times\nabla$（算子）。
\item $\nabla\cdot(\nabla\times\vecA) \equiv 0$ 和 $\nabla\times(\nabla\psi) \equiv \vect{0}$ 是恒等式，但 $\nabla\cdot\nabla\psi = \lap\psi$ 不是恒等式——它定义了拉普拉斯算子。
\item 高斯定理要求曲面\textbf{闭合}；斯托克斯定理要求曲线\textbf{闭合}。
\item $\delta$ 函数不是普通函数，是广义函数（分布）。严格说它只在积分号下有意义——脱离积分号直接做"$\delta(x) = 0$、$\delta^2(x)$"之类的等式运算属于形式简写，没有严格数学意义。本书习题中大量对 $\delta$ 直接做代数运算，都是这种简写；正确做法是两边同乘检验函数并积分。
\end{enumerate}
\end{warning}

\begin{history}{矢量分析的发展脉络}{vector-history-2}
现代矢量分析的符号体系主要在1880-1890年代由吉布斯和亥维赛建立。在此之前，麦克斯韦使用"四元数"记号（哈密顿发明），虽然优雅但计算繁琐。有趣的是，麦克斯韦本人其实不太喜欢四元数——他在信中写道："四元数的算子方法就像一种语言，你必须学会用它思考。"矢量分析的诞生，正是这种"让数学更贴近物理直觉"的需求驱动。今天，当我们写下 $\curl\vecE$ 时，很少有人意识到这个简洁的记号背后有一个世纪的数学演进。
\end{history}


% ============================================================

\begin{center}\small \hyperlink{ch2-goal}{下一章：第2章}\end{center}

% 第2章：静电学基础（优化版）
% ============================================================
\chapter{静电学基础}
\epigraph{给我一个支点，我就能撬动地球。}{--- 阿基米德（静电力也是一种"支点"）}

\hypertarget{ch2-goal}{}\begin{insight}{本章核心目标}{ch2-goal-v2}
从库仑定律出发，建立静电学的完整理论框架：电场、电势、高斯定理、环路定理，最终得到泊松方程。掌握唯一性定理——它是镜像法和分离变量法的理论基石。本章的每个推导都应达到"自己能独立复现"的程度。
\end{insight}

静电学研究\textbf{静止电荷}产生的电场及其性质。这是整个电磁理论的起点——不仅因为静电现象最容易观测，更因为静电学中建立的概念和方法（如标量势、高斯定理、唯一性定理）贯穿整个电动力学。

\begin{history}{静电学的发展}{electrostatics-history}
摩擦起电早在公元前 600 年就被古希腊人发现（希腊语"琥珀"elektron 正是"电"的词源），但定量研究始于 18 世纪：\textbf{库仑} 1785 年用扭秤实验确立平方反比律；\textbf{泊松}与\textbf{高斯}在 1813 年分别给出泊松方程与散度定理；\textbf{格林} 1828 年引入势函数方法。法拉第 1830 年代的"力线"概念虽未数学形式化，却为麦克斯韦的场论提供了关键物理图像。
\end{history}

\section{库仑定律与电场}

\subsection{库仑定律}

\begin{theorem}{库仑定律}{coulomb-law-v2}
真空中，两个静止点电荷 $q_1$ 和 $q_2$ 之间的作用力为：
\[
\vect{F} = \frac{1}{4\pi\eps_0} \frac{q_1 q_2}{r^2} \unitvec{r}
\]
其中 $r = |\vect{r}_1 - \vect{r}_2|$ 是两电荷间的距离，$\unitvec{r} = (\vect{r}_1 - \vect{r}_2)/r$ 是从 $q_2$ 指向 $q_1$ 的单位矢量。$\eps_0 = 8.854 \times 10^{-12} \, \si{C^2/(N.m^2)}$ 是\textbf{真空介电常数}。同性电荷相斥，异性电荷相吸。
\end{theorem}

\begin{insight}{$1/4\pi\eps_0$ 的物理意义}{coulomb-constant-v2}
常数 $1/4\pi\eps_0 \approx 8.99 \times 10^9 \, \si{N.m^2/C^2}$。在国际单位制（SI）中，这是电学量的基本常数。选择 $4\pi$ 出现在分母而非分子中，是所谓"\textbf{合理化单位制}"（rationalized units）的选择——它使得球面积分的结果（$4\pi r^2$）能与分母的 $4\pi$ 抵消，简化后续公式。虽然这看起来只是数学技巧，但在电磁学中这种"合理化"贯穿始终，使得 Maxwell 方程组以最简洁的形式出现。
\end{insight}

\subsection{电场}

\begin{definition}{电场}{electric-field-v2}
电场 $\vecE(\vect{r})$ 定义为：单位正试探电荷在空间中某点所受的力：
\[
\vecE(\vect{r}) = \frac{\vect{F}}{q_0} = \frac{1}{4\pi\eps_0} \frac{q}{|\vect{r} - \vect{r}'|^2} \unitvec{R}
\]
其中 $\vect{R} = \vect{r} - \vect{r}'$，$\unitvec{R} = \vect{R}/R$。电场的单位是伏特/米（V/m）或牛顿/库仑（N/C）。
\end{definition}

\begin{insight}{电场的物理本质}{field-concept-v2}
电场概念的引入是物理学史上的革命性一步。在牛顿力学中，力的作用是"超距"的——两个物体通过某种神秘的"超距作用"瞬间相互影响。法拉第和麦克斯韦提出的"场"概念改变了这一图景：电荷在其周围空间中\textbf{产生}了一种物理实体——电场，其他电荷通过感受这个场而受到力。电场是局域的、连续的，以有限速度（光速）传播。这个观念的转变是现代物理学的基础——从引力场到规范场，"场"已经成为描述自然界的基本语言。
\end{insight}

\subsection{叠加原理}

\begin{theorem}{叠加原理}{superposition-v2}
对于多个点电荷产生的电场，总电场等于各电荷产生电场的\textbf{矢量和}：
\[
\vecE(\vect{r}) = \frac{1}{4\pi\eps_0} \sum_{i=1}^N \frac{q_i}{|\vect{r} - \vect{r}_i|^2} \unitvec{R}_i
\]
对于连续电荷分布：
\[
\vecE(\vect{r}) = \frac{1}{4\pi\eps_0} \int \frac{\rho(\vect{r}')}{|\vect{r} - \vect{r}'|^2} \unitvec{R} \, \diff^3r'
\]
\end{theorem}

\begin{insight}{叠加原理的深层含义}{superposition-deep-v2}
叠加原理并非理所当然。在广义相对论中，引力场不满足线性叠加——质量产生的时空弯曲相互影响，导致非线性方程。电磁学中的叠加原理源于麦克斯韦方程组的线性性，而线性性又与光子的无质量性和电磁相互作用的阿贝尔规范对称性（U(1)）密切相关。在强场极限下（如超强激光与等离子体相互作用），即使是电磁场也会出现非线性效应（如相对论自聚焦），但在经典电动力学的通常范围内，叠加原理精确成立。
\end{insight}

\begin{example}{电偶极子的电场}{dipole-field-ex-v2}
\difficulty{提高}
两个等量异号电荷 $+q$ 和 $-q$ 相距 $d$，构成\textbf{电偶极子}，电偶极矩 $\vect{p} = q\vect{d}$（$\vect{d}$ 从 $-q$ 指向 $+q$）。求远处（$r \gg d$）的电场。
\tcblower
\textbf{解答：}

取坐标原点在偶极子中心，$+q$ 在 $\vect{d}/2$，$-q$ 在 $-\vect{d}/2$。

$\vecE(\vect{r}) = \dfrac{q}{4\pi\eps_0} \left[ \dfrac{\vect{r} - \vect{d}/2}{|\vect{r} - \vect{d}/2|^3} - \dfrac{\vect{r} + \vect{d}/2}{|\vect{r} + \vect{d}/2|^3} \right]$

对 $d/r$ 展开到一阶：

$|\vect{r} \mp \vect{d}/2|^{-3} \approx r^{-3}\left(1 \pm \dfrac{3\vect{r}\cdot\vect{d}}{2r^2}\right)$

经过代数运算（展开后线性项保留，高阶项略去），得到：
\[
\vecE(\vect{r}) = \frac{1}{4\pi\eps_0 r^3} \left[ 3(\vect{p}\cdot\unitvec{r})\unitvec{r} - \vect{p} \right]
\]

这是电偶极子远场的一般公式，在球坐标中：
\begin{align*}
E_r &= \frac{2p\cos\theta}{4\pi\eps_0 r^3} \\
E_\theta &= \frac{p\sin\theta}{4\pi\eps_0 r^3} \\
E_\phi &= 0
\end{align*}

注意：电偶极子场按 $1/r^3$ 衰减，比点电荷的 $1/r^2$ 更快——因为正负电荷的贡献在大距离上几乎抵消，只剩下一阶修正项。
\end{example}

\section{高斯定理}

\subsection{电通量}

\begin{definition}{电通量}{electric-flux-v2}
电场穿过曲面 $S$ 的\textbf{电通量}定义为：
\[
\Phi_E = \int_S \vecE \cdot \diff\vect{S}
\]
对于闭合曲面，通量表示从曲面内\textbf{净流出}的电场线数目。
\end{definition}

\subsection{高斯定理的积分形式}

\begin{theorem}{高斯定理（积分形式）}{gauss-law-integral-v2}
\[
\oint_S \vecE \cdot \diff\vect{S} = \frac{Q_{\text{enc}}}{\eps_0} = \frac{1}{\eps_0} \int_V \rho(\vect{r}) \, \diff^3r
\]
其中 $Q_{\text{enc}}$ 是闭合曲面 $S$ 所包围的净电荷，$V$ 是 $S$ 所围的体积。
\end{theorem}

\begin{insight}{高斯定理的物理图像}{gauss-insight-2-v2}
高斯定理告诉我们：穿过任意闭合曲面的电通量只与曲面内的净电荷有关，与电荷如何分布、曲面形状无关。这是静电学中最强大的工具之一——当电荷分布具有高度对称性（球对称、柱对称、平面对称）时，可以用它快速求出电场，而无需做复杂的积分。

为什么通量只与内部电荷有关？因为电场线从正电荷发出、终止于负电荷（或延伸至无穷远）。对于闭合曲面外的电荷，其发出的电场线必然同时穿入和穿出曲面，净贡献为零。
\end{insight}

\subsection{高斯定理的微分形式}

由高斯散度定理：$\displaystyle\oint_S \vecE \cdot \diff\vect{S} = \int_V \divg\vecE \, \diff V$

结合高斯定理：$\displaystyle\int_V \divg\vecE \, \diff V = \frac{1}{\eps_0} \int_V \rho \, \diff V$

由于这对任意体积 $V$ 成立，故被积函数相等：

\begin{theorem}{高斯定理（微分形式）}{gauss-law-diff-v2}
\[
\divg\vecE = \frac{\rho}{\eps_0}
\]
这是麦克斯韦方程组的第一个方程（在静电学极限下）。
\end{theorem}

\begin{example}{利用高斯定理求电场}{gauss-app-ex-v2}
\difficulty{基础}
求均匀带电球面（半径 $R$，总电荷 $Q$）产生的电场。
\tcblower
\textbf{解答：}

由球对称性，电场只沿径向且大小仅依赖于 $r$。

\textbf{区域一：$r > R$}（球外）

取半径为 $r$ 的同心球面为高斯面。由高斯定理：
\[
\oint \vecE \cdot \diff\vect{S} = E(r) \cdot 4\pi r^2 = \frac{Q}{\eps_0}
\]

$E(r) = \dfrac{Q}{4\pi\eps_0 r^2}$

与点电荷电场相同！——球对称电荷分布在外部等效于点电荷。

\textbf{区域二：$r < R$}（球内）

高斯面内无电荷：$Q_{\text{enc}} = 0$

$E(r) = 0$！——均匀带电球面内部电场为零。

\textbf{结论}：
\[
E(r) = \begin{cases} 0, & r < R \\ \dfrac{Q}{4\pi\eps_0 r^2}, & r > R \end{cases}
\]

在 $r = R$ 处，电场从0跳跃到有限值，这是面电荷分布的特征。
\end{example}

\begin{example}{无限长带电直线的电场}{line-charge-ex-v2}
\difficulty{基础}
求电荷线密度为 $\lambda$ 的无限长带电直线的电场。
\tcblower
\textbf{解答：}

由柱对称性，电场只有径向分量，大小只与到直线的距离 $\rho$ 有关。

取同轴圆柱面（半径 $\rho$，高度 $h$）为高斯面。侧面贡献通量，上下底面 $\vecE \perp \diff\vect{S}$ 无贡献。

$\displaystyle\oint \vecE \cdot \diff\vect{S} = E(\rho) \cdot 2\pi\rho h = \frac{\lambda h}{\eps_0}$

$E(\rho) = \dfrac{\lambda}{2\pi\eps_0 \rho}$

方向沿径向向外（$\lambda > 0$）或向内（$\lambda < 0$）。
\end{example}

\begin{warning}{高斯定理使用要点}{gauss-warning-v2}
\begin{enumerate}
\item 高斯定理对\textbf{任何}闭合曲面都成立，但要"好用"，必须选择\textbf{对称性匹配}的高斯面：
  \begin{itemize}
  \item 球对称 $\to$ 同心球面
  \item 柱对称 $\to$ 同轴圆柱面
  \item 平面对称 $\to$ "药盒"形（pillbox）
  \end{itemize}
\item 高斯定理只能直接求出\textbf{电场大小}（通过对称性确定方向后）。如果电场方向不能先验确定，仅靠高斯定理不够。
\item 非对称电荷分布（如电偶极子）不能用高斯定理直接求电场——需要叠加原理或其他方法。
\end{enumerate}
\end{warning}

\section{环路定理与电势}

\subsection{静电场的无旋性}

\begin{theorem}{静电场环路定理}{loop-theorem-v2}
\[
\oint_C \vecE \cdot \diff\vect{l} = 0
\]
对任意闭合路径 $C$ 成立。
\end{theorem}

\begin{insight}{环路定理的物理含义}{loop-meaning-v2}
环路定理说：静电场沿任何闭合路径的\textbf{功为零}。这意味着将试探电荷从一点移动到另一点，电场力做的功与路径无关——静电场是\textbf{保守场}。保守场的一个等价描述是：可以定义一个标量势函数，使得 $\vecE = -\grad\phi$。

为什么静电场是无旋的？从本质上说，静电场由静止电荷产生，而静止电荷的电场线从正电荷发出、终止于负电荷，不会形成闭合回路。这与感应电场（由变化磁场产生）形成对比——后者可以是有旋的（法拉第定律）。
\end{insight}

由斯托克斯定理：$\displaystyle\oint_C \vecE \cdot \diff\vect{l} = \int_S (\curl\vecE) \cdot \diff\vect{S} = 0$，对任意曲面 $S$ 成立，故：

\begin{theorem}{静电场无旋性的微分形式}{curl-e-zero-v2}
\[
\curl\vecE = \vect{0}
\]
这是麦克斯韦方程组中法拉第定律在静电学极限（$\partial\vecB/\partial t = \vect{0}$）下的形式。
\end{theorem}

\subsection{电势}

\begin{definition}{电势}{electric-potential-v2}
由于静电场是保守场，可定义\textbf{电势} $\phi(\vect{r})$ 满足：
\[
\vecE = -\grad\phi
\]
电势的差值具有物理意义：$\phi(\vect{b}) - \phi(\vect{a}) = -\displaystyle\int_{\vect{a}}^{\vect{b}} \vecE \cdot \diff\vect{l}$

点电荷的电势（取无穷远为零势点）：
\[
\phi(\vect{r}) = \frac{q}{4\pi\eps_0 r}
\]
\end{definition}

\begin{insight}{电势与电场的关系}{potential-field-rel-v2}
$\vecE = -\grad\phi$ 是静电学中最基本的关系之一。从物理上看：
\begin{itemize}
\item 电场方向是电势\textbf{下降最快}的方向（负号！）
\item 等势面（$\phi = \text{常数}$）处处与电场线\textbf{垂直}
\item 电场越强的地方，等势面越密集
\end{itemize}

从数学上看，$\vecE = -\grad\phi$ 是\textbf{矢量场到标量场的约化}——将三个分量的信息压缩到一个标量函数中。这是可能的，正是因为静电场满足 $\curl\vecE = 0$（根据矢量分析，无旋场必可表为标量势的梯度）。
\end{insight}

\begin{example}{均匀带电圆环轴线上的电势}{ring-potential-ex-v2}
\difficulty{基础}
半径为 $R$ 的圆环均匀带电，总电荷 $Q$。求轴线上距离圆心 $z$ 处的电势和电场。
\tcblower
\textbf{解答：}

圆环上电荷元 $\diff q = \lambda R\diff\theta = \dfrac{Q}{2\pi} \diff\theta$

电荷元到轴上点的距离：$r = \sqrt{R^2 + z^2}$（对所有电荷元相同！）

电势（标量叠加，无需矢量积分）：
\[
\phi(z) = \frac{1}{4\pi\eps_0} \oint \frac{\diff q}{r} = \frac{1}{4\pi\eps_0} \frac{Q}{\sqrt{R^2+z^2}}
\]

由对称性，电场只有 $z$ 分量：
\[
E_z = -\frac{\partial\phi}{\partial z} = \frac{Qz}{4\pi\eps_0 (R^2+z^2)^{3/2}}
\]

在 $z = 0$（圆心）处，$\phi = \dfrac{Q}{4\pi\eps_0 R}$，$E_z = 0$（对称性预期）。

在 $z \gg R$ 处，$\phi \approx \dfrac{Q}{4\pi\eps_0 z}$，$E_z \approx \dfrac{Q}{4\pi\eps_0 z^2}$——退化为点电荷。
\end{example}

\section{泊松方程与拉普拉斯方程}

\subsection{从基本方程到泊松方程}

静电学的两个基本方程为：
\begin{align}
\divg\vecE &= \frac{\rho}{\eps_0} \label{eq:gauss-static} \\
\curl\vecE &= \vect{0} \label{eq:curl-static}
\end{align}

由式\eqref{eq:curl-static}，$\vecE = -\grad\phi$。代入式\eqref{eq:gauss-static}：
\[
\divg(-\grad\phi) = -\lap\phi = \frac{\rho}{\eps_0}
\]

\begin{theorem}{泊松方程}{poisson-eq-v2}
\begin{impEqbox}{静电学核心方程}
\begin{equation}
\lap\phi = -\frac{\rho}{\eps_0}\label{eq:poisson}
\end{equation}
在无电荷区域（$\rho = 0$），退化为\textbf{拉普拉斯方程}：
\[
\lap\phi = 0
\]
\end{impEqbox}
\end{theorem}

\begin{insight}{泊松方程的物理意义}{poisson-insight-v2}
泊松方程 $\lap\phi = -\rho/\eps_0$ 是静电学的核心方程。它告诉我们：
\begin{itemize}
\item 正电荷所在处（$\rho > 0$），$\lap\phi < 0$——电势函数在此"向下凹"（像山谷）
\item 负电荷所在处（$\rho < 0$），$\lap\phi > 0$——电势函数在此"向上凸"（像山峰）
\item 无电荷区域，$\lap\phi = 0$——电势函数满足"平均值性质"：任一点的电势等于其周围球面上电势的平均值
\end{itemize}

这个方程的解（给定适当的边界条件）唯一确定了空间中的电势分布，从而通过 $\vecE = -\grad\phi$ 确定了电场。
\end{insight}

\subsection{泊松方程的通解}

\begin{theorem}{电势的积分公式}{potential-integral-v2}
泊松方程的解可写为：
\[
\phi(\vect{r}) = \frac{1}{4\pi\eps_0} \int \frac{\rho(\vect{r}')}{|\vect{r} - \vect{r}'|} \, \diff^3r'
\]
这表示：总电势是所有电荷元产生的 $1/r$ 势的叠加。
\end{theorem}

\textbf{证明}：对积分表达式作用拉普拉斯算子：
\[
\lap\phi(\vect{r}) = \frac{1}{4\pi\eps_0} \int \rho(\vect{r}') \lap\frac{1}{|\vect{r} - \vect{r}'|} \, \diff^3r'
\]

利用 $\lap\dfrac{1}{|\vect{r} - \vect{r}'|} = -4\pi\delta^{(3)}(\vect{r} - \vect{r}')$（第1章已证）：
\[
\lap\phi = \frac{1}{4\pi\eps_0} \int \rho(\vect{r}') (-4\pi) \delta^{(3)}(\vect{r} - \vect{r}') \, \diff^3r' = -\frac{\rho(\vect{r})}{\eps_0}
\]

证毕。

\section{唯一性定理}

唯一性定理是边值问题求解的理论基石。它告诉我们：给定适当的边界条件，泊松/拉普拉斯方程的解是唯一的。

\begin{theorem}{唯一性定理}{uniqueness-theorem-v2}
设区域 $V$ 内的电荷分布 $\rho(\vect{r})$ 给定。若边界 $S$ 上满足以下两类边界条件之一：
\begin{enumerate}
\item \textbf{Dirichlet边界条件}：$\phi|_S = \phi_0(\vect{r})$（边界电势给定）
\item \textbf{Neumann边界条件}：$\partial\phi/\partial n|_S = g(\vect{r})$（边界法向导数给定）
\end{enumerate}
则区域 $V$ 内的电势 $\phi(\vect{r})$ 被唯一确定（Neumann条件下允许差一个任意常数）。
\end{theorem}

\begin{tip}{通俗类比：为什么镜像法靠的就是唯一性定理}{uniqueness-analogy}
唯一性定理听起来抽象，物理含义一句话就能说清：\textbf{"边条件定了，解就定了——不管你怎么找出来的"}。

\textbf{前提条件不要丢}：唯一性定理要求\textbf{求解区域内的电荷分布 $\rho(\vect{r})$ 全部给定}。如果区域内部还有未知的电荷（如未标注的感应电荷分布），仅有边界条件不能唯一确定场——这正是镜像法必须把感应电荷"换成"已知镜像电荷的原因。初学者做题时若发现"边界条件够了解却不唯一"，先检查是否漏给了内部电荷。

一个类比：给定"山谷的地形边界"，山谷里的水系（等高线）就完全确定。你可以用任何方法猜出一张水系图，只要它满足边界条件，它就是\textbf{唯一正确的那张}。

\textbf{这正是镜像法的理论根基}：镜像电荷根本不存在，只是我们\textbf{猜}出来的一个虚构电荷；用它算出的电势恰好满足真实边界（导体表面等势）和泊松方程（区域内的真实电荷）。由唯一性定理，这个"凑出来的解"就是\textbf{唯一正确}的解——虚构的图像、真实的答案。没有唯一性定理，镜像法只是碰运气；有了它，猜对边界就赢了。

\textbf{做题启示}：分离变量法同理——级数解只要满足边界条件，唯一性定理保证它就是答案，不用怀疑"是不是还有别的解"。
\end{tip}

\textbf{证明}（Dirichlet情况）：

设 $\phi_1$ 和 $\phi_2$ 都满足泊松方程和相同的边界条件。令 $\psi = \phi_1 - \phi_2$，则：
\begin{itemize}
\item $\lap\psi = 0$ 在 $V$ 内（因为两者都满足泊松方程）
\item $\psi|_S = 0$（相同的Dirichlet边界条件）
\end{itemize}

考虑恒等式：$\divg(\psi\grad\psi) = |\grad\psi|^2 + \psi\lap\psi = |\grad\psi|^2$（因为 $\lap\psi = 0$）

对 $V$ 积分并用高斯定理：
\[
\int_V |\grad\psi|^2 \, \diff V = \oint_S \psi\grad\psi \cdot \diff\vect{S} = 0
\]
（因为 $\psi = 0$ 在 $S$ 上）

由于 $|\grad\psi|^2 \geq 0$，积分等于零意味着 $\grad\psi = 0$ 在 $V$ 内处处成立，即 $\psi = \text{常数}$。又 $\psi = 0$ 在边界上，故 $\psi \equiv 0$，即 $\phi_1 = \phi_2$。

\textbf{证明}（Neumann情况）：

设 $\phi_1$ 和 $\phi_2$ 都满足泊松方程和相同的Neumann边界条件 $\partial\phi/\partial n|_S = g(\vect{r})$。令 $\psi = \phi_1 - \phi_2$，则：
\begin{itemize}
\item $\lap\psi = 0$ 在 $V$ 内（因为两者都满足泊松方程）
\item $\partial\psi/\partial n|_S = 0$（相同的Neumann边界条件）
\end{itemize}

考虑恒等式：$\divg(\psi\grad\psi) = |\grad\psi|^2 + \psi\lap\psi = |\grad\psi|^2$（因为 $\lap\psi = 0$）

对 $V$ 积分并用高斯定理：
\[
\int_V |\grad\psi|^2 \, \diff V = \oint_S \psi\grad\psi \cdot \diff\vect{S}
= \oint_S \psi \frac{\partial\psi}{\partial n} \, \diff S = 0
\]
（因为 $\partial\psi/\partial n = 0$ 在 $S$ 上）

由于 $|\grad\psi|^2 \geq 0$，积分等于零意味着 $\grad\psi = 0$ 在 $V$ 内处处成立，即 $\psi = \text{常数}$。

因此 $\phi_1 = \phi_2 + \text{常数}$。Neumann边界条件下，电势被确定到相差一个任意常数。

\begin{insight}{Neumann条件下的常数不确定性}{neumann-constant-v2}
Neumann边界条件下解差一个常数，其物理根源在于：电势的\textbf{零点可以任意选取}。只有电势差 $\phi(\vect{b}) - \phi(\vect{a})$ 才有物理意义（对应电场力做的功），而电势的绝对值本身不可测量。

这与Dirichlet条件形成对照：Dirichlet条件通过固定边界上的电势值（通常取接地 $\phi=0$），消除了这种不确定性。在Neumann条件下，如果我们人为选定某一点的电势值（例如 $\phi(\vect{r}_0) = 0$），则解就完全确定了。

从物理上看，Neumann条件给定的是边界上的电场法向分量（即面电荷密度，因为 $E_n = -\partial\phi/\partial n$），它只确定了电荷的分布规律，但没有指定电势的绝对水平。正如在海拔高度问题中，知道了各点的坡度并不能确定绝对海拔，还需要一个基准点。
\end{insight}

\begin{insight}{唯一性定理的重要性}{uniqueness-importance-v2}
唯一性定理的重要性怎么强调都不为过。它告诉我们：
\begin{enumerate}
\item 只要找到一个满足方程和边界条件的解，它就是\textbf{唯一正确的解}——无需担心"还有没有其他解"。
\item \textbf{镜像法}之所以有效，正是因为"镜像电荷 + 真实电荷"产生的势满足相同的边界条件。
\item 数值计算（有限差分、有限元）求解泊松方程时，必须有适当的边界条件——否则解不唯一。
\end{enumerate}

唯一性定理还揭示了静电学的深刻结构：给定区域内的电荷分布和边界条件，完全确定了区域内的场。边界条件不是"额外的物理假设"，而是物理系统的必要组成部分——它们代表了区域外部对内部场的影响。
\end{insight}

\begin{review}{第2章要点回顾}{ch2-review}
\begin{center}
\begin{tabular}{@{}ll@{}}
\toprule
\textbf{方程} & \textbf{物理内容} \\
\midrule
$\vect{F} = \dfrac{1}{4\pi\eps_0}\dfrac{q_1 q_2}{r^2}\unitvec{r}$ & 库仑定律 \\
$\vecE = -\grad\phi$ & 电场与电势的关系 \\
$\displaystyle\oint_S \vecE\cdot\diff\vect{S} = \dfrac{Q_{\text{enc}}}{\eps_0}$ & 高斯定理（积分） \\
$\divg\vecE = \rho/\eps_0$ & 高斯定理（微分） \\
$\curl\vecE = \vect{0}$ & 环路定理（微分） \\
$\lap\phi = -\rho/\eps_0$ & 泊松方程 \\
$\lap\phi = 0$ & 拉普拉斯方程 \\
\bottomrule
\end{tabular}
\end{center}
\end{review}

\begin{tip}{本章核心框架}{ch2-summary}
\textbf{一句话总结}：静电场的有源无旋特性通过高斯定理和环路定理完全刻画。\\
\textbf{核心方程}：$\divg\vecE=\rho/\eps_0$（有源）、$\curl\vecE=\vect{0}$（无旋，保守场）、$\lap\phi=-\rho/\eps_0$（泊松方程）。\\
\textbf{核心方法}：库仑叠加（点电荷）、高斯定理（对称性）、电势法（标量比矢量好算）、唯一性定理（边值问题定解）。\\
\textbf{衔接}：第3章引入电介质，将$\eps_0$替换为$\eps$，新增极化强度$\vect{P}$和电位移$\vect{D}$。
\end{tip}

% ---------- 第2章随堂自测 ----------
\section*{第2章随堂自测}
\addcontentsline{toc}{section}{第2章随堂自测}

\begin{miniquiz}
\textbf{一、选择题}
\begin{enumerate}[label=\arabic*.]
\item 静电场的高斯定理 $\divg\vecE = \rho/\eps_0$ 说明：
  \begin{enumerate}[label=(\Alph*),nosep]
  \item 静电场总是无旋的
  \item 电荷是电场的源
  \item 电场线总是闭合的
  \item 静电场可以被屏蔽
  \end{enumerate}
\begin{ocg}{答案}{ocg答案2}{off}
  \textit{答案：B。散度非零意味着存在源。}
\end{ocg}

\item 导体接地后置于外电场中，导体内部：
  \begin{enumerate}[label=(\Alph*),nosep]
  \item $\vecE=0$, $\rho=0$
  \item $\vecE=0$, $\rho\neq 0$
  \item $\vecE\neq 0$, $\rho=0$
  \item $\vecE\neq 0$, $\rho\neq 0$
  \end{enumerate}
\begin{ocg}{答案}{ocg答案2}{off}
  \textit{答案：A。导体内部电场为零，电荷只能分布在表面。}
\end{ocg}

\item 点电荷的标量势 $\phi$ 随距离的衰减规律是：
  \begin{enumerate}[label=(\Alph*),nosep]
  \item $\propto r$
  \item $\propto 1/r$
  \item $\propto 1/r^2$
  \item $\propto \ln r$
  \end{enumerate}
\begin{ocg}{答案}{ocg答案2}{off}
  \textit{答案：B。$\phi = q/(4\pi\eps_0 r)$。}
\end{ocg}
\end{enumerate}

\vspace{0.5em}
\textbf{二、填空题}
\begin{enumerate}[label=\arabic*.,resume]
\item 静电场的环路定理等价于微分形式 $\curl\vecE=$ \underline{$\vect{0}$}。
\item 引入标量势的定义是 $\vecE = $ \underline{$-\grad\phi$}。
\item 泊松方程的形式为 \underline{$\lap\phi = -\rho/\eps_0$}。
\item 导体表面附近的电场强度 $E = $ \underline{$\sigma/\eps_0$}（法向）。
\end{enumerate}
\end{miniquiz}
\toggleButton{显示答案}{ocg答案2}

\begin{chapterreview}
\begin{itemize}[nosep]
\item \textbf{核心方程}：高斯定理、环路定理、泊松方程、拉普拉斯方程。
\item \textbf{核心方法}：库仑定律叠加法、高斯定理对称性分析、电势积分定义、唯一性定理。
\item \textbf{重要概念}：导体静电平衡条件、镜像法原理。
\item \textbf{自检验}：能否推导均匀带电球体内外的电场？
\item \textbf{衔接预告}：第3章将引入电介质，讨论极化、电位移矢量$\vect{D}$和边界条件。
\end{itemize}
\end{chapterreview}

\begin{tip}{本章自检清单}{ch2-checklist}
学完本章后对照下列条目逐条自测，全部通过即可进入下一章；有未通过项时，回到对应小节复习：
\begin{itemize}[label=$\square$, nosep]
\item 由库仑定律导出电场，并计算典型电荷分布的 $\vecE$ 与 $\phi$
\item 用高斯定理处理球/柱/面对称问题
\item 写出泊松方程并用叠加原理验证特解
\item 说出导体静电平衡的三条性质
\end{itemize}
\end{tip}

% ---------- 第2章综合练习题 ----------
\section*{第2章综合练习题}
\addcontentsline{toc}{section}{第2章综合练习题}
\begin{enumerate}[label=\textbf{\arabic*.}]
\item \textbf{基础} 半径为 $R$ 的均匀带电球体（总电荷 $Q$），用高斯定理求球内外电场分布。
\item \textbf{基础} 求无限大均匀带电平面（面密度 $\sigma$）的电场，并证明两侧电场跃变 $\sigma/\eps_0$。
\item \textbf{提高} 证明：若 $\rho(\vect{r})$ 已知且满足 $r^3\rho(\vect{r})\to 0$（当 $r\to\infty$），则泊松方程的解唯一。
\item \textbf{提高} 一个电偶极子 $\vect{p}$ 位于均匀外电场 $\vecE_0$ 中，求其受力、力矩和势能。
\item \textbf{挑战} 用格林定理证明静电学的唯一性定理（同时覆盖Dirichlet和Neumann边界条件）。
\end{enumerate}
\newpage
\section*{习题解答}
\addcontentsline{toc}{section}{习题解答}

\textbf{题1解析：}

[分析] 本题利用高斯定理的对称性分析，关键是识别球对称性并选择合适的高斯面。

[解答] 由球对称性，$\vecE$ 沿径向，大小仅依赖于 $r$。取半径为 $r$ 的同心球面为高斯面。

\textbf{区域一：$r > R$}（球外）
\[
E \cdot 4\pi r^2 = \frac{Q}{\eps_0} \quad\Rightarrow\quad E(r) = \frac{Q}{4\pi\eps_0 r^2}
\]

\textbf{区域二：$r < R$}（球内）
电荷体密度 $\rho = Q/(\frac{4}{3}\pi R^3)$，$Q_{\text{enc}} = \rho \cdot \frac{4}{3}\pi r^3 = Qr^3/R^3$。
\[
E \cdot 4\pi r^2 = \frac{Qr^3}{\eps_0 R^3} \quad\Rightarrow\quad E(r) = \frac{Qr}{4\pi\eps_0 R^3}
\]

[量纲校验] $[E] = [MLT^{-3}I^{-1}]$，球外 $[Q/(\eps_0 r^2)] = [C]/([F/m][L^2]) = [MLT^{-3}I^{-1}]$；球内 $[Qr/(\eps_0 R^3)]$ 量纲相同，一致。\qedsymbol

\textbf{题2解析：}

[分析] 本题利用平面对称性和高斯定理求无限大带电平面的电场，并验证边界条件。

[解答] 取一个垂直于平面的"药盒"形高斯面，上下底面积各为 $A$，距平面等距。由平面对称性，$\vecE$ 垂直于平面，两侧大小相等方向相反。

上底：$E \cdot A$（外法向与 $E$ 同向）；下底：$E \cdot A$；侧面：无通量。
\[
2EA = \frac{\sigma A}{\eps_0} \quad\Rightarrow\quad E = \frac{\sigma}{2\eps_0}
\]

两侧电场：一侧 $E_+ = \sigma/(2\eps_0)$（远离平面），另一侧 $E_- = -\sigma/(2\eps_0)$（指向平面）。

跃变：$\Delta E = E_+ - E_- = \sigma/(2\eps_0) - (-\sigma/(2\eps_0)) = \sigma/\eps_0$。\qedsymbol

[量纲校验] $[\sigma/\eps_0] = [C/m^2]/[F/m] = [MLT^{-3}I^{-1}] = [E]$，一致。\qedsymbol

\textbf{题3解析：}

[分析] 本题要求证明泊松方程解的唯一性，需要利用格林第二恒等式和边界条件。

[解答] 设 $\phi_1$ 和 $\phi_2$ 都是满足泊松方程和相同边界条件的解。令 $\psi = \phi_1 - \phi_2$，则 $\lap\psi = 0$（拉普拉斯方程），且边界上 $\psi|_S = 0$（Dirichlet）或 $\partial\psi/\partial n|_S = 0$（Neumann）。

格林第一恒等式：$\int_V (\psi\lap\psi + |\grad\psi|^2)\diff V = \oint_S \psi\frac{\partial\psi}{\partial n}\diff S$。

因 $\lap\psi = 0$：$\int_V |\grad\psi|^2\diff V = \oint_S \psi\frac{\partial\psi}{\partial n}\diff S$。

Dirichlet条件：$\psi|_S = 0$，右端为0，故 $\int_V |\grad\psi|^2\diff V = 0$，$\grad\psi = \vect{0}$，$\psi = \text{常数}$。Dirichlet条件下常数必须为0，故 $\phi_1 = \phi_2$。

Neumann条件：$\partial\psi/\partial n|_S = 0$，右端同样为0，故 $\grad\psi = \vect{0}$，$\psi = \text{常数}$。Neumann条件下允许解相差一个常数，电场 $\vecE = -\grad\phi$ 唯一确定。\qedsymbol

[量纲校验] $[\int|\grad\psi|^2\diff V] = [L^{-2}][V^2][L^3] = [V^2L]$，右端 $[\oint\psi\partial\psi/\partial n\,\diff S] = [V][V/L][L^2] = [V^2L]$，一致。\qedsymbol

\textbf{题4解析：}

[分析] 本题考查电偶极子在外电场中的受力、力矩和势能，关键在于理解偶极子与均匀场的相互作用特性。

[解答] 设偶极子由 $-q$（在 $-\vect{d}/2$）和 $+q$（在 $\vect{d}/2$）组成，$\vect{p} = q\vect{d}$。

\textbf{受力}：$\vect{F} = q\vecE_0(-\vect{d}/2) + (-q)\vecE_0(\vect{d}/2) = q(\vecE_0 - \vecE_0) = \vect{0}$（均匀场中受力为零）。

\textbf{力矩}：$\vect{\tau} = \vect{p} \times \vecE_0$。

\textbf{势能}：$U = q\phi(\vect{d}/2) - q\phi(-\vect{d}/2) = -q\vecE_0\cdot\vect{d} = -\vect{p}\cdot\vecE_0$。\qedsymbol

[量纲校验] $[\tau] = [ML^2T^{-2}]$，$[pE] = [CL][MLT^{-3}I^{-1}] = [ML^2T^{-2}]$；$[U] = [ML^2T^{-2}]$，一致。\qedsymbol

\textbf{题5解析：}

[分析] 本题用格林定理证明静电学唯一性定理，同时覆盖Dirichlet和Neumann条件。

[解答] 设 $\phi_1$ 和 $\phi_2$ 满足 $\lap\phi = -\rho/\eps_0$。令 $\psi = \phi_1 - \phi_2$，则 $\lap\psi = 0$。

考虑格林第二恒等式：$\int_V (\psi\lap\varphi - \varphi\lap\psi)\diff V = \oint_S \left(\psi\frac{\partial\varphi}{\partial n} - \varphi\frac{\partial\psi}{\partial n}\right)\diff S$（$\varphi$ 为任一光滑函数）。

更直接地，用格林第一恒等式：$\int_V |\grad\psi|^2\diff V = \oint_S \psi\partial\psi/\partial n\,\diff S$。

\textbf{Dirichlet}：$\psi|_S = 0$，$\partial\psi/\partial n$ 任意，右端为0，$\grad\psi = \vect{0}$，$\psi = 0$。

\textbf{Neumann}：$\partial\psi/\partial n|_S = 0$，$\psi$ 任意，右端为0，$\grad\psi = \vect{0}$，$\psi = \text{常数}$。\qedsymbol

[量纲校验] 同上题，一致。\qedsymbol



\begin{warning}{第2章常见误区}{ch2-warnings}
\begin{enumerate}
\item $\vecE = -\grad\phi$ 中的负号！电场指向电势\textbf{下降}方向。
\item 高斯定理要求闭合曲面；计算通量时注意 $\diff\vect{S}$ 的方向（外法向）。
\item 电偶极子场按 $1/r^3$ 衰减，不是 $1/r^2$。
\item 导体内部 $E=0$（静电平衡），但电势可以不为零——$E=0$ 只意味着 $\phi = \text{常数}$。
\item 唯一性定理的Neumann条件要求边界上给定的是$\partial\phi/\partial n$（法向导数），不是$\partial\phi/\partial x$等。
\end{enumerate}
\end{warning}


% ============================================================

\begin{center}\small \hyperlink{ch1-goal}{上一章：第1章} \quad|\quad \hyperlink{ch3-goal}{下一章：第3章}\end{center}

% 第3章：电介质与静电边界条件（新增章节，优化版）
% ============================================================
\chapter{电介质与静电边界条件}
\epigraph{物质的结构决定了它的性质。}{--- 德谟克利特}

\hypertarget{ch3-goal}{}\begin{insight}{本章核心目标}{ch3-goal-v2}
理解电介质在电场中的响应机制——极化。引入电位移矢量 $\vect{D}$ 简化方程，建立不同介质界面处的边界条件。这些边界条件是求解"分区均匀"介质中电场的关键工具，也是理解电容、电磁波反射等现象的基础。
\end{insight}

前几章讨论的是真空（或空气近似）中的静电场。然而，实际的电场几乎总是在某种物质中存在。当电场作用于物质时，物质中的电荷会重新分布，反过来又影响电场。本章研究这种相互作用的最简单形式——\textbf{电介质}中的静电场。

\begin{tip}{本章新符号}{ch3-new-symbols}
\begin{tabular}{@{}ll@{}}
$\vect{P}$ & 极化强度（偶极矩密度，单位 \si{C/m^2}） \\
$\vect{D}$ & 电位移矢量（$\vect{D} = \eps_0\vecE + \vect{P}$，单位 \si{C/m^2}） \\
$\chi_e$ & 电极化率（无量纲，$\eps_r = 1 + \chi_e$） \\
$\eps_r,\ \eps$ & 相对/绝对介电常数（$\eps = \eps_0\eps_r$） \\
$\sigma_f,\ \sigma_b$ & 自由/束缚面电荷密度 \\
\end{tabular}
\end{tip}

\begin{history}{电介质研究的历史}{dielectric-history}
电介质的系统研究始于法拉第。他在1837年发现，插入绝缘材料（如玻璃或硫磺）后，电容器的电容会增大。他将这种现象归因于材料"极化"——内部电荷在电场作用下发生微小位移，产生与外加电场相反的极化场。\textbf{法拉第}引入了"介电常数"的概念，但数学描述要等到\textbf{麦克斯韦}才完善。现代对电介质的理解基于原子尺度的极化机制，包括电子极化、离子极化和取向极化。
\end{history}

\section{电介质的极化}

\subsection{极化的物理机制}

电介质由原子或分子组成。在外加电场作用下，这些微观结构发生响应，产生净的极化效应。主要有三种机制：

\begin{enumerate}
\item \textbf{电子极化}：外电场使原子中的电子云相对于原子核发生微小位移，形成感应电偶极矩。这种极化在所有电介质中都存在，响应时间极短（$\sim 10^{-15}$秒）。

\item \textbf{离子极化}：在离子晶体（如NaCl）中，外电场使正负离子沿相反方向发生微小位移，产生感应电偶极矩。响应时间约为 $10^{-13}$秒。

\item \textbf{取向极化}：具有永久电偶极矩的分子（如H$_2$O）在外电场作用下，原本随机取向的偶极子趋向于沿电场方向排列。这种极化与温度有关——热运动会干扰有序排列。响应时间较慢（$\sim 10^{-12}$至$10^{-6}$秒）。
\end{enumerate}

\subsection{极化强度}

\begin{definition}{极化强度}{polarization-v2}
\textbf{极化强度} $\vect{P}$ 定义为单位体积内的电偶极矩：
\[
\vect{P} = \lim_{\Delta V \to 0} \frac{\sum_i \vect{p}_i}{\Delta V}
\]
其中 $\vect{p}_i$ 是体积元 $\Delta V$ 内第 $i$ 个分子的电偶极矩。$\vect{P}$ 的单位是 $\si{C/m^2}$（与面电荷密度单位相同，这并非巧合）。
\end{definition}

\begin{insight}{极化强度与面束缚电荷}{polarization-charge-v2}
极化电介质内部的电偶极子虽然整体中性，但在电介质\textbf{表面}会出现\textbf{束缚电荷}（bound charge）。考虑一个均匀极化的电介质块：内部的正、负电荷中心相互抵消，但在上表面露出负电荷，下表面露出正电荷。

定量地，极化产生的\textbf{体束缚电荷密度}和\textbf{面束缚电荷密度}分别为：
\begin{align}
\rho_b &= -\divg\vect{P} \label{eq:rho-bound} \\
\sigma_b &= \vect{P} \cdot \unitvec{n} \label{eq:sigma-bound}
\end{align}
其中 $\unitvec{n}$ 是电介质表面的外法向单位矢量。
\end{insight}

\textbf{式\eqref{eq:rho-bound}的推导}：考虑体积元 $\Delta V$ 内的总偶极矩为 $\vect{P}\Delta V$。用宏观观点，这等价于在体积元边界上分布的面电荷。通过高斯定理类型的分析可以证明，体束缚电荷密度正是 $-\divg\vect{P}$。如果极化是均匀的（$\vect{P} = \text{常数}$），则 $\divg\vect{P} = 0$，内部无体束缚电荷——只有表面有束缚电荷。

\section{电位移矢量}

\begin{tip}{为什么要引入 $\vect{D}$？}{why-D-tip}
明明有 $\vecE$，为什么还要定义 $\vect{D}$？动机是一个"打包"操作：

介质中的高斯定理 $\divg\vecE = (\rho_f + \rho_b)/\eps_0$ 用起来很别扭——$\rho_f$（自由电荷，通常已知可控）和 $\rho_b$（束缚电荷，取决于介质响应）混在\textbf{一个方程}里。做题时你往往只知道 $\rho_f$，却要先算出 $\rho_b$ 才能解 $\vecE$。

$\vect{D} \equiv \eps_0\vecE + \vect{P}$ 的作用是\textbf{把束缚电荷打包进定义}：$\divg\vect{D} = \rho_f$——源只剩自由电荷。求解流程变成：先用已知的 $\rho_f$ 解 $\vect{D}$，再由 $\vecE = \vect{D}/\eps$（线性介质）回代得 $\vecE$。

\textbf{代价}：$\vect{D}$ 不是真实力场——电荷受力是 $q\vecE$，\textbf{不是} $q\vect{D}$（见下方"易错点"）。$\vect{D}$ 只是让方程更好解的辅助量。
\end{tip}

极化电荷与自由电荷一样产生电场。因此，高斯定理应写为：
\[
\divg\vecE = \frac{\rho_f + \rho_b}{\eps_0} = \frac{\rho_f - \divg\vect{P}}{\eps_0}
\]
其中 $\rho_f$ 是\textbf{自由电荷密度}（可由外部控制的电荷，如导体上的电荷、掺杂的离子等），$\rho_b$ 是\textbf{束缚电荷密度}。

整理得：
\[
\divg(\eps_0\vecE + \vect{P}) = \rho_f
\]

\begin{definition}{电位移矢量}{d-field-v2}
定义\textbf{电位移矢量}：
\[
\vect{D} = \eps_0\vecE + \vect{P}
\]
则高斯定理变为：
\[
\divg\vect{D} = \rho_f
\]
\textbf{关键性质}：$\vect{D}$ 的源仅是自由电荷，与束缚电荷无关！
\end{definition}

\begin{insight}{理解D的物理意义}{d-meaning-v2}
$\vect{D}$ 是一个\textbf{辅助场量}，它的引入是为了简化方程——用自由电荷（通常已知且可控）作为源，而"隐藏"了复杂的极化响应。但要注意：
\begin{itemize}
\item $\vect{D}$ 本身\textbf{不是}基本物理量，它是为了数学方便而定义的。
\item 真正有力学效应的是 $\vecE$（电场对电荷的作用力 $q\vecE$）。
\item 在不同介质的边界上，$\vect{D}$ 的边界条件比 $\vecE$ 更简单。
\end{itemize}

可以这样理解：总电场 $\vecE$ 是"自由电荷产生的电场"加上"束缚电荷产生的电场"。$\vect{D}$ 则"提前计入"了极化效应，使其方程形式上只与自由电荷有关。
\end{insight}

\begin{tip}{E、D、P 三场完整对比表}{edp-compare-table}
\begin{center}
\small
\begin{tabular}{@{}>{\raggedright\arraybackslash}p{2.2cm}>{\raggedright\arraybackslash}p{3.6cm}>{\raggedright\arraybackslash}p{3.2cm}>{\raggedright\arraybackslash}p{3.6cm}@{}}
\toprule
 & $\vecE$（电场强度） & $\vect{D}$（电位移矢量） & $\vect{P}$（极化强度） \\
\midrule
定义来源 & 基本场（库仑力 $q\vecE$ 定义） & 辅助定义：$\vect{D} = \eps_0\vecE + \vect{P}$ & 偶极矩密度 $\vect{P} = \sum\vect{p}_i/\Delta V$ \\
\midrule
方程源 & $\divg\vecE = (\rho_f+\rho_b)/\eps_0$（全部电荷） & $\divg\vect{D} = \rho_f$（仅自由电荷） & $\rho_b = -\divg\vect{P}$（其散度即束缚电荷） \\
\midrule
单位 & \si{V/m} & \si{C/m^2} & \si{C/m^2} \\
\midrule
可测量性 & 可直接测量（试探电荷受力） & 不可直接测量 & 不可直接测量 \\
\midrule
线性介质中 & $\vecE = \vect{D}/\eps$ & $\vect{D} = \eps\vecE$ & $\vect{P} = \eps_0(\eps_r-1)\vecE$ \\
\midrule
边界条件 & 切向连续：$E_{2t}=E_{1t}$ & 法向：$D_{2n}-D_{1n}=\sigma_f$ & 法向：$P_{2n}-P_{1n}=-\sigma_b$ \\
\midrule
物理角色 & 决定受力与能量 & 简化方程，绕开束缚电荷 & 描述介质响应本身 \\
\bottomrule
\end{tabular}
\end{center}

\textbf{易错点}：$\vect{D}$ 和 $\vect{P}$ 单位相同（\si{C/m^2}）但物理含义完全不同——$\vect{D}$ 包含真空部分 $\eps_0\vecE$，$\vect{P}$ 只描述介质响应。在真空中 $\vect{P}=0$ 但 $\vect{D}=\eps_0\vecE\neq 0$。判断介质分界面上的物理问题时，先分清问的是哪一种"面电荷"：自由电荷用 $\vect{D}$ 的法向跃变，束缚电荷用 $\vect{P}$ 的法向跃变。
\end{tip}

\subsection{线性各向同性电介质}

对于\textbf{线性各向同性}电介质，实验表明：
\[
\vect{P} = \eps_0\chi_e \vecE
\]
其中 $\chi_e$ 是\textbf{电极化率}（electric susceptibility），无量纲。

于是：
\[
\vect{D} = \eps_0\vecE + \eps_0\chi_e\vecE = \eps_0(1+\chi_e)\vecE = \eps_0\eps_r\vecE = \eps\vecE
\]

\begin{definition}{介电常数}{permittivity-v2}
\textbf{相对介电常数}（或介电常数）：$\eps_r = 1 + \chi_e$

\textbf{绝对介电常数}：$\eps = \eps_0\eps_r$

常见材料的相对介电常数：真空 $\eps_r = 1$，空气 $\eps_r \approx 1.0006$，水 $\eps_r \approx 80$，玻璃 $\eps_r \approx 4$-$10$，陶瓷 $\eps_r \approx 6$-$1000$。
\end{definition}

\begin{example}{平行板电容器中的电介质}{capacitor-dielectric-ex-v2}
\difficulty{基础}
平行板电容器板间距为 $d$，面积为 $A$，两板间充满相对介电常数为 $\eps_r$ 的电介质。极板上自由电荷面密度为 $\sigma_f$。求：(a) $\vect{D}$、$\vecE$、$\vect{P}$；(b) 束缚电荷面密度；(c) 电容。
\tcblower
\textbf{解答：}

(a) 用高斯定理求 $\vect{D}$：取底面在极板内、顶面在两板之间的"药盒"形高斯面。由于极板内 $E=0$（导体），只有顶面贡献：
\[
D \cdot A = \sigma_f A \quad\Rightarrow\quad D = \sigma_f
\]
方向从正极板指向负极板。

$E = D/\eps = \sigma_f/(\eps_0\eps_r)$

$P = \eps_0\chi_e E = \eps_0(\eps_r - 1)\dfrac{\sigma_f}{\eps_0\eps_r} = \dfrac{(\eps_r - 1)\sigma_f}{\eps_r}$

(b) 束缚电荷面密度：

在上表面（靠近正极板）：$\sigma_b = \vect{P}\cdot\unitvec{n} = -P = -\dfrac{(\eps_r-1)\sigma_f}{\eps_r}$（负号因为 $\vect{P}$ 指向负极板，而表面外法向指向正极板）

在下表面：$\sigma_b = +\dfrac{(\eps_r-1)\sigma_f}{\eps_r}$

(c) 电容：$C = \dfrac{Q}{V} = \dfrac{\sigma_f A}{E d} = \dfrac{\sigma_f A}{\sigma_f d/(\eps_0\eps_r)} = \dfrac{\eps_0\eps_r A}{d} = \eps_r C_0$

其中 $C_0 = \eps_0 A/d$ 是无电介质时的电容。插入电介质后，电容增大了 $\eps_r$ 倍——这正是法拉第1837年的发现！
\end{example}

\section{静电边界条件}

不同介质的界面（或导体与介质的界面）上，电场会发生"跳跃"。边界条件描述了这种跳跃的规律，是求解分区均匀介质中电场的基本工具。

\subsection{电场的法向分量}

在界面处取一个极薄的"药盒"形高斯面，上下底面分别在两种介质中：

\begin{theorem}{D法向分量的边界条件}{d-normal-bc-v2}
\[
D_{2n} - D_{1n} = \sigma_f
\]
或矢量形式：
\[
\unitvec{n} \cdot (\vect{D}_2 - \vect{D}_1) = \sigma_f
\]
其中 $\unitvec{n}$ 是从介质1指向介质2的单位法向矢量，$\sigma_f$ 是界面上的\textbf{自由电荷面密度}。

若界面上无自由电荷（$\sigma_f = 0$）：
\[
D_{1n} = D_{2n}\text{（法向 } \vect{D} \text{ 连续）}
\]
\end{theorem}

\textbf{推导}：对药盒应用高斯定理 $\oint\vect{D}\cdot\diff\vect{S} = Q_{f,\text{enc}}$。当药盒厚度趋于零时，侧面贡献消失，只剩下两底面。$Q_{f,\text{enc}} = \sigma_f A$，$\vect{D}_2\cdot\unitvec{n}A - \vect{D}_1\cdot\unitvec{n}A = \sigma_f A$，即得结果。

\begin{warning}{初学者易错：跃变只来自自由电荷}{d-normal-bc-error}
\textbf{最容易踩的坑}：$D_{2n} - D_{1n} = \sigma_f$ 中的 $\sigma_f$ 是\textbf{自由面电荷}，不是总面电荷！极化束缚电荷 $\sigma_b = \vect{P}\cdot\unitvec{n}$ \textbf{不会}出现在 $\vect{D}$ 的边界条件里——它已经被打包进 $\vect{D}$ 的定义（$\vect{D} = \eps_0\vecE + \vect{P}$）。

\textit{反例}：介质-真空界面无自由电荷（$\sigma_f = 0$）时，$\vect{D}$ 法向连续，但 $\vecE$ 法向\textbf{不}连续（$E_{2n} - E_{1n} = \sigma_b/\eps_0$）——束缚电荷在 $\vecE$ 的边界条件里体现。初学者常把 $\sigma_b$ 代入 $\vect{D}$ 的公式，得到 $\vect{D}$ 法向不连续的错误结论。

\textbf{记忆口诀}：$\vect{D}$ 的跃变问自由电荷，$\vecE$ 的跃变问总电荷（自由 + 束缚）。
\end{warning}

\subsection{电场的切向分量}

在界面处取一个极窄的矩形回路，长边分别位于两种介质中、平行于界面：

\begin{theorem}{E切向分量的边界条件}{e-tangential-bc-v2}
\[
E_{2t} - E_{1t} = 0
\]
或矢量形式：
\[
\unitvec{n} \times (\vecE_2 - \vecE_1) = \vect{0}
\]
即\textbf{切向电场总是连续的}，无论界面上是否有电荷。
\end{theorem}

\textbf{推导}：对矩形回路应用静电场的环路定理 $\oint\vecE\cdot\diff\vect{l} = 0$。当矩形宽度趋于零时，短边贡献消失，长边贡献为 $E_{1t}l - E_{2t}l = 0$。

\begin{insight}{边界条件的物理图像}{bc-intuition-v2}
\begin{itemize}
\item \textbf{法向D跳跃}：反映了界面上的自由面电荷。这些自由电荷是电位移场的"源"，就像 $\divg\vect{D} = \rho_f$ 在局部的表现。
\item \textbf{切向E连续}：反映了静电场的无旋性。如果切向电场不连续，意味着沿界面存在无穷大的旋度——这与 $\curl\vecE = 0$ 矛盾。
\end{itemize}

边界条件的本质是：将微分方程在\textbf{不连续处}的"极限"行为提取出来。它们是连接不同区域的"桥梁"，没有它们，分区求解就无从谈起。
\end{insight}

\begin{example}{电介质-真空界面}{dielectric-vacuum-bc-ex-v2}
\difficulty{提高}
区域1（$z<0$）为真空，区域2（$z>0$）为相对介电常数 $\eps_r$ 的电介质。已知真空中电场为 $\vecE_1 = E_0(\sin\theta\,\unitvec{x} + \cos\theta\,\unitvec{z})$，界面无自由电荷。求电介质中的电场 $\vecE_2$ 和界面上的束缚电荷面密度。
\tcblower
\textbf{解答：}

界面法向为 $\unitvec{n} = \unitvec{z}$（从1指向2）。

由 $D_{1n} = D_{2n}$（无自由电荷）：
\[
\eps_0 E_{1z} = \eps_0\eps_r E_{2z} \quad\Rightarrow\quad E_{2z} = \frac{E_{1z}}{\eps_r} = \frac{E_0\cos\theta}{\eps_r}
\]

由 $E_{1t} = E_{2t}$：
\[
E_{2x} = E_{1x} = E_0\sin\theta, \quad E_{2y} = E_{1y} = 0
\]

因此：
\[
\vecE_2 = E_0\sin\theta\,\unitvec{x} + \frac{E_0\cos\theta}{\eps_r}\unitvec{z}
\]

$\vect{D}_2 = \eps_0\eps_r\vecE_2 = \eps_0\eps_r E_0\sin\theta\,\unitvec{x} + \eps_0 E_0\cos\theta\,\unitvec{z}$

极化强度：$\vect{P} = \vect{D}_2 - \eps_0\vecE_2 = \eps_0(\eps_r-1)E_0\sin\theta\,\unitvec{x} + \eps_0(\eps_r-1)\dfrac{E_0\cos\theta}{\eps_r}\unitvec{z}$

界面束缚电荷面密度：
\[
\sigma_b = \vect{P}\cdot\unitvec{n} = -P_z = -\eps_0(\eps_r-1)\frac{E_0\cos\theta}{\eps_r} \qquad (\unitvec{n} = -\unitvec{z}\ \text{为介质的外法向})
\]

其中 $\unitvec{n}$ 取\textbf{介质的外法向}——介质占据 $z>0$，界面 $z=0$ 处的外法向是 $-\unitvec{z}$（由介质指向真空），而不是 $+\unitvec{z}$。因此当 $E_0\cos\theta>0$ 时，界面上出现的是\textbf{负}束缚电荷：$\sigma_b = -P_z < 0$。

\textbf{物理解读}：$\vecE_1$ 的 $z$ 分量为正意味着真空侧场线指向 $+z$。介质中分子偶极沿场排列（偶极正端朝 $+z$），而界面 $z=0$ 是介质朝向真空的那一侧表面，露出的正是偶极的\textbf{负端}，故 $\sigma_b = \vect{P}\cdot(-\unitvec{z}) = -P_z < 0$。这与介质内场被削弱（$E_{2z} = E_{1z}/\eps_r < E_{1z}$）完全自洽：界面上的负束缚电荷在介质内产生与 $\vecE_1$ 反向的场。

\textbf{易混淆点}：判断 $\sigma_b$ 符号的可靠方法永远是 $\sigma_b = \vect{P}\cdot\unitvec{n}$，且 $\unitvec{n}$ 必须取\textbf{该表面的介质外法向}——本例中介质在 $z>0$、界面处外法向为 $-\unitvec{z}$；若误取 $+\unitvec{z}$（那是真空侧的外法向），符号就会弄反。一个可靠的自检：束缚电荷总要\textbf{削弱介质内的场}——本题 $E_{2z} = E_{1z}/\eps_r < E_{1z}$，故界面必带负束缚电荷。
\end{example}

\begin{warning}{边界条件的符号约定}{bc-sign-warning}
使用边界条件时，法向 $\unitvec{n}$ 的方向至关重要。通常约定 $\unitvec{n}$ 从介质1指向介质2。则：
\begin{itemize}
\item $D_{2n} - D_{1n} = \sigma_f$ 中，$D_{1n} = \vect{D}_1\cdot\unitvec{n}$，$D_{2n} = \vect{D}_2\cdot\unitvec{n}$
\item $\sigma_b = \vect{P}\cdot\unitvec{n}$ 给出的是从介质指向外的束缚电荷
\end{itemize}
务必在计算开始时明确法向方向，并始终一致。
\end{warning}

\begin{example}{介质包裹导体球}{coated-conductor-sphere-ex-v2}
\difficulty{提高}
半径为 $a$ 的导体球，外部包裹一层外半径为 $b$、介电常数为 $\eps$ 的均匀电介质，导体球带电 $Q$。求：(a) 各区域电场分布；(b) 导体球电势（取无穷远为零势点）；(c) 电介质内的束缚电荷。

\begin{center}
\begin{tikzpicture}[scale=1.2]
% 真空区（最外层，白色）
\draw[black,thick] (0,0) circle (3);
% 电介质层（蓝色环）
\filldraw[fill=blue!15,draw=black,thick] (0,0) circle (2);
% 导体球（灰色实心）
\filldraw[fill=gray!35,draw=black,thick] (0,0) circle (1);
% 半径标注
\draw[dashed] (0,0) -- (0.707,-0.707) node[midway,right] {$a$};
\draw[dashed] (0,0) -- (1.414,-1.414) node[midway,right] {$b$};
% 区域标签
\node[font=\small] at (0,0.35) {导体};
\node[font=\small] at (0,-0.45) {$\sigma_f$};
\node[font=\small,blue!70!black] at (0,1.5) {电介质 $\eps$};
\node[font=\small] at (2.45,0) {真空 $\eps_0$};
% 电场方向（介质内径向）
\draw[->,thick,red] (0.9,0.35) -- (1.75,0.68) node[above,midway,font=\small] {$\vecE_{II}$};
\draw[->,thick,red] (2.1,-0.6) -- (2.8,-0.8) node[below,right,font=\small] {$\vecE_{III}$};
\end{tikzpicture}
\end{center}
\tcblower
\textbf{解答：}

由球对称性，电场只有径向分量。将空间分为三个区域：

\textbf{区域 I：$r < a$（导体内部）}

导体静电平衡：$\vecE_I = \vect{0}$

\textbf{区域 II：$a < r < b$（电介质中）}

取半径为 $r$ 的同心球面为高斯面，用 $\vect{D}$ 的高斯定理：
\[
\oint \vect{D}\cdot\diff\vect{S} = D_{II}(r) \cdot 4\pi r^2 = Q
\]
\[
D_{II}(r) = \frac{Q}{4\pi r^2}, \quad \vect{D}_{II} = \frac{Q}{4\pi r^2}\unitvec{r}
\]

由 $\vect{D} = \eps\vecE$：
\[
\vecE_{II} = \frac{Q}{4\pi\eps r^2}\unitvec{r}
\]

\textbf{区域 III：$r > b$（真空中）}

同理：
\[
\vect{D}_{III} = \frac{Q}{4\pi r^2}\unitvec{r}, \quad \vecE_{III} = \frac{Q}{4\pi\eps_0 r^2}\unitvec{r}
\]

\textbf{(b) 导体球电势}

导体是等势体，电势为：
\[
\phi(a) = -\int_\infty^a \vecE\cdot\diff\vect{r}
= -\int_\infty^b E_{II}\,\diff r - \int_b^a E_{I}\,\diff r
\]

\[
= -\int_\infty^b \frac{Q}{4\pi\eps_0 r^2}\,\diff r - \int_b^a \frac{Q}{4\pi\eps r^2}\,\diff r
\]

\[
= \frac{Q}{4\pi\eps_0 b} + \frac{Q}{4\pi\eps}\left(\frac{1}{a} - \frac{1}{b}\right)
\]

\textbf{(c) 电介质内的束缚电荷}

电介质内的极化强度：
\[
\vect{P} = \vect{D}_{II} - \eps_0\vecE_{II}
= \frac{Q}{4\pi r^2}\unitvec{r} - \frac{\eps_0 Q}{4\pi\eps r^2}\unitvec{r}
= \frac{(\eps - \eps_0)Q}{4\pi\eps r^2}\unitvec{r}
\]

\textbf{体束缚电荷}：由于 $\vect{P}$ 的径向依赖为 $1/r^2$，
\[
\rho_b = -\divg\vect{P} = -\frac{1}{r^2}\frac{\partial}{\partial r}(r^2 P_r)
= -\frac{1}{r^2}\frac{\partial}{\partial r}\left(\frac{(\eps-\eps_0)Q}{4\pi\eps}\right) = 0
\]
体束缚电荷为零（均匀极化的结果）。

\textbf{面束缚电荷}：

在 $r = a$ 处（电介质内表面，法向指向电介质内部即 $-\unitvec{r}$）：
\[
\sigma_{b,a} = -\vect{P}\cdot\unitvec{r}\big|_{r=a} = -\frac{(\eps-\eps_0)Q}{4\pi\eps a^2}
\]

在 $r = b$ 处（电介质外表面，法向指向外即 $+\unitvec{r}$）：
\[
\sigma_{b,b} = \vect{P}\cdot\unitvec{r}\big|_{r=b} = \frac{(\eps-\eps_0)Q}{4\pi\eps b^2}
\]

\textbf{验证}：总束缚电荷
\[
Q_b = \sigma_{b,a}\cdot 4\pi a^2 + \sigma_{b,b}\cdot 4\pi b^2
= -\frac{(\eps-\eps_0)Q}{\eps} + \frac{(\eps-\eps_0)Q}{\eps} = 0
\]
符合预期——束缚电荷总量为零。

\textbf{边界条件验证}：
\begin{itemize}
\item 在 $r = a$：$D_{II,r} - D_{I,r} = Q/(4\pi a^2) - 0 = \sigma_f$（自由面电荷密度）$\checkmark$
\item 在 $r = a$：$E_{II,t} = E_{I,t} = 0$（无切向分量）$\checkmark$
\item 在 $r = b$：$D_{III,r} = D_{II,r}$（无自由电荷）$\checkmark$
\item 在 $r = b$：$E_{III,t} = E_{II,t} = 0$（无切向分量）$\checkmark$
\end{itemize}
\end{example}

\section{导体表面的边界条件}

导体是电介质的一个极端情形：内部存在大量自由电子，可以任意移动。静电平衡时：
\begin{enumerate}
\item 导体内部 $\vecE = 0$（否则自由电荷会继续移动）
\item 导体内部 $\phi = \text{常数}$（整个导体是等势体）
\item 导体表面 $\vecE$ 垂直于表面（否则切向分量会驱动电荷移动）
\item 电荷只分布在导体表面
\end{enumerate}

\begin{theorem}{导体表面的边界条件}{conductor-bc-v2}
设导体为介质1，外部介质2的介电常数为 $\eps$。则：
\begin{align}
\vecE_{\text{内}} &= \vect{0} \\
E_{\text{外},n} &= \frac{\sigma}{\eps} \\
E_{\text{外},t} &= 0
\end{align}
或用电势表示：$\phi|_{\text{导体表面}} = \text{常数}$
\end{theorem}

\begin{review}{第3章要点回顾}{ch3-review}
\begin{center}
\begin{tabular}{@{}ll@{}}
\toprule
\textbf{公式/概念} & \textbf{内容} \\
\midrule
$\vect{P}$ & 极化强度 = 单位体积偶极矩 \\
$\rho_b = -\divg\vect{P}$ & 体束缚电荷密度 \\
$\sigma_b = \vect{P}\cdot\unitvec{n}$ & 面束缚电荷密度 \\
$\vect{D} = \eps_0\vecE + \vect{P}$ & 电位移矢量 \\
$\divg\vect{D} = \rho_f$ & D的高斯定理 \\
$\vect{D} = \eps\vecE$（线性介质） & 本构关系 \\
$\eps = \eps_0\eps_r$，$\eps_r = 1+\chi_e$ & 介电常数 \\
$D_{2n} - D_{1n} = \sigma_f$ & D法向边界条件 \\
$E_{2t} = E_{1t}$ & E切向边界条件 \\
$\phi|_{\text{导体}} = \text{常数}$ & 导体边界条件 \\
\bottomrule
\end{tabular}
\end{center}
\end{review}

\begin{tip}{本章核心框架}{ch3-summary}
\textbf{一句话总结}：电介质通过极化响应外电场，产生束缚电荷，改变有效介电常数。\\
\textbf{核心关系}：$\vect{D}=\eps_0\vecE+\vect{P}=\eps\vecE$（LIH介质）、$\rho_p=-\divg\vect{P}$（束缚电荷）。\\
\textbf{边界条件}：$E_{1t}=E_{2t}$（切向连续）、$D_{1n}=D_{2n}$（法向连续，无自由面电荷）。\\
\textbf{衔接}：第4章将用分离变量法、镜像法等解析方法求解含电介质的边值问题。
\end{tip}

% ---------- 第3章随堂自测 ----------
\section*{第3章随堂自测}
\addcontentsline{toc}{section}{第3章随堂自测}

\begin{miniquiz}
\textbf{一、选择题}
\begin{enumerate}[label=\arabic*.]
\item 电位移矢量 $\vect{D}$ 的定义是：
  \begin{enumerate}[label=(\Alph*),nosep]
  \item $\vect{D} = \eps_0\vecE$
  \item $\vect{D} = \eps_0\vecE + \vect{P}$
  \item $\vect{D} = \eps\vecE$
  \item $\vect{D} = \vecE/\eps_0$
  \end{enumerate}
\begin{ocg}{答案}{ocg答案3}{off}
  \textit{答案：B。这是$\vect{D}$的普适定义，与介质性质无关。}
\end{ocg}

\item 在电介质界面上，电场法向分量的边界条件是：
  \begin{enumerate}[label=(\Alph*),nosep]
  \item $E_{1n} = E_{2n}$
  \item $\eps_1 E_{1n} = \eps_2 E_{2n}$
  \item $D_{1n} = D_{2n}$（无自由面电荷时）
  \item $\vecE_1 = \vecE_2$
  \end{enumerate}
\begin{ocg}{答案}{ocg答案3}{off}
  \textit{答案：C。$\vect{D}$的法向分量在无自由面电荷时连续。}
\end{ocg}

\item 极化强度 $\vect{P}$ 与极化电荷密度的关系是：
  \begin{enumerate}[label=(\Alph*),nosep]
  \item $\rho_p = \divg\vect{P}$
  \item $\rho_p = -\divg\vect{P}$
  \item $\rho_p = \curl\vect{P}$
  \item $\rho_p = 0$
  \end{enumerate}
\begin{ocg}{答案}{ocg答案3}{off}
  \textit{答案：B。极化电荷是束缚电荷，满足$\rho_p = -\divg\vect{P}$。}
\end{ocg}
\end{enumerate}

\vspace{0.5em}
\textbf{二、判断题}
\begin{enumerate}[label=\arabic*.,resume]
\item 线性各向同性均匀介质中，$\vect{D} = \eps\vecE = \eps_0\eps_r\vecE$。
\begin{ocg}{答案}{ocg答案3}{off}
  \textit{答案：正确。这是LIH介质的本构关系。}
\end{ocg}

\item 在导体与电介质的界面上，$\vect{D}$ 的切向分量连续。
\begin{ocg}{答案}{ocg答案3}{off}
  \textit{答案：错误。切向分量连续的是$\vecE$，不是$\vect{D}$。}
\end{ocg}
\end{enumerate}
\end{miniquiz}
\toggleButton{显示答案}{ocg答案3}

\begin{chapterreview}
\begin{itemize}[nosep]
\item \textbf{核心概念}：极化强度$\vect{P}$、极化电荷、电位移矢量$\vect{D}$、介电常数$\eps_r$。
\item \textbf{边界条件}：$\vecE$切向连续、$\vect{D}$法向连续（无自由面电荷）；$\vecB$法向连续、$\vect{H}$切向连续（无自由面电流）。
\item \textbf{自检验}：能否从$\vect{D}$的定义出发推导$\divg\vect{D}=\rho_f$？
\item \textbf{衔接预告}：第4章将用分离变量法、镜像法等解析方法求解边值问题。
\end{itemize}
\end{chapterreview}

\begin{tip}{本章自检清单}{ch3-checklist}
学完本章后对照下列条目逐条自测，全部通过即可进入下一章；有未通过项时，回到对应小节复习：
\begin{itemize}[label=$\square$, nosep]
\item 由极化机制推导 $\rho_b = -\divg\vect{P}$
\item 区分 $\vecE$、$\vect{D}$、$\vect{P}$ 的源与单位
\item 写出电介质界面的法向/切向边界条件
\item 计算含介质的平行板电容器
\end{itemize}
\end{tip}

% ---------- 第3章综合练习题 ----------
\section*{第3章综合练习题}
\addcontentsline{toc}{section}{第3章综合练习题}
\begin{enumerate}[label=\textbf{\arabic*.}]
\item \textbf{基础} 平行板电容器充满介电常数为 $\eps_r$ 的介质，极板电荷面密度 $\pm\sigma$，求 $D$、$E$、$P$ 及束缚电荷。
\item \textbf{基础} 证明：在均匀极化的介质球表面，束缚面电荷密度 $\sigma_b = P\cos\theta$（$\theta$ 为极角）。
\item \textbf{提高} 半径为 $a$ 的导体球带电 $Q$，外部包裹外半径 $b$、介电常数 $\eps$ 的电介质。求各区域电场、导体球电势及电介质内束缚电荷。
\item \textbf{提高} 两种各向异性介质的分界面上，$\vect{D}$ 和 $\vecE$ 的方向如何关系？给出折射定律的各向异性推广。
\item \textbf{挑战} 利用边界条件证明，在导体表面外侧，电场总是垂直于导体表面，且 $E=\sigma/\eps_0$。
\end{enumerate}
\newpage
\section*{习题解答}
\addcontentsline{toc}{section}{习题解答}

\textbf{题1解析：}

[分析] 本题是平行板电容器充满电介质的基本计算，要求熟练运用 $\vect{D}$、$\vecE$、$\vect{P}$ 之间的关系。

[解答] (a) 取药盒形高斯面，底面在导体极板内（$E=0$），顶面在介质中：
\[
D \cdot A = \sigma A \quad\Rightarrow\quad D = \sigma
\]

(b) $E = D/(\eps_0\eps_r) = \sigma/(\eps_0\eps_r)$。

(c) $P = \eps_0\chi_e E = \eps_0(\eps_r - 1)\dfrac{\sigma}{\eps_0\eps_r} = \dfrac{(\eps_r - 1)\sigma}{\eps_r}$。

(d) 束缚电荷面密度：
\begin{itemize}
\item 靠近正极板表面：$\sigma_b = \vect{P}\cdot\unitvec{n} = -P = -\dfrac{(\eps_r - 1)\sigma}{\eps_r}$（$\unitvec{n}$ 指向正极板，$\vect{P}$ 指向负极板）
\item 靠近负极板表面：$\sigma_b = +P = +\dfrac{(\eps_r - 1)\sigma}{\eps_r}$
\end{itemize}

[量纲校验] $[D] = [C/m^2] = [\sigma]$，$[E] = [C/m^2]/[F/m] = [MLT^{-3}I^{-1}]$，$[P] = [C/m^2]$，一致。\qedsymbol

\textbf{题2解析：}

[分析] 本题要求从极化强度的定义出发，推导均匀极化介质球表面的束缚面电荷分布。

[解答] 由边界条件，面束缚电荷密度 $\sigma_b = \vect{P}\cdot\unitvec{n}$。对于均匀极化球，取极化方向为 $z$ 轴，则 $\vect{P} = P\unitvec{z}$。

球面上外法向 $\unitvec{n} = \sin\theta\cos\phi\,\unitvec{x} + \sin\theta\sin\phi\,\unitvec{y} + \cos\theta\,\unitvec{z}$。

$\sigma_b = \vect{P}\cdot\unitvec{n} = P\cos\theta$。

在北极（$\theta = 0$）：$\sigma_b = +P$（最大正束缚电荷）；
在赤道（$\theta = \pi/2$）：$\sigma_b = 0$；
在南极（$\theta = \pi$）：$\sigma_b = -P$（最大负束缚电荷）。\qedsymbol

[量纲校验] $[\sigma_b] = [C/m^2] = [P]$，一致。\qedsymbol

\textbf{题3解析：}

[分析] 本题是导体--电介质分层球体的典型问题，需要分区域应用高斯定理。

[解答] 设 $r < a$ 为导体内部，$a < r < b$ 为电介质，$r > b$ 为真空。

\textbf{区域一：$r < a$}（导体内部）
$\vecE = \vect{0}$，导体是等势体。

\textbf{区域二：$a < r < b$}（电介质中）
高斯定理：$D \cdot 4\pi r^2 = Q$，$D = Q/(4\pi r^2)$。
$E = D/\eps = Q/(4\pi\eps r^2)$。

\textbf{区域三：$r > b$}（真空中）
$E = Q/(4\pi\eps_0 r^2)$。

导体球电势：$\phi_a = \int_a^b E\,\diff r + \int_b^\infty E\,\diff r = \dfrac{Q}{4\pi\eps}\left(\dfrac{1}{a} - \dfrac{1}{b}\right) + \dfrac{Q}{4\pi\eps_0 b}$。

电介质内束缚电荷：体束缚电荷 $\rho_b = -\divg\vect{P} = 0$（均匀介质），面束缚电荷在 $r=a$ 和 $r=b$ 处。

[量纲校验] $[\phi] = [ML^2T^{-3}I^{-1}]$，$[Q/(\eps r)] = [C]/([F/m][L]) = [ML^2T^{-3}I^{-1}]$，一致。\qedsymbol

\textbf{题4解析：}

[分析] 本题考查各向异性介质中 $\vect{D}$ 与 $\vecE$ 的关系，需要将各向同性结论推广到介电张量。

[解答] 在各向异性介质中，$D_i = \eps_{ij}E_j$（$\eps_{ij}$ 为对称正定介电张量）。

在边界处，$D_{1n} = D_{2n}$（法向 $D$ 连续，无自由面电荷）和 $E_{1t} = E_{2t}$（切向 $E$ 连续）。

设界面法向为 $z$ 方向，则 $D_{1z} = D_{2z}$，$E_{1x} = E_{2x}$，$E_{1y} = E_{2y}$。

在各向异性介质中，$\vect{D}$ 与 $\vecE$ 方向不同。折射定律推广：能量传播方向（Poynting矢量）与波矢方向不同，需要用到介电张量的本征值分析。\qedsymbol

[量纲校验] $[D_i] = [\eps_{ij}][E_j]$，$[C/m^2] = [F/m][V/m]$，一致。\qedsymbol

\textbf{题5解析：}

[分析] 本题利用静电学边界条件证明导体表面的电场性质，是导体静电平衡的基本结论。

[解答] 导体内部 $E = 0$。由 $E$ 切向分量的连续性：$E_{t,\text{外}} = E_{t,\text{内}} = 0$，故 $\vecE$ 只能有垂直于表面的法向分量。

取药盒形高斯面，一底面在导体内（$E=0$），一底面在外（$E_n$），$D$ 法向跃变：$D_{\text{外},n} - D_{\text{内},n} = \sigma_f$。

因导体内 $E=0$，$D_{\text{内}} = 0$（真空或介质中），故 $D_{\text{外},n} = \sigma_f$。
$\eps_0 E = \sigma_f$（真空外），$E = \sigma_f/\eps_0$。\qedsymbol

[量纲校验] $[\sigma/\eps_0] = [C/m^2]/[F/m] = [V/m] = [E]$，一致。\qedsymbol



\begin{warning}{第3章常见误区}{ch3-warnings}
\begin{enumerate}
\item $\vect{D}$ 不是物理力场，$\vecE$ 才是。带电粒子受力是 $q\vecE$，不是 $q\vect{D}$。
\item 只在\textbf{线性各向同性}均匀介质中才有 $\vect{D} = \eps\vecE$。在各向异性介质（如晶体）中，$\vect{D}$ 和 $\vecE$ 方向不同，需要用介电张量 $\varepsilon_{ij}$。
\item 边界条件中的 $\sigma_f$ 仅指自由电荷。束缚电荷已"隐藏"在 $\vect{P}$ 中。
\item 导体是等势体，但不同导体可以有不同电势。
\item 电介质内部可以有电场（$\vecE \neq 0$），只是被削弱了；导体内部 $E=0$。
\end{enumerate}
\end{warning}

% ============================================================

\begin{center}\small \hyperlink{ch2-goal}{上一章：第2章} \quad|\quad \hyperlink{ch4-goal}{下一章：第4章}\end{center}

% 第4章：静电学边值问题（优化版）
% ============================================================
\chapter{静电学边值问题}
\epigraph{简单是终极的复杂。}{--- 列奥纳多·达·芬奇}

\hypertarget{ch4-goal}{}\begin{insight}{本章核心目标}{ch4-goal-v2}
掌握求解泊松方程/拉普拉斯方程的系统方法：镜像法、分离变量法和多极展开。
\tcblower
\textbf{本章学完你应该能够：}
\begin{itemize}[label=$\square$]
\item 用分离变量法求解球/柱坐标下的拉普拉斯方程
\item 解释为什么勒让德多项式出现在球坐标解中（角向方程的本征函数）
\item 用多极展开近似计算远场电势，并判断哪一阶矩主导
\item 用格林函数法理解点源产生的场
\end{itemize}

\textbf{前置知识}：第1章\S1.3（曲线坐标系）、\S1.5（$\delta$函数）、第2章（唯一性定理）\\
\textbf{难点预告}：勒让德多项式和贝塞尔函数的性质第一次用到时会觉得陌生——\textbf{可以先接受它们是"柱/球坐标下的正弦/余弦"，具体性质用到时再查}。
\end{insight}

边值问题是电动力学的核心计算技能。给定电荷分布和边界条件，求空间中的电势和电场——这就是静电学的基本问题。本章介绍三种经典解法。

\section{镜像法}

\subsection{基本原理}

\begin{theorem}{镜像法的原理}{image-method-v2}
镜像法的核心思想是：用假想的"镜像电荷"代替导体表面或介质界面的影响，使得在新的电荷配置下，边界条件得到满足。由唯一性定理，这样得到的解就是唯一正确的解。
\end{theorem}

\begin{insight}{为什么镜像法有效}{image-why-v2}
镜像法的有效性根植于唯一性定理。导体表面是等势面——如果我们能在导体"背后"放置适当的镜像电荷，使得导体表面恰好成为等势面，那么由唯一性定理，这个解就是正确的。镜像电荷不是"真实的"，它们只是数学技巧——真正物理的电荷仍然只分布在导体表面（作为感应电荷）。

镜像法的本质是：将边界条件对场的影响，等效为虚构的电荷分布。这种方法只适用于具有高度对称性的边界（平面、球面、圆柱面等），但其物理图像清晰，计算简洁。
\end{insight}

\subsection{点电荷与无限大导体平面}

\begin{example}{点电荷与接地导体平面}{image-plane-ex-v2}
\difficulty{基础}
在 $z > 0$ 的半空间，距接地导体平面（$z=0$）$d$ 处有一点电荷 $q$（位于 $(0,0,d)$）。求 $z>0$ 区域的电势。
\tcblower
\textbf{解题思路}：无限大接地平面的对称性极好——在平面另一侧对称位置放一个镜像电荷，平面就成了一个天然的"对称面"。要猜的镜像电荷大小和位置只需满足一个条件：\textbf{它在平面上产生的电势与原电荷等值反号}，使平面电势为零（接地条件）。坐标系选直角坐标（平面对称），$z$ 轴垂直平面。\textbf{已知}：$q$ 在 $(0,0,d)$，$z=0$ 接地；\textbf{求}：$z>0$ 区域 $\phi(\vect{r})$。

\textbf{解答：}

镜像电荷 $q' = -q$ 置于 $(0,0,-d)$。

验证：在 $z=0$ 平面上，任意点 $(x,y,0)$ 到 $q$ 和 $q'$ 的距离相等：
\[
r_+ = \sqrt{x^2+y^2+d^2} = r_-
\]
电势：
\[
\phi(x,y,0) = \frac{1}{4\pi\eps_0}\left(\frac{q}{r_+} + \frac{-q}{r_-}\right) = 0
\]
满足接地条件！

因此，$z>0$ 区域的电势为：
\[
\phi(\vect{r}) = \frac{q}{4\pi\eps_0}\left(\frac{1}{|\vect{r} - d\unitvec{z}|} - \frac{1}{|\vect{r} + d\unitvec{z}|}\right)
\]

导体表面的感应电荷面密度：
\[
\sigma = -\eps_0 \frac{\partial\phi}{\partial z}\bigg|_{z=0}
\]

计算（留作练习）得：
\[
\sigma(x,y) = -\frac{qd}{2\pi(x^2+y^2+d^2)^{3/2}}
\]

总感应电荷：$\int\sigma\,\diff S = -q$（用极坐标积分验证）——感应电荷总量等于镜像电荷，符合电荷守恒。

\textbf{结果解读}：(1) 电势只对 $z>0$ 有效——镜像电荷在"镜子"后面，物理空间里不存在；(2) 极限校验 $d\to\infty$ 时 $\sigma\to 0$（电荷离平面越远感应越弱），$d\to 0$ 时 $\sigma \to -\infty\cdot$点状分布（点电荷贴近导体）；(3) 受力验证：$q$ 受到指向平面的吸引力 $F = -q^2/(16\pi\eps_0 d^2)$——与相距 $2d$ 的异号电荷间作用力相同，这正是"镜像"名字的由来。这道题考察\textbf{唯一性定理}的应用：猜对边界条件就赢了。
\end{example}

% ---------- 镜像电荷示意图 ----------
\begin{center}
\begin{tikzpicture}[scale=1.0]
  % conductor plane (z=0 horizontal)
  \fill[gray!30] (-4,0) rectangle (4,0.12);
  \draw[thick] (-4,0.12) -- (4,0.12);
  \node[font=\small] at (3.3,0.35) {接地导体平面};
  \node[font=\scriptsize] at (-3.3,0.35) {$\phi = 0$};

  % real charge q at (0,0,d) -> draw at (0,1.8)
  \fill[red] (0,1.8) circle (0.12);
  \node[right, font=\small] at (0.2,1.8) {$q$（真实电荷）};
  \draw[->, thick] (0,1.8) -- (1.4,2.6) node[right, font=\scriptsize] {真实空间 $z>0$};

  % image charge at (0,-1.8) (virtual, behind plane)
  \fill[blue] (0,-1.8) circle (0.12);
  \node[right, font=\small] at (0.2,-1.8) {$q' = -q$（镜像）};
  \node[font=\scriptsize, blue!60!black] at (0,-2.3) {虚像区 $z<0$（不存在）};

  % dashed connection
  \draw[dashed, gray] (0,1.8) -- (0,-1.8);
  \node[font=\scriptsize] at (0.55,0.6) {$d$};
  \node[font=\scriptsize] at (0.55,-0.6) {$d$};

  % field lines from q curving down to plane
  \foreach \x in {-2.4,-1.5,0,1.5,2.4} {
    \draw[->, red!60, thin] (\x*0.45,1.65) .. controls (\x*0.75,0.9) .. (\x*0.62,0.14);
  }

  % induced sigma annotation
  \node[font=\scriptsize] at (-2.2,-0.45) {感应电荷 $\sigma(x,y) = -\dfrac{qd}{2\pi(x^2+y^2+d^2)^{3/2}}$};
\end{tikzpicture}
\end{center}

\textbf{图解}：镜像电荷 $q'=-q$ 放在平面\textbf{另一侧}对称位置（虚像区），不是真实电荷——它只是让平面电势恰好为零的数学技巧。真实空间（$z>0$）的电场线从 $q$ 出发、垂直终止于导体表面（表面感应出负电荷 $\sigma$）。由唯一性定理，这个"虚构电荷 + 真实电荷"算出的场就是真实场的唯一解。

\begin{example}{点电荷与导体球}{image-sphere-ex-v2}
\difficulty{提高}
半径为 $R$ 的接地导体球，球外距球心 $a$ 处有一点电荷 $q$（$a > R$）。求球外电势。
\tcblower
\textbf{解答：}

镜像电荷 $q'$ 置于球内距球心 $b$ 处。由球对称性，$q'$ 应在 $q$ 与球心的连线上。

由球面 $\phi = 0$ 的条件，对球面上任意点：
\[
\frac{q}{r_1} + \frac{q'}{r_2} = 0
\]
其中 $r_1, r_2$ 是到 $q$ 和 $q'$ 的距离。

取球面上两点：最近点（距 $q$ 为 $a-R$）和最远点（距 $q$ 为 $a+R$），可得：
\[
q' = -\frac{R}{a}q, \quad b = \frac{R^2}{a}
\]

验证：对球面上任意点，由相似三角形可证 $r_2/r_1 = R/a = |q'/q|$，故 $\phi = 0$。

球外电势：
\[
\phi = \frac{1}{4\pi\eps_0}\left(\frac{q}{|\vect{r} - a\unitvec{z}|} - \frac{Rq/a}{|\vect{r} - (R^2/a)\unitvec{z}|}\right)
\]

[量纲校验] $[q/|\vect{r}-\vect{r}'|] = \mathrm{C}/\mathrm{m}$，乘 $[1/(4\pi\eps_0)] = \mathrm{N\cdot m^2/C^2}$ 得 $\mathrm{N\cdot m/C} = \mathrm{V}$，电势量纲一致。\checkmark（极限校验：$a\to\infty$ 时 $q' \to 0$，退化为点电荷电势）
\end{example}

\toggleButton{挑战题}{ocg挑战2}
\begin{ocg}{挑战}{ocg挑战2}{off}
\begin{example}{点电荷与电介质球——边值问题视角}{dielectric-sphere-ex-v2}
\difficulty{挑战}
半径为 $R$、介电常数为 $\eps$ 的均匀电介质球，置于均匀外电场 $\vecE_0 = E_0\unitvec{z}$ 中。用分离变量法求球内外的电势分布。

\tcblower
\textbf{解答：}

这是从边值问题角度处理导体-介质系统的经典例题。球外为真空（$\eps_0$），球内为电介质（$\eps$）。无自由电荷，故处处满足拉普拉斯方程 $\lap\phi = 0$。

由轴对称性，选取球坐标 $(r,\theta)$，$z$ 轴沿 $\vecE_0$ 方向。

\textbf{边界条件}：
\begin{enumerate}
\item $r \to \infty$：$\phi_{\text{外}} \to -E_0 r\cos\theta$（恢复均匀外场）
\item $r = 0$：$\phi_{\text{内}}$ 有限
\item $r = R$：$\phi_{\text{外}} = \phi_{\text{内}}$（电势连续）
\item $r = R$：$\eps_0 E_{\text{外},n} = \eps E_{\text{内},n}$，即 $\eps_0\dfrac{\partial\phi_{\text{外}}}{\partial r} = \eps\dfrac{\partial\phi_{\text{内}}}{\partial r}$（法向 $\vect{D}$ 连续，无自由面电荷）
\end{enumerate}

由球坐标分离变量（详见第4.2节），轴对称拉普拉斯方程的通解为：
\[
\phi(r,\theta) = \sum_{l=0}^{\infty}(A_l r^l + B_l r^{-(l+1)})P_l(\cos\theta)
\]

\textbf{球内}（$r < R$）：要求 $r=0$ 处有限，故所有 $B_l = 0$：
\[
\phi_{\text{内}} = \sum_{l=0}^{\infty} A_l r^l P_l(\cos\theta)
\]

\textbf{球外}（$r > R$）：将外场写成 $-E_0 r P_1(\cos\theta)$，再加上感应项：
\[
\phi_{\text{外}} = -E_0 r P_1(\cos\theta) + \sum_{l=0}^{\infty} C_l r^{-(l+1)} P_l(\cos\theta)
\]

由边界条件分析，只有 $l=1$ 项非零（因为外场只有 $P_1$ 分量）。设：
\begin{align*}
\phi_{\text{内}} &= -A r\cos\theta \\
\phi_{\text{外}} &= -E_0 r\cos\theta + \frac{B}{r^2}\cos\theta
\end{align*}

在 $r = R$ 处：
\begin{align*}
\text{电势连续：} &\quad -AR = -E_0 R + \frac{B}{R^2} \\
\text{D法向连续：} &\quad \eps_0\left(E_0 + \frac{2B}{R^3}\right) = \eps A
\end{align*}

解得：
\[
A = \frac{3\eps_0}{\eps + 2\eps_0}E_0, \quad B = \frac{\eps - \eps_0}{\eps + 2\eps_0}R^3 E_0
\]

因此：
\[
\phi_{\text{内}} = -\frac{3\eps_0}{\eps + 2\eps_0}E_0 r\cos\theta
\]
\[
\phi_{\text{外}} = -E_0 r\cos\theta + \frac{\eps - \eps_0}{\eps + 2\eps_0}\frac{R^3 E_0}{r^2}\cos\theta
\]

\textbf{物理图像}：
\begin{itemize}
\item 球内为\textbf{均匀电场}：$\vecE_{\text{内}} = \dfrac{3\eps_0}{\eps + 2\eps_0}\vecE_0$，方向与外场相同，但被削弱了（$\eps > \eps_0$ 时 $E_{\text{内}} < E_0$）
\item 球外相当于外场加上一个位于原点的\textbf{感应电偶极子}，偶极矩 $p = 4\pi\eps_0\dfrac{\eps-\eps_0}{\eps+2\eps_0}R^3 E_0$
\item 当 $\eps \gg \eps_0$（强极化介质），$E_{\text{内}} \approx 0$，球内近似等势——趋近于导体球的行为
\end{itemize}
\end{example}
\end{ocg}

\section{分离变量法}

\subsection{直角坐标系中的分离变量}

对于拉普拉斯方程 $\lap\phi = 0$ 在直角坐标系中：
\[
\frac{\partial^2\phi}{\partial x^2} + \frac{\partial^2\phi}{\partial y^2} + \frac{\partial^2\phi}{\partial z^2} = 0
\]

设 $\phi(x,y,z) = X(x)Y(y)Z(z)$，代入后除以 $\phi$：
\[
\frac{X''}{X} + \frac{Y''}{Y} + \frac{Z''}{Z} = 0
\]

每项只依赖于一个变量，故必须各自为常数：
\[
\frac{X''}{X} = -k_x^2, \quad \frac{Y''}{Y} = -k_y^2, \quad \frac{Z''}{Z} = k_x^2 + k_y^2 = k_z^2
\]

\begin{example}{矩形区域内的电势}{separation-rect-ex-v2}
\difficulty{提高}
一个边长为 $a \times b$ 的矩形区域，边界条件为：$\phi(0,y) = \phi(a,y) = 0$，$\phi(x,0) = 0$，$\phi(x,b) = V_0\sin(\pi x/a)$。求区域内的电势。
\tcblower
\textbf{解答：}

设 $\phi(x,y) = X(x)Y(y)$。$X''/X = -k^2$，$Y''/Y = k^2$。

由 $X(0) = X(a) = 0$：$X_n(x) = \sin(n\pi x/a)$，$k_n = n\pi/a$。

$Y(0) = 0$：$Y_n(y) = \sinh(n\pi y/a)$。

通解：$\phi = \sum_n A_n \sin(n\pi x/a)\sinh(n\pi y/a)$

由 $y=b$ 处的边界条件：$\sum_n A_n\sin(n\pi x/a)\sinh(n\pi b/a) = V_0\sin(\pi x/a)$

比较系数，只有 $n=1$ 项非零：
\[
A_1\sinh(\pi b/a) = V_0 \quad\Rightarrow\quad A_1 = \frac{V_0}{\sinh(\pi b/a)}
\]

因此：
\[
\phi(x,y) = V_0 \frac{\sin(\pi x/a)\sinh(\pi y/a)}{\sinh(\pi b/a)}
\]
\end{example}

\subsection{球坐标系中的分离变量}

球坐标下拉普拉斯方程已在第1章给出：
\[
\lap\phi = \frac{1}{r^2}\frac{\partial}{\partial r}\left(r^2\frac{\partial\phi}{\partial r}\right)
+ \frac{1}{r^2\sin\theta}\frac{\partial}{\partial\theta}\left(\sin\theta\frac{\partial\phi}{\partial\theta}\right)
+ \frac{1}{r^2\sin^2\theta}\frac{\partial^2\phi}{\partial\varphi^2} = 0
\]

对于\textbf{轴对称问题}（与 $\varphi$ 无关，即绕 $z$ 轴旋转对称），方程简化为：
\[
\frac{1}{r^2}\frac{\partial}{\partial r}\left(r^2\frac{\partial\phi}{\partial r}\right)
+ \frac{1}{r^2\sin\theta}\frac{\partial}{\partial\theta}\left(\sin\theta\frac{\partial\phi}{\partial\theta}\right) = 0
\]

\begin{theorem}{球坐标轴对称拉普拉斯方程的通解}{spherical-separation-v2}
设 $\phi(r,\theta)$ 与方位角 $\varphi$ 无关。通过分离变量 $\phi(r,\theta) = R(r)\Theta(\theta)$，可得轴对称拉普拉斯方程的通解：
\[
\phi(r,\theta) = \sum_{l=0}^{\infty}\left(A_l r^l + B_l r^{-(l+1)}\right)P_l(\cos\theta)
\]
其中 $P_l(\cos\theta)$ 是\textbf{勒让德多项式}，$A_l$ 和 $B_l$ 由边界条件确定。
\end{theorem}

\begin{insight}{勒让德多项式的物理意义}{legendre-meaning-v2}
勒让德多项式 $P_l(\cos\theta)$ 描述了角向分布的"标准模式"，类似于振动弦上的驻波模式。最低两阶的显式形式为：
\begin{align*}
P_0(\cos\theta) &= 1 \quad \text{（各向同性，球对称）} \\
P_1(\cos\theta) &= \cos\theta \quad \text{（偶极型，关于赤道面反对称）}
\end{align*}

$l=0$ 项对应\textbf{点电荷}或\textbf{单极子}的角向分布——完全球对称。
$l=1$ 项对应\textbf{偶极子}的角向分布——正北极、负南极（或反之）。
$l=2$ 项对应\textbf{四极子}，描述电荷分布偏离球对称的更精细结构。

在量子力学中，$l$ 对应角动量量子数；在经典电动力学中，$l$ 表征多极展开的阶数。不必深究量子力学的诠释，只需记住：每一项 $l$ 对应一种特定的角向"图案"。

\textbf{正交归一}：$P_l(\cos\theta)$ 在 $[-1,1]$ 上满足正交归一条件
\[
\int_{-1}^{1} P_l(\mu)\,P_{l'}(\mu)\,\diff\mu = \frac{2}{2l+1}\,\delta_{ll'}
\]
这就是"级数系数怎么算"的来源：把边界条件两边乘以 $P_{l'}(\cos\theta)\sin\theta\,\diff\theta$（即 $\diff\mu$）并对 $\mu\in[-1,1]$ 积分，正交性使求和只剩 $l = l'$ 一项，系数即可解出。第 4 章均匀外场导体球、电介质球两例正是这样定出 $B_1 = E_0R^3$、球内 $A_1$ 系数的——做题时若边界条件含 $\cos\theta$，只有 $l=1$ 项非零，可直接比较系数（等价于利用正交性）。
\end{insight}

\begin{example}{均匀外场中的导体球}{conductor-sphere-uniform-field-v2}
\difficulty{提高}
半径为 $R$ 的导体球置于均匀外电场 $\vecE_0 = E_0\unitvec{z}$ 中。用球坐标分离变量法求球外电势分布。

\tcblower
\textbf{解答：}

导体球静电平衡时为等势体，设电势为 $\phi_0$（常数，取 $\phi_0 = 0$ 即接地）。

\textbf{边界条件}：
\begin{enumerate}
\item $r \to \infty$：$\phi \to -E_0 r\cos\theta = -E_0 r P_1(\cos\theta)$
\item $r = R$：$\phi = 0$
\end{enumerate}

球外（$r > R$）满足拉普拉斯方程，通解为：
\[
\phi(r,\theta) = \sum_{l=0}^{\infty}\left(A_l r^l + B_l r^{-(l+1)}\right)P_l(\cos\theta)
\]

由边界条件1，$r \to \infty$ 时 $\phi \sim -E_0 r P_1(\cos\theta)$，故只有 $l=1$ 项有非零的 $A_1 = -E_0$，其他 $A_l = 0$（$l \neq 1$）。

因此：
\[
\phi(r,\theta) = -E_0 r P_1(\cos\theta) + \sum_{l=0}^{\infty} B_l r^{-(l+1)}P_l(\cos\theta)
\]

由边界条件2，$r = R$ 处 $\phi = 0$：
\[
-E_0 R P_1(\cos\theta) + \sum_{l=0}^{\infty} B_l R^{-(l+1)}P_l(\cos\theta) = 0
\]

由于不同 $l$ 的 $P_l(\cos\theta)$ 线性无关，比较系数：
\begin{align*}
l=1: &\quad -E_0 R + B_1 R^{-2} = 0 \quad\Rightarrow\quad B_1 = E_0 R^3 \\
l \neq 1: &\quad B_l = 0
\end{align*}

因此：
\[
\phi(r,\theta) = -E_0 r\cos\theta + \frac{E_0 R^3}{r^2}\cos\theta
\]

球外电场：
\[
\vecE = -\grad\phi = E_0\cos\theta\left(1+\frac{2R^3}{r^3}\right)\unitvec{r}
- E_0\sin\theta\left(1-\frac{R^3}{r^3}\right)\unitvec{\theta}
\]

\textbf{物理图像}：
\begin{itemize}
\item 第一项 $-E_0 r\cos\theta$ 是外场贡献
\item 第二项 $\dfrac{E_0 R^3}{r^2}\cos\theta$ 是导体球表面感应电荷产生的偶极子势，等价于一个位于球心的电偶极子 $p = 4\pi\eps_0 R^3 E_0$
\item 在导体表面（$r = R$），$E_\theta = 0$（切向电场为零），$E_r = 3E_0\cos\theta$（法向电场，正比于面电荷密度 $\sigma = 3\eps_0 E_0\cos\theta$）
\end{itemize}
\end{example}

\subsection{柱坐标系中的轴对称分离变量}

柱坐标 $(\rho,\varphi,z)$ 下拉普拉斯方程为：
\[
\lap\phi = \frac{1}{\rho}\frac{\partial}{\partial\rho}\left(\rho\frac{\partial\phi}{\partial\rho}\right)
+ \frac{1}{\rho^2}\frac{\partial^2\phi}{\partial\varphi^2}
+ \frac{\partial^2\phi}{\partial z^2} = 0
\]

对于\textbf{轴对称问题}（与 $\varphi$ 无关），方程简化为：
\[
\frac{1}{\rho}\frac{\partial}{\partial\rho}\left(\rho\frac{\partial\phi}{\partial\rho}\right)
+ \frac{\partial^2\phi}{\partial z^2} = 0
\]

\begin{theorem}{柱坐标轴对称拉普拉斯方程的通解}{cylindrical-separation-v2}
设 $\phi(\rho,z)$ 与方位角 $\varphi$ 无关。通过分离变量，径向部分满足\textbf{贝塞尔方程}，其解涉及贝塞尔函数 $J_n(k\rho)$ 和 $Y_n(k\rho)$。对于要求 $\rho=0$ 处有限的解，零阶贝塞尔函数 $J_0(k\rho)$ 是关键函数。

零阶贝塞尔函数 $J_0(x)$ 是下列方程的解：
\[
x^2\frac{\diff^2 J_0}{\diff x^2} + x\frac{\diff J_0}{\diff x} + x^2 J_0 = 0
\]

$J_0(x)$ 具有振荡衰减特性（类似衰减的余弦函数），在 $x=0$ 处取值为1，并随着 $x$ 增大而振荡衰减。它是处理圆柱形边界问题的自然函数。
\end{theorem}

\toggleButton{挑战题}{ocg挑战3}
\begin{ocg}{挑战}{ocg挑战3}{off}
\begin{example}{无限长圆柱面的电势问题}{cylinder-potential-ex-v2}
\difficulty{提高}
半径为 $R$ 的无限长圆柱面，其表面电势分布为 $\phi(R,\varphi) = V_0\cos\varphi$（与 $z$ 无关）。求柱面外（$\rho > R$）的电势分布。

\tcblower
\textbf{解答：}

这是一个二维问题（与 $z$ 无关），在柱坐标下拉普拉斯方程简化为：
\[
\frac{1}{\rho}\frac{\partial}{\partial\rho}\left(\rho\frac{\partial\phi}{\partial\rho}\right)
+ \frac{1}{\rho^2}\frac{\partial^2\phi}{\partial\varphi^2} = 0
\]

设 $\phi(\rho,\varphi) = R(\rho)\Phi(\varphi)$，分离变量后：
\[
\rho^2\frac{R''}{R} + \rho\frac{R'}{R} = -\frac{\Phi''}{\Phi} = m^2
\]

角向方程 $\Phi'' + m^2\Phi = 0$ 的解为 $\Phi_m(\varphi) = \cos(m\varphi)$ 或 $\sin(m\varphi)$（$m$ 为整数，保证周期性）。

径向方程：
\[
\rho^2 R'' + \rho R' - m^2 R = 0
\]

这是欧拉方程，解为 $R_m(\rho) = A_m\rho^m + B_m\rho^{-m}$（$m \neq 0$），或 $R_0 = A_0 + B_0\ln\rho$（$m = 0$）。

由边界条件 $\phi \to 0$（或有限）当 $\rho \to \infty$，需要 $A_m = 0$（对所有 $m \geq 0$）。因此：
\[
\phi(\rho,\varphi) = \sum_{m=0}^{\infty} \rho^{-m}\left(C_m\cos m\varphi + D_m\sin m\varphi\right)
\]

由柱面边界条件 $\phi(R,\varphi) = V_0\cos\varphi$：
\[
\sum_{m=0}^{\infty} R^{-m}\left(C_m\cos m\varphi + D_m\sin m\varphi\right) = V_0\cos\varphi
\]

比较系数：$C_1 R^{-1} = V_0$，其他系数均为零。故 $C_1 = V_0 R$，$D_m = 0$。

因此：
\[
\phi(\rho,\varphi) = \frac{V_0 R}{\rho}\cos\varphi = \frac{V_0 R}{\rho^2}\,\rho\cos\varphi
= \frac{V_0 R}{\rho^2}\,x
\]

这等价于一个位于原点的二维电偶极子势！事实上，$\phi = V_0 R\cos\varphi/\rho$ 可改写为 $\phi = \vect{p}\cdot\vect{\rho}/\rho^2$ 的形式，其中二维偶极矩 $p = V_0 R$。

柱面外电场：
\[
E_\rho = -\frac{\partial\phi}{\partial\rho} = \frac{V_0 R}{\rho^2}\cos\varphi, \quad
E_\varphi = -\frac{1}{\rho}\frac{\partial\phi}{\partial\varphi} = \frac{V_0 R}{\rho^2}\sin\varphi
\]
\end{example}
\end{ocg}

\section{多极展开}

对于远处的电势，可将电荷分布展开为多极矩的贡献：

\begin{theorem}{多极展开}{multipole-expansion-v2}
\[
\phi(\vect{r}) = \frac{1}{4\pi\eps_0}\left(\frac{Q}{r} + \frac{\vect{p}\cdot\unitvec{r}}{r^2} + \frac{1}{2r^3}\sum_{i,j} Q_{ij}\hat{r}_i\hat{r}_j + \cdots\right)
\]
其中：
\begin{itemize}
\item $Q = \int\rho\,\diff^3r'$ —— 总电荷（单极矩）
\item $\vect{p} = \int\vect{r}'\rho(\vect{r}')\,\diff^3r'$ —— 电偶极矩
\item $Q_{ij} = \int(3x'_i x'_j - r'^2\delta_{ij})\rho\,\diff^3r'$ —— 电四极矩张量
\end{itemize}
\end{theorem}

\begin{insight}{多极展开的物理意义}{multipole-phy-v2}
多极展开将复杂的电荷分布分解为一系列"标准模式"：
\begin{itemize}
\item \textbf{单极项}（$1/r$）：总电荷的贡献，球对称，所有电荷分布都有
\item \textbf{偶极项}（$1/r^2$）：反映电荷分布的非对称性，正负电荷中心不重合时出现
\item \textbf{四极项}（$1/r^3$）：更精细的结构信息
\end{itemize}
对于中性物体（$Q=0$），偶极矩主导；对于具有中心对称的物体，偶极矩为零，四极矩主导。在原子核物理中，四极矩是研究核形状的重要工具。
\end{insight}

\begin{example}{均匀带电圆盘的多极展开}{disk-multipole-ex-v2}
\difficulty{提高}
半径为 $R$ 的圆盘均匀带电，总电荷 $Q$。求圆盘轴线上距离圆心 $z$ 处的电势，并将其展开至多极矩形式（保留到四极矩项）。

\tcblower
\textbf{解答：}

圆盘面电荷密度 $\sigma = Q/(\pi R^2)$。

\textbf{直接计算电势}：

轴线上电势（用柱坐标）：
\[
\phi(z) = \frac{1}{4\pi\eps_0}\int_0^R \frac{\sigma \cdot 2\pi\rho\,\diff\rho}{\sqrt{\rho^2+z^2}}
= \frac{\sigma}{2\eps_0}\left[\sqrt{R^2+z^2} - |z|\right]
\]

对 $z > 0$：
\[
\phi(z) = \frac{Q}{2\pi\eps_0 R^2}\left(\sqrt{R^2+z^2} - z\right)
\]

\textbf{多极展开}（$z \gg R$ 的远场）：

将 $\sqrt{R^2+z^2} = z\sqrt{1+R^2/z^2}$ 对 $R/z$ 展开：
\[
\sqrt{R^2+z^2} = z\left(1 + \frac{R^2}{2z^2} - \frac{R^4}{8z^4} + \cdots\right)
\]

因此：
\[
\phi(z) = \frac{Q}{2\pi\eps_0 R^2}\left(z + \frac{R^2}{2z} - \frac{R^4}{8z^3} + \cdots - z\right)
= \frac{Q}{4\pi\eps_0 z} - \frac{Q R^2}{16\pi\eps_0 z^3} + O(z^{-5})
\]

\textbf{与标准多极展开对照}：
\begin{align*}
\text{单极项：} &\quad \frac{Q}{4\pi\eps_0 z} \quad\checkmark \\
\text{偶极项：} &\quad 0 \quad\text{（圆盘关于中心对称，偶极矩为零）} \\
\text{四极项：} &\quad -\frac{QR^2}{16\pi\eps_0 z^3} = \frac{1}{4\pi\eps_0}\frac{1}{2z^3}Q_{zz}
\end{align*}

其中圆盘的电四极矩张量（取 $z$ 轴为对称轴）：
\[
Q_{zz} = \int(3z'^2 - r'^2)\sigma\,\diff S = -\int\rho'^2\sigma\,\diff S
= -\sigma \int_0^{2\pi}\diff\varphi\int_0^R \rho'^3\,\diff\rho'
= -\frac{Q}{\pi R^2}\cdot 2\pi\cdot\frac{R^4}{4} = -\frac{QR^2}{2}
\]

两种方法完全自洽：直接展开给出 $-QR^2/(16\pi\eps_0 z^3)$；标准四极势公式在轴线上（$\hat{r}_z = 1$，$Q_{xx} = Q_{yy} = -Q_{zz}/2$）为 $\dfrac{1}{4\pi\eps_0}\dfrac{Q_{zz}}{2z^3} = -\dfrac{QR^2}{16\pi\eps_0 z^3}$。负号来自 $Q_{zz} < 0$——圆盘电荷分布在垂直于 $z$ 的平面内"压扁"，沿轴方向的四极势因此为负。
\end{example}

\toggleButton{挑战题}{ocg挑战4}
\begin{ocg}{挑战}{ocg挑战4}{off}
\begin{example}{线性四极子——四极矩张量计算}{quadrupole-tensor-ex-v2}
\difficulty{挑战}
三个点电荷沿 $y$ 轴排列：$+q$ 在 $(0,d,0)$，$-2q$ 在 $(0,0,0)$，$+q$ 在 $(0,-d,0)$。计算该电荷分布的电四极矩张量。

\tcblower
\textbf{解答：}

四极矩张量定义为：
\[
Q_{ij} = \sum_k q_k(3x_{k,i}x_{k,j} - r_k^2\delta_{ij})
\]

各电荷贡献：

\textbf{电荷1}：$+q$ 在 $(0,d,0)$，$r^2 = d^2$
\begin{align*}
Q_{xx}^{(1)} &= q(0 - d^2) = -qd^2 \\
Q_{yy}^{(1)} &= q(3d^2 - d^2) = 2qd^2 \\
Q_{zz}^{(1)} &= q(0 - d^2) = -qd^2 \\
Q_{ij}^{(1)} &= 0 \quad (i \neq j)
\end{align*}

\textbf{电荷2}：$-2q$ 在 $(0,0,0)$，$r^2 = 0$
\[
Q_{ij}^{(2)} = 0 \quad \text{（所有分量）}
\]

\textbf{电荷3}：$+q$ 在 $(0,-d,0)$，$r^2 = d^2$
\begin{align*}
Q_{xx}^{(3)} &= q(0 - d^2) = -qd^2 \\
Q_{yy}^{(3)} &= q(3d^2 - d^2) = 2qd^2 \\
Q_{zz}^{(3)} &= q(0 - d^2) = -qd^2 \\
Q_{ij}^{(3)} &= 0 \quad (i \neq j)
\end{align*}

\textbf{总和}：
\begin{align*}
Q_{xx} &= -qd^2 + 0 - qd^2 = -2qd^2 \\
Q_{yy} &= 2qd^2 + 0 + 2qd^2 = 4qd^2 \\
Q_{zz} &= -qd^2 + 0 - qd^2 = -2qd^2 \\
Q_{xy} &= Q_{xz} = Q_{yz} = 0
\end{align*}

验证：$\text{Tr}(Q) = Q_{xx} + Q_{yy} + Q_{zz} = -2qd^2 + 4qd^2 - 2qd^2 = 0$。\checkmark

四极矩张量可写成：
\[
Q_{ij} = 2qd^2\begin{pmatrix} -1 & 0 & 0 \\ 0 & 2 & 0 \\ 0 & 0 & -1 \end{pmatrix}
\]

这是一个\textbf{无迹对称张量}（traceless symmetric tensor），描述电荷分布沿 $y$ 轴方向的伸长（$Q_{yy} > 0$）和在垂直方向的压缩（$Q_{xx}, Q_{zz} < 0$）。

\textbf{远场电势}（在 $y$ 轴上，$r = y$）：
\[
\phi_{quad} = \frac{1}{4\pi\eps_0}\frac{1}{2r^3}Q_{yy} = \frac{1}{4\pi\eps_0}\frac{4qd^2}{2y^3} = \frac{qd^2}{2\pi\eps_0 y^3}
\]
按 $1/y^3$ 衰减，比偶极子更快。
\end{example}
\end{ocg}

\begin{tip}{四极矩定义差异对照表}{quadrupole-def-table}
不同教材对四极矩张量的定义可能有系数差异，阅读文献时需注意（下表第三种的迹为 $\int r'^2\rho\,\diff^3r'$，一般不为零，此即未减迹之意）：
\begin{center}
\begin{tabular}{lll}
\toprule
\textbf{教材} & \textbf{定义} & \textbf{说明} \\
\midrule
本书 & $Q_{ij}=\int(3x'_ix'_j-r'^2\delta_{ij})\rho\,\diff^3r'$ & Jackson 4th ed. 标准定义 \\
Griffiths & $Q_{ij}=\int(3x'_ix'_j-r'^2\delta_{ij})\rho\,\diff^3r'$ & 与 Jackson 一致 \\
部分教材 & $\tilde{Q}_{ij}=\int x'_ix'_j\rho\,\diff^3r'$（未减迹形式） & 需额外减去迹项 \\
\bottomrule
\end{tabular}
\end{center}
\textbf{关键}：按本书定义（$Q_{ij}=\int(3x'_ix'_j-r'^2\delta_{ij})\rho\,\diff^3r'$），四极势为 $4\pi\eps_0\phi = \dfrac{1}{2r^3}Q_{ij}\hat{r}_i\hat{r}_j$，即 $4\pi\eps_0\phi\approx Q_{ij}\hat{r}_i\hat{r}_j/(2r^3)$——系数与定义方式严格配对。若采用下表第三种（未减迹）定义，展开系数需相应调整。
\end{tip}


\begin{insight}{四极矩的物理应用}{quadrupole-applications-v2}
电四极矩在现代物理中有广泛应用：

\textbf{原子核物理}：原子核的四极矩 $Q$ 是表征核偏离球对称程度的关键参数。球形核（如 $^{16}$O）$Q = 0$；长椭球形核（如 $^{238}$U）$Q > 0$；扁椭球形核 $Q < 0$。核四极矩通过超精细结构影响原子光谱，是核结构研究的重要探针。

\textbf{分子光谱}：多原子分子的电荷分布通常具有非零四极矩。例如 CO$_2$ 分子的线性对称结构使其四极矩不为零，这直接影响分子间的范德华力和红外光谱的选择定则。分子四极矩的测量为确定分子构型提供了重要信息。

\textbf{引力波天文学}：质量分布的四极矩变化是引力波辐射的主要来源（类比于电磁辐射中偶极矩变化产生电磁波）。LIGO 探测到的双黑洞并合信号，本质上是追踪系统质量四极矩随时间的演化。
\end{insight}

\begin{tip}{拓展：用 Python 数值求解泊松方程}{python-poisson-tip}
镜像法和分离变量法只适用于对称边界；对一般几何，\textbf{数值方法}是标准工具。最简单的入门方案是\textbf{有限差分 + 迭代}：把求解区域切成网格，将 $\lap\phi = 0$ 离散为"每点电势 = 四邻平均"，反复迭代直到收敛（雅可比/高斯-塞德尔迭代），边界格子固定为已知的边界条件值。十几行 Python（NumPy 数组 + 双重 for 循环或向量化切片更新）即可算出任意二维截面的电势分布，再用 \texttt{matplotlib} 画等势线直观验证。想深入可查 \texttt{FEniCS} 有限元库或 SciPy 的稀疏线性求解器。本书不展开代码细节，但理解"泊松方程 + 边界条件 = 唯一解"后，数值求解的思想是自然延伸。
\end{tip}

\begin{review}{第4章要点回顾}{ch4-review}
\begin{center}
\begin{tabular}{@{}lll@{}}
\toprule
\textbf{方法} & \textbf{适用条件} & \textbf{核心思想} \\
\midrule
镜像法 & 高度对称边界 & 虚构镜像电荷满足边界条件 \\
分离变量 & 规则边界 & 将偏微分方程化为常微分方程 \\
多极展开 & 远场近似 & 将电荷分布展开为标准矩 \\
\bottomrule
\end{tabular}
\end{center}
\end{review}

\begin{tip}{本章核心框架}{ch4-summary}
\textbf{一句话总结}：静电学边值问题可通过分离变量、镜像法、格林函数三大解析方法求解。\\
\textbf{核心工具}：勒让德多项式（球坐标）、贝塞尔函数（柱坐标）、镜像电荷（导体边界）、格林函数（$\delta$源响应）。\\
\textbf{关键定理}：唯一性定理（给定边界条件则解唯一）、多极展开（远场近似）。\\
\textbf{衔接}：第5章进入静磁学，类比地引入矢势$\vecA$和毕奥-萨伐尔定律。
\end{tip}

% ---------- 第4章随堂自测 ----------
\section*{第4章随堂自测}
\addcontentsline{toc}{section}{第4章随堂自测}

\begin{miniquiz}
\textbf{一、选择题}
\begin{enumerate}[label=\arabic*.]
\item 拉普拉斯方程 $\lap\phi = 0$ 在球坐标下的分离变量解中，角向部分由什么函数描述？
  \begin{enumerate}[label=(\Alph*),nosep]
  \item 贝塞尔函数
  \item 球谐函数
  \item 勒让德多项式
  \item 厄米多项式
  \end{enumerate}
\begin{ocg}{答案}{ocg答案4}{off}
  \textit{答案：C。勒让德多项式（及连带勒让德函数）是球坐标角向方程的解。}
\end{ocg}

\item 接地导体球的镜像法中，像电荷的位置在：
  \begin{enumerate}[label=(\Alph*),nosep]
  \item 球心
  \item 球外与源电荷对称的位置
  \item 球内某一点
  \item 无穷远处
  \end{enumerate}
\begin{ocg}{答案}{ocg答案4}{off}
  \textit{答案：C。像电荷位于球内，距球心距离为 $R^2/d$。}
\end{ocg}

\item 格林函数法求解泊松方程的核心思想是：
  \begin{enumerate}[label=(\Alph*),nosep]
  \item 直接积分
  \item 叠加原理与$\delta$函数源
  \item 变量替换
  \item 傅里叶变换
  \end{enumerate}
\begin{ocg}{答案}{ocg答案4}{off}
  \textit{答案：B。格林函数是点源产生的场，任意分布可看作点源的叠加。}
\end{ocg}
\end{enumerate}

\vspace{0.5em}
\textbf{二、填空题}
\begin{enumerate}[label=\arabic*.,resume]
\item 唯一性定理告诉我们：给定边界上的\underline{\hspace{2cm}}或\underline{\hspace{2cm}}，区域内的场唯一确定。
\begin{ocg}{答案}{ocg答案4}{off}
  \textit{答案：电势（Dirichlet）；电势的法向导数（Neumann）。}
\end{ocg}
\item 柱坐标下的分离变量中，径向方程的解是\underline{\hspace{2cm}}函数。
\begin{ocg}{答案}{ocg答案4}{off}
  \textit{答案：贝塞尔。}
\end{ocg}
\end{enumerate}
\end{miniquiz}
\toggleButton{显示答案}{ocg答案4}

\begin{warning}{第4章常见误区}{ch4-warnings}
\begin{enumerate}
\item 镜像电荷\textbf{不在求解区域内}——接地导体球问题中镜像电荷在球内，计算球外电势时才能使用；误把镜像电荷算进区域内会重复计数。
\item 镜像电荷的大小\textbf{不等于} $-q$（球面问题中 $q' = -qR/a$），只有无限大平面才是 $q' = -q$。
\item 分离变量的级数系数由\textbf{边界条件}（而非初始条件）确定——忘记归一化或漏掉正交性积分是最常见错误。
\item 多极展开的适用条件是 $r \gg r'$（远场）——近距离使用多极展开会得到错误结论。
\item 唯一性定理只保证"解唯一"，不提供解本身——它验证猜测，不产生猜测（镜像电荷仍要靠对称性去猜）。
\item 格林函数法的完整推导见第 1 章 $\delta$ 函数与格林函数部分；本章直接使用其结论。
\end{enumerate}
\end{warning}

\begin{chapterreview}
\begin{itemize}[nosep]
\item \textbf{核心方法}：分离变量法（直角/柱/球）、镜像法、格林函数法、多极展开。
\item \textbf{关键函数}：勒让德多项式、球谐函数、贝塞尔函数。
\item \textbf{核心定理}：唯一性定理。
\item \textbf{自检验}：能否推导点电荷在接地导体球外的镜像电荷大小和位置？
\item \textbf{衔接预告}：第5章进入静磁学——电流、毕奥-萨伐尔定律、安培环路定理。
\end{itemize}
\end{chapterreview}

\begin{tip}{本章自检清单}{ch4-checklist}
学完本章后对照下列条目逐条自测，全部通过即可进入下一章；有未通过项时，回到对应小节复习：
\begin{itemize}[label=$\square$, nosep]
\item 用镜像法处理接地导体球与无限大导体平板
\item 用分离变量法求解直角/球坐标拉普拉斯方程
\item 做多极展开并判断主导阶
\item 说出唯一性定理的内容及其在镜像法中的角色
\end{itemize}
\end{tip}

% ---------- 第4章综合练习题 ----------
\section*{第4章综合练习题}
\addcontentsline{toc}{section}{第4章综合练习题}
\begin{enumerate}[label=\textbf{\arabic*.}]
\item \textbf{基础} 接地导体球（半径 $R$）外距球心 $a$ 处有一点电荷 $q$（$a>R$），用镜像法求球外电势和感应电荷。
\item \textbf{基础} 直角坐标中，拉普拉斯方程 $\lap\phi=0$ 在矩形区域 $0<x<a$、$0<y<b$ 内，边界条件为四边分别给定电势。写出分离变量后的通解形式。
\item \textbf{提高} 均匀外场 $\vecE_0$ 中的接地导体球，用球坐标分离变量求感应电荷分布和总感应电荷。
\item \textbf{提高} 均匀带电圆盘（半径 $a$，面密度 $\sigma$）在轴线上的电势，展开到四极矩项，并与直接积分结果比较。
\item \textbf{挑战} 一个点电荷位于两个互相垂直的接地导体平面之间，用镜像法求所有镜像电荷的位置和大小，并讨论解的唯一性。
\end{enumerate}
\newpage
\section*{习题解答}
\addcontentsline{toc}{section}{习题解答}

\textbf{题1解析：}

[分析] 本题是导体球镜像法的经典问题，关键是利用球面电势为零的条件确定镜像电荷的大小和位置。

[解答] 设镜像电荷 $q'$ 位于球内距球心 $b$ 处（在 $q$ 与球心连线上）。球面上任一点满足 $\phi = 0$：
\[
\frac{q}{r_1} + \frac{q'}{r_2} = 0
\]
取球面上最近点和最远点列方程：
\begin{align*}
\frac{q}{a-R} + \frac{q'}{R-b} &= 0 \\
\frac{q}{a+R} + \frac{q'}{R+b} &= 0
\end{align*}
解得：$q' = -\dfrac{R}{a}q$，$b = \dfrac{R^2}{a}$。

球外电势：$\phi(\vect{r}) = \dfrac{1}{4\pi\eps_0}\left(\dfrac{q}{|\vect{r}-a\unitvec{z}|} - \dfrac{Rq/a}{|\vect{r}-(R^2/a)\unitvec{z}|}\right)$。

感应电荷：由导体表面 $\sigma = -\eps_0\partial\phi/\partial n|_{r=R}$，总感应电荷 $Q_{\text{ind}} = \int\sigma\,\diff S = -Rq/a = q'$。\qedsymbol

[量纲校验] $[q'] = [C]$，$[Rq/a] = [L][C]/[L] = [C]$，一致。\qedsymbol

\textbf{题2解析：}

[分析] 本题要求写出矩形区域内拉普拉斯方程分离变量后的通解形式，是分离变量法的基础。

[解答] 设 $\phi(x,y) = X(x)Y(y)$，代入 $\lap\phi = 0$：
\[
\frac{X''}{X} + \frac{Y''}{Y} = 0
\]
令 $\dfrac{X''}{X} = -\alpha^2$，$\dfrac{Y''}{Y} = \alpha^2$（$\alpha$ 为分离常数）。

解：$X(x) = A\sinh(\alpha x) + B\cosh(\alpha x)$，$Y(y) = C\sin(\alpha y) + D\cos(\alpha y)$。

考虑到边界条件 $y=0$ 和 $y=b$ 处给定电势，设 $Y(0) = Y(b) = 0$（齐次边界），则 $Y_n(y) = \sin(n\pi y/b)$，$\alpha_n = n\pi/b$。

通解：$\phi(x,y) = \sum_{n=1}^{\infty}\left[A_n\sinh\left(\frac{n\pi x}{b}\right) + B_n\cosh\left(\frac{n\pi x}{b}\right)\right]\sin\left(\frac{n\pi y}{b}\right)$。

若 $x=0$ 和 $x=a$ 处也有边界条件，则系数由边界条件确定。\qedsymbol

[量纲校验] $[\alpha] = [L^{-1}]$，$[\sinh(\alpha x)]$ 无量纲，$[\phi]$ 为电势量纲，级数各项量纲一致。\qedsymbol

\textbf{题3解析：}

[分析] 本题用球坐标分离变量法求解均匀外场中的导体球问题，是电动力学经典边值问题。

[解答] 球坐标下轴对称拉普拉斯方程通解：$\phi(r,\theta) = \sum_{l=0}^{\infty}(A_l r^l + B_l r^{-(l+1)})P_l(\cos\theta)$。

边界条件：
\begin{enumerate}
\item $r \to \infty$：$\phi \to -E_0 r\cos\theta$（恢复均匀外场）
\item $r = R$：$\phi = 0$（接地导体）
\end{enumerate}

由条件1，只有 $l=1$ 项：$\phi = -E_0 r\cos\theta + B_1 r^{-2}\cos\theta$。

由条件2：$-E_0 R + B_1/R^2 = 0$，$B_1 = E_0 R^3$。

故：$\phi(r,\theta) = -E_0\left(r - \frac{R^3}{r^2}\right)\cos\theta$。

感应电荷面密度：$\sigma = -\eps_0\left.\frac{\partial\phi}{\partial r}\right|_{r=R} = -\eps_0\left(-E_0 - 2E_0\right)\cos\theta = 3\eps_0 E_0\cos\theta$。

总感应电荷：$Q_{\text{ind}} = \int_0^{2\pi}\diff\phi\int_0^\pi 3\eps_0 E_0\cos\theta \cdot R^2\sin\theta\,\diff\theta = 0$（正负电荷抵消）。\qedsymbol

[量纲校验] $[\phi] = [V]$，$[E_0 r] = [V/m][L] = [V]$；$[\sigma] = [\eps_0 E_0] = [F/m][V/m] = [C/m^2]$，一致。\qedsymbol

\textbf{题4解析：}

[分析] 本题要求将带电圆盘的电势展开到多极矩，并与直接积分结果比较。

[解答] 直接积分：圆盘在轴线上 $z$ 处的电势为
\[
\phi(z) = \frac{1}{4\pi\eps_0}\int_0^a \frac{\sigma \cdot 2\pi\rho\,\diff\rho}{\sqrt{\rho^2+z^2}} = \frac{\sigma}{2\eps_0}\left(\sqrt{a^2+z^2} - |z|\right)
\]

对 $z \gg a$ 展开：$\sqrt{a^2+z^2} = |z|\sqrt{1+a^2/z^2} \approx |z|\left(1 + \frac{a^2}{2z^2} - \frac{a^4}{8z^4} + \cdots\right)$。

$\phi(z) \approx \dfrac{\sigma}{2\eps_0}\left(|z| + \dfrac{a^2}{2|z|} - \dfrac{a^4}{8|z|^3} - |z|\right) = \dfrac{\sigma a^2}{4\eps_0|z|} - \dfrac{\sigma a^4}{16\eps_0|z|^3}$。

第一项：$Q/(4\pi\eps_0|z|)$，$Q = \pi a^2\sigma$（单极项）；
第二项：四极矩贡献。\qedsymbol

[量纲校验] $[\phi] = [V]$，$[\sigma a^2/(\eps_0 z)] = [C/m^2][L^2]/([F/m][L]) = [V]$，一致。\qedsymbol

\textbf{题5解析：}

[分析] 本题是两个垂直接地导体平面间的点电荷镜像法问题，需要构造无限镜像序列。

[解答] 设两平面为 $x=0$ 和 $y=0$（第一象限），点电荷 $q$ 在 $(a,b)$。

对 $x=0$ 镜像：$q_1 = -q$ 在 $(-a,b)$；对 $y=0$ 镜像：$q_2 = -q$ 在 $(a,-b)$；
$q_1$ 对 $y=0$ 镜像：$q_3 = +q$ 在 $(-a,-b)$；$q_2$ 对 $x=0$ 镜像：$q_4 = +q$ 在 $(-a,-b)$（与 $q_3$ 重合）。

实际上，求解区域（第一象限）之外只需 \textbf{3} 个镜像电荷：$q' = -q$ 在 $(-a,b)$，$q'' = -q$ 在 $(a,-b)$，$q''' = +q$ 在 $(-a,-b)$。

验证：在 $x=0$ 上，$q$ 与 $q'$ 抵消，$q''$ 与 $q'''$ 抵消；在 $y=0$ 上，$q$ 与 $q''$ 抵消，$q'$ 与 $q'''$ 抵消。

总电势：$\phi = \dfrac{q}{4\pi\eps_0}\left(\dfrac{1}{r_0} - \dfrac{1}{r_1} - \dfrac{1}{r_2} + \dfrac{1}{r_3}\right)$。\qedsymbol

[量纲校验] 电势由点电荷叠加，量纲一致。\qedsymbol




% ============================================================

\begin{center}\small \hyperlink{ch3-goal}{上一章：第3章} \quad|\quad \hyperlink{ch5-goal}{下一章：第5章}\end{center}

% 第5章：静磁学（优化版）
% ============================================================
\chapter{静磁学}
\epigraph{想象力比知识更重要。}{--- 阿尔伯特·爱因斯坦}

\hypertarget{ch5-goal}{}\begin{insight}{本章核心目标}{ch5-goal-v2}
建立静磁学的完整框架：毕奥-萨伐尔定律、安培环路定理、磁矢势，以及磁偶极子的性质。注意静磁学与静电学的深刻类比和本质区别——磁场是有旋无源场，不存在磁单极子。
\end{insight}

静磁学研究\textbf{恒定电流}产生的磁场。虽然静磁学在结构上与静电学有许多相似之处，但磁场的有旋性（$\curl\vecB \neq 0$）使其数学性质与静电场截然不同。

\begin{tip}{本章新符号}{ch5-new-symbols}
\begin{tabular}{@{}ll@{}}
$\vecA$ & 磁矢势（$\vecB = \curl\vecA$，单位 \si{T.m} = \si{Wb/m}） \\
$\vect{\mu}$ & 磁偶极矩（$\vect{\mu} = I\vect{S}$，单位 \si{A.m^2}） \\
\end{tabular}
\end{tip}

\begin{history}{静磁学的发展}{magnetostatics-history}
磁现象的发现远早于电——天然磁石在古代中国和希腊就被制成指南针。定量研究始于 1820 年：\textbf{奥斯特}发现电流使磁针偏转，首次揭示电与磁的联系；\textbf{毕奥}和\textbf{萨伐尔}随后给出电流元磁场的定量规律；\textbf{安培}在 1820--1825 年间建立电流间作用力的完整理论并提出"分子电流"假说。1831 年\textbf{法拉第}发现电磁感应，为麦克斯韦的统一理论铺平了道路。
\end{history}

\section{毕奥-萨伐尔定律}

\begin{theorem}{毕奥-萨伐尔定律}{biot-savart-v2}
电流元 $I\diff\vect{l}'$ 在 $\vect{r}$ 处产生的磁场为：
\[
\diff\vecB = \frac{\muz}{4\pi} \frac{I\diff\vect{l}' \times \unitvec{R}}{R^2}
\]
其中 $\vect{R} = \vect{r} - \vect{r}'$，$R = |\vect{R}|$。对连续电流分布：
\[
\vecB(\vect{r}) = \frac{\muz}{4\pi} \int \frac{\vecJ(\vect{r}') \times \unitvec{R}}{R^2} \, \diff^3r'
\]
其中 $\muz = 4\pi \times 10^{-7} \, \si{T.m/A}$ 是\textbf{真空磁导率}。
\end{theorem}

\begin{insight}{毕奥-萨伐尔定律的物理图像}{biot-savart-insight-v2}
毕奥-萨伐尔定律与库仑定律有深刻的相似性和区别：
\begin{itemize}
\item 相似性：都遵循 $1/r^2$ 的距离依赖，都是叠加原理的基础
\item 区别：库仑定律中源是标量（电荷），场是径向的；毕奥-萨伐尔定律中源是矢量（电流元），场是\textbf{环绕}电流的（由叉积决定）
\end{itemize}

叉积 $\diff\vect{l}' \times \unitvec{R}$ 的方向垂直于 $\diff\vect{l}'$ 和 $\vect{R}$ 所在平面，由右手定则确定。这意味着：电流产生的磁场线总是\textbf{环绕}电流的闭合回路——没有起点和终点。
\end{insight}

\begin{example}{无限长直导线的磁场}{straight-wire-b-v2}
\difficulty{基础}
求电流为 $I$ 的无限长直导线在距离 $\rho$ 处的磁场。
\tcblower
\textbf{解答：}

设导线沿 $z$ 轴。电流元 $I\diff z'\,\unitvec{z}$ 到场点的距离矢量为 $\vect{R} = \rho\unitvec{\rho} - z'\unitvec{z}$。

$\diff\vect{l}' \times \vect{R} = \diff z'\,\unitvec{z} \times (\rho\unitvec{\rho} - z'\unitvec{z}) = \rho\diff z'\,\unitvec{\phi}$

（因为 $\unitvec{z} \times \unitvec{\rho} = \unitvec{\phi}$）

$\diff B = \dfrac{\muz I}{4\pi} \dfrac{\rho\diff z'}{(\rho^2 + z'^2)^{3/2}}$

积分：$B = \dfrac{\muz I\rho}{4\pi} \int_{-\infty}^{\infty} \dfrac{\diff z'}{(\rho^2+z'^2)^{3/2}}$

令 $z' = \rho\tan\theta$，$\diff z' = \rho\sec^2\theta\,\diff\theta$：

$\displaystyle\int \dfrac{\rho\sec^2\theta\,\diff\theta}{\rho^3\sec^3\theta} = \frac{1}{\rho^2}\int\cos\theta\,\diff\theta = \frac{\sin\theta}{\rho^2} = \frac{z'}{\rho^2\sqrt{\rho^2+z'^2}}$

代入上下限：$\left[\dfrac{z'}{\rho^2\sqrt{\rho^2+z'^2}}\right]_{-\infty}^{\infty} = \dfrac{1}{\rho^2}(1 - (-1)) = \dfrac{2}{\rho^2}$

因此：
\[
B = \frac{\muz I\rho}{4\pi} \cdot \frac{2}{\rho^2} = \frac{\muz I}{2\pi\rho}
\]
方向沿 $\unitvec{\phi}$（环绕导线）。
\end{example}

\begin{example}{圆形电流环轴线上的磁场}{loop-b-v2}
\difficulty{提高}
半径为 $R$ 的圆形电流环，电流为 $I$。求轴线上距圆心 $z$ 处的磁场。
\tcblower
\textbf{解答：}

由对称性，只有 $z$ 分量非零。电流元产生的磁场 $d\vecB$ 垂直于 $\diff\vect{l}'$ 和 $\vect{R}$。

设源点到场点的矢量 $\vect{R} = z\unitvec{z} - R\unitvec{\rho}$，其模长 $R_{\!d} = \sqrt{R^2+z^2}$。电流元 $\diff\vect{l}' = \diff l'\,\unitvec{\phi}$ 沿切向，与 $\vect{R}$ \textbf{垂直}（$\unitvec{\phi}\cdot\vect{R} = 0$），故叉积的模长：

$|\diff\vect{l}' \times \vect{R}| = \diff l'\, R_{\!d}$

毕奥--萨伐尔定律给出电流元磁场的大小：
\[
dB = \dfrac{\muz I}{4\pi} \dfrac{|\diff\vect{l}' \times \vect{R}|}{R_{\!d}^3} = \dfrac{\muz I}{4\pi} \dfrac{\diff l'}{R^2+z^2}
\]

$d\vecB$ 的方向垂直于 $\unitvec{\phi}$ 与 $\vect{R}$ 张成的平面，与 $z$ 轴夹角等于 $\vect{R}$ 与 $\unitvec{\rho}$ 的夹角——投影到 $z$ 轴的因子为 $\cos\angle(\vect{R},\unitvec{\rho}) = R/R_{\!d}$：

$dB_z = dB \cdot \dfrac{R}{R_{\!d}} = \dfrac{\muz I R\,\diff l'}{4\pi(R^2+z^2)^{3/2}}$

积分一圈（$\oint\diff l' = 2\pi R$）：
\[
B_z = \frac{\muz I R}{4\pi(R^2+z^2)^{3/2}} \cdot 2\pi R = \frac{\muz I R^2}{2(R^2+z^2)^{3/2}}
\]

在圆心处（$z=0$）：$B = \dfrac{\muz I}{2R}$

远处（$z \gg R$）：$B \approx \dfrac{\muz I R^2}{2z^3} = \dfrac{\muz}{4\pi} \dfrac{2m}{z^3}$，其中 $m = I\pi R^2$ 是磁偶极矩。

[量纲校验] $[\muz I R^2/R_{\!d}^3] = (\mathrm{H/m})(\mathrm{A})(\mathrm{m^2})/\mathrm{m^3} = \mathrm{H\cdot A/m^2} = \mathrm{T}$，与磁场量纲一致。\checkmark
\end{example}

\section{安培环路定理}

\subsection{积分形式}

\begin{theorem}{安培环路定理（积分形式）}{ampere-law-int-v2}
\[
\oint_C \vecB \cdot \diff\vect{l} = \muz I_{\text{enc}}
\]
其中 $I_{\text{enc}}$ 是穿过以 $C$ 为边界的任意曲面的净电流。
\end{theorem}

\begin{insight}{安培定理与高斯定理的对比}{ampere-vs-gauss-v2}
\begin{center}
\begin{tabular}{@{}lll@{}}
\toprule
& \textbf{高斯定理} & \textbf{安培定理} \\
\midrule
场源 & 电荷（标量） & 电流（矢量） \\
积分类型 & 面积分 & 线积分 \\
对称性要求 & 球/柱/平面对称 & 柱/平面对称 \\
结果 & 场的大小 & 场的环流 \\
\bottomrule
\end{tabular}
\end{center}

安培定理比高斯定理更"微妙"，因为电流的方向很重要——只有与回路"链接"的电流才贡献。
\end{insight}

\subsection{微分形式}

由斯托克斯定理：$\displaystyle\oint_C \vecB\cdot\diff\vect{l} = \int_S (\curl\vecB)\cdot\diff\vect{S} = \muz \int_S \vecJ\cdot\diff\vect{S}$

由于这对任意曲面 $S$ 成立：

\begin{theorem}{安培定理（微分形式）}{ampere-law-diff-v2}
\[
\curl\vecB = \muz\vecJ
\]
这是静磁学的核心微分方程。
\end{theorem}

\begin{example}{利用安培定理求磁场}{ampere-app-ex-v2}
\difficulty{基础}
求无限长圆柱形导线（半径 $R$，均匀电流密度 $j$）内外的磁场。
\tcblower
\textbf{解答：}

由柱对称性，$\vecB$ 只有 $\phi$ 分量，大小只与 $\rho$ 有关。

\textbf{区域一：$\rho > R$}（导线外）

取半径为 $\rho$ 的圆形安培回路：
\[
B \cdot 2\pi\rho = \muz I = \muz j\pi R^2
\]
\[
B = \frac{\muz j R^2}{2\rho} = \frac{\muz I}{2\pi\rho}
\]
与线电流结果一致。

\textbf{区域二：$\rho < R$}（导线内）

$I_{\text{enc}} = j\pi\rho^2 = I\rho^2/R^2$

$B \cdot 2\pi\rho = \muz I\rho^2/R^2$
\[
B = \frac{\muz I\rho}{2\pi R^2}
\]
磁场在中心为零，线性增长到表面，然后按 $1/\rho$ 衰减。
\end{example}

\section{静磁镜像法}

第4章介绍的静电镜像法可以平行地推广到静磁学。其核心思想是：用假想的\textbf{镜像电流}替代磁介质边界的影响，使得在求解区域内磁场方程不变，同时满足边界条件。

\subsection{铁磁平面边界}

考虑一个磁导率 $\mu \to \infty$ 的理想铁磁平面（如软铁表面），真空区域中有一电流分布。铁磁表面要求 $\vecB$ 线几乎\textbf{垂直}入射（见第6章例题\ref{ex:ferromagnet-bc-ex-v2}）。

对于\textbf{无限长直导线}平行于铁磁平面（距离平面 $d$），镜像法给出：
\begin{itemize}
\item 镜像电流位于平面另一侧对称位置（距离 $d$）
\item 镜像电流方向与\textbf{原电流相同}
\item 镜像电流大小 $I' = I$（对 $\mu \to \infty$ 的铁磁平面）
\end{itemize}

这与导体平面的静电镜像（异号电荷）形成对照——铁磁平面的镜像电流\textbf{同向}，因为铁磁体要\textbf"增强"而不是抵消磁场。

\begin{tip}{对比记忆}{magnetic-image-tip}
\begin{center}
\begin{tabular}{@{}lll@{}}
\toprule
\textbf{问题} & \textbf{镜像位置} & \textbf{镜像强度/方向} \\
\midrule
导体平面+点电荷 & 对称点 & $q' = -q$（异号） \\
铁磁平面+平行电流 & 对称点 & $I' = +I$（同向） \\
\bottomrule
\end{tabular}
\end{center}
\end{tip}

\subsection*{超导迈斯纳效应的镜像（拓展）}

超导体具有\textbf{完全抗磁性}（迈斯纳效应，$\vecB = 0$ 在体内）。当电流靠近超导体表面时，表面会感应出\textbf{屏蔽电流}，使得体内 $\vecB = 0$。用镜像法描述时，镜像电流位于对称位置但方向\textbf{相反}（$I' = -I$），以保证超导体表面处 $\vecB_\parallel = 0$。这与导体静电镜像（异号电荷）更为相似。

\begin{example}{无限长直导线平行于铁磁平面}{ferro-image-ex-v2}
\difficulty{提高}
电流为 $I$ 的无限长直导线平行于理想铁磁平面（$\mu \to \infty$），距离平面 $d$。求真空侧磁场分布及导线单位长度所受的力。
\tcblower
\textbf{解答：}

设铁磁平面为 $y=0$，导线位于 $y=d$ 沿 $z$ 轴方向。镜像电流 $I'=I$ 位于 $y=-d$。

真空侧任意点 $(x,y)$（$y>0$）的磁场由两导线叠加：
\[
\vecB = \frac{\muz I}{2\pi r_1}\unitvec{\phi}_1 + \frac{\muz I'}{2\pi r_2}\unitvec{\phi}_2
\]
其中 $r_1 = \sqrt{x^2+(y-d)^2}$，$r_2 = \sqrt{x^2+(y+d)^2}$。

在铁磁表面 $y=0$ 上，镜像电流的选取应使切向磁场为零：$\vect{H}$ \textbf{只有法向分量}（即 $H_{1t}=H_{2t}=0$，因 $\mu\to\infty$ 时铁磁体内 $\vect{H}\to\vect{0}$）。注意这里\textbf{法向性属于 $\vect{H}$ 而非 $\vecB$}——$\vecB$ 的法向分量 $B_y$ 由 $B_{1n}=B_{2n}$ 定出；而在真空侧（$y>0$）的\textbf{内部}，$\vecB$ 一般同时含有 $x$、$y$ 两个分量（只有表面 $y=0$ 这一层上 $B_x=0$）。

\textbf{导线受力}：镜像电流对原导线的作用力。两导线间距 $2d$，电流同向相吸：
\[
F = \frac{\muz I^2}{2\pi(2d)} = \frac{\muz I^2}{4\pi d}
\]
方向指向铁磁平面（$y$ 减小方向）。这说明铁磁体\textbf{吸引}载流导线。
\end{example}

\section{磁矢势}

\subsection{引入磁矢势}

静磁学的两个基本方程：
\begin{align}
\divg\vecB &= 0 \label{eq:div-b-zero} \\
\curl\vecB &= \muz\vecJ \label{eq:curl-b}
\end{align}

由式\eqref{eq:div-b-zero}，$\vecB$ 可表为某矢量的旋度：

\begin{tip}{为什么要引入 $\vecA$？}{why-vector-potential-tip}
明明有 $\vecB$，为什么还要引入矢势 $\vecA$？动机有三层：

\textbf{1. 数学结构}：$\divg\vecB = 0$ 说磁场无源。而第 1 章证明过：\textbf{无散场一定是某个矢量的旋度}（旋度的散度恒为零）。所以 $\vecB$ 必然可以写成 $\vecB = \curl\vecA$——这不是假设，是数学必然。

\textbf{2. 简化求解}：直接求 $\vecB$ 要同时满足 2 个方程（$\divg\vecB=0$ 和 $\curl\vecB=\muz\vecJ$）；引入 $\vecA$ 后，$\divg\vecB=0$ 被\textbf{自动满足}（$\divg(\curl\vecA)\equiv 0$），只剩一个方程 $\lap\vecA = -\muz\vecJ$——和静电泊松方程完全同形，第 2 章的所有求解技巧（格林函数、镜像法、分离变量）直接照搬。

\textbf{3. 代价}：$\vecA$ 不唯一（$\vecA' = \vecA + \grad\chi$ 给出同一个 $\vecB$）——这就是\textbf{规范自由度}，见下文。此外第 7 章会看到：时变场中 $\vecA$ 不只是数学工具，Aharonov--Bohm 效应表明它有可观测的物理效应。

\textbf{一句话}：引入 $\vecA$ = 把"两个约束的方程"换成"一个自动满足 + 一个同形泊松方程"，用规范自由度做代价换求解效率。
\end{tip}

\begin{definition}{磁矢势}{vector-potential-v2}
\[
\vecB = \curl\vecA
\]
$\vecA$ 称为\textbf{磁矢势}。由矢量分析恒等式 $\divg(\curl\vecA) = 0$，自动满足 $\divg\vecB = 0$。
\end{definition}

\begin{insight}{磁矢势的物理意义}{a-meaning-v2}
与静电学中 $\vecE = -\grad\phi$ 类似，磁矢势是磁场的"势"。但有本质区别：
\begin{itemize}
\item 电势 $\phi$ 是标量，降维了（3个分量 $\to$ 1个标量）
\item 磁矢势 $\vecA$ 仍是矢量，没有降维——这反映了磁场本质上是\textbf{有旋场}
\item 电势有直接的物理意义（电势能），磁矢势的"直接"物理意义更微妙（Aharonov-Bohm效应表明 $\vecA$ 具有可观测效应）
\end{itemize}

磁矢势的\textbf{规范自由度}：若 $\vecA' = \vecA + \grad\chi$，则 $\curl\vecA' = \curl\vecA + \curl(\grad\chi) = \curl\vecA = \vecB$。即不同的 $\vecA$ 可以给出相同的 $\vecB$。通常选择\textbf{库仑规范} $\divg\vecA = 0$ 来固定这种自由度。
\end{insight}

\subsection{磁矢势的方程}

将 $\vecB = \curl\vecA$ 代入式\eqref{eq:curl-b}：
\[
\curl(\curl\vecA) = \grad(\divg\vecA) - \lap\vecA = \muz\vecJ
\]

选择库仑规范 $\divg\vecA = 0$，得：

\begin{theorem}{磁矢势的泊松方程}{a-poisson-v2}
\[
\lap\vecA = -\muz\vecJ
\]
在各分量上：$\lap A_i = -\muz j_i$。这与标量泊松方程形式相同！
\end{theorem}

解为：
\[
\vecA(\vect{r}) = \frac{\muz}{4\pi} \int \frac{\vecJ(\vect{r}')}{|\vect{r} - \vect{r}'|} \, \diff^3r'
\]

\subsection{规范变换}

\begin{theorem}{规范变换}{gauge-transform-v2}
磁矢势 $\vecA$ 的\textbf{规范变换}定义为：
\[
\vecA' = \vecA + \grad\chi
\]
其中 $\chi(\vect{r})$ 是任意标量函数。规范变换不改变磁场 $\vecB = \curl\vecA$。
\end{theorem}

\textbf{证明} $\vecB$ 不变：
\[
\vecB' = \curl\vecA' = \curl(\vecA + \grad\chi) = \curl\vecA + \curl(\grad\chi) = \curl\vecA = \vecB
\]
其中用到恒等式 $\curl(\grad\chi) = \vect{0}$（第1章定理\ref{thm:grad-curl-zero-v2}）。

\textbf{库仑规范的推导}：给定任意 $\vecA$，总可以选取适当的 $\chi$ 使得新势满足 $\divg\vecA' = 0$。要求：
\[
\divg\vecA' = \divg\vecA + \divg(\grad\chi) = \divg\vecA + \lap\chi = 0
\]
即需要解泊松方程：
\[
\lap\chi = -\divg\vecA
\]
该方程在数学上恒有解（对应第4章的格林函数方法），因此\textbf{库仑规范 $\divg\vecA = 0$ 总可以实现}。

\begin{tip}{学习预警：规范选择初学可略}{coulomb-gauge-skip-tip}
\textbf{第一次阅读可以跳过"库仑规范"相关段落}。在静磁学中，直接令 $\divg\vecA = 0$ 只是众多可能选择中最方便的一种——它的唯一作用是把矢势方程 $\lap\vecA = -\muz\vecJ$ 化成与静电泊松方程\textbf{完全同形}的形式，从而可以直接套用第 2 章的全部求解技巧（格林函数、镜像法、分离变量）。规范自由度真正的物理讨论要到\textbf{第 7 章}（时变场中的库仑规范 vs 洛伦兹规范完整对比表）和\textbf{第 10 章}（协变形式 $\partial_\mu A^\mu = 0$）才展开。初学时只需记住一句话：$\vecA$ 不唯一，$\vecB$ 才是物理。
\end{tip}

\begin{insight}{规范自由度的物理意义}{gauge-meaning-v2}
规范自由度意味着 $\vecA$ 本身\textbf{不可直接测量}——只有 $\vecB = \curl\vecA$ 才有直接的物理意义。这与静电学中电势零点的任意性类似：$\phi$ 可以整体加减常数而不改变 $\vecE = -\grad\phi$。

两者的对比：
\begin{itemize}
\item 静电学：$\phi' = \phi + C$（常数）$\Rightarrow$ $\vecE$ 不变
\item 静磁学：$\vecA' = \vecA + \grad\chi$（函数）$\Rightarrow$ $\vecB$ 不变
\end{itemize}
磁学的规范自由度\textbf{更大}（任意函数 vs 任意常数），这反映了矢量势比标量势多一个自由度。

\textbf{前瞻}：第7章研究时变电磁场时，规范变换将扩展为包含标量势 $\phi$ 的联合变换；第10章相对论电动力学中，$\phi$ 与 $\vecA$ 将组合成\textbf{四维电磁势} $A^\mu$，规范变换获得更优美的协变形式。
\end{insight}

\section{磁偶极子}

\begin{definition}{磁偶极矩}{magnetic-dipole-v2}
对于平面电流环，\textbf{磁偶极矩}定义为：
\[
\vect{\mu} = I \vect{S}
\]
其中 $\vect{S}$ 是回路面积矢量（方向由右手定则确定）。
\end{definition}

磁偶极子在远处的磁场：
\[
\vecB(\vect{r}) = \frac{\muz}{4\pi r^3} \left[ 3(\vect{\mu}\cdot\unitvec{r})\unitvec{r} - \vect{\mu} \right]
\]

这与电偶极子的电场公式形式完全相同！

\begin{example}{磁偶极子在外磁场中的受力}{dipole-force-v2}
\difficulty{提高}
磁偶极子 $\vect{\mu}$ 在外磁场 $\vecB$ 中的势能为 $U = -\vect{\mu}\cdot\vecB$。求其所受的力和力矩。
\tcblower
\textbf{解答：}

\textbf{力}：$\vect{F} = -\grad U = \grad(\vect{\mu}\cdot\vecB)$

对于恒定 $\vect{\mu}$，利用矢量恒等式：
\[
\vect{F} = (\vect{\mu}\cdot\grad)\vecB
\]

在非均匀磁场中，磁偶极子受到指向磁场增强方向的力。

\textbf{力矩}：$\vect{\tau} = \vect{\mu} \times \vecB$

这解释了为什么指南针会沿地磁场方向排列。
\end{example}

\begin{review}{第5章要点回顾}{ch5-review}
\begin{center}
\begin{tabular}{@{}ll@{}}
\toprule
\textbf{方程/公式} & \textbf{内容} \\
\midrule
$\vecB = \dfrac{\muz}{4\pi}\displaystyle\int\dfrac{\vecJ\times\unitvec{R}}{R^2}\diff^3r'$ & 毕奥-萨伐尔定律 \\
$\oint\vecB\cdot\diff\vect{l} = \muz I_{\text{enc}}$ & 安培环路定理（积分） \\
$\curl\vecB = \muz\vecJ$ & 安培环路定理（微分） \\
$\divg\vecB = 0$ & 磁场无源 \\
$\vecB = \curl\vecA$ & 磁矢势定义 \\
$\lap\vecA = -\muz\vecJ$（库仑规范） & 矢势方程 \\
$\vect{\mu} = I\vect{S}$ & 磁偶极矩 \\
\bottomrule
\end{tabular}
\end{center}
\end{review}

\begin{tip}{本章核心框架}{ch5-summary}
\textbf{一句话总结}：静磁场由电流产生，无源有旋，需要用矢势$\vecA$描述。\\
\textbf{核心方程}：$\divg\vecB=0$（无源）、$\curl\vecB=\mu_0\vecJ$（安培定理，恒定电流）、$\lap\vecA=-\mu_0\vecJ$（库仑规范下）。\\
\textbf{核心方法}：毕奥-萨伐尔定律（积分叠加）、安培环路定理（对称性）、矢势法（自动满足$\divg\vecB=0$）。\\
\textbf{衔接}：第6章引入磁介质，类比电介质引入磁化强度$\vect{M}$和$\vect{H}$场。
\end{tip}

% ---------- 第5章随堂自测 ----------
\section*{第5章随堂自测}
\addcontentsline{toc}{section}{第5章随堂自测}

\begin{miniquiz}
\textbf{一、选择题}
\begin{enumerate}[label=\arabic*.]
\item 安培环路定理 $\oint\vecB\cdot\diff\vect{l} = \muz I_{enc}$ 适用于：
  \begin{enumerate}[label=(\Alph*),nosep]
  \item 只有恒定电流
  \item 任何电流分布
  \item 只有对称性高的电流分布
  \item 只有真空中的电流
  \end{enumerate}
\begin{ocg}{答案}{ocg答案5}{off}
  \textit{答案：A。原始安培定理仅适用于恒定电流（静磁学）。}
\end{ocg}

\item 长直螺线管内部的磁场：
  \begin{enumerate}[label=(\Alph*),nosep]
  \item 为零
  \item 均匀且平行于轴线
  \item 随半径增大而增大
  \item 随半径增大而减小
  \end{enumerate}
\begin{ocg}{答案}{ocg答案5}{off}
  \textit{答案：B。理想长直螺线管内部磁场均匀，$B = \muz n I$。}
\end{ocg}

\item 矢势 $\vecA$ 的物理意义是：
  \begin{enumerate}[label=(\Alph*),nosep]
  \item 磁场的能量
  \item 使 $\vecB = \curl\vecA$
  \item 磁场的源
  \item 磁标势
  \end{enumerate}
\begin{ocg}{答案}{ocg答案5}{off}
  \textit{答案：B。矢势的定义就是使$\vecB=\curl\vecA$，且自动满足$\divg\vecB=0$。}
\end{ocg}
\end{enumerate}

\vspace{0.5em}
\textbf{二、判断题}
\begin{enumerate}[label=\arabic*.,resume]
\item 磁场$\vecB$的散度恒为零意味着不存在磁单极子。
\begin{ocg}{答案}{ocg答案5}{off}
  \textit{答案：正确。$\divg\vecB=0$等价于"磁荷密度恒为零"。}
\end{ocg}

\item 静磁学中$\vecB$和$\vect{H}$在任何情况下都平行。
\begin{ocg}{答案}{ocg答案5}{off}
  \textit{答案：错误。只有在各向同性介质中才有$\vecB = \mu\vect{H}$。各向异性介质中二者方向可以不同。}
\end{ocg}
\end{enumerate}
\end{miniquiz}
\toggleButton{显示答案}{ocg答案5}

\begin{chapterreview}
\begin{itemize}[nosep]
\item \textbf{核心定律}：毕奥-萨伐尔定律、安培力定律、安培环路定理。
\item \textbf{核心概念}：矢势$\vecA$、库仑规范$\divg\vecA=0$。
\item \textbf{自检验}：能否用毕奥-萨伐尔定律推导无限长直导线的磁场？
\item \textbf{衔接预告}：第6章将讨论磁介质、磁化强度$\vect{M}$和磁场强度$\vect{H}$。
\end{itemize}
\end{chapterreview}

\begin{tip}{本章自检清单}{ch5-checklist}
学完本章后对照下列条目逐条自测，全部通过即可进入下一章；有未通过项时，回到对应小节复习：
\begin{itemize}[label=$\square$, nosep]
\item 用毕奥--萨伐尔定律计算直导线与电流环的磁场
\item 用安培环路定理处理对称电流分布
\item 理解矢势 $\vecA$ 与规范自由度
\item 计算磁偶极矩及其在外场中的受力
\end{itemize}
\end{tip}

% ---------- 第5章综合练习题 ----------
\section*{第5章综合练习题}
\addcontentsline{toc}{section}{第5章综合练习题}
\begin{enumerate}[label=\textbf{\arabic*.}]
\item \textbf{基础} 半径为 $R$ 的无限长圆柱导体，通有均匀电流 $I$，用安培环路定理求柱内外磁场分布。
\item \textbf{基础} 证明：对任意规范函数 $\chi$，规范变换 $\vecA'=\vecA+\nabla\chi$ 不改变磁场 $\vecB$。
\item \textbf{提高} 在库仑规范下，证明静磁学中 $\vecA$ 满足泊松方程 $\lap\vecA=-\muz\vecJ$。
\item \textbf{提高} 无限长直导线（电流 $I$）平行于铁磁平面（磁导率 $\mu\to\infty$），距离平面 $d$。求镜像电流的大小和方向，并给出空间中磁场分布。
\item \textbf{挑战} 证明：对于任意闭合电流回路，磁矢势在远处可展开为 $\vecA(\vect{r})=\frac{\muz}{4\pi}\frac{\vect{m}\times\vect{r}}{r^3}+\cdots$，其中 $\vect{m}$ 为磁偶极矩。
\end{enumerate}
\newpage
\section*{习题解答}
\addcontentsline{toc}{section}{习题解答}

\textbf{题1解析：}

[分析] 本题利用安培环路定理和对称性分析求圆柱导体内外的磁场分布，与静电学中均匀带电球体问题完全类比。

[解答] 由柱对称性，$\vecB$ 只有 $\phi$ 分量，大小只依赖于 $\rho$。取半径为 $\rho$ 的圆形安培回路。

\textbf{区域一：$\rho > R$}（导线外）
\[
B \cdot 2\pi\rho = \muz I \quad\Rightarrow\quad B = \frac{\muz I}{2\pi\rho}
\]

\textbf{区域二：$\rho < R$}（导线内）
电流密度 $j = I/(\pi R^2)$，$I_{\text{enc}} = j\pi\rho^2 = I\rho^2/R^2$。
\[
B \cdot 2\pi\rho = \muz I\rho^2/R^2 \quad\Rightarrow\quad B = \frac{\muz I\rho}{2\pi R^2}
\]

磁场在中心为零，线性增长到表面，然后按 $1/\rho$ 衰减。\qedsymbol

[量纲校验] $[B] = [MT^{-2}I^{-1}]$，$[\muz I/\rho] = [H/m][A]/[L] = [MT^{-2}I^{-1}]$，一致。\qedsymbol

\textbf{题2解析：}

[分析] 本题证明规范不变性，是理解磁矢势规范自由度的核心。

[解答] 规范变换：$\vecA' = \vecA + \grad\chi$。

计算新磁场：
\[
\vecB' = \curl\vecA' = \curl(\vecA + \grad\chi) = \curl\vecA + \curl(\grad\chi) = \curl\vecA + \vect{0} = \vecB
\]

这里用到矢量恒等式 $\curl(\grad\chi) = \vect{0}$（梯度的旋度恒为零）。因此磁场在任意规范变换下保持不变。\qedsymbol

[量纲校验] $[\grad\chi] = [L^{-1}][\chi]$，若 $[A] = [MLT^{-2}I^{-1}]$，则 $[\chi] = [ML^2T^{-2}I^{-1}]$，$[\grad\chi] = [MLT^{-2}I^{-1}] = [A]$，一致。\qedsymbol

\textbf{题3解析：}

[分析] 本题在库仑规范下推导磁矢势的泊松方程，是静磁学中计算矢势的基础方程。

[解答] 库仑规范：$\divg\vecA = 0$。

磁矢势满足：$\curl\vecB = \curl(\curl\vecA) = \grad(\divg\vecA) - \lap\vecA = \muz\vecJ$。

在库仑规范下 $\divg\vecA = 0$，故：
\[
-\lap\vecA = \muz\vecJ \quad\Rightarrow\quad \lap\vecA = -\muz\vecJ
\]

这是每个分量 $A_i$ 的泊松方程，可用格林函数求解。\qedsymbol

[量纲校验] $[\lap A] = [L^{-2}][MLT^{-2}I^{-1}] = [MT^{-3}I^{-1}]$，$[\muz j] = [H/m][A/L^2] = [MT^{-3}I^{-1}]$，一致。\qedsymbol

\textbf{题4解析：}

[分析] 本题考查铁磁平面的磁镜像法，关键是理解铁磁平面与导体平面的镜像区别。

[解答] 对于磁导率 $\mu \to \infty$ 的理想铁磁平面，边界条件要求 $\vecB$ 线几乎垂直于表面射出。

无限长直导线平行于铁磁平面，距离 $d$：
\begin{itemize}
\item 镜像电流位于平面另一侧对称位置（距离 $d$）
\item 镜像电流方向与原电流\textbf{相同}（$I' = I$）
\item 这与导体平面的静电镜像（异号电荷）形成鲜明对照——铁磁平面要"增强"磁场
\end{itemize}

空间中任一点的磁场：$\vecB = \dfrac{\muz I}{2\pi}\left(\dfrac{1}{\rho_1}\unitvec{\phi}_1 + \dfrac{1}{\rho_2}\unitvec{\phi}_2\right)$，其中 $\rho_1, \rho_2$ 到场点和镜像点的距离。\qedsymbol

[量纲校验] $[I'] = [I]$，$[\muz I/\rho] = [B]$，一致。\qedsymbol

\textbf{题5解析：}

[分析] 本题推导磁矢势的多极展开，重点是识别磁偶极矩项。

[解答] 磁矢势：$\vecA(\vect{r}) = \dfrac{\muz}{4\pi}\displaystyle\int\dfrac{\vecJ(\vect{r}')}{|\vect{r} - \vect{r}'|}\,\diff^3r'$。

对 $r' \ll r$ 展开：$\dfrac{1}{|\vect{r} - \vect{r}'|} \approx \dfrac{1}{r} + \dfrac{\vect{r}\cdot\vect{r}'}{r^3} + \cdots$。

单极项：$\vecA_0 = \dfrac{\muz}{4\pi r}\int\vecJ(\vect{r}')\,\diff^3r' = \vect{0}$（稳恒电流 $\divg\vecJ = 0$ 意味着 $\int\vecJ\,\diff^3r' = 0$）。

偶极项：$\vecA_1 = \dfrac{\muz}{4\pi r^3}\int\vecJ(\vect{r}')(\vect{r}\cdot\vect{r}')\,\diff^3r'$。

利用矢量恒等式：$\int\vecJ(\vect{r}\cdot\vect{r}')\,\diff^3r' = -\vect{r}\times\int\dfrac{1}{2}(\vect{r}'\times\vecJ)\,\diff^3r' = \vect{m}\times\vect{r}$。

其中 $\vect{m} = \dfrac{1}{2}\int\vect{r}'\times\vecJ\,\diff^3r'$ 是磁偶极矩。

因此：$\vecA(\vect{r}) = \dfrac{\muz}{4\pi}\dfrac{\vect{m}\times\vect{r}}{r^3} + \cdots$。\qedsymbol

[量纲校验] $[A] = [MLT^{-2}I^{-1}]$，$[\muz m/r^2] = [H/m][A\cdot L^2]/[L^2] = [MLT^{-2}I^{-1}]$，一致。\qedsymbol



\begin{warning}{第5章常见误区}{ch5-warnings}
\begin{enumerate}
\item 毕奥-萨伐尔定律中 $\diff\vecB$ 的方向由叉积决定——先用右手定则判断方向再计算大小。
\item 安培定理中的 $I_{\text{enc}}$ 是\textbf{净}电流，要考虑方向（右手定则确定正负）。
\item $\divg\vecB = 0$ 意味着不存在磁单极子——磁场线永远是闭合回路。
\item 磁矢势 $\vecA$ 不是唯一的（规范自由度），但 $\vecB = \curl\vecA$ 是唯一的。
\item 磁偶极子场公式与电偶极子场公式形式相同，但物理来源完全不同。
\end{enumerate}
\end{warning}

% ============================================================

\begin{center}\small \hyperlink{ch4-goal}{上一章：第4章} \quad|\quad \hyperlink{ch6-goal}{下一章：第6章}\end{center}

% 第6章：磁介质（新增章节，优化版）
% ============================================================
\chapter{磁介质}
\epigraph{没有磁单极子，是这个世界的缺憾，也是它的美妙之处。}{--- 保罗·狄拉克（意译）}

\hypertarget{ch6-goal}{}\begin{insight}{本章核心目标}{ch6-goal-v2}
理解磁介质在磁场中的响应机制——磁化。引入磁场强度 $\vect{H}$ 简化方程，建立磁学边界条件。区分顺磁性、抗磁性和铁磁性三种磁性机制。
\end{insight}

磁介质与电介质类似，但磁现象更为复杂。除了与电极化类似的\textbf{磁化}效应外，磁介质还表现出顺磁性、抗磁性和铁磁性等不同的响应行为。

\begin{tip}{本章新符号}{ch6-new-symbols}
\begin{tabular}{@{}ll@{}}
$\vect{M}$ & 磁化强度（磁偶极矩密度，单位 \si{A/m}） \\
$\vect{H}$ & 磁场强度（$\vect{H} = \vecB/\muz - \vect{M}$，单位 \si{A/m}） \\
$\chi_m$ & 磁化率（无量纲，$\mu_r = 1 + \chi_m$） \\
$\mu_r,\ \mu$ & 相对/绝对磁导率（$\mu = \muz\mu_r$） \\
$\vecJ_M,\ \vect{K}_M$ & 体/面磁化电流密度 \\
\end{tabular}
\end{tip}

\section{磁介质的磁化}

\subsection{磁化的物理机制}

\begin{enumerate}
\item \textbf{抗磁性}：所有物质都具有的抗磁性源于电磁感应——外磁场在原子中感应出反向的环流，产生与外场相反的磁矩。这是\textbf{楞次定律}在原子尺度上的表现。抗磁性很弱，但在超导体中表现为完全抗磁性（迈斯纳效应）。

\item \textbf{顺磁性}：具有永久磁矩的原子/分子在外磁场中趋向于沿磁场方向排列。但热运动会干扰这种排列，因此顺磁效应随温度升高而减弱（居里定律 $\chi \propto 1/T$）。

\item \textbf{铁磁性}：某些材料（如铁、钴、镍）中存在\textbf{交换相互作用}，使相邻原子的磁矩自发平行排列，形成\textbf{磁畴}。即使在零外场下，这些材料也可以具有宏观磁矩。铁磁性是量子力学效应，无法用经典物理完全解释。
\end{enumerate}

\subsection{磁化强度}

\begin{definition}{磁化强度}{magnetization-v2}
\textbf{磁化强度} $\vect{M}$ 定义为单位体积内的磁偶极矩：
\[
\vect{M} = \lim_{\Delta V \to 0} \frac{\sum_i \vect{\mu}_i}{\Delta V}
\]
单位：$\si{A/m}$（与 $\vect{H}$ 相同）。
\end{definition}

磁化电流：磁化介质中存在\textbf{体磁化电流密度}和\textbf{面磁化电流密度}：
\begin{align}
\vecJ_M &= \curl\vect{M} \\
\vect{K}_M &= \vect{M} \times \unitvec{n}
\end{align}

\section{磁场强度}

磁化电流与自由电流一样产生磁场。安培定理应写为：
\[
\curl\vecB = \muz(\vecJ_f + \vecJ_M) = \muz(\vecJ_f + \curl\vect{M})
\]

整理：
\[
\curl\left(\frac{\vecB}{\muz} - \vect{M}\right) = \vecJ_f
\]

\begin{definition}{磁场强度}{h-field-v2}
定义\textbf{磁场强度}：
\[
\vect{H} = \frac{\vecB}{\muz} - \vect{M}
\]
则安培定理变为：
\[
\curl\vect{H} = \vecJ_f
\]
\textbf{关键性质}：$\vect{H}$ 的旋度源仅是自由电流。
\end{definition}

对于线性各向同性磁介质：
\[
\vect{M} = \chi_m \vect{H}
\]
其中 $\chi_m$ 是\textbf{磁化率}。

于是：
\[
\vecB = \muz(\vect{H} + \vect{M}) = \muz(1+\chi_m)\vect{H} = \muz \mu_r\vect{H} = \mu\vect{H}
\]

其中 $\mu_r = 1 + \chi_m$ 是\textbf{相对磁导率}，$\mu = \muz \mu_r$ 是\textbf{绝对磁导率}。

\begin{insight}{理解B、H、M的关系}{bhm-relation-v2}
可以类比电场中的 $\vect{D}, \vecE, \vect{P}$：
\begin{center}
\begin{tabular}{@{}lll@{}}
\toprule
\textbf{电场} & & \textbf{磁场} \\
\midrule
$\vect{D} = \eps_0\vecE + \vect{P}$ & $\longleftrightarrow$ & $\vecB = \muz(\vect{H} + \vect{M})$ \\
$\divg\vect{D} = \rho_f$ & $\longleftrightarrow$ & $\curl\vect{H} = \vecJ_f$ \\
$\vect{P} = \eps_0\chi_e\vecE$ & $\longleftrightarrow$ & $\vect{M} = \chi_m\vect{H}$ \\
$\vect{D} = \eps\vecE$ & $\longleftrightarrow$ & $\vecB = \mu\vect{H}$ \\
\bottomrule
\end{tabular}
\end{center}

\textbf{三个磁学量的完整对比}：
\begin{center}
\small
\begin{tabular}{@{}>{\raggedright\arraybackslash}p{2.2cm}>{\raggedright\arraybackslash}p{3.6cm}>{\raggedright\arraybackslash}p{3.2cm}>{\raggedright\arraybackslash}p{3.6cm}@{}}
\toprule
 & $\vecB$（磁感应强度） & $\vect{H}$（磁场强度） & $\vect{M}$（磁化强度） \\
\midrule
定义来源 & 基本场（电流元受力 $I\diff\vect{l}\times\vecB$） & 辅助定义：$\vect{H} = \vecB/\muz - \vect{M}$ & 偶极矩密度 $\vect{M} = \sum\vect{\mu}_i/\Delta V$ \\
\midrule
方程源 & $\curl\vecB = \muz(\vecJ_f+\vecJ_M)$（全部电流） & $\curl\vect{H} = \vecJ_f$（仅自由电流） & $\vecJ_M = \curl\vect{M}$（其旋度即磁化电流） \\
\midrule
单位 & \si{T}（特斯拉） & \si{A/m} & \si{A/m} \\
\midrule
线性介质中 & $\vecB = \mu\vect{H}$ & $\vect{H} = \vecB/\mu$ & $\vect{M} = \chi_m\vect{H}$ \\
\midrule
边界条件 & 法向连续（无磁单极） & 切向：$\unitvec{n}\times(\vect{H}_2-\vect{H}_1)=\vect{K}_f$ & 法向：$M_{2n}-M_{1n}=-K_{Mb}$ \\
\midrule
物理角色 & 决定洛伦兹力与能量 & 简化方程，绕开磁化电流 & 描述介质磁响应本身 \\
\bottomrule
\end{tabular}
\end{center}

但要注意一个关键区别：$\vecE$ 是基本场（电荷受力 $q\vecE$），而 $\vecB$ 是基本场（电流元受力 $I\diff\vect{l}\times\vecB$）。对应关系中，$\vecE \leftrightarrow \vecB$（基本场），$\vect{D} \leftrightarrow \vect{H}$（辅助场）。

另一个常见困惑：$\vect{H}$ 与 $\vect{M}$ 单位相同（\si{A/m}）但含义完全不同——$\vect{H}$ 的源是\textbf{自由电流}（外部可控），$\vect{M}$ 的源是介质内的\textbf{分子环流}（介质响应）。真空中 $\vect{M}=0$ 但 $\vect{H}\neq 0$。磁介质分界面上的"电流"同理要分清：自由电流出现在 $\vect{H}$ 的切向边界条件中，磁化面电流出现在 $\vect{M}$ 的法向跃变中。
\end{insight}

\begin{warning}{$\vect{P}$ 与 $\vect{M}$ 的本质不对称}{p-m-asymmetry-v2}
初学者常对以下对比感到困惑：
\begin{itemize}
\item 极化偶极密度 $\vect{P}$ 产生\textbf{标量}束缚电荷：$\rho_b = -\divg\vect{P}$
\item 磁化偶极密度 $\vect{M}$ 产生\textbf{矢量}束缚电流：$\vecJ_M = \curl\vect{M}$
\end{itemize}

为什么一个是散度、一个是旋度？为什么一个是标量、一个是矢量？

\textbf{根源在于自然界不存在磁单极子}。电场可以有源（电荷），因此极化用\textbf{散度}描述"源"的分布；而磁场无源（$\divg\vecB=0$），磁化只能产生\textbf{涡旋}（电流），用旋度描述。

更深层的电磁不对称：
\begin{itemize}
\item 电：$\divg\vecE = \rho/\eps_0$ $\Rightarrow$ $\rho_b = -\divg\vect{P}$（标量方程对标量方程）
\item 磁：$\curl\vecB = \muz\vecJ$ $\Rightarrow$ $\vecJ_M = \curl\vect{M}$（矢量方程对矢量方程）
\end{itemize}

这种不对称性将在第7章被"修补"：麦克斯韦引入位移电流后，方程组展现出更优美的对称结构——但磁单极子的缺失仍然是一个深刻的物理事实。
\end{warning}

\section{磁学边界条件}

\begin{theorem}{磁学边界条件}{magnetic-bc-v2}
在两种磁介质的界面处：
\begin{align}
B_{1n} &= B_{2n} \quad\text{（法向 } \vecB \text{ 连续）} \\
H_{2t} - H_{1t} &= K_f \quad\text{（切向 } \vect{H} \text{ 跳跃 = 自由面电流）}
\end{align}
若界面上无自由面电流（$K_f = 0$）：$H_{1t} = H_{2t}$（切向 $\vect{H}$ 连续）。
\end{theorem}

\textbf{推导}：
\begin{itemize}
\item 法向 $\vecB$ 连续：对药盒形高斯面，$\oint\vecB\cdot\diff\vect{S} = 0$（因为 $\divg\vecB = 0$ 处处成立），故 $B_{2n} = B_{1n}$
\item 切向 $\vect{H}$ 跳跃：对矩形回路，$\oint\vect{H}\cdot\diff\vect{l} = I_{f,\text{enc}}$，当回路跨界面且高度趋于零时，$H_{2t}l - H_{1t}l = K_f l$
\end{itemize}

\begin{example}{铁磁材料表面的磁场}{ferromagnet-bc-ex-v2}
\difficulty{提高}
铁磁材料（$\mu \gg \muz$）表面外为真空。若铁磁材料内部的 $\vecB$ 与法向成 $\theta_1$ 角，求外部的角度 $\theta_2$。
\tcblower
\textbf{解答：}

设介质1为铁磁材料，介质2为真空。法向从1指向2。

$B_{1n} = B_{2n}$：$B_1\cos\theta_1 = B_2\cos\theta_2$

$H_{1t} = H_{2t}$（无自由面电流）：$\dfrac{B_1\sin\theta_1}{\mu} = \dfrac{B_2\sin\theta_2}{\muz}$

两式相除：
\[
\frac{\tan\theta_1}{\mu/\muz} = \tan\theta_2
\]

由于 $\mu/\muz \gg 1$，$\tan\theta_2 \ll \tan\theta_1$，即 $\theta_2 \approx 0$。

\textbf{结论}：铁磁材料表面的磁场线几乎\textbf{垂直}于表面射出——这与导体表面电场线垂直于表面类似！这一性质被广泛用于磁路设计和磁屏蔽。
\end{example}

\begin{example}{铁磁--真空分界面的数值变式}{ferro-vacuum-bc-ex-v2}
\difficulty{提高}
上例取具体数值：铁磁材料 $\mu = 2000\muz$，$\vect{H}_1$ 大小 $H_1 = \SI{100}{A/m}$、与法向成 $\theta_1 = 60^\circ$。求真空侧的 $H_{2n}$、$H_{2t}$、$H_2$ 和角度 $\theta_2$。
\tcblower
\textbf{解答：}

由边界条件（沿用上例记号）：
\begin{align*}
H_{2t} &= H_{1t} = H_1\sin\theta_1 = 100\sin 60^\circ = \SI{86.6}{A/m} \\
H_{2n} &= \frac{\mu}{\muz}H_1\cos\theta_1 = 2000 \times 100 \times \cos 60^\circ = \SI{1.0e5}{A/m}
\end{align*}

合成：
\[
H_2 = \sqrt{H_{2n}^2 + H_{2t}^2} = \sqrt{(1.0\times10^5)^2 + (86.6)^2} \approx \SI{1.0e5}{A/m}
\]

\[
\tan\theta_2 = \frac{H_{2t}}{H_{2n}} = \frac{86.6}{1.0\times10^5} = 8.66\times10^{-4} \quad\Rightarrow\quad \theta_2 \approx 0.05^\circ
\]

\textbf{结论}：真空侧磁场几乎完全垂直于表面（偏角仅 $0.05^\circ$）。数值验证了铁磁界面"磁场线垂直射出"的性质——$\mu/\muz$ 越大偏角越小，磁路设计中正是利用这一点把磁通"约束"在铁芯内。
\end{example}

\subsection{磁多极展开}

类比静电学的多极展开，矢势的远场近似也可以展开为多极矩的叠加。

\begin{theorem}{磁偶极矩的矢势}{magnetic-dipole-potential}
对于局域电流分布，远场矢势的主导项为：
\[
\vecA(\vect{r}) = \frac{\mu_0}{4\pi}\frac{\vect{m}\times\vect{r}}{r^3}
\]
其中磁偶极矩 $\vect{m} = \frac{1}{2}\int \vect{r}'\times\vecJ(\vect{r}')\,\diff^3r'$。
\end{theorem}

% ---------- 图：磁偶极子的磁场线 ----------
\begin{center}
\begin{tikzpicture}[scale=0.9]
% 偶极子（小电流环，位于原点，轴沿 z）
\draw[thick,blue] (-0.5,0) arc (180:360:0.5 and 0.18);
\draw[thick,blue,dashed] (-0.5,0) arc (180:0:0.5 and 0.18);
\node[blue,font=\small,anchor=west] at (0.6,0) {$I$};

% 磁场线：闭合回线从环的上方 (0, 0.2w) 出发，向两侧张开再弯回下方
% 控制点在水平方向远处、垂直方向近处，形成"饱满"的回线
\foreach \w in {1.2, 2.0, 2.8} {
  \draw[thick,red!70] (0,{0.22*\w}) .. controls ({1.1*\w},{0.5*\w}) and ({1.1*\w},{-0.5*\w}) .. (0,{-0.22*\w});
  \draw[thick,red!70] (0,{-0.22*\w}) .. controls ({-1.1*\w},{-0.5*\w}) and ({-1.1*\w},{0.5*\w}) .. (0,{0.22*\w});
}
% 方向箭头（左右最外线上，指示外部 B 从上（N）出、下（S）回）
\draw[->,thick,red!70] ({2.6},0.9) -- ({2.65},0.55);
\draw[->,thick,red!70] ({-2.6},-0.9) -- ({-2.65},-0.55);
% 磁矩方向箭头（轴向上，从环心向上）
\draw[->,thick,blue] (0,0.15) -- (0,1.3) node[above] {$\vect{m}$};
\node[red,font=\small] at (3.1,1.6) {$\vecB$};

\node[align=center,font=\small] at (0,-2.6) {磁偶极子磁场线\\外部似条形磁铁，内部闭合};
\end{tikzpicture}
\end{center}


\begin{insight}{电多极与磁多极的对称性}{electric-magnetic-multipole}
电多极展开：$\phi = \phi_{(0)} + \phi_{(1)} + \phi_{(2)} + \cdots$（单极、偶极、四极）\\
磁多极展开：$\vecA = \vecA_{(1)} + \vecA_{(2)} + \cdots$（无磁单极！）\\
关键差异：磁多极展开从偶极开始，因为不存在磁单极。这是麦克斯韦方程$\divg\vecB=0$的直接后果。
\end{insight}

% ---------- 图：铁磁介质的磁畴结构 ----------
\begin{center}
\begin{tikzpicture}[scale=0.8]
% 未磁化：每个磁畴是带框的小方块，磁化方向随机
\begin{scope}[shift={(-3,0)}]
\draw[thick] (-2,-1.3) rectangle (2,1.3);
\foreach \x/\a in {-1.5/30, -0.5/150, 0.5/80, 1.5/210} {
  \foreach \y in {-0.75, 0, 0.75} {
    \draw[gray!60] (\x-0.42,\y-0.28) rectangle (\x+0.42,\y+0.28);
    \draw[->,thick,blue,rotate around={\a:(\x,\y)}] (\x-0.25,\y) -- (\x+0.25,\y);
  }
}
\node at (0,-2.0) {\textbf{(a) 未磁化}};
\node[align=center,font=\small] at (0,-2.8) {磁畴随机取向\\$\langle\vect{M}\rangle=0$};
\end{scope}

% 磁化饱和：所有磁畴方向一致（向右），外场 H
\begin{scope}[shift={(3,0)}]
\draw[thick] (-2,-1.3) rectangle (2,1.3);
\foreach \x in {-1.5,-0.5,0.5,1.5} {
  \foreach \y in {-0.75, 0, 0.75} {
    \draw[gray!60] (\x-0.42,\y-0.28) rectangle (\x+0.42,\y+0.28);
    \draw[->,thick,blue] (\x-0.25,\y) -- (\x+0.25,\y);
  }
}
\draw[->,thick,green!60!black] (-2,1.9) -- (2,1.9) node[right] {$\vect{H}$};
\node at (0,-2.0) {\textbf{(b) 磁化饱和}};
\node[align=center,font=\small] at (0,-2.8) {外场对齐磁畴\\$\langle\vect{M}\rangle\neq 0$};
\end{scope}
\end{tikzpicture}
\end{center}



\begin{review}{第6章要点回顾}{ch6-review}
\begin{center}
\begin{tabular}{@{}ll@{}}
\toprule
\textbf{公式/概念} & \textbf{内容} \\
\midrule
$\vecJ_M = \curl\vect{M}$，$\vect{K}_M = \vect{M}\times\unitvec{n}$ & 磁化电流 \\
$\vect{H} = \vecB/\muz - \vect{M}$ & 磁场强度 \\
$\curl\vect{H} = \vecJ_f$ & H的安培定理 \\
$\vecB = \mu\vect{H}$（线性介质） & 本构关系 \\
$B_{1n} = B_{2n}$ & B法向连续 \\
$H_{2t} - H_{1t} = K_f$ & H切向跳跃 \\
\bottomrule
\end{tabular}
\end{center}
\end{review}

\begin{tip}{本章核心框架}{ch6-summary}
\textbf{一句话总结}：磁介质通过磁化响应磁场，分类为抗磁、顺磁、铁磁，改变有效磁导率。\\
\textbf{核心关系}：$\vecB=\mu_0(\vect{H}+\vect{M})=\mu\vect{H}$（LIH介质）、$\vecJ_m=\curl\vect{M}$（磁化电流）。\\
\textbf{边界条件}：$B_{1n}=B_{2n}$（法向连续）、$H_{1t}=H_{2t}$（切向连续，无自由面电流）。\\
\textbf{衔接}：第7章将统一电磁学，安培定理修正为麦克斯韦-安培方程，引入位移电流。
\end{tip}

% ---------- 第6章随堂自测 ----------
\section*{第6章随堂自测}
\addcontentsline{toc}{section}{第6章随堂自测}

\begin{miniquiz}
\textbf{一、选择题}
\begin{enumerate}[label=\arabic*.]
\item 磁场强度 $\vect{H}$ 的定义是：
  \begin{enumerate}[label=(\Alph*),nosep]
  \item $\vect{H} = \vecB/\muz$
  \item $\vect{H} = \vecB/\muz - \vect{M}$
  \item $\vect{H} = \vecB/\mu$
  \item $\vect{H} = \vect{M}/\chi_m$
  \end{enumerate}
\begin{ocg}{答案}{ocg答案6}{off}
  \textit{答案：B。这是$\vect{H}$的普适定义，$\vect{H}$的环路定理只含自由电流。}
\end{ocg}

\item 抗磁性材料的磁化率 $\chi_m$：
  \begin{enumerate}[label=(\Alph*),nosep]
  \item $\chi_m > 0$
  \item $\chi_m < 0$
  \item $\chi_m = 0$
  \item $\chi_m \gg 1$
  \end{enumerate}
\begin{ocg}{答案}{ocg答案6}{off}
  \textit{答案：B。抗磁性材料的$\chi_m$为很小的负数。}
\end{ocg}

\item 在磁介质界面上，$\vecB$ 的法向分量的边界条件是：
  \begin{enumerate}[label=(\Alph*),nosep]
  \item $B_{1n} = B_{2n}$
  \item $B_{1n} - B_{2n} = \mu_0 K$
  \item $H_{1n} = H_{2n}$
  \item $B_{1t} = B_{2t}$
  \end{enumerate}
\begin{ocg}{答案}{ocg答案6}{off}
  \textit{答案：A。$\divg\vecB=0$保证$\vecB$法向分量始终连续。}
\end{ocg}
\end{enumerate}

\vspace{0.5em}
\textbf{二、填空题}
\begin{enumerate}[label=\arabic*.,resume]
\item 磁化强度$\vect{M}$与磁化电流密度的关系：$\vecJ_m = $ \underline{$\curl\vect{M}$}。
\item 相对磁导率$\mu_r$与磁化率的关系：$\mu_r = $ \underline{$1 + \chi_m$}。
\end{enumerate}
\end{miniquiz}
\toggleButton{显示答案}{ocg答案6}

\begin{chapterreview}
\begin{itemize}[nosep]
\item \textbf{核心概念}：磁化强度$\vect{M}$、磁场强度$\vect{H}$、磁化电流、束缚电流。
\item \textbf{材料分类}：抗磁体（$\chi_m<0$）、顺磁体（$\chi_m>0$很小）、铁磁体（$\chi_m\gg 1$，有磁滞）。
\item \textbf{自检验}：能否从$\vect{H}$的定义推导$\curl\vect{H}=\vecJ_f$？
\item \textbf{衔接预告}：第7章将统一电磁学——麦克斯韦方程组登场。
\end{itemize}
\end{chapterreview}

\begin{tip}{本章自检清单}{ch6-checklist}
学完本章后对照下列条目逐条自测，全部通过即可进入下一章；有未通过项时，回到对应小节复习：
\begin{itemize}[label=$\square$, nosep]
\item 由分子环流图像推导 $\vecJ_M = \curl\vect{M}$
\item 区分 $\vecB$、$\vect{H}$、$\vect{M}$ 的源与单位
\item 写出磁介质界面的边界条件
\item 分类抗磁/顺磁/铁磁并解释磁滞
\end{itemize}
\end{tip}


% ---------- 第6章综合练习题 ----------
\section*{第6章综合练习题}
\addcontentsline{toc}{section}{第6章综合练习题}
\begin{enumerate}[label=\textbf{\arabic*.}]
\item \textbf{基础} 一根长直螺线管（单位长度匝数 $n$，电流 $I$）内部充满磁导率为 $\mu_r$ 的磁介质，求内部的 $B$、$H$、$M$。
\item \textbf{基础} 铁磁体（$\mu_r=1000$）与真空分界面，铁磁体内 $\vect{H}_1$ 与法线成 $5°$ 角。求真空侧 $\vecB_2$ 和 $\vect{H}_2$ 的方向。
\item \textbf{提高} 证明：在均匀磁化的介质球中，球内磁场是均匀的，且 $\vecB_{\text{in}}=\frac{2}{3}\muz\vect{M}$、$\vect{H}_{\text{in}}=-\frac{1}{3}\vect{M}$。
\item \textbf{提高} 一个超导体球置于均匀外磁场 $\vect{H}_0$ 中，利用迈斯纳效应（内部 $\vecB=0$）和磁镜像法，求球外磁场分布。
\item \textbf{挑战} 论证：由于不存在磁单极，$\nabla\cdot\vecB=0$ 在任何介质中严格成立，因此磁场的边界条件总是 $B_{1\perp}=B_{2\perp}$，而 $H$ 的切向条件才受表面电流影响。
\end{enumerate}
\newpage
\section*{习题解答}
\addcontentsline{toc}{section}{习题解答}

\textbf{题1解析：}

[分析] 本题是螺线管填充磁介质后的基本场量计算，要求区分 $B$、$H$、$M$ 的物理意义。

[解答] 对理想长直螺线管，安培环路定理给出：$H = nI$（沿轴向）。

$\vecB = \muz\mu_r\vect{H} = \muz\mu_r nI\,\unitvec{z}$。

$\vect{M} = \chi_m\vect{H} = (\mu_r - 1)nI\,\unitvec{z}$。

验证：$\vecB = \muz(\vect{H} + \vect{M}) = \muz(nI + (\mu_r - 1)nI) = \muz\mu_r nI$，一致。\qedsymbol

[量纲校验] $[H] = [A/m]$，$[B] = [MT^{-2}I^{-1}]$，$[\muz H] = [H/m][A/m] = [MT^{-2}I^{-1}]$；$[M] = [A/m]$，一致。\qedsymbol

\textbf{题2解析：}

[分析] 本题利用磁学边界条件求折射角，是磁介质界面行为的重要应用。

[解答] 边界条件：$B_{1n} = B_{2n}$，$H_{1t} = H_{2t}$（无自由面电流）。

$\tan\theta_1 = H_{1t}/H_{1n}$，$\tan\theta_2 = H_{2t}/H_{2n}$。

由 $B_{1n} = B_{2n}$：$\mu_1 H_{1n} = \mu_2 H_{2n}$，$H_{2n} = (\mu_1/\mu_2)H_{1n}$。

由 $H_{1t} = H_{2t}$：$\tan\theta_2 = H_{1t}/((\mu_1/\mu_2)H_{1n}) = (\mu_2/\mu_1)\tan\theta_1$。

代入 $\mu_r = 1000$：$\tan\theta_2 = (1/1000)\tan 5° \approx 0.001 \times 0.0875 = 8.75 \times 10^{-5}$。
$\theta_2 \approx 0.005° \approx 0$。真空侧磁场几乎垂直于表面。\qedsymbol

[量纲校验] $[\tan\theta]$ 无量纲，$(\mu_2/\mu_1)$ 无量纲，一致。\qedsymbol

\textbf{题3解析：}

[分析] 本题推导均匀磁化介质球内部的磁场，需要利用磁标势方法或磁化电流计算。

[解答] 均匀磁化球 $\vect{M} = M\unitvec{z}$。等效磁化电流：体电流 $\vecJ_M = \curl\vect{M} = \vect{0}$（均匀），面电流 $\vect{K}_M = \vect{M}\times\unitvec{n} = M\sin\theta\,\unitvec{\phi}$。

面电流在球内产生的磁场（类比旋转带电球壳）：
\[
\vecB_{\text{in}} = \frac{2}{3}\muz\vect{M}
\]

$\vect{H}_{\text{in}} = \vecB_{\text{in}}/\muz - \vect{M} = \dfrac{2}{3}\vect{M} - \vect{M} = -\dfrac{1}{3}\vect{M}$。

[量纲校验] $[B] = [\muz M] = [H/m][A/m] = [MT^{-2}I^{-1}]$，$[H] = [A/m] = [M]$，一致。\qedsymbol

\textbf{题4解析：}

[分析] 本题利用迈斯纳效应和磁镜像法求解超导体球在外磁场中的磁场分布。

[解答] 迈斯纳效应要求超导体内部 $\vecB = 0$。利用磁镜像法，等效于在球内放置一个磁偶极子镜像，使得球面上 $\vecB_n = 0$。

均匀外场 $\vect{H}_0 = H_0\unitvec{z}$ 对应的磁标势：$\phi_m = -H_0 r\cos\theta$。

为满足球面 $B_r = 0$，需要感应偶极子项：$\phi_m = -H_0\left(r + \dfrac{R^3}{2r^2}\right)\cos\theta$。

球外磁场：$\vecB = \muz\vect{H}_0 - \dfrac{\muz R^3}{2r^3}[3(\vect{H}_0\cdot\unitvec{r})\unitvec{r} - \vect{H}_0]$（感应磁偶极矩 $\vect{m} = -2\pi R^3\vect{H}_0$；截面校验：赤道面上 $r=R$ 处 $B = \frac{3}{2}\muz H_0$）。\qedsymbol

[量纲校验] $[\phi_m] = [A]$，$[Hr] = [A/m][L] = [A]$，一致。\qedsymbol

\textbf{题5解析：}

[分析] 本题从麦克斯韦方程 $\divg\vecB = 0$ 出发，论证磁单极不存在以及法向边界条件。

[解答] 麦克斯韦方程 $\divg\vecB = 0$ 对任意体积 $V$ 都成立：
\[
\int_V \divg\vecB\,\diff V = \oint_S \vecB\cdot\diff\vect{S} = 0
\]

若存在磁单极（磁荷 $\rho_g$），则 $\divg\vecB = \mu_0\rho_g$，高斯定理给出 $\oint_S \vecB\cdot\diff\vect{S} = \mu_0 Q_g$，磁荷可以作为磁场的源。但实验上从未观测到磁单极，因此 $\divg\vecB = 0$ 严格成立。

在两种介质界面上，对药盒形高斯面：
\[
B_{2n}A - B_{1n}A = 0 \quad\Rightarrow\quad B_{1n} = B_{2n}
\]

法向 $\vecB$ 总是连续的，与介质性质无关。切向 $\vect{H}$ 的跳跃仅来自自由面电流 $K_f$。\qedsymbol

[量纲校验] $[\divg B] = [L^{-1}][MT^{-2}I^{-1}] = [ML^{-1}T^{-2}I^{-1}]$，$[\mu_0\rho_g]$ 若存在，$[H/m][C/m^3]$ 需要磁荷量纲匹配，但 $\divg B = 0$ 无需额外量纲。\qedsymbol




% ============================================================

\begin{center}\small \hyperlink{ch5-goal}{上一章：第5章} \quad|\quad \hyperlink{ch7-goal}{下一章：第7章}\end{center}

% 第7章：麦克斯韦方程组（优化版）
% ============================================================
\chapter{麦克斯韦方程组}
\epigraph{光就是电磁场的扰动。}{--- 詹姆斯·克拉克·麦克斯韦}

\hypertarget{ch7-goal}{}\begin{insight}{本章核心目标}{ch7-goal-v2}
将静电学和静磁学的方程推广到随时间变化的一般情形，引入位移电流和法拉第电磁感应定律，完成经典电磁理论的统一框架。理解麦克斯韦方程组的物理内涵、对称性和完备性。
\end{insight}

麦克斯韦方程组是经典物理学的巅峰成就之一。它不仅统一了电学和磁学，预言了电磁波的存在，揭示了光的电磁本质，更为狭义相对论的诞生埋下了种子。本章将系统推导这组方程，并讨论其物理意义。

\begin{tip}{本章新符号}{ch7-new-symbols}
\begin{tabular}{@{}ll@{}}
$\vecJ_D$ & 位移电流密度（$\vecJ_D = \partial\vect{D}/\partial t$） \\
$\vect{S}$ & 坡印廷矢量（能流密度，$\vect{S} = \vecE \times \vect{H}$，单位 \si{W/m^2}） \\
$u$ & 电磁场能量密度（$u = \frac{1}{2}(\eps_0 E^2 + B^2/\muz)$，单位 \si{J/m^3}） \\
$\chi$ & 规范函数（$\vecA' = \vecA + \grad\chi$） \\
\end{tabular}
\end{tip}

\begin{history}{麦克斯韦与电磁理论统一}{maxwell-history}
\textbf{麦克斯韦}（James Clerk Maxwell，1831--1879）在 1865 年的《电磁场的动力学理论》中引入"位移电流"，把电磁学定律统一为自洽的方程组，并算出电磁波速恰等于光速——"光是一种电磁振动"。23 年后\textbf{赫兹}在实验室中产生并检测到电磁波，证实了这一预言。
\end{history}

\section{法拉第电磁感应定律}

\subsection{实验基础}

1831年，法拉第发现：当穿过闭合回路的磁通量发生变化时，回路中会产生感应电动势。实验表明：
\begin{enumerate}
\item 感应电动势的大小与磁通量变化率成正比
\item 感应电流的方向总是阻碍磁通量的变化（楞次定律）
\end{enumerate}

\subsection{积分形式}

\begin{theorem}{法拉第电磁感应定律（积分形式）}{faraday-law-int-v2}
\[
\oint_C \vecE \cdot \diff\vect{l} = -\frac{\diff}{\diff t} \int_S \vecB \cdot \diff\vect{S}
\]
其中 $C$ 是闭合回路，$S$ 是以 $C$ 为边界的任意曲面。
\end{theorem}

\begin{insight}{感应电场的本质}{induced-e-field-v2}
法拉第定律揭示了一种全新的电场——\textbf{感应电场}。与静电场不同，感应电场：
\begin{itemize}
\item 由变化磁场产生，而非静止电荷
\item 是\textbf{有旋场}（$\curl\vecE \neq 0$），不能用电势描述
\item 电场线形成闭合回路，没有起点和终点
\end{itemize}

这意味着电场不再仅仅是保守场——当磁场随时间变化时，电场获得了"涡旋"成分。这是电磁理论从静态到动态的质变。
\end{insight}

\subsection{微分形式}

由斯托克斯定理，左边 $= \int_S (\curl\vecE)\cdot\diff\vect{S}$。若回路静止（无动生电动势）：
\[
\int_S (\curl\vecE)\cdot\diff\vect{S} = -\int_S \frac{\partial\vecB}{\partial t}\cdot\diff\vect{S}
\]

由于 $S$ 任意：

\begin{theorem}{法拉第定律（微分形式）}{faraday-law-diff-v2}
\begin{equation}\curl\vecE = -\frac{\partial\vecB}{\partial t}\label{eq:faraday}\end{equation}
这是麦克斯韦方程组的第三个方程。
\end{theorem}

\subsection{时变电磁势}

法拉第定律的微分形式 $\curl\vecE = -\partial\vecB/\partial t$ 揭示了一个深刻事实：当磁场随时间变化时，电场不再是保守场。为了在这种一般情形下仍然引入标势 $\phi$ 和矢势 $\vecA$，我们需要从磁矢势的定义出发进行系统推导。

由 $\divg\vecB = 0$ 可知磁场可表示为矢势的旋度：
\[
\vecB = \curl\vecA
\]
将这一关系代入法拉第定律：
\[
\curl\vecE = -\frac{\partial}{\partial t}(\curl\vecA) = -\curl\left(\frac{\partial\vecA}{\partial t}\right)
\]
移项得：
\[
\curl\left(\vecE + \frac{\partial\vecA}{\partial t}\right) = \vect{0}
\]
由于旋度为零的矢量场必可表示为某个标量场的梯度，因此存在一个标量场 $\phi$（称为\textbf{标量势}），使得：
\[
\vecE + \frac{\partial\vecA}{\partial t} = -\grad\phi
\]

\begin{theorem}{时变电磁场的标势与矢势}{time-varying-potentials-v2}
在随时间变化的电磁场中，电场和磁场可以用标势 $\phi$ 和矢势 $\vecA$ 表示为：
\begin{align}
\vecB &= \curl\vecA \\
\vecE &= -\grad\phi - \frac{\partial\vecA}{\partial t}
\end{align}
其中 $-\grad\phi$ 是电场的保守成分（由电荷产生），$-\partial\vecA/\partial t$ 是电场的涡旋成分（由变化磁场产生）。
\end{theorem}

\begin{insight}{感应电场的矢势起源}{induced-e-from-A-v2}
时变电磁场中电场的分解 $\vecE = -\grad\phi - \partial\vecA/\partial t$ 具有深刻的物理内涵：
\begin{itemize}
\item \textbf{保守部分} $-\grad\phi$：与静电场相同，由电荷分布决定，满足 $\curl(-\grad\phi) = \vect{0}$
\item \textbf{感应部分} $-\partial\vecA/\partial t$：完全来自矢势的时间变化率，这是法拉第电磁感应的微观本质——变化的磁场通过其矢势的时间导数直接激发涡旋电场
\item 在静电学极限 $\partial\vecA/\partial t = \vect{0}$ 下，自然退化回 $\vecE = -\grad\phi$
\item 这一分解为后续引入规范变换、推导达朗贝尔方程奠定了数学基础
\end{itemize}
\end{insight}

\begin{warning}{静磁学的$\curl\vecE=0$还成立吗？}{curl-e-change}
不成立了！$\curl\vecE = 0$ 仅在静电学（$\partial\vecB/\partial t = 0$）中成立。一般情况下，电场同时具有保守成分（由电荷产生）和涡旋成分（由变化磁场产生）：
\[
\vecE = -\grad\phi - \frac{\partial\vecA}{\partial t}
\]
其中 $-\grad\phi$ 是保守部分，$-\partial\vecA/\partial t$ 是涡旋部分。
\end{warning}

\subsection{规范条件与达朗贝尔方程}

在时变电磁场中，标势 $\phi$ 和矢势 $\vecA$ 的引入虽然恢复了用电势描述电磁场的可能性，但也带来了一个本质性的问题：四个势分量（$\phi$ 和 $\vecA$ 的三个分量）满足的只是三个独立的麦克斯韦方程（$\divg\vecB=0$ 已自动满足），因此存在冗余自由度。为了消除这种冗余，必须对势函数施加额外的约束条件，称为\textbf{规范条件}（gauge condition）。本节讨论两种最重要的规范条件，并在此基础上导出时变电磁势满足的达朗贝尔方程。

\noindent\textbf{本节要点}：洛伦兹规范的定义、库仑规范与洛伦兹规范的对比，以及\textbf{达朗贝尔方程}——这三项务必掌握。其下的逐步推导属补充内容，第一次阅读可先跳过。

\toggleButton{展开详细推导}{ocg推导}
\begin{ocg}{推导}{ocg推导}{off}

\subsubsection{库仑规范与洛伦兹规范}

在静磁学中，我们引入了\textbf{库仑规范}（Coulomb gauge）：
\[
\nabla\cdot\vecA = 0
\]
在此规范下，标势满足泊松方程 $-\lap\phi = \rho/\eps_0$，可以"瞬时"地由电荷分布确定。然而，当电磁场随时间变化时，库仑规范并不能将矢势方程完全解耦，且标势的"瞬时"响应与狭义相对论的光速极限原理相矛盾。

为了建立相对论协变的电磁理论，需要采用另一种规范条件。将 $\vecE = -\grad\phi - \partial\vecA/\partial t$ 代入高斯定理 $\divg\vecE = \rho/\eps_0$，得：
\[
-\lap\phi - \frac{\partial}{\partial t}(\nabla\cdot\vecA) = \frac{\rho}{\eps_0}
\]
若将矢势的散度取为
\[
\nabla\cdot\vecA + \muz\eps_0\frac{\partial\phi}{\partial t} = 0
\]
则上式左边的时间导数项恰好被抵消，标势方程与矢势方程将被完全解耦。这就是\textbf{洛伦兹规范}（Lorenz gauge\footnote{注意两人之别：洛伦兹规范以丹麦物理学家 Ludvig Lorenz（1829--1891）命名，他早于（且独立于）Hendrik Lorentz 提出该规范；而\textbf{洛伦兹变换、洛伦兹力}以荷兰物理学家 Hendrik Lorentz（1853--1928）命名。第8章的"洛伦兹色散模型"（Lorentz model）也是 Hendrik Lorentz 的工作。中文都译作"洛伦兹"，英文文献中 Lorenz 与 Lorentz 是两个人，极易混淆。}）。

\begin{definition}{洛伦兹规范}{lorentz-gauge-def}
\textbf{洛伦兹规范}定义为：
\[
\nabla\cdot\vecA + \muz\eps_0\frac{\partial\phi}{\partial t} = 0
\]
它同时约束了标势和矢势的时间演化，将标势与矢势的耦合方程彻底解耦为两个独立的非齐次波动方程。
\end{definition}
\subsubsection{库仑规范与洛伦兹规范对比}

\begin{center}
\begin{tabular}{lll}
\toprule
\textbf{特性} & \textbf{库仑规范} & \textbf{洛伦兹规范} \\
\midrule
条件 & $\nabla\cdot\vecA=0$ & $\nabla\cdot\vecA+\mu_0\eps_0\partial\phi/\partial t=0$ \\
标势方程 & $-\lap\phi = \rho/\eps_0$（泊松方程） & $\Box\phi = -\rho/\eps_0$（波动方程） \\
矢势方程 & $\lap\vecA-\mu_0\eps_0\partial^2\vecA/\partial t^2 = -\mu_0\vecJ+\mu_0\eps_0\grad(\partial\phi/\partial t)$ & $\Box\vecA = -\mu_0\vecJ$ \\
相对论协变性 & 否（$\phi$瞬时响应） & 是（推迟效应自然包含） \\
适用场景 & 非相对论、静场问题 & 电磁波、相对论框架 \\
\bottomrule
\end{tabular}
\end{center}

\begin{insight}{为什么洛伦兹规范更自然？}{lorentz-natural}
库仑规范中$\phi$由电荷"瞬时"确定（泊松方程），似乎违反相对论的光速极限。但实际上，物理可观测量$\vecE$和$\vecB$仍然满足推迟效应——因为$\vecA$的方程补偿了$\phi$的瞬时性。洛伦兹规范将推迟效应显式地包含在$\phi$和$\vecA$的方程中，使数学结构更简洁，理论更自洽。
\end{insight}


\begin{insight}{两种规范的物理意义}{gauge-meaning-ch7-v2}
库仑规范与洛伦兹规范代表了两种截然不同的物理图像：
\begin{itemize}
\item \textbf{库仑规范}：$\nabla\cdot\vecA=0$ 是"瞬时"的。标势 $\phi$ 由电荷分布通过泊松方程瞬时确定，不包含推迟效应。这种规范在静场问题和非相对论近似下非常方便，但在严格相对论框架中会导致标势与矢势的非协变耦合。
\item \textbf{洛伦兹规范}：$\nabla\cdot\vecA + \muz\eps_0\partial\phi/\partial t = 0$ 是"相对论协变"的。标势和矢势都满足具有相同波动算符的方程，推迟效应被自然地包含在解中。它保证了电磁势的变换规律与狭义相对论完全兼容。
\end{itemize}

值得注意的是，在\textbf{第10章}的四维形式下，洛伦兹规范可以写为极其简洁的协变形式：
\[
\partial_\mu A^\mu = 0
\]
其中 $A^\mu = (\phi/c, \vecA)$ 是四维电磁势。这个简洁的表达式充分揭示了洛伦兹规范的深刻几何意义。
\end{insight}

\subsubsection{达朗贝尔方程的推导}

在洛伦兹规范下，我们可以从麦克斯韦方程组出发，系统地推导标势和矢势所满足的微分方程。

\textbf{第一步：从 $\divg\vecE = \rho/\eps_0$ 推导标势方程}

将 $\vecE = -\grad\phi - \partial\vecA/\partial t$ 代入高斯定理：
\begin{align*}
\divg\left(-\grad\phi - \frac{\partial\vecA}{\partial t}\right) &= \frac{\rho}{\eps_0} \\
-\lap\phi - \frac{\partial}{\partial t}(\nabla\cdot\vecA) &= \frac{\rho}{\eps_0}
\end{align*}
应用洛伦兹规范 $\nabla\cdot\vecA = -\muz\eps_0\,\partial\phi/\partial t$，代入上式：
\begin{align*}
-\lap\phi - \frac{\partial}{\partial t}\left(-\muz\eps_0\frac{\partial\phi}{\partial t}\right) &= \frac{\rho}{\eps_0} \\
\lap\phi - \muz\eps_0\frac{\partial^2\phi}{\partial t^2} &= -\frac{\rho}{\eps_0}
\end{align*}

\textbf{第二步：从 $\curl\vecB = \muz\vecJ + \muz\eps_0\partial\vecE/\partial t$ 推导矢势方程}

由 $\vecB = \curl\vecA$，有 $\curl\vecB = \curl(\curl\vecA) = \grad(\divg\vecA) - \lap\vecA$。将修正的安培-麦克斯韦定律写为：
\begin{align*}
\grad(\nabla\cdot\vecA) - \lap\vecA &= \muz\vecJ + \muz\eps_0\frac{\partial}{\partial t}\left(-\grad\phi - \frac{\partial\vecA}{\partial t}\right) \\
&= \muz\vecJ - \muz\eps_0\grad\frac{\partial\phi}{\partial t} - \muz\eps_0\frac{\partial^2\vecA}{\partial t^2}
\end{align*}
整理并应用洛伦兹规范 $\nabla\cdot\vecA = -\muz\eps_0\,\partial\phi/\partial t$，即 $\grad(\nabla\cdot\vecA) = -\muz\eps_0\grad(\partial\phi/\partial t)$，上式左边第一项恰好与右边第二项抵消：
\begin{align*}
-\lap\vecA + \muz\eps_0\frac{\partial^2\vecA}{\partial t^2} &= \muz\vecJ \\
\lap\vecA - \muz\eps_0\frac{\partial^2\vecA}{\partial t^2} &= -\muz\vecJ
\end{align*}

\begin{theorem}{达朗贝尔方程}{d-alembert-eqs}
\begin{impEqbox}{时变电磁场的核心方程}
在洛伦兹规范下，标势 $\phi$ 和矢势 $\vecA$ 满足\textbf{达朗贝尔方程}（非齐次波动方程）：
\begin{align}
\lap\phi - \muz\eps_0\frac{\partial^2\phi}{\partial t^2} &= -\frac{\rho}{\eps_0} \label{eq:d-alembert-phi} \\
\lap\vecA - \muz\eps_0\frac{\partial^2\vecA}{\partial t^2} &= -\muz\vecJ \label{eq:d-alembert-A}
\end{align}
两式具有相同的微分算符结构——\textbf{达朗贝尔算符} 全书统一取 $\Box \equiv \lap - \muz\eps_0\,\dfrac{\partial^2}{\partial t^2} = -\partial_\mu\partial^\mu$（因 $\muz\eps_0 = 1/c^2$；另一常见写法 $\dfrac{1}{c^2}\dfrac{\partial^2}{\partial t^2} - \lap$ 与本书 $\Box$ \textbf{相差一个整体负号}，见下方提醒），这保证了电磁势在四维时空中具有统一的协变描述。

\textbf{跨教材核对提醒}：本书的 $\Box$ 定义为 $\nabla^2 - \frac{1}{c^2}\partial_t^2 = -\partial_\mu\partial^\mu$；部分教材（如 Jackson 某些章节、朗道《场论》）把达朗贝尔算符定义为 $\frac{1}{c^2}\partial_t^2 - \nabla^2 = \partial_\mu\partial^\mu$——\textbf{两者差一个负号}，相应波动方程右边符号也随之相反。参考外文教材做题时，先核对该书 $\Box$ 的定义，不要直接照搬方程两边的符号。
\end{impEqbox}
\end{theorem}
\end{ocg}

\begin{example}{真空中平面波解的验证}{plane-wave-dalembert-ex}
\difficulty{基础}
在真空中（$\rho = 0$，$\vecJ = \vect{0}$），验证平面波解
\[
\phi = 0, \qquad \vecA = \vecA_0 e^{i(kz-\omega t)}\unitvec{x}
\]
满足洛伦兹规范和达朗贝尔方程。其中 $\vecA_0$ 为常数振幅，$k$ 和 $\omega$ 满足色散关系 $\omega = ck$（$c = 1/\sqrt{\muz\eps_0}$）。
\tcblower
\textbf{解答：}

\textbf{步骤1：验证洛伦兹规范}

由 $\phi = 0$，洛伦兹规范要求 $\nabla\cdot\vecA = 0$。计算：
\[
\nabla\cdot\vecA = \frac{\partial A_x}{\partial x} + \frac{\partial A_y}{\partial y} + \frac{\partial A_z}{\partial z}
\]
由于 $A_x = A_0 e^{i(kz-\omega t)}$，$A_y = A_z = 0$，且 $A_x$ 不依赖于 $x$ 或 $y$，故：
\[
\nabla\cdot\vecA = \frac{\partial A_x}{\partial x} = 0
\]
洛伦兹规范满足。

\textbf{步骤2：验证达朗贝尔方程}

对矢势 $\vecA$ 计算各项：
\begin{align*}
\lap\vecA &= \frac{\partial^2\vecA}{\partial z^2} = -k^2 \vecA_0 e^{i(kz-\omega t)}\unitvec{x} = -k^2\vecA \\
\frac{\partial^2\vecA}{\partial t^2} &= -\omega^2 \vecA_0 e^{i(kz-\omega t)}\unitvec{x} = -\omega^2\vecA
\end{align*}
代入达朗贝尔方程（$\rho=0$ 时齐次方程）：
\begin{align*}
\lap\vecA - \muz\eps_0\frac{\partial^2\vecA}{\partial t^2} &= -k^2\vecA - \muz\eps_0(-\omega^2\vecA) \\
&= (-k^2 + \muz\eps_0\omega^2)\vecA
\end{align*}
由色散关系 $\omega = ck$ 且 $c = 1/\sqrt{\muz\eps_0}$，即 $k^2 = \muz\eps_0\omega^2$，故：
\[
-k^2 + \muz\eps_0\omega^2 = -\muz\eps_0\omega^2 + \muz\eps_0\omega^2 = 0
\]
因此 $\lap\vecA - \muz\eps_0\partial^2\vecA/\partial t^2 = \vect{0}$，矢势的达朗贝尔方程满足。

标势 $\phi = 0$ 在 $\rho = 0$ 时自然满足齐次达朗贝尔方程。证毕。
\end{example}

\section{位移电流}

\subsection{安培定理的危机}

静磁学中的安培定理 $\curl\vecB = \muz\vecJ$ 在时变情形下遇到严重问题。对其取散度：
\[
\divg(\curl\vecB) = 0 = \muz\divg\vecJ
\]

但电荷守恒要求连续性方程：
\[
\divg\vecJ + \frac{\partial\rho}{\partial t} = 0
\]

若 $\partial\rho/\partial t \neq 0$（时变情形），则 $\divg\vecJ \neq 0$，与安培定理矛盾！

\subsection{麦克斯韦的修正}

麦克斯韦的天才之处在于：他注意到高斯定理 $\divg\vecE = \rho/\eps_0$ 意味着：
\[
\divg\vecJ = -\frac{\partial\rho}{\partial t} = -\eps_0\frac{\partial}{\partial t}(\divg\vecE) = -\divg\left(\eps_0\frac{\partial\vecE}{\partial t}\right)
\]

因此，若将安培定理修正为：
\[
\curl\vecB = \muz\vecJ + \muz\eps_0\frac{\partial\vecE}{\partial t}
\]

则取散度：
\[
0 = \muz\divg\vecJ + \muz\eps_0\frac{\partial}{\partial t}(\divg\vecE) = \muz\left(\divg\vecJ + \frac{\partial\rho}{\partial t}\right) = 0
\]

与电荷守恒自洽！


\begin{insight}{电容器困境：为什么需要位移电流？}{capacitor-dilemma}
想象一个正在充电的电容器：导线中有电流$I$流入上极板，导线中有电流$I$流出下极板，但\textbf{两极板之间是真空/电介质，没有传导电流}。

现在，用安培环路定理 $\oint\vecB\cdot\diff\vect{l}=\mu_0 I_{enc}$ 计算电容器极板之间的磁场。问题来了：
\begin{itemize}
\item 如果取一个穿过导线的高斯面（面1）：$I_{enc}=I$，磁场不为零。
\item 如果取一个穿过极板之间的高斯面（面2）：$I_{enc}=0$，安培定理说磁场为零。
\end{itemize}

但同一个环路，磁场应该只有一个值！\textbf{矛盾出现了}。

麦克斯韦的天才在于：他注意到极板之间的电场在变化（充电时$\vecE$增大）。如果把 $\eps_0\partial\vecE/\partial t$ 也当作一种"电流"加到安培定理中，两个面的结果就一致了。

这就是\textbf{位移电流}的物理起源——它不是真的电荷在流动，而是变化的电场"等效地"产生了磁场。
\end{insight}

\begin{definition}{位移电流}{displacement-current-v2}
\textbf{位移电流密度}定义为：
\[
\vecJ_D = \eps_0\frac{\partial\vecE}{\partial t}
\]
它不是真实的电荷流动，而是变化电场等效的"电流"。
\end{definition}

\begin{insight}{位移电流的物理意义}{displacement-meaning-v2}
位移电流的引入是麦克斯韦最深刻的洞见之一。它表明：
\begin{itemize}
\item 变化的电场像电流一样产生磁场
\item 在电容器充电过程中，虽然传导电流在极板间中断了，但位移电流"接上"了——保证了电流的连续性
\item 电场和磁场可以\textbf{相互激发}，形成自维持的波动——这就是\textbf{电磁波}
\end{itemize}

从对称性角度看，法拉第定律说"变化磁场产生电场"，位移电流则说"变化电场产生磁场"。这种对称性是电磁波存在的基础。
\end{insight}

\begin{example}{平行板电容器中的位移电流}{capacitor-displacement-ex-v2}
\difficulty{基础}
圆形平行板电容器，板半径 $R$，板间距 $d \ll R$，正在充电，传导电流为 $I$。求：(a) 极板间的位移电流；(b) 距轴线 $r$ 处的磁场。
\tcblower
\textbf{解答：}

(a) 极板间电场：$E = Q/(\eps_0\pi R^2)$（均匀，边缘效应忽略）

位移电流密度：$j_D = \eps_0 \partial E/\partial t = (1/\pi R^2)\, \diff Q/\diff t = I/(\pi R^2)$

总位移电流：$I_D = j_D \cdot \pi R^2 = I$——等于传导电流！

(b) 由修正的安培定理（轴对称性），取半径为 $r$ 的圆形回路：

$\oint\vecB\cdot\diff\vect{l} = \muz(I_{\text{enc}} + I_{D,\text{enc}})$

极板间无传导电流：$I_{\text{enc}} = 0$

$I_{D,\text{enc}} = j_D \cdot \pi r^2 = I r^2/R^2$

$B \cdot 2\pi r = \muz I r^2/R^2$

$B = \dfrac{\muz I r}{2\pi R^2}$（$r < R$）

在 $r = R$ 处，$B = \muz I/(2\pi R)$，与导线中的磁场连续衔接。
\end{example}

% ---------- 位移电流示意图 ----------
\begin{center}
\begin{tikzpicture}[scale=1.2]
  % Capacitor plates（左正右负）
  \fill[blue!10] (-2,0) rectangle (-1.8,3);
  \fill[red!10] (1.8,0) rectangle (2,3);
  \node[blue!70!black,anchor=east] at (-2.15,2.7) {$+Q$};
  \node[red!70!black,anchor=west] at (2.15,2.7) {$-Q$};

  % Wires and current（导线在极板外侧，电流标注在箭头上方错开）
  \draw[thick] (-1.9,4.2) -- (-1.9,3);
  \draw[->, thick] (-1.9,3.9) -- (-1.9,3.4);
  \node at (-1.55,4.0) {$I$};
  \draw[thick] (1.9,3) -- (1.9,4.2);
  \draw[->, thick] (1.9,3.4) -- (1.9,3.9);
  \node at (2.25,4.0) {$I$};

  % Electric field between plates（红箭头，避开放大区）
  \foreach \y in {0.4,0.8,2.2,2.6} {
    \draw[->, red, thick] (-1.5,\y) -- (1.5,\y);
  }
  \node[red] at (0,2.85) {$\vecE$};

  % Displacement current（橙色虚线，单独一行居中，上下留白）
  \draw[->, orange, very thick, dashed] (-1.2,1.5) -- (1.2,1.5);
  \node[orange,anchor=west] at (1.35,1.5) {$\vecJ_D$};
  \node[font=\small,orange!80!black] at (0,1.05) {$\eps_0\dfrac{\partial\vecE}{\partial t}$};

  % Caption
  \node at (0,-0.6) {\small 电容器极板间：位移电流接上传导电流的断裂};
\end{tikzpicture}
\end{center}


\section{麦克斯韦方程组的完整形式}

\begin{theorem}{麦克斯韦方程组}{maxwell-eqs-v2}
\begin{impEqbox}{全书最重要的方程（SI 单位制）}
\textbf{积分形式}：
\begin{align}
\oint_S \vect{D} \cdot \diff\vect{S} &= Q_{f,\text{enc}} \label{eq:M1} \\
\oint_C \vecE \cdot \diff\vect{l} &= -\frac{\diff}{\diff t}\int_S \vecB\cdot\diff\vect{S} \label{eq:M2} \\
\oint_S \vecB \cdot \diff\vect{S} &= 0 \label{eq:M3} \\
\oint_C \vect{H} \cdot \diff\vect{l} &= I_{f,\text{enc}} + \frac{\diff}{\diff t}\int_S \vect{D}\cdot\diff\vect{S} \label{eq:M4}
\end{align}

\textbf{微分形式}：
\begin{align}
\divg\vect{D} &= \rho_f \label{eq:m1} \\
\curl\vecE &= -\frac{\partial\vecB}{\partial t} \label{eq:m2} \\
\divg\vecB &= 0 \label{eq:m3} \\
\curl\vect{H} &= \vecJ_f + \frac{\partial\vect{D}}{\partial t} \label{eq:m4}
\end{align}
\end{impEqbox}
\end{theorem}

\begin{insight}{麦克斯韦方程组的对称性与不对称性}{maxwell-symmetry-v2}
\textbf{对称性}：
\begin{itemize}
\item 方程\eqref{eq:m2}和\eqref{eq:m4}形成优美的对称：$\curl\vecE \propto -\partial\vecB/\partial t$，$\curl\vect{H} \propto +\partial\vect{D}/\partial t$（加上源项）
\item 这种对称暗示了电场和磁场的统一性
\end{itemize}

\textbf{不对称性}：
\begin{itemize}
\item $\divg\vect{D} = \rho_f$（有电荷源），但 $\divg\vecB = 0$（无磁荷源）
\item 如果磁单极子存在，方程将完全对称：$\divg\vecB = \mu_0\rho_m$，$\curl\vecE = -\partial\vecB/\partial t - \mu_0\vecJ_m$
\item 狄拉克（1931）从量子力学角度论证了磁单极子的存在可能性，但实验上至今未找到确凿证据
\end{itemize}
\end{insight}

\begin{history}{磁单极子的当代实验搜寻}{monopole-search-v2}
磁单极子是粒子物理学和宇宙学中最重要的未解之谜之一。当代实验对磁单极子施加了极为严格的限制：

\begin{itemize}
\item \textbf{加速器实验}：大型强子对撞机（LHC）的MoEDAL实验专门搜寻磁单极子，目前对质量低于约 \SI{3}{\tera\eV} 的磁单极子已排除一定耦合强度范围
\item \textbf{宇宙线实验}：南极冰立方（IceCube）和 MACRO 地下实验通过寻找穿过地球的磁单极子轨迹来限制其流强，当前上限约为 $10^{-16}\,\si{cm^{-2}s^{-1}sr^{-1}}$（对相对论性磁单极子）
\item \textbf{天体物理约束}：如果磁单极子大量存在，它们会催化质子衰变并显著影响恒星冷却速率。观测到的中子星表面磁场与标准磁流体动力学一致，不支持大量磁单极子的存在
\item \textbf{理论动机}：大统一理论（GUT）预言磁单极子质量约为 $10^{16}\,\si{GeV}$，远超当前加速器能量；宇宙暴胀可能稀释了原初磁单极子密度，解释了为什么我们在宇宙中极少看到它们
\end{itemize}

尽管实验上尚未发现磁单极子，但它仍是连接经典电动力学、量子力学和粒子物理前沿的重要概念。
\end{history}

\section{电磁场能量与坡印廷矢量}

\subsection{能量密度}

电磁场携带能量。能量密度为：
\[
u = \frac{1}{2}\vecE\cdot\vect{D} + \frac{1}{2}\vecB\cdot\vect{H} = \frac{1}{2}\eps E^2 + \frac{1}{2\mu}B^2
\]
（线性各向同性介质）

\subsection{坡印廷矢量与能量守恒}

\begin{theorem}{坡印廷定理}{poynting-theorem-v2}
定义\textbf{坡印廷矢量}：
\[
\vect{S} = \vecE \times \vect{H}
\]
它表示单位面积上的电磁能流（功率/面积）。

能量守恒方程：
\[
\frac{\partial u}{\partial t} + \divg\vect{S} = -\vecJ_f\cdot\vecE
\]
右边 $-\vecJ_f\cdot\vecE$ 是电磁场对自由电荷做的功（焦耳热等）。
\end{theorem}

\begin{example}{同轴电缆中的能量传输}{coaxial-poynting-ex-v2}
\difficulty{提高}
同轴电缆内导体半径为 $a$，外导体内半径为 $b$，中间填充空气。电缆传输直流功率，内导体上电流为 $I$（沿 $+z$ 方向），内外导体间电压为 $V$。证明能量主要在导体之间的空间中传播，而非在导线内部。
\tcblower
\textbf{解答：}

\textbf{步骤1：分析场分布}

由对称性，电场沿径向，磁场沿角向。设内导体单位长度带电荷 $\lambda$，则内外导体间电场由高斯定理：
\[
E \cdot 2\pi r \cdot l = \frac{\lambda l}{\eps_0} \quad\Rightarrow\quad E = \frac{\lambda}{2\pi\eps_0 r} = \frac{V}{r\ln(b/a)} \quad (a < r < b)
\]

由安培环路定理，磁场：
\[
B \cdot 2\pi r = \muz I \quad\Rightarrow\quad B = \frac{\muz I}{2\pi r} \quad (a < r < b)
\]

\textbf{步骤2：计算坡印廷矢量}

坡印廷矢量 $\vect{S} = \vecE \times \vect{H} = \dfrac{1}{\muz}\vecE \times \vecB$：

$\vecE$ 沿径向向外，$\vecB$ 沿角向（$\hat{\phi}$ 方向），因此 $\vect{S}$ 沿 $+z$ 方向：
\[
S = \frac{EB}{\muz} = \frac{V}{r\ln(b/a)} \cdot \frac{I}{2\pi r} = \frac{VI}{2\pi r^2\ln(b/a)}
\]

\textbf{步骤3：验证总功率}

总功率通过对两导体间截面的坡印廷矢量积分：
\[
P = \int_a^b S \cdot 2\pi r\,\diff r = \int_a^b \frac{VI}{2\pi r^2\ln(b/a)} \cdot 2\pi r\,\diff r = \frac{VI}{\ln(b/a)} \int_a^b \frac{\diff r}{r} = VI
\]

这正是电路理论中的 $P = VI$！

\textbf{步骤4：导线内部的能流}

在导体内部（良导体），电场 $E_{\text{内}} = j/\sigma = I/(\pi a^2\sigma)$（沿 $z$ 方向），而表面磁场 $B = \muz I/(2\pi a)$。坡印廷矢量指向导体内部（径向向内），表示能量从电磁场流入导体，以焦耳热的形式耗散：
\[
S_{\text{内}} = E_{\text{内}} \cdot H = \frac{I}{\pi a^2\sigma} \cdot \frac{I}{2\pi a} = \frac{I^2}{2\pi^2 a^3\sigma}
\]
其量级远小于导体间的坡印廷矢量，且方向是径向的（不是沿传输方向）。

\textbf{结论}：能量主要在导体之间的空间中沿 $+z$ 方向传播，导线本身只是"引导"电磁场的边界条件。这一结论彻底颠覆了"电流携带能量"的直观图像——在电磁学中，能量存储在电磁场中，通过坡印廷矢量传输。
\end{example}

\begin{insight}{坡印廷矢量的物理意义}{poynting-meaning-v2}
坡印廷矢量 $\vect{S}$ 的方向表示能量流动的方向，大小表示能流密度。例如：
\begin{itemize}
\item 平行板电容器充电时，能量不是沿导线"流入"电容器，而是从电容器侧面（导线与极板之间的空间）以坡印廷矢量的形式流入
\item 直流电路中，能量主要在导线\textbf{外部}的电磁场中传输，导线只是"引导"了场的分布
\item 电磁波中，$\vect{S}$ 指向波的传播方向
\end{itemize}
\end{insight}

\subsection{电磁场动量与麦克斯韦应力张量}

坡印廷矢量描述了电磁场的能量流，而电磁场同样携带动量。为了完整地描述电磁场的力学性质，需要引入动量密度和麦克斯韦应力张量。

\subsubsection{动量密度}

从相对论的能量-动量关系可以预期，电磁场的能量流 $\vect{S}$ 与其动量密度之间存在简单的比例关系。定义\textbf{电磁场动量密度}：
\[
\vect{g} = \muz\eps_0\vect{S} = \muz\eps_0\vecE\times\vect{H}
\]
在真空中，利用 $\vecB = \muz\vect{H}$，也可写为 $\vect{g} = \eps_0\vecE\times\vecB$。

\subsubsection{麦克斯韦应力张量}

电磁场对带电体施加洛伦兹力，其力密度为 $\vect{f} = \rho\vecE + \vecJ\times\vecB$。从洛伦兹力密度出发，结合麦克斯韦方程组，可以导出动量守恒方程。详细推导如下：

将电荷密度和电流密度用场量表示：
\begin{align*}
\rho &= \eps_0\nabla\cdot\vecE \\
\vecJ &= \frac{1}{\muz}\curl\vecB - \eps_0\frac{\partial\vecE}{\partial t}
\end{align*}
代入洛伦兹力密度：
\begin{align*}
f_i &= \rho E_i + (\vecJ\times\vecB)_i \\
&= \eps_0(\nabla\cdot\vecE)E_i + \left[\left(\frac{1}{\muz}\curl\vecB - \eps_0\frac{\partial\vecE}{\partial t}\right)\times\vecB\right]_i
\end{align*}

经过矢量恒等式的系统整理（利用 $\divg\vecB=0$ 和法拉第定律），可以得到：
\begin{theorem}{电磁场动量守恒方程}{maxwell-stress-theorem}
电磁场与带电物质系统的总动量守恒可表述为
\[
\frac{\partial g_i}{\partial t} + \partial_j T_{ij} = -f_i
\]
其中 $\vect{g} = \muz\eps_0\vecE\times\vect{H}$ 是电磁场动量密度，$T_{ij}$ 是\textbf{麦克斯韦应力张量}：
\[
T_{ij} = \eps_0\left(E_i E_j - \frac{1}{2}\delta_{ij}E^2\right) + \frac{1}{\muz}\left(B_i B_j - \frac{1}{2}\delta_{ij}B^2\right)
\]
\end{theorem}

\begin{insight}{麦克斯韦应力张量的物理意义}{stress-tensor-meaning}
麦克斯韦应力张量 $T_{ij}$ 是第1章中介绍的并矢（dyad）概念在电磁学中的核心应用。$T_{ij}$ 表示电磁场通过单位面积表面传递的动量流，即电磁场对物体表面的\textbf{应力}：
\begin{itemize}
\item 对角元 $T_{ii}$ 表示沿 $i$ 方向的法向应力（压强或张力），由电场和磁场的"自压"与"背景压"组成。
\item 非对角元 $T_{ij}$（$i\neq j$）表示切向应力，描述电磁场沿 $j$ 方向通过垂直于 $i$ 轴的表面所传递的动量流。
\item 在真空中平面电磁波的情形下，$T_{ij}$ 具有特别简单的形式：沿传播方向存在辐射压强，而垂直于传播方向的应力分量为零。
\end{itemize}
麦克斯韦应力张量使我们能够直接计算电磁场对宏观物体施加的辐射压力和辐射力，是激光捕获、太阳帆等应用的理论基础。
\end{insight}

\begin{example}{平面电磁波的麦克斯韦应力张量}{plane-wave-stress-ex}
\difficulty{提高}
考虑沿 $+z$ 方向传播的线偏振平面电磁波，电场沿 $x$ 方向：$\vecE = E_0\cos(kz-\omega t)\unitvec{x}$，磁场沿 $y$ 方向：$\vecB = B_0\cos(kz-\omega t)\unitvec{y}$，其中 $B_0 = E_0/c$。求该电磁波的麦克斯韦应力张量。
\tcblower
\textbf{解答：}

\textbf{步骤1：计算各场量分量}

场的非零分量为 $E_x = E_0\cos(kz-\omega t)$，$B_y = B_0\cos(kz-\omega t)$，其余分量为零。且有 $E^2 = E_x^2$，$B^2 = B_y^2$。

\textbf{步骤2：计算电场对应部分}

对 $i=j=x$：
\[
T_{xx}^{(E)} = \eps_0\left(E_x^2 - \frac{1}{2}E^2\right) = \eps_0\left(E_x^2 - \frac{1}{2}E_x^2\right) = \frac{1}{2}\eps_0 E_x^2
\]

对 $i=j=y$：$E_y=0$，故 $T_{yy}^{(E)} = -\frac{1}{2}\eps_0 E_x^2$。

对 $i=j=z$：$E_z=0$，故 $T_{zz}^{(E)} = -\frac{1}{2}\eps_0 E_x^2$。

非对角元：由于只有一个电场分量 $E_x$ 非零，所有 $i\neq j$ 的 $T_{ij}^{(E)} = 0$。

\textbf{步骤3：计算磁场对应部分}

类似地，对 $i=j=y$：
\[
T_{yy}^{(B)} = \frac{1}{\muz}\left(B_y^2 - \frac{1}{2}B^2\right) = \frac{1}{2\muz}B_y^2
\]

对 $i=j=x$：$B_x=0$，故 $T_{xx}^{(B)} = -\frac{1}{2\muz}B_y^2$。

对 $i=j=z$：$B_z=0$，故 $T_{zz}^{(B)} = -\frac{1}{2\muz}B_y^2$。

非对角元：$T_{ij}^{(B)} = 0$（$i\neq j$）。

\textbf{步骤4：合成总应力张量}

利用 $B_0 = E_0/c$ 和 $c^2 = 1/(\muz\eps_0)$，有 $B_0^2/\muz = \eps_0 E_0^2$。因此：
\begin{align*}
T_{xx} &= T_{xx}^{(E)} + T_{xx}^{(B)} = \frac{1}{2}\eps_0 E_x^2 - \frac{1}{2\muz}B_y^2 = 0 \\
T_{yy} &= T_{yy}^{(E)} + T_{yy}^{(B)} = -\frac{1}{2}\eps_0 E_x^2 + \frac{1}{2\muz}B_y^2 = 0 \\
T_{zz} &= T_{zz}^{(E)} + T_{zz}^{(B)} = -\frac{1}{2}\eps_0 E_x^2 - \frac{1}{2\muz}B_y^2 = -\eps_0 E_x^2 = -\frac{1}{\muz}B_y^2
\end{align*}

所有非对角元均为零。取时间平均，利用 $\langle\cos^2(kz-\omega t)\rangle = 1/2$，得：
\[
\langle T_{zz}\rangle = -\frac{1}{2}\eps_0 E_0^2 = -\langle u\rangle
\]
其中 $\langle u\rangle = \frac{1}{2}\eps_0 E_0^2$ 是\textbf{时间平均}能量密度。平面波的能量密度瞬时值 $u = \eps_0 E^2$（电场 $\frac12\eps_0 E^2$ 与磁场 $\frac{B^2}{2\muz} = \frac{\eps_0 E^2}{2}$ 之和，对行波恒成立），时间平均 $\langle u\rangle = \eps_0\langle E^2\rangle = \frac12\eps_0 E_0^2$——与 $\langle T_{zz}\rangle$ 恰好差一个负号，即时间平均的动量流大小等于能量密度。负号表示沿传播方向 $+z$ 存在\textbf{辐射压强}（压力），这正是光压的微观起源。
\end{example}

\begin{review}{第7章要点回顾}{ch7-review}
\begin{center}
\begin{tabular}{@{}ll@{}}
\toprule
\textbf{方程} & \textbf{名称/含义} \\
\midrule
$\divg\vect{D} = \rho_f$ & 高斯定理（电场有源） \\
$\curl\vecE = -\partial\vecB/\partial t$ & 法拉第定律（变化磁场生电场） \\
$\divg\vecB = 0$ & 磁场无源（无磁单极子） \\
$\curl\vect{H} = \vecJ_f + \partial\vect{D}/\partial t$ & 安培-麦克斯韦定律 \\
$\vect{S} = \vecE \times \vect{H}$ & 坡印廷矢量（能流密度） \\
\bottomrule
\end{tabular}
\end{center}
\end{review}

\begin{tip}{本章核心框架}{ch7-summary}
\textbf{一句话总结}：麦克斯韦方程组统一了电磁学，预言了电磁波，奠定了狭义相对论的基础。\\
\textbf{核心方程}：四个方程（高斯、法拉第、安培-麦克斯韦、无磁单极），两个势方程（达朗贝尔方程），能量守恒（坡印廷定理）。\\
\textbf{核心方法}：规范变换（库仑规范vs洛伦兹规范）、推迟势（$t_r=t-R/c$）。\\
\textbf{衔接}：第8章将研究电磁波在介质中的传播——反射、折射、波导、色散。
\end{tip}

% ---------- 第7章随堂自测 ----------
\section*{第7章随堂自测}
\addcontentsline{toc}{section}{第7章随堂自测}

\begin{miniquiz}
\textbf{一、选择题}
\begin{enumerate}[label=\arabic*.]
\item 麦克斯韦引入位移电流 $\vecJ_D = \eps_0\partial\vecE/\partial t$ 的目的是：
  \begin{enumerate}[label=(\Alph*),nosep]
  \item 解释电磁感应
  \item 修正安培定理，使其在非恒定情况下也成立
  \item 引入电磁波
  \item 统一电场和磁场
  \end{enumerate}
\begin{ocg}{答案}{ocg答案7}{off}
  \textit{答案：B。位移电流保证了电荷守恒在非恒定情况下仍然成立。}
\end{ocg}

\item 洛伦兹规范的条件是：
  \begin{enumerate}[label=(\Alph*),nosep]
  \item $\divg\vecA = 0$
  \item $\divg\vecA + \muz\eps_0\partial\phi/\partial t = 0$
  \item $\phi = 0$
  \item $\vecA = 0$
  \end{enumerate}
\begin{ocg}{答案}{ocg答案7}{off}
  \textit{答案：B。洛伦兹规范使$\phi$和$\vecA$满足协变的波动方程。}
\end{ocg}

\item 真空中的电磁波速度 $c$ 满足：
  \begin{enumerate}[label=(\Alph*),nosep]
  \item $c = 1/\sqrt{\muz\eps_0}$
  \item $c = \sqrt{\muz/\eps_0}$
  \item $c = \muz\eps_0$
  \item $c = 1/(\muz\eps_0)$
  \end{enumerate}
\begin{ocg}{答案}{ocg答案7}{off}
  \textit{答案：A。这是麦克斯韦理论最惊人的预言之一。}
\end{ocg}
\end{enumerate}

\vspace{0.5em}
\textbf{二、判断题}
\begin{enumerate}[label=\arabic*.,resume]
\item 真空中的麦克斯韦方程组是线性的，满足叠加原理。
\begin{ocg}{答案}{ocg答案7}{off}
  \textit{答案：正确。线性是波动方程可解的关键。}
\end{ocg}

\item 库仑规范在相对论框架下比洛伦兹规范更自然。
\begin{ocg}{答案}{ocg答案7}{off}
  \textit{答案：错误。洛伦兹规范具有相对论协变性，在相对论框架下更自然。}
\end{ocg}
\end{enumerate}
\end{miniquiz}
\toggleButton{显示答案}{ocg答案7}

\begin{chapterreview}
\begin{itemize}[nosep]
\item \textbf{麦克斯韦方程组}（微分形式）：\begin{align*}\divg\vecE&=\rho/\eps_0 & \curl\vecE&=-\partial\vecB/\partial t\\ \divg\vecB&=0 & \curl\vecB&=\muz\vecJ+\muz\eps_0\partial\vecE/\partial t\end{align*}
\item \textbf{核心推导}：位移电流的引入、能量守恒（坡印廷定理）、洛伦兹规范下达朗贝尔方程。
\item \textbf{自检验}：能否从麦克斯韦方程组导出真空中$\vecE$和$\vecB$的波动方程？
\item \textbf{衔接预告}：第8章将研究电磁波在介质中的传播——反射、折射、波导。
\end{itemize}
\end{chapterreview}

\begin{tip}{本章自检清单}{ch7-checklist}
学完本章后对照下列条目逐条自测，全部通过即可进入下一章；有未通过项时，回到对应小节复习：
\begin{itemize}[label=$\square$, nosep]
\item 完整写出麦克斯韦方程组并说明每条的实验来源
\item 解释位移电流为什么必须引入
\item 用坡印廷矢量验证能量守恒
\item 推导达朗贝尔方程并说明规范条件的作用
\end{itemize}
\end{tip}

% ---------- 第7章综合练习题 ----------
\section*{第7章综合练习题}
\addcontentsline{toc}{section}{第7章综合练习题}
\begin{enumerate}[label=\textbf{\arabic*.}]
\item \textbf{基础} 从法拉第定律微分形式 $\curl\vecE=-\partial\vecB/\partial t$ 和 $\vecB=\curl\vecA$ 出发，完整推导 $\vecE=-\nabla\phi-\partial\vecA/\partial t$。
\item \textbf{基础} 同轴电缆内导体半径 $a$，外导体内半径 $b$，通直流电流 $I$。求导体间坡印廷矢量 $\vect{S}$，并证明传输功率 $P=VI$ 与电路理论一致。
\item \textbf{提高} 证明：在洛伦兹规范 $\nabla\cdot\vecA+\muz\eps_0\partial\phi/\partial t=0$ 下，达朗贝尔方程可简化为非齐次波动方程。
\item \textbf{提高} 一平行板电容器（极板面积 $A$，间距 $d$）接交流电源 $V=V_0\cos\omega t$。求位移电流、磁场分布和流入电容器的坡印廷矢量功率。
\item \textbf{挑战} 假设磁单极 $g$ 存在，写出修改后的麦克斯韦方程组，并证明其与电荷守恒定律的一致性需要磁荷守恒 $\partial\rho_g/\partial t+\nabla\cdot\vecJ_g=0$。
\end{enumerate}
\newpage
\section*{习题解答}
\addcontentsline{toc}{section}{习题解答}

\textbf{题1解析：}

[分析] 本题从法拉第定律和磁矢势定义出发，严格推导时变电磁场中标量势与矢势的表达式。

[解答] 已知 $\curl\vecE = -\partial\vecB/\partial t$ 和 $\vecB = \curl\vecA$。

代入：$\curl\vecE = -\partial(\curl\vecA)/\partial t = -\curl(\partial\vecA/\partial t)$。

移项：$\curl(\vecE + \partial\vecA/\partial t) = \vect{0}$。

旋度为零的矢量场必可表示为某个标量场的梯度：
\[
\vecE + \frac{\partial\vecA}{\partial t} = -\grad\phi
\]

即：$\vecE = -\grad\phi - \partial\vecA/\partial t$。\qedsymbol

[量纲校验] $[\partial\vecA/\partial t] = [MLT^{-2}I^{-1}]/[T] = [MLT^{-3}I^{-1}] = [E]$；

$[\grad\phi] = [V/m] = [MLT^{-3}I^{-1}]$，两者一致。\qedsymbol

\textbf{题2解析：}

[分析] 本题计算同轴电缆的坡印廷矢量，证明电磁场能量传输的物理图像。

[解答] 由高斯定理：$E = \lambda/(2\pi\eps_0\rho)$（$a < \rho < b$）。由安培定理：$B = \muz I/(2\pi\rho)$。

坡印廷矢量 $\vect{S} = \vecE\times\vect{H} = \dfrac{1}{\muz}\vecE\times\vecB$：
$\vecE$ 沿径向，$\vecB$ 沿 $\unitvec{\phi}$，故 $\vect{S}$ 沿 $+z$ 方向。

大小：$S = EB/\muz = \dfrac{\lambda I}{4\pi^2\eps_0\rho^2}$。

传输功率：$P = \int_a^b S \cdot 2\pi\rho\,\diff\rho = \dfrac{\lambda I}{2\pi\eps_0}\int_a^b \dfrac{\diff\rho}{\rho} = \dfrac{\lambda I}{2\pi\eps_0}\ln\dfrac{b}{a}$。

又 $V = \int_a^b E\,\diff\rho = \dfrac{\lambda}{2\pi\eps_0}\ln\dfrac{b}{a}$，故 $P = VI$。\qedsymbol

[量纲校验] $[S] = [V/m][A/m] = [VA/m^2] = [W/m^2]$，$[P] = [VI] = [W]$，一致。\qedsymbol

\textbf{题3解析：}

[分析] 本题证明洛伦兹规范下达朗贝尔方程的简化，是电磁势理论的关键步骤。

[解答] 达朗贝尔方程（一般形式）：
\begin{align*}
\lap\vecA - \muz\eps_0\partial_t^2\vecA &= -\muz\vecJ + \grad(\divg\vecA + \muz\eps_0\partial_t\phi) \\
\lap\phi - \muz\eps_0\partial_t^2\phi &= -\frac{\rho}{\eps_0} - \partial_t(\divg\vecA + \muz\eps_0\partial_t\phi)
\end{align*}

洛伦兹规范：$\divg\vecA + \muz\eps_0\partial_t\phi = 0$。

代入后耦合项消失：
\begin{align*}
\lap\vecA - \frac{1}{c^2}\partial_t^2\vecA &= -\muz\vecJ \\
\lap\phi - \frac{1}{c^2}\partial_t^2\phi &= -\frac{\rho}{\eps_0}
\end{align*}

两个独立的非齐次波动方程。\qedsymbol

[量纲校验] $[\lap A] = [L^{-2}A]$，$[\muz j] = [H/m][A/L^2] = [L^{-2}A]$，一致。\qedsymbol

\textbf{题4解析：}

[分析] 本题计算平行板电容器接交流电源时的位移电流和磁场，是理解时变电磁场耦合的经典例题。

[解答] 极板间电场：$E = V/d = V_0\cos\omega t/d$。

位移电流密度：
\[
j_D = \eps_0\partial E/\partial t = -\frac{\eps_0\omega V_0\sin\omega t}{d}
\]
总位移电流：
\[
I_D = j_D A = -\frac{\eps_0 A\omega V_0\sin\omega t}{d}
\]

由修正安培定理（轴对称），取半径为 $\rho$ 的圆形回路：
\[
B \cdot 2\pi\rho = \muz I_{D,\text{enc}} = \muz j_D \pi\rho^2
\]
$I_{D,\text{enc}}$ 是位移电流在半径 $\rho$ 内的部分（$\rho < R$ 时，$R$ 为极板半径）。

$B = -\muz\eps_0\omega V_0\rho\sin\omega t/(2d)$（$\rho < R$；负号由 $j_D$ 的负号带入，它决定能流方向，不可略去）。

坡印廷矢量（径向分量）：$S = EB/\muz = \eps_0\omega V_0^2\rho\cos\omega t\sin\omega t/(2d^2)$——正值表示此刻能流由极板间沿 $+\unitvec{\rho}$ 向外流出。
穿出柱侧面的功率（侧面积 $2\pi R d$）：$P = S(R)\cdot 2\pi R d = \dfrac{\pi\eps_0\omega V_0^2R^2\sin\omega t\cos\omega t}{d}$。\qedsymbol

\textbf{能量守恒校验}：电容器储能
\[
U = \frac{1}{2}CV^2 = \dfrac{\eps_0\pi R^2V_0^2\cos^2\omega t}{2d}
\]
故 $\dfrac{\diff U}{\diff t} = -\dfrac{\pi\eps_0\omega R^2V_0^2\sin\omega t\cos\omega t}{d} = -P$，即流出的功率恰等于储能的减少率。

[量纲校验] $[B] = [\muz\eps_0\omega V\rho/d] = [H/m][F/m][T^{-1}][V][L]/[L] = [MT^{-2}I^{-1}]$，一致。\qedsymbol

\textbf{题5解析：}

[分析] 本题假设磁单极存在，修改麦克斯韦方程组，并讨论磁荷守恒的必要性。

[解答] 修改后的麦克斯韦方程组：
\begin{align*}
\divg\vecE &= \frac{\rho}{\eps_0} \\
\divg\vecB &= \mu_0\rho_g \\
\curl\vecE &= -\frac{\partial\vecB}{\partial t} - \mu_0\vecJ_g \\
\curl\vecB &= \muz\vecJ + \muz\eps_0\frac{\partial\vecE}{\partial t}
\end{align*}

对 $\divg\vecB = \mu_0\rho_g$ 取时间导数：$\partial_t(\divg\vecB) = \mu_0\partial_t\rho_g$。

对 $\curl\vecE = -\partial_t\vecB - \mu_0\vecJ_g$ 取散度：$0 = -\partial_t(\divg\vecB) - \mu_0\divg\vecJ_g$。

联立：$\mu_0\partial_t\rho_g + \mu_0\divg\vecJ_g = 0$，即 $\partial_t\rho_g + \divg\vecJ_g = 0$（磁荷守恒）。\qedsymbol

[量纲校验] $[\rho_g] = [Wb/m^3]$，$[j_g] = [Wb/(m^2\cdot s)]$，$[\partial_t\rho_g] = [Wb/(m^3\cdot s)] = [\divg j_g]$，一致。\qedsymbol



\begin{warning}{第7章常见误区}{ch7-warnings}
\begin{enumerate}
\item 位移电流不是真实电荷的流动——它是变化电场的等效效果。

  \textit{反例}：充电平板电容器两极板之间是\textbf{真空}，没有任何电荷移动，但 $\partial\vect{D}/\partial t \neq 0$ 使得安培环路定理在极板间依然成立、磁场依然存在。若位移电流是"真实电流"，就要求真空中有电荷在流动——显然矛盾。位移电流的正确理解：它是 $\partial \vecE/\partial t$ 的\textbf{等效效应}，与传导电流共同构成"全电流"，两者各自都能产生磁场，但物理本质完全不同。

\item 法拉第定律中感应电动势的方向（负号）——总是阻碍磁通量变化（楞次定律）。

  \textit{反例}：磁铁 N 极插入线圈时，线圈中感应电流产生的磁场在插入端为 N 极，\textbf{排斥}磁铁——若方向相反（助磁），感应电流会加速磁铁插入、产生更大电流，能量凭空增加，违反能量守恒。负号正是能量守恒的体现。

\item $\curl\vecE \neq 0$ 时，不能定义全局电势——但局部仍可定义。

  \textit{反例}：围绕通电螺线管的闭合回路上 $\oint\vecE\cdot\diff\vect{l} = -\diff\Phi/\diff t \neq 0$（Aharonov--Bohm 实验配置），电势的单值定义失效——同一空间点沿不同路径积分会得到不同的"势差"。但在不含时变磁通的任意小区域内，仍然可以局部定义 $\phi = -\int\vecE\cdot\diff\vect{l}$。

\item 麦克斯韦方程组是\textbf{完备的}：给定初始条件和边界条件，电磁场的演化被完全确定。
\end{enumerate}
\end{warning}

% ============================================================

\begin{center}\small \hyperlink{ch6-goal}{上一章：第6章} \quad|\quad \hyperlink{ch8-goal}{下一章：第8章}\end{center}

% 第8章：电磁波的传播（优化版）
% ============================================================
\chapter{电磁波的传播}
\epigraph{波是自然界传递信息的通用语言。}{--- 佚名}

\hypertarget{ch8-goal}{}\begin{insight}{本章核心目标}{ch8-goal-v2}
从麦克斯韦方程组导出电磁波方程，研究平面波在真空、介质和导体中的传播特性，掌握反射、折射和波导等基本现象。理解电磁波是电磁场的自维持振荡——电场和磁场相互激发、共同传播。
\end{insight}

电磁波是麦克斯韦方程组最辉煌的预言。光、无线电波、X射线、微波——它们本质上都是电磁波，只是频率不同。

\section{波动方程的推导}

\subsection{真空中的波动方程}

在无源（$\rho = 0$，$\vecJ = 0$）真空中，麦克斯韦方程组简化为：
\begin{align}
\divg\vecE &= 0 \label{eq:vac1} \\
\curl\vecE &= -\frac{\partial\vecB}{\partial t} \label{eq:vac2} \\
\divg\vecB &= 0 \label{eq:vac3} \\
\curl\vecB &= \muz\eps_0\frac{\partial\vecE}{\partial t} \label{eq:vac4}
\end{align}

对式\eqref{eq:vac2}取旋度：
\[
\curl(\curl\vecE) = \grad(\divg\vecE) - \lap\vecE = -\frac{\partial}{\partial t}(\curl\vecB)
\]

由式\eqref{eq:vac1}，$\divg\vecE = 0$；由式\eqref{eq:vac4}，$\curl\vecB = \muz\eps_0\partial\vecE/\partial t$。

代入得：
\[
-\lap\vecE = -\muz\eps_0\frac{\partial^2\vecE}{\partial t^2}
\]

\begin{theorem}{真空中的电磁波方程}{wave-eq-v2}
\[
\lap\vecE - \muz\eps_0\frac{\partial^2\vecE}{\partial t^2} = \vect{0}
\]
同理：
\[
\lap\vecB - \muz\eps_0\frac{\partial^2\vecB}{\partial t^2} = \vect{0}
\]
\end{theorem}

\begin{insight}{波动方程的深刻含义}{wave-meaning-v2}
这是一个标准的\textbf{波动方程}，形式为 $\lap\psi - (1/c^2)\partial^2\psi/\partial t^2 = 0$。对比可得：
\[
c = \frac{1}{\sqrt{\muz\eps_0}} \approx 2.998 \times 10^8 \, \si{m/s}
\]
这就是\textbf{光速}！麦克斯韦由此得出结论：光是一种电磁波。

更深刻的是：电磁波的速度完全由真空磁导率和介电常数决定——这两个常数原本来自静磁学和静电学的独立实验。它们的组合恰好给出光速，说明光、电、磁是同一物理实体的不同表现。
\end{insight}

\section{平面电磁波}

\subsection{单色平面波的解}

波动方程的最简单解是\textbf{单色平面波}：
\[
\vecE(\vect{r},t) = \vecE_0 \ee^{\ii(\vecK\cdot\vect{r} - \omega t)}
\]

代入波动方程，得\textbf{色散关系}：
\[
k^2 = \muz\eps_0\omega^2 = \frac{\omega^2}{c^2}
\]
即 $\omega = ck$，$c = \omega/k$ 是相速度。

\subsection{横波性质}

由 $\divg\vecE = 0$：
\[
\vecK \cdot \vecE_0 = 0
\]
即 $\vecE$ \textbf{垂直}于传播方向——电磁波是\textbf{横波}。

由 $\curl\vecE = -\partial\vecB/\partial t$，得：
\[
\vecB = \frac{1}{\omega}\vecK \times \vecE = \frac{1}{c}\unitvec{k} \times \vecE
\]

\begin{insight}{电磁波的结构}{em-wave-structure-v2}
平面电磁波具有如下结构：
\begin{itemize}
\item $\vecE$、$\vecB$、$\vecK$ 三者互相垂直，形成右手系：$\vecE \times \vecB \parallel \vecK$
\item $|\vecB| = |\vecE|/c$
\item $\vecE$ 和 $\vecB$ \textbf{同相位}振荡
\item 能量密度：$u = \eps_0 E^2$（因为 $B^2/\muz = \eps_0 E^2$）
\item 能流密度（坡印廷矢量）：$S = EH = E^2/(c\muz) = c\eps_0 E^2 = cu$
\end{itemize}

对时间平均的强度（辐照度）：
\[
I = \langle S \rangle = \frac{1}{2}c\eps_0 E_0^2
\]
\end{insight}

% ---------- 平面波 E/B/k 三色图 ----------
\begin{center}
\begin{tikzpicture}[scale=1.3]
  % Axes
  \draw[->, thick, green!60!black] (-4,0) -- (4.5,0) node[right] {$\vecK$};

  % E field (red sine wave) in y-direction
  \draw[red, very thick, domain=-4:4, samples=100] plot(\x, {1.2*sin(2*\x r)});
  \foreach \x in {-3.5,-2.5,-1.5,-0.5,0.5,1.5,2.5,3.5} {
    \pgfmathsetmacro{\ey}{1.2*sin(2*\x r)}
    \draw[->, red!70!black, thick] (\x,0) -- (\x,\ey);
  }
  \node[red, font=\Large] at (0,1.8) {$\vecE$};

  % B field: ⊙ (out of page) above axis, ⊗ (into page) below — B oscillates in z
  % E 与 B 同相位: B_z ∝ sin(2x), E 峰值处 B 最大
  \foreach \x in {-3.5,-2.5,-1.5,-0.5,0.5,1.5,2.5,3.5} {
    \pgfmathsetmacro{\ey}{1.2*sin(2*\x r)}
    \pgfmathsetmacro{\bsep}{0.34*abs(\ey)}
    \ifdim\ey pt>0pt
      % E 朝上：B 朝 +z（出纸 ⊙ 在下方，入纸 ⊗ 在更下方）
      \draw[blue!70!black, thick] (\x,{-0.6}) circle (0.14);
      \fill[blue!70!black] (\x,{-0.6}) circle (0.035);
      \draw[blue!70!black, thick] (\x,{-1.0-\bsep}) circle (0.14);
      \draw[blue!70!black, thick] (\x-0.098,{-1.0-\bsep-0.098}) -- (\x+0.098,{-1.0-\bsep+0.098});
      \draw[blue!70!black, thick] (\x-0.098,{-1.0-\bsep+0.098}) -- (\x+0.098,{-1.0-\bsep-0.098});
    \else
      % E 朝下：B 朝 -z（入纸 ⊗ 在下方，出纸 ⊙ 在更下方）
      \draw[blue!70!black, thick] (\x,{-0.6}) circle (0.14);
      \draw[blue!70!black, thick] (\x-0.098,{-0.6-0.098}) -- (\x+0.098,{-0.6+0.098});
      \draw[blue!70!black, thick] (\x-0.098,{-0.6+0.098}) -- (\x+0.098,{-0.6-0.098});
      \draw[blue!70!black, thick] (\x,{-1.0-\bsep}) circle (0.14);
      \fill[blue!70!black] (\x,{-1.0-\bsep}) circle (0.035);
    \fi
  }
  \node[blue, font=\Large] at (0,-1.9) {$\vecB$（$\odot$ 出纸 / $\otimes$ 入纸）};

  % Labels and annotations
  \node[green!60!black, font=\large] at (4.2,-0.4) {传播方向};
  \node at (0,-2.8) {$\vecE\perp\vecB\perp\vecK$ \quad (横波)};
  \draw[<->, thick] (2.5,0) -- (2.5,1.2) node[midway, right] {$E_0$};
\end{tikzpicture}
\end{center}


\section{介质中的电磁波}

在线性各向同性介质中，$\vect{D} = \eps\vecE$，$\vecB = \mu\vect{H}$。波动方程变为：
\[
\lap\vecE - \mu\eps\frac{\partial^2\vecE}{\partial t^2} = \vect{0}
\]

波速：
\[
v = \frac{1}{\sqrt{\mu\eps}} = \frac{c}{\sqrt{\mu_r\eps_r}} = \frac{c}{n}
\]

\begin{definition}{折射率}{refractive-index-v2}
\textbf{折射率}：
\[
n = \sqrt{\mu_r\eps_r}
\]
对非磁性介质（$\mu_r \approx 1$）：$n \approx \sqrt{\eps_r}$。例如水的 $\eps_r \approx 80$，$n \approx 1.33$（注意：静态 $\eps_r$ 和光频 $\eps_r$ 不同！光频下水的 $\eps_r \approx 1.77$，$n \approx 1.33$）。
\end{definition}

\section{电磁波在导体中的传播}

在导体中，欧姆定律 $\vecJ = \sigma\vecE$ 成立。麦克斯韦方程组变为：
\[
\lap\vecE - \mu\eps\frac{\partial^2\vecE}{\partial t^2} - \mu\sigma\frac{\partial\vecE}{\partial t} = \vect{0}
\]

对单色波 $\ee^{-\ii\omega t}$：
\[
\lap\vecE + k^2\vecE = \vect{0}, \quad k^2 = \mu\eps\omega^2 + \ii\mu\sigma\omega
\]

\begin{insight}{良导体的近似}{good-conductor-v2}
对良导体（$\sigma/\eps\omega \gg 1$）：
\[
k^2 \approx \ii\mu\sigma\omega \quad\Rightarrow\quad k \approx \sqrt{\ii\mu\sigma\omega} = (1+\ii)\sqrt{\frac{\mu\sigma\omega}{2}}
\]

令 $k = \beta + \ii\alpha$，则：
\[
\beta = \alpha = \sqrt{\frac{\mu\sigma\omega}{2}}
\]

电场：$E \propto \ee^{-\alpha z}\ee^{\ii(\beta z - \omega t)}$

\begin{itemize}
\item $\ee^{-\alpha z}$：指数衰减——电磁波不能深入导体
\item \textbf{趋肤深度}：$\delta = 1/\alpha = \sqrt{2/(\mu\sigma\omega)}$
\item 对铜（$\sigma = 5.8\times 10^7 \, \si{S/m}$）在1 GHz：$\delta \approx 2\,\si{\mu m}$
\end{itemize}

这就是微波炉不能用金属容器、手机信号在电梯里变差的原因。
\end{insight}

\begin{tip}{工程联系：趋肤深度决定了屏蔽与布线设计}{skin-depth-engineering-v2}
\begin{itemize}[nosep]
\item \textbf{电磁屏蔽}：屏蔽罩厚度只需几个趋肤深度即可将泄露场衰减 $e^{-5}\approx 150$ 倍以上。电子设备的金属外壳对 1 GHz 干扰只需 \SI{10}{\mu m} 量级的铜层——这就是为什么"薄薄一层铝箔"就能屏蔽大部分射频干扰。
\item \textbf{高频布线}：高频电流只流经导线表面 $\delta$ 厚的薄层，有效截面积变小、交流电阻增大。射频连接器和高频电感常采用\textbf{镀银}（$\sigma$ 更高）而非加粗导线。
\item \textbf{电力传输}：工频 50 Hz 下铜的 $\delta \approx \SI{9.3}{mm}$，粗导线中心部分几乎不载流——大电流母线有时做成空心管或采用多股绝缘绞线（利兹线）。
\item \textbf{感应加热}：工业感应加热利用 $\delta \propto 1/\sqrt{\omega}$ 把能量集中注入金属表层，频率越高加热层越浅，用于表面淬火。
\end{itemize}
\end{tip}

% ---------- 趋肤效应示意图 ----------
\begin{center}
\begin{tikzpicture}[scale=1.0]
  % conductor cross-section (semicircle cross section of a wire)
  \fill[gray!25] (0,0) circle (2);
  \draw[thick] (0,0) circle (2);
  \node[font=\small] at (0,-2.45) {导体横截面};

  % skin depth layer (annulus near surface)
  \draw[red, very thick] (0,0) circle (2);
  \draw[red, dashed, thick] (0,0) circle (1.55);
  \node[font=\scriptsize, red] at (1.05,1.35) {载流层};
  \draw[<->, red] (1.415,1.03) -- (1.78,1.29) node[midway, right, font=\scriptsize] {$\delta$};

  % current density arrows (strong near surface, weak inside)
  \foreach \a in {0,45,90,135,180,225,270,315} {
    \draw[->, red, thick] ({\a-90}:{1.85}) -- ({\a-90}:{1.6});
    \draw[->, red!40, thin] ({\a-90}:{1.0}) -- ({\a-90}:{0.92});
  }

  % annotations
  \node[font=\scriptsize, align=center] at (0,0) {中心：电流密度\\几乎为零};
  \node[font=\scriptsize, align=center] at (4.1,0.9) {$J(z) = J_0\ee^{-z/\delta}$\\表面最强\\指数衰减};
  \draw[->, gray] (3.3,0.9) -- (2.0,0.9);
\end{tikzpicture}
\end{center}

\textbf{图解}：交流电流不是均匀流过导体截面，而是集中在表面 $\delta$ 厚的薄层——表面电流密度最强，向内指数衰减 $J = J_0\ee^{-z/\delta}$。频率越高，$\delta$ 越薄，中心"闲置"的导体越多——这就是为什么高频导线要镀银而不是加粗。

\section{波导基础}

电磁波在自由空间中传播时，TEM波（横电磁波）可以存在。但在波导（中空的金属管）中，边界条件要求电磁波必须以特定的\textbf{模式}传播。

\subsection{矩形波导的TE/TM模式}

考虑横截面为 $a \times b$（$a > b$）的矩形金属波导，内壁为理想导体。电磁波在波导中沿 $z$ 方向传播，假设场量具有形式 $\ee^{\ii(k_z z - \omega t)}$。

由麦克斯韦方程组和边界条件（切向电场在导体壁为零），可以证明：
\begin{itemize}
\item \textbf{TE模式}（横电模式）：$E_z = 0$，$H_z \neq 0$
\item \textbf{TM模式}（横磁模式）：$H_z = 0$，$E_z \neq 0$
\item TEM模式在空心波导中\textbf{不能存在}（因为没有内导体提供纵向电流）
\end{itemize}

对TE和TM模式，波动方程的分离变量解要求横向波数满足：
\[
k_x = \frac{m\pi}{a}, \quad k_y = \frac{n\pi}{b}
\]
其中 $m, n$ 为非负整数（对TE模式不同时为零，对TM模式均不为零）。

\subsection{截止频率}

纵向波数 $k_z$ 与频率 $\omega$ 的关系为：
\[
k_z^2 = \frac{\omega^2}{c^2} - \left(\frac{m\pi}{a}\right)^2 - \left(\frac{n\pi}{b}\right)^2
\]

\begin{theorem}{矩形波导的截止频率}{cutoff-freq-v2}
矩形波导中 $\text{TE}_{mn}$ 或 $\text{TM}_{mn}$ 模式的截止角频率为：
\[
\omega_c = c\pi\sqrt{\left(\frac{m}{a}\right)^2 + \left(\frac{n}{b}\right)^2}
\]
只有当工作频率 $\omega > \omega_c$（即 $k_z$ 为实数）时，该模式才能在波导中传播；当 $\omega < \omega_c$ 时，$k_z$ 为纯虚数，波呈指数衰减（\textbf{截止}）。

[量纲校验] $[c/a] = (\mathrm{m/s})/\mathrm{m} = \mathrm{s}^{-1}$，$\omega_c$ 单位为 \si{rad/s}。\checkmark
\end{theorem}

最低截止频率对应\textbf{主模} $\text{TE}_{10}$（$m=1, n=0$）：
\[
f_{c,10} = \frac{c}{2a}
\]
通常波导设计为只让 $\text{TE}_{10}$ 单模工作，避免多模传播引起的信号畸变。

\begin{insight}{截止频率的物理图像}{cutoff-physical-v2}
截止频率的出现是电磁波在波导壁之间多次反射的干涉结果。可以将波导中的传播模式想象成平面波以某个角度在波导壁之间来回反射并前进。当频率太低（波长太长）时，入射角太小，波在反射后"折返"而无法有效前进——这就是截止。只有当波长足够短（频率足够高），波才能以足够陡的角度反射，形成有效的纵向传播。

这与光纤中的全反射导光、以及量子力学中势阱中的束缚态有深刻的数学同构性。
\end{insight}

\subsection{TE\texorpdfstring{$_{10}$}{10}主模的场分布}

$\text{TE}_{10}$ 模式的电磁场分量为：
\begin{align}
E_y &= E_0 \sin\left(\frac{\pi x}{a}\right) \ee^{\ii(k_z z - \omega t)} \\
H_x &= -\frac{k_z}{\omega\muz} E_0 \sin\left(\frac{\pi x}{a}\right) \ee^{\ii(k_z z - \omega t)} \\
H_z &= -\frac{\ii\pi}{\omega\muz a} E_0 \cos\left(\frac{\pi x}{a}\right) \ee^{\ii(k_z z - \omega t)}
\end{align}

特点：
\begin{itemize}
\item 电场只有 $y$ 分量，沿 $x$ 方向呈半正弦分布，在波导宽边中央最强
\item 电场在窄边（$y$ 方向）上均匀——与 $y$ 无关
\item 磁场有 $x$ 和 $z$ 两个分量
\item 表面电流在宽边中央呈\textbf{纵向}（沿 $z$），向窄边方向连续转为\textbf{横向}并接到窄边上——宽边中央的电流并不为零（$x=a/2$ 处 $H_x$ 取极大）
\end{itemize}

% ---------- 波导模式示意图（横截面场分布） ----------
\begin{center}
\begin{tikzpicture}[scale=1.0]
  % waveguide cross-section (a wide, b tall)
  \draw[very thick] (0,0) rectangle (6,3);
  \node[below] at (3,-0.15) {$a$（宽边）};
  \node[left] at (-0.15,1.5) {$b$（窄边）};

  % TE10 E_y field: half-sine along x, uniform along y (drawn as vertical arrows)
  \foreach \x in {0.6,1.2,1.8,2.4,3.0,3.6,4.2,4.8,5.4} {
    \pgfmathsetmacro{\ee}{1.1*sin(30*\x)}
    \ifdim\ee pt>0pt
      \draw[->, red, thick] (\x,1.5) -- (\x,{1.5+\ee});
    \else
      \draw[->, red, thick] (\x,1.5) -- (\x,{1.5+\ee});
    \fi
  }
  \node[red] at (3,2.9) {$\vecE$（$y$ 方向，半正弦分布）};

  % H field: closed loops around E maxima (drawn as small ellipses)
  \draw[blue!70!black, thick] (3.0,1.5) ellipse (0.8 and 0.45);
  \draw[blue!70!black, thick] (3.0,1.5) ellipse (0.45 and 0.25);
  \node[blue!70!black] at (3.0,0.45) {$\vect{H}$（绕 $E$ 闭环）};

  % labels
  \node at (0.9,0.35) {\small 导体壁};
  \draw[gray, dashed] (0,0.7) -- (6,0.7);
  \node[right, gray] at (6.1,0.7) {\small 中线};
\end{tikzpicture}
\end{center}

\begin{example}{矩形波导尺寸设计}{waveguide-design-ex-v2}
\difficulty{提高}
设计一个矩形波导，要求在 \SI{3}{\giga\hertz} 到 \SI{5}{\giga\hertz} 频段内单模（$\text{TE}_{10}$）工作。已知宽边 $a = \SI{5}{\cm}$，求窄边 $b$ 的取值范围。
\tcblower
\textbf{解题思路}：单模工作 = 主模 $\text{TE}_{10}$ 传播 + \textbf{所有高阶模截止}。审题分三步：(1) 算出 $\text{TE}_{10}$ 截止频率（由宽边 $a$ 决定），确认它在频段下限之下；(2) 算出最邻近的高阶模 $\text{TE}_{20}$（由 $a$）和 $\text{TE}_{01}$（由 $b$）的截止频率，确认它们在频段上限之上；(3) 由 $\text{TE}_{01}$ 的条件反解 $b$ 的范围。\textbf{已知}：频段 3--5 GHz，$a = 5$ cm；\textbf{求}：$b$ 取值范围。

\textbf{解答：}

单模工作的条件是：$\text{TE}_{10}$ 模式可以传播，而 $\text{TE}_{20}$ 和 $\text{TE}_{01}$ 模式被截止。

$\text{TE}_{10}$ 截止频率：
\[
f_{c,10} = \frac{c}{2a} = \frac{3\times 10^8}{2\times 0.05} = \SI{3}{\giga\hertz}
\]

$\text{TE}_{20}$ 截止频率：
\[
f_{c,20} = \frac{c}{a} = \SI{6}{\giga\hertz} > \SI{5}{\giga\hertz}
\]
在频段内被截止，满足要求。

$\text{TE}_{01}$ 截止频率：
\[
f_{c,01} = \frac{c}{2b}
\]
要求 $f_{c,01} > \SI{5}{\giga\hertz}$，即：
\[
\frac{c}{2b} > 5\times 10^9 \quad\Rightarrow\quad b < \frac{3\times 10^8}{2\times 5\times 10^9} = \SI{0.03}{\m} = \SI{3}{\cm}
\]

同时为保证单模工作稳定性，通常取 $b < a/2$，且 $b$ 不能太小（功率容量和损耗考虑）。

\textbf{结论}：窄边应满足 $b < \SI{3}{\cm}$，工程上常取 $b = \SI{2.5}{\cm}$ 左右。

\textbf{结果解读}：(1) 结论中 $b$ 只设上限不设下限，但工程上 $b$ 太小会降低功率容量、增大壁损耗——所以实际取 $b\approx a/2$；(2) 极限校验 $b\to a/2 = 2.5$ cm 时 $\text{TE}_{01}$ 与 $\text{TE}_{10}$ 简并（$f_{c,01} = f_{c,10}$），多模风险最大，这就是工程上避开 $b = a/2$ 的原因；(3) 本题考察\textbf{截止频率公式的反向使用}——由"哪个模必须截止"反推几何尺寸，这是微波工程设计的标准思路。

\begin{insight}{为什么空心金属波导不能传TEM波}{tem-no-propagation-v2}
上面的设计中我们默认只有 TE/TM 模式——为什么 TEM 波在空心波导中不存在？

\textbf{论证}：TEM 波要求 $E_z = H_z = 0$，即电场和磁场都完全横向。由 $\curl\vecE = 0$（对横向场而言），横向电场必须可以写成某个横向标量势的梯度：$\vecE_t = -\grad_t\phi$。代入横向的 $\divg\vect{D} = 0$ 得 $\lap_t\phi = 0$。而波导壁是等势体（理想导体），边界条件要求 $\phi$ 在整个横截面边界上为常数——由拉普拉斯方程解的唯一性，$\phi$ 在整个截面上恒为常数，于是 $\vecE_t = \vect{0}$，波不存在。

\textbf{物理图像}：TEM 波本质上是"双导体传输线"模式（如同轴线），电场线必须从一个导体出发、终止于另一个导体。空心波导只有\textbf{一个}导体边界，电场线无处安放。同轴线之所以能传 TEM 波，正是因为它有内、外两个导体。

\textbf{工程意义}：这也是低频下同轴电缆优于空心波导的原因——TEM 模\textbf{无截止频率}，可以传输任意低频；而空心波导的 $\text{TE}_{10}$ 主模有截止频率 $f_c = c/2a$，只适合微波频段（雷达、卫星通信）。
\end{insight}
\end{example}

\section{色散与洛伦兹模型}

前文将介电常数 $\eps$ 和折射率 $n$ 视为与频率无关的常数。然而，实验表明折射率强烈依赖于频率——这就是\textbf{色散}现象。棱镜分光、彩虹、光纤中的脉冲展宽都是色散的表现。

\subsection{电子受迫振动模型}

洛伦兹将介质中的束缚电子建模为受电场驱动的简谐振子。设电子偏离平衡位置的位移为 $x$，在交变电场 $E = E_0\ee^{-\ii\omega t}$ 作用下，电子运动方程为：
\[
m\ddot{x} + m\gamma\dot{x} + m\omega_0^2 x = -e E_0\ee^{-\ii\omega t}
\]
其中 $\omega_0$ 是束缚电子的固有频率，$\gamma$ 是阻尼系数（反映辐射阻尼和碰撞损耗）。

设稳态解 $x = x_0\ee^{-\ii\omega t}$，代入得：
\[
x_0 = \frac{-eE_0/m}{\omega_0^2 - \omega^2 - \ii\gamma\omega}
\]

单位体积内的电极化强度：
\[
\vect{P} = N(-e)\vect{x} = \frac{Ne^2/m}{\omega_0^2 - \omega^2 - \ii\gamma\omega}\vecE
\]

由 $\vect{D} = \eps_0\vecE + \vect{P} = \eps(\omega)\vecE$，得复介电常数：
\[
\eps(\omega) = \eps_0 + \frac{Ne^2/m}{\omega_0^2 - \omega^2 - \ii\gamma\omega}
\]

\begin{theorem}{洛伦兹色散模型}{lorentz-model-v2}
介质的复折射率平方满足洛伦兹模型：
\[
n^2(\omega) = 1 + \frac{Ne^2}{\eps_0 m}\frac{1}{\omega_0^2 - \omega^2 - \ii\gamma\omega}
\]
对低密度气体（$n \approx 1$），实部近似为：
\[
n(\omega) \approx 1 + \frac{Ne^2}{2\eps_0 m}\frac{\omega_0^2 - \omega^2}{(\omega_0^2 - \omega^2)^2 + \gamma^2\omega^2}
\]
\end{theorem}

\subsection{正常色散与反常色散}

由洛伦兹模型，折射率的频率依赖呈现两种典型行为：

\textbf{正常色散}（远离共振频率 $\omega_0$）：
\begin{itemize}
\item $n$ 随 $\omega$ 增大而缓慢增大
\item 相速度 $v_p = c/n$ 随 $\omega$ 增大而减小
\item 典型例子：可见光在玻璃中的色散（紫光折射率大于红光）
\end{itemize}

\textbf{反常色散}（共振频率附近 $|\omega - \omega_0| \lesssim \gamma$）：
\begin{itemize}
\item $n$ 随 $\omega$ 增大而急剧减小
\item 相速度可以超过 $c$（但信号传播速度——群速度——仍小于 $c$）
\item 伴随强烈的吸收（$\eps$ 的虚部达到峰值）
\item 经典例子：钠蒸气在D线（\SI{589}{\nano\m}）附近的反常色散
\end{itemize}

\begin{insight}{色散的物理本质}{dispersion-physical-v2}
色散的本质是介质对电磁场的\textbf{延迟响应}。当电场频率接近电子固有频率 $\omega_0$ 时，电子发生共振，大量吸收电磁能量并以特定相位重新辐射——这种共振行为导致折射率的剧烈变化。

从更根本的层面看，色散告诉我们"真空"本身也是具有量子结构的介质。在量子电动力学中，真空涨落（虚电子-正电子对的产生和湮灭）导致光在真空中传播时也存在极微弱的色散——这在极端强场（如脉冲星附近）或极高精度测量中才可能显现。洛伦兹经典模型虽然简化，却抓住了色散的物理核心：介质中的带电粒子在电磁场作用下的受迫振动。
\end{insight}

\section{反射与折射}

\subsection{菲涅尔公式}

电磁波在介质界面处发生反射和折射。设入射面为 $xz$ 平面，$\vecK$ 在 $xz$ 平面内。

\begin{theorem}{菲涅尔公式}{fresnel-v2}
定义振幅反射系数 $r$ 和透射系数 $t$：

\textbf{s偏振}（$\vecE$ 垂直于入射面）：
\[
r_s = \frac{n_1\cos\theta_i - n_2\cos\theta_t}{n_1\cos\theta_i + n_2\cos\theta_t}, \quad t_s = \frac{2n_1\cos\theta_i}{n_1\cos\theta_i + n_2\cos\theta_t}
\]

\textbf{p偏振}（$\vecE$ 在入射面内）：
\[
r_p = \frac{n_2\cos\theta_i - n_1\cos\theta_t}{n_2\cos\theta_i + n_1\cos\theta_t}, \quad t_p = \frac{2n_1\cos\theta_i}{n_2\cos\theta_i + n_1\cos\theta_t}
\]
\end{theorem}

\subsection{全反射与倏逝波}

当电磁波从光密介质（$n_1$）入射到光疏介质（$n_2 < n_1$）时，存在\textbf{临界角}：
\[
\theta_c = \arcsin\frac{n_2}{n_1}
\]

当入射角 $\theta_i > \theta_c$ 时，由斯涅尔定律 $n_1\sin\theta_i = n_2\sin\theta_t$，$\sin\theta_t > 1$，不存在实数解的折射角——此时发生\textbf{全反射}。

\begin{theorem}{全反射的反射系数}{total-reflection-v2}
当 $\theta_i > \theta_c$（$n_1 > n_2$）时，s偏振和p偏振的反射系数模均为1：
\[
|r_s| = |r_p| = 1
\]
即入射能量全部反射，无能量透入第二种介质。
\end{theorem}

\begin{insight}{倏逝波的物理图像}{evanescent-wave-v2}
全反射时，第二种介质中并非完全无场。在界面附近存在\textbf{倏逝波}（evanescent wave）：

透射波的法向波数变为纯虚数：$k_{tz} = \ii\kappa$，其中 $\kappa > 0$。透射电场具有形式：
\[
\vecE_t \propto \ee^{-\kappa z}\ee^{\ii(k_x x - \omega t)}
\]

关键特征：
\begin{itemize}
\item 振幅沿法向（$z$ 方向）呈\textbf{指数衰减}，衰减长度约为波长量级
\item 沿切向（$x$ 方向）仍呈行波形式传播
\item 平均能流密度的法向分量为零——无净能量流入第二种介质
\item 如果在倏逝波区域放置第三块介质（间距小于衰减长度），光可以"隧穿"过去——这就是\textbf{受抑全反射}（frustrated total internal reflection），是光纤分束器和近场光学显微镜的物理基础
\end{itemize}

倏逝波的存在表明：全反射不是场在界面处的简单"反弹"，而是电磁场以非辐射形式短暂渗入第二种介质、再返回第一种介质的量子力学式过程。它在光纤通信、棱镜耦合器和近场光学中都有重要应用。
\end{insight}

\begin{insight}{布儒斯特角}{brewster-v2}
当p偏振光以\textbf{布儒斯特角}入射时，反射系数为零：
\[
\tan\theta_B = \frac{n_2}{n_1}
\]

物理原因：此时反射波和折射波垂直，p偏振的振荡方向恰好与反射方向平行，无法激发反射波。利用布儒斯特角可以制作偏振器。
\end{insight}

\begin{review}{第8章要点回顾}{ch8-review}
\begin{center}
\begin{tabular}{@{}ll@{}}
\toprule
\textbf{公式/概念} & \textbf{内容} \\
\midrule
$c = 1/\sqrt{\muz\eps_0}$ & 真空中光速 \\
$\vecE$, $\vecB$, $\vecK$ 互相垂直 & 横波性质 \\
$|\vecB| = |\vecE|/c$ & 场振幅关系 \\
$I = \frac{1}{2}c\eps_0 E_0^2$ & 光强 \\
$n = \sqrt{\mu_r\eps_r}$ & 折射率 \\
$\delta = \sqrt{2/(\mu\sigma\omega)}$ & 趋肤深度 \\
$\tan\theta_B = n_2/n_1$ & 布儒斯特角 \\
\bottomrule
\end{tabular}
\end{center}
\end{review}

\begin{tip}{本章核心框架}{ch8-summary}
\textbf{一句话总结}：电磁波是横波，在介质中传播遵循反射/折射定律，在波导中形成离散模式。\\
\textbf{核心内容}：平面波（TEM）、反射折射（菲涅尔公式、布儒斯特角、全反射）、波导（TE/TM模式、截止频率）、色散（相速/群速）。\\
\textbf{关键公式}：$v=c/n$、$n=\sqrt{\mu_r\eps_r}$、$\omega_{mn}=c\pi\sqrt{(m/a)^2+(n/b)^2}$。\\
\textbf{衔接}：第9章将讨论电磁波的产生——加速电荷的辐射、天线辐射场。
\end{tip}

% ---------- 第8章随堂自测 ----------
\section*{第8章随堂自测}
\addcontentsline{toc}{section}{第8章随堂自测}

\begin{miniquiz}
\textbf{一、选择题}
\begin{enumerate}[label=\arabic*.]
\item 平面电磁波在真空中传播时，$\vecE$ 和 $\vecB$ 的关系是：
  \begin{enumerate}[label=(\Alph*),nosep]
  \item $\vecE \parallel \vecB$
  \item $\vecE \perp \vecB$，且都垂直于传播方向
  \item $\vecE$和$\vecB$的夹角为$45^\circ$
  \item $\vecE$沿传播方向，$\vecB$垂直于传播方向
  \end{enumerate}
\begin{ocg}{答案}{ocg答案8}{off}
  \textit{答案：B。电磁波是横波，$\vecE\perp\vecB\perp\vecK$。}
\end{ocg}

\item 电磁波在介质中的相速度 $v = c/n$，其中折射率 $n$：
  \begin{enumerate}[label=(\Alph*),nosep]
  \item $n = \sqrt{\mu_r\eps_r}$
  \item $n = \mu_r\eps_r$
  \item $n = \sqrt{\mu_r/\eps_r}$
  \item $n = 1/\sqrt{\mu_r\eps_r}$
  \end{enumerate}
\begin{ocg}{答案}{ocg答案8}{off}
  \textit{答案：A。$n=\sqrt{\mu_r\eps_r}$，对非磁性介质$\mu_r\approx 1$，则$n\approx\sqrt{\eps_r}$。}
\end{ocg}

\item 全反射发生在：
  \begin{enumerate}[label=(\Alph*),nosep]
  \item 光从光疏介质射向光密介质
  \item 光从光密介质射向光疏介质，入射角大于临界角
  \item 任何界面
  \item 只有金属表面
  \end{enumerate}
\begin{ocg}{答案}{ocg答案8}{off}
  \textit{答案：B。全反射需要$n_1>n_2$且$\theta_i>\theta_c$。}
\end{ocg}
\end{enumerate}

\vspace{0.5em}
\textbf{二、填空题}
\begin{enumerate}[label=\arabic*.,resume]
\item 真空中电磁波的能流密度由\underline{\hspace{2cm}}矢量描述：$\vect{S} = $ \underline{$\vecE\times\vecB/\muz$}。
\item TE波是指电场\underline{\hspace{2cm}}于传播方向的波；TM波是指磁场\underline{\hspace{2cm}}于传播方向的波。
\begin{ocg}{答案}{ocg答案8}{off}
  \textit{答案：垂直；垂直。}
\end{ocg}
\end{enumerate}
\end{miniquiz}
\toggleButton{显示答案}{ocg答案8}

\begin{chapterreview}
\begin{itemize}[nosep]
\item \textbf{核心内容}：平面波解、偏振（线/圆/椭圆）、反射与折射（菲涅尔公式）、波导模式（TE/TM/TEM）。
\item \textbf{关键公式}：菲涅尔公式、布儒斯特角$\theta_B$、临界角$\theta_c$、截止频率$\omega_{mn}$。
\item \textbf{自检验}：能否推导TE波的截止频率？能否解释布儒斯特角的物理起源？
\item \textbf{衔接预告}：第9章将讨论电磁波的辐射——天线、偶极辐射、推迟势。
\end{itemize}
\end{chapterreview}

\begin{tip}{本章自检清单}{ch8-checklist}
学完本章后对照下列条目逐条自测，全部通过即可进入下一章；有未通过项时，回到对应小节复习：
\begin{itemize}[label=$\square$, nosep]
\item 由麦克斯韦方程组导出真空波动方程
\item 说明平面波的横波结构与 $\vecE$、$\vecB$、$\vecK$ 关系
\item 计算介质与导体中的波速、衰减与趋肤深度
\item 求矩形波导截止频率并解释 TEM 波不存在的机理
\end{itemize}
\end{tip}

% ---------- 第8章综合练习题 ----------
\section*{第8章综合练习题}
\addcontentsline{toc}{section}{第8章综合练习题}
\begin{enumerate}[label=\textbf{\arabic*.}]
\item \textbf{基础} 证明：在均匀无损耗介质中，平面电磁波的 $|\vecE|/|\vecB|=1/\sqrt{\muz\eps}=\omega/k$。
\item \textbf{基础} 光从空气（$n_1=1$）入射到玻璃（$n_2=1.5$），入射角 $30°$。求反射角、折射角和反射率（s偏振和p偏振）。
\item \textbf{提高} 设计一个矩形波导（宽 $a=2$ cm，高 $b=1$ cm），求TE$_{10}$模式的截止频率，并计算 $\omega=2\pi\times 10$ GHz 时的波导波长。
\item \textbf{提高} 光从光密介质（$n_1=1.5$）到光疏介质（$n_2=1$），入射角 $50°$。求全反射时的倏逝波透射深度（定义为振幅降至 $1/e$ 的距离）。
\item \textbf{挑战} 利用洛伦兹模型推导折射率 $n(\omega)$ 的色散关系，证明在共振频率附近出现反常色散（$dn/d\omega<0$），并讨论其物理意义。
\end{enumerate}
\newpage
\section*{习题解答}
\addcontentsline{toc}{section}{习题解答}

\textbf{题1解析：}

[分析] 本题从麦克斯韦方程组出发证明平面电磁波中电场与磁场振幅的关系。

[解答] 对单色平面波 $\vecE = \vecE_0 e^{i(\vecK\cdot\vect{r} - \omega t)}$，代入 $\curl\vecE = -\partial\vecB/\partial t$：

$\curl\vecE = i\vecK\times\vecE = i\omega\vecB$。

故 $\vecB = \dfrac{1}{\omega}\vecK\times\vecE$，$|\vecB| = \dfrac{k}{\omega}|\vecE| = \dfrac{1}{v}|\vecE|$。

在介质中 $v = 1/\sqrt{\mu\eps}$，$\omega/k = v$。因此：
\[
\frac{|\vecE|}{|\vecB|} = \frac{\omega}{k} = v = \frac{1}{\sqrt{\mu\eps}}
\]

真空中 $\mu = \mu_0$，$\eps = \eps_0$，$v = c$。\qedsymbol

[量纲校验] $[E/B] = [V/m]/[T] = [LT^{-1}] = [v]$，一致。\qedsymbol

\textbf{题2解析：}

[分析] 本题应用菲涅尔公式计算反射率和折射率，是电磁波界面问题的基本计算。

[解答] 反射角：$\theta_r = \theta_i = 30°$。

折射角（斯涅尔定律）：$n_1\sin\theta_i = n_2\sin\theta_t$，$\sin\theta_t = \sin 30°/1.5 = 1/3$，$\theta_t = \arcsin(1/3) \approx 19.5°$。

菲涅尔反射系数：
\begin{align*}
r_s &= \frac{n_1\cos\theta_i - n_2\cos\theta_t}{n_1\cos\theta_i + n_2\cos\theta_t} = \frac{1\cdot\cos 30° - 1.5\cdot\cos 19.5°}{1\cdot\cos 30° + 1.5\cdot\cos 19.5°} \approx -0.240 \\
r_p &= \frac{n_2\cos\theta_i - n_1\cos\theta_t}{n_2\cos\theta_i + n_1\cos\theta_t} = \frac{1.5\cdot\cos 30° - 1\cdot\cos 19.5°}{1.5\cdot\cos 30° + 1\cdot\cos 19.5°} \approx 0.159
\end{align*}

反射率：$R_s = |r_s|^2 \approx 0.058$（5.8\%），$R_p = |r_p|^2 \approx 0.025$（2.5\%）。\qedsymbol

[量纲校验] $r_s, r_p$ 无量纲，$R_s, R_p$ 无量纲，一致。\qedsymbol

\textbf{题3解析：}

[分析] 本题计算矩形波导的截止频率和波导波长，是波导理论的基础。

[解答] TE$_{10}$模式（$m=1, n=0$）：

截止频率：$\omega_c = c\pi/a$，$f_c = c/(2a) = 3\times 10^8/(2\times 0.02) = 7.5$ GHz。

截止波数：$k_c = \pi/a = \pi/0.02 = 50\pi$ rad/m。

工作频率 $\omega = 2\pi\times 10$ GHz，$k = \omega/c = 2\pi\times 10^{10}/(3\times 10^8) = 200\pi/3$ rad/m。

波导波数：$k_g = \sqrt{k^2 - k_c^2} = \sqrt{(200\pi/3)^2 - (50\pi)^2} = 50\pi\sqrt{16/9 - 1} = 50\pi\sqrt{7}/3$。

波导波长：$\lambda_g = 2\pi/k_g = 6/(50\sqrt{7}) \approx 0.045$ m $= 4.5$ cm。\qedsymbol

[量纲校验] $[f_c] = [T^{-1}]$，$[c/a] = [LT^{-1}]/[L] = [T^{-1}]$；$[\lambda_g] = [L]$，一致。\qedsymbol

\textbf{题4解析：}

[分析] 本题计算全反射时的倏逝波透射深度，是理解全反射物理机制的关键。

[解答] 临界角：$\theta_c = \arcsin(n_2/n_1) = \arcsin(1/1.5) = \arcsin(2/3) \approx 41.8°$。

$\theta_i = 50° > \theta_c$，发生全反射。

由斯涅尔定律：$n_1\sin\theta_i = n_2\sin\theta_t$，$\sin\theta_t = 1.5\sin 50° > 1$。

设 $\cos\theta_t = i\sqrt{\sin^2\theta_t - 1}$（虚数）。透射波沿 $z$ 方向（界面法向）传播常数：
\[
\kappa = k_1\sqrt{\sin^2\theta_i - (n_2/n_1)^2} = \frac{2\pi n_1}{\lambda_0}\sqrt{\sin^2 50° - (2/3)^2}
\]

取 $\lambda_0 = 500$ nm（可见光），$\kappa \approx \dfrac{2\pi\times 1.5}{500\times 10^{-9}}\sqrt{0.587 - 0.444} \approx 1.885\times 10^{7}\times 0.377 \approx 7.1\times 10^{6}$ m$^{-1}$。

透射深度（振幅降至 $1/e$）：$\delta = 1/\kappa \approx 1.4\times 10^{-7}$ m $= 0.14$ $\mu$m（$\kappa$ 与 $k_1$ 同量级但更小，符合预期）。\qedsymbol

[量纲校验] $[\kappa] = [L^{-1}]$，$[k] = [L^{-1}]$，一致；$[\delta] = [L]$，一致。\qedsymbol

\textbf{题5解析：}

[分析] 本题推导洛伦兹色散模型，解释正常色散与反常色散的物理起源。

[解答] 洛伦兹模型：电子受弹性恢复力 $-m\omega_0^2 x$、阻尼力 $-m\gamma\dot{x}$ 和驱动力 $-eE$。

运动方程：$m\ddot{x} + m\gamma\dot{x} + m\omega_0^2 x = -eE_0 e^{-i\omega t}$。

稳态解：$x = -eE_0/(m(\omega_0^2 - \omega^2 - i\gamma\omega))$。

极化强度：$P = -Nex = Ne^2E_0/(m(\omega_0^2 - \omega^2 - i\gamma\omega))$。

$\chi_e = P/(\eps_0 E) = \dfrac{Ne^2}{\eps_0 m(\omega_0^2 - \omega^2 - i\gamma\omega)}$。

折射率：$n^2 = 1 + \chi_e$，$n(\omega) \approx 1 + \dfrac{Ne^2}{2\eps_0 m(\omega_0^2 - \omega^2 - i\gamma\omega)}$。

$\dfrac{\diff n}{\diff\omega} = \dfrac{Ne^2}{2\eps_0 m}\cdot\dfrac{2\omega\left[(\omega_0^2 - \omega^2)^2 - \gamma^2\omega_0^2\right]}{\left((\omega_0^2 - \omega^2)^2 + \gamma^2\omega^2\right)^2}$（实部近似）。

\textbf{关键在分子方括号里的 $-\gamma^2\omega_0^2$}：正是它使 $dn/d\omega$ 可以在共振附近变号。若像无阻尼近似那样只保留 $2\omega$（见下文"注意"），$dn/d\omega$ 将处处为正，无法得到反常色散。

远离共振（$\omega\ll\omega_0$ 或 $\omega\gg\omega_0$）：$dn/d\omega > 0$，\textbf{正常色散}；
共振附近（$|\omega-\omega_0|\lesssim\gamma$，且必须保留阻尼项）：$dn/d\omega < 0$，\textbf{反常色散}。\textbf{注意}：分母 $((\omega_0^2-\omega^2)^2+\gamma^2\omega^2)^2$ 恒正、因子 $2\omega>0$，故符号完全由分子方括号决定——当 $(\omega_0^2-\omega^2)^2 < \gamma^2\omega_0^2$（即 $|\omega-\omega_0| \lesssim \gamma/2$ 的共振邻域）时方括号为负，$dn/d\omega<0$；离开该邻域后方括号转正，$dn/d\omega>0$。若取无阻尼极限 $\gamma\to0$，方括号退化为 $(\omega_0^2-\omega^2)^2\geq0$，$dn/d\omega$ 除共振点外处处为正——这正是"反常色散必须保留阻尼项"的原因（与 §8.6.2 正文一致）。

反常色散意味着群速度可以超过相速度，但信号速度不超过 $c$（因果性限制）。\qedsymbol

[量纲校验] $[\chi_e]$ 无量纲，$[Ne^2/(\eps_0 m\omega^2)] = [L^{-3}][C^2]/([F/m][M][T^{-2}]) = 1$（无量纲），一致。\qedsymbol




% ============================================================

\begin{center}\small \hyperlink{ch7-goal}{上一章：第7章} \quad|\quad \hyperlink{ch9-goal}{下一章：第9章}\end{center}

% 第9章：电磁波的辐射（优化版）
% ============================================================
\chapter{电磁波的辐射}
\epigraph{加速的电荷会辐射——这是电动力学对宇宙最慷慨的馈赠。}{--- 本书作者}

\hypertarget{ch9-goal}{}\begin{insight}{本章核心目标}{ch9-goal-v2}
理解加速运动的电荷如何产生电磁辐射。掌握推迟势的概念，学会计算电偶极辐射的场和功率。理解天线辐射的基本原理。
\end{insight}

静止电荷产生静电场，匀速运动电荷产生静磁场——但都不辐射电磁波。只有当电荷\textbf{加速运动}时，才会产生脱离源而独立传播的电磁辐射。

\begin{tip}{想一想}{think-why-radiate}
不急着看下面的数学，先自己想一想：\textbf{为什么静止电荷不辐射、加速电荷才辐射？}

提示一：辐射是"脱离源独立传播"的波——它携带能量飞向远方。如果场完全由电荷当前位置决定（瞬时作用），电荷一停，远方的场就要立刻消失——能量凭空消失，违反能量守恒。

提示二：场的扰动以光速 $c$ 传播（推迟势）。电荷加速时，它"通知"远处场变化的消息需要时间 $r/c$ 在路上跑——这个正在路上跑的"消息"就是辐射波。

读完推迟势一节后回头对照：辐射正是"推迟消息"在电荷停止加速后仍在路上传播的部分。
\end{tip}

\section{推迟势}

\subsection{非齐次波动方程}

含源的麦克斯韦方程组导出的势方程为：
\begin{align}
\lap\phi - \muz\eps_0\frac{\partial^2\phi}{\partial t^2} &= -\frac{\rho}{\eps_0} \\
\lap\vecA - \muz\eps_0\frac{\partial^2\vecA}{\partial t^2} &= -\muz\vecJ
\end{align}

这是\textbf{非齐次波动方程}（达朗贝尔方程）。

\subsection{推迟解}

\begin{theorem}{推迟势}{retarded-potentials-v2}
达朗贝尔方程的解为：
\begin{align}
\phi(\vect{r},t) &= \frac{1}{4\pi\eps_0} \int \frac{\rho(\vect{r}', t_r)}{|\vect{r} - \vect{r}'|} \, \diff^3r' \\
\vecA(\vect{r},t) &= \frac{\muz}{4\pi} \int \frac{\vecJ(\vect{r}', t_r)}{|\vect{r} - \vect{r}'|} \, \diff^3r'
\end{align}
其中 $t_r = t - |\vect{r} - \vect{r}'|/c$ 是\textbf{推迟时间}。符号逐项说明：$t_r$ = 推迟时刻（源"发出"影响的时刻）；$|\vect{r}-\vect{r}'|$ = 源点到场点的距离；$c$ = 光速；$\rho(\vect{r}',t_r)$ = 在推迟时刻的源电荷密度；$\vecJ(\vect{r}',t_r)$ = 在推迟时刻的源电流密度。

[量纲校验] $[\rho\,\diff^3r'/R] = (\mathrm{C/m^3})(\mathrm{m^3})/\mathrm{m} = \mathrm{C/m}$，乘 $[1/(4\pi\eps_0)] = \mathrm{N\cdot m^2/C^2}$ 得 $\mathrm{N\cdot m/C} = \mathrm{V}$ = 电势单位。\checkmark；$[\vecJ\,\diff^3r'/R] = (\mathrm{A/m^2})(\mathrm{m^3})/\mathrm{m} = \mathrm{A\cdot m}$，乘 $[\muz/4\pi] = \mathrm{H/m}$ 得 $\mathrm{H\cdot A} = \mathrm{Wb/m} = \mathrm{T\cdot m}$ = 矢势单位。\checkmark
\end{theorem}

\begin{insight}{推迟势与格林函数}{retarded-green-v2}
回顾第2章，泊松方程的格林函数满足：
\[
\lap \frac{1}{|\vect{r} - \vect{r}'|} = -4\pi \, \delta^{(3)}(\vect{r} - \vect{r}')
\]
这描述了静态点源的势分布。在含时情形下，推迟势可以看作\textbf{含时格林函数}的积分：
\[
G(\vect{r}, t) = \frac{\delta(t - r/c)}{4\pi r}
\]
其中 $\delta(t - r/c)$ 保证了因果性——只有早于当前时刻 $t_r = t - r/c$ 的源才能影响场点。从静态的 $1/r$ 到含时的 $\delta(t-r/c)/r$，数学结构一脉相承：格林函数从拉普拉斯算子的逆推广到了达朗贝尔算子的逆。
\end{insight}

\begin{insight}{推迟时间的物理意义}{retarded-time-v2}
推迟时间反映了电磁相互作用的\textbf{有限传播速度}。空间中某点在时刻 $t$ 的电势，不是由同一时刻的电荷分布决定的，而是由\textbf{更早时刻} $t_r = t - R/c$ 的电荷分布决定的——电磁波以光速 $c$ 传播，需要时间 $R/c$ 才能从源点到达场点。

这是相对论因果性的直接要求：如果场能瞬时传播，就违背了"信息不能超光速"的基本原理。推迟势的存在表明，电磁相互作用不是超距的，而是通过场以光速传播的。
\end{insight}

% ---------- 推迟势因果示意图 ----------
\begin{center}
\begin{tikzpicture}[scale=1.0]
  % timeline axis
  \draw[->, thick] (-0.5,0) -- (9.5,0) node[right, font=\small] {$t$};

  % source event (t_r)
  \fill[red] (2,0) circle (0.1);
  \node[above, font=\small, red] at (2,0.15) {源"发出"影响：$t_r$};
  \node[below, font=\scriptsize] at (2,-0.1) {电荷此刻在 $\vect{r}'$};

  % propagation wave packet
  \draw[->, thick, green!60!black] (2.3,0.5) -- (5.7,0.5) node[midway, above, font=\scriptsize] {电磁波以 $c$ 传播（在路上跑 $R/c$）};
  \draw[decorate, decoration={snake, amplitude=1mm, segment length=3mm}, green!60!black, thick] (2.3,0.5) -- (5.7,0.5);

  % observation event (t)
  \fill[blue] (8,0) circle (0.1);
  \node[above, font=\small, blue] at (8,0.15) {场点"收到"影响：$t$};
  \node[below, font=\scriptsize] at (8,-0.1) {场点此刻在 $\vect{r}$};

  % time interval R/c
  \draw[<->, thick] (2,-0.9) -- (8,-0.9) node[midway, below, font=\small] {$t - t_r = R/c$};
  \draw[dashed, gray] (2,-0.15) -- (2,-1.05);
  \draw[dashed, gray] (8,-0.15) -- (8,-1.05);
\end{tikzpicture}
\end{center}

\textbf{图解}：场点在时刻 $t$ "感受"到的电势，是源在\textbf{更早}时刻 $t_r = t - R/c$ 发出的影响——消息以光速在路上跑了 $R/c$ 才到达。类比声音：你听到的雷声是闪电\textbf{几秒前}发生的事，看到闪电听到雷声的时间差就是"推迟"。电磁场同理，只是"消息"跑得是光速。

\section{电偶极辐射}

\subsection{电偶极辐射场}

对于局域在很小区域内的时变电荷分布，远处的辐射场由电偶极矩 $\vect{p}(t)$ 主导。

\begin{theorem}{电偶极辐射}{dipole-radiation-v2}
时变电偶极矩 $\vect{p}(t) = p_0\ee^{-\ii\omega t}\unitvec{z}$ 产生的辐射场（远场区 $r \gg \lambda$）：
\begin{align}
\vecE &= \frac{k^2}{4\pi\eps_0} \frac{\sin\theta}{r} \, p_0 \ee^{\ii(kr - \omega t)} \, \unitvec{\theta} \\
\vecB &= \frac{\muz\omega^2}{4\pi c} \frac{\sin\theta}{r} \, p_0 \ee^{\ii(kr - \omega t)} \, \unitvec{\phi} = \frac{\vecE}{c}
\end{align}
其中 $\theta$ 是 $\vect{r}$ 与偶极轴（$z$ 轴）的夹角，$k = \omega/c$。
\end{theorem}

\begin{insight}{电偶极辐射的特征}{dipole-rad-features-v2}
\begin{itemize}
\item 场按 $1/r$ 衰减——比静电场（$1/r^2$）慢得多，因为辐射能量要覆盖越来越大的球面
\item $\sin\theta$ 角分布：在偶极轴方向（$\theta = 0, \pi$）辐射为零，在赤道面（$\theta = \pi/2$）最强
\item $\vecE$ 沿 $\unitvec{\theta}$ 方向，$\vecB$ 沿 $\unitvec{\phi}$ 方向，两者与传播方向 $\unitvec{r}$ 互相垂直
\item 坡印廷矢量 $S \propto \sin^2\theta/r^2$，对球面积分后与 $r$ 无关——能量守恒
\end{itemize}
\end{insight}

% ---------- 电偶极辐射角分布图 ----------
\begin{center}
\begin{tikzpicture}[scale=1.5]
  % Polar grid
  \draw[gray!30] (0,0) circle (1);
  \draw[gray!30] (0,0) circle (2);
  \draw[gray!30] (0,0) -- (2.5,0);
  \draw[gray!30] (0,0) -- (0,2.5);
  \draw[gray!30] (0,0) -- (-2.5,0);
  \draw[gray!30] (0,0) -- (0,-2.5);

  % sin^2(theta) 花生形：极坐标 r = 2 sin^2(theta)，x = r sin(theta)（对水平轴的夹角），y = r cos(theta)
  % 完整 0-360 度，上下对称双瓣
  \draw[red, very thick, domain=0:360, samples=240, smooth] plot ({2*sin(\x)*sin(\x)*sin(\x)}, {2*sin(\x)*sin(\x)*cos(\x)});

  % Dipole axis（z 轴向上，p 标注在轴右侧错开）
  \draw[->, thick, blue] (0,-2.5) -- (0,2.5) node[above] {$z$（偶极轴）};
  \node[blue] at (0.35,2.0) {$\vect{p}$};

  % Labels
  \node[red] at (1.35,0.25) {$\sin^2\theta$};
  \node at (0,-2.85) {\small 赤道面最强，轴方向为零};

  % Angular labels（θ=0 移到轴左侧，避开 z 标签）
  \node[anchor=west] at (2.35,0.12) {$\theta=90°$};
  \node[anchor=east,font=\small] at (-0.15,2.62) {$\theta=0$};
\end{tikzpicture}
\end{center}


\subsection{磁偶极辐射}

对于时变电流分布，除电偶极辐射外，还存在磁偶极辐射。设磁偶极矩为 $\vect{m}(t) = m_0 \ee^{-\ii\omega t} \unitvec{z}$（例如小电流环，$\vect{m} = I\vect{S}$）。

\begin{theorem}{磁偶极辐射场}{mag-dipole-rad-v2}
磁偶极矩 $\vect{m}(t)$ 产生的辐射场（远场区）：
\begin{align}
\vecE &= \frac{\muz k\omega}{4\pi r} \, (\unitvec{r} \times \vect{m}) \, \ee^{\ii(kr - \omega t)} = \frac{\muz\omega^2}{4\pi c r} \, (\unitvec{r} \times \vect{m}) \, \ee^{\ii(kr - \omega t)} \\
\vecB &= \frac{1}{c}\,\unitvec{r} \times \vecE = \frac{\muz k^2}{4\pi r} \, \unitvec{r} \times (\unitvec{r} \times \vect{m}) \, \ee^{\ii(kr - \omega t)}
\end{align}
其中用到 $k = \omega/c$，保证 $\vecE = c\vecB$（真空辐射场关系）。
\end{theorem}

\begin{insight}{电偶极与磁偶极辐射的对比}{dipole-compare-v2}
\begin{center}
\small
\begin{tabular}{@{}p{3.0cm}p{4.6cm}p{4.6cm}@{}}
\toprule
 & \textbf{电偶极辐射} & \textbf{磁偶极辐射} \\
\midrule
辐射源 & 时变电偶极矩 $\vect{p}(t)$（如振荡电荷对） & 时变磁偶极矩 $\vect{m}(t)$（如小电流环） \\
\midrule
远场 $\vecE$ & 沿 $\unitvec{\theta}$ 方向 & 沿 $\unitvec{\phi}$ 方向 \\
远场 $\vecB$ & 沿 $\unitvec{\phi}$ 方向 & 沿 $\unitvec{\theta}$ 方向 \\
\midrule
角分布 & $\dfrac{\diff P}{\diff\Omega}\propto\sin^2\theta$ & $\dfrac{\diff P}{\diff\Omega}\propto\sin^2\theta$（相同） \\
\midrule
总功率 & $P = \dfrac{\muz p_0^2\omega^4}{12\pi c}$ & $P = \dfrac{\muz m_0^2\omega^4}{12\pi c^3}$ \\
\midrule
极化平面 & $\vecE$ 在含偶极轴的平面内振动 & $\vecE$ 在垂直于偶极轴的平面内振动 \\
\midrule
典型物理源 & 半波天线、原子电偶极跃迁 & 小电流环天线、原子磁偶极跃迁 \\
\bottomrule
\end{tabular}
\end{center}

\textbf{易错点}：两者角分布相同但\textbf{极化方向旋转了 $90^\circ$}。强度上磁偶极辐射通常比同尺寸的电偶极辐射弱 $(v/c)^2$ 量级——这就是为什么实用天线几乎都是电偶极型（如半波天线），磁偶极天线只在小尺寸（如 NFC 线圈）或特殊场合使用。

\textbf{适用条件}：$P_M/P_E \sim (v/c)^2$ 这一比例\textbf{只对相同几何尺度的源}成立（推导中用到 $m_0 \sim qva$、$p_0 \sim qa$，即电流环与电荷对尺度相同）。如果电流环很大或电荷分布尺度不同，该比例不再成立——例如大环天线配合谐振匹配可以有很高的效率。记忆结论时务必连同适用条件一起记。

\textbf{学习预警}：本小节公式推导用到推迟势展开（$\ee^{\ii kr}$ 的多极展开），第一次阅读可以只记结论表，推导细节留到复习时再看。
\end{insight}

\toggleButton{挑战题}{ocg挑战5}
\begin{ocg}{挑战}{ocg挑战5}{off}
\begin{example}{小电流环的辐射}{small-loop-rad-v2}
\difficulty{挑战}
一个半径为 $a$ 的小电流环，电流 $I(t) = I_0 \ee^{-\ii\omega t}$，且 $a \ll \lambda$。求其磁偶极辐射的总功率。
\tcblower
\textbf{解答：}

磁偶极矩振幅：$m_0 = I_0 \pi a^2$。

代入磁偶极辐射功率公式：
\[
P = \frac{\muz m_0^2 \omega^4}{12\pi c^3} = \frac{\muz (I_0 \pi a^2)^2 \omega^4}{12\pi c^3} = \frac{\muz \pi I_0^2 a^4 \omega^4}{12 c^3}
\]
\end{example}
\end{ocg}

\subsection{辐射功率}

\begin{theorem}{拉莫尔公式}{larmor-v2}
加速运动的点电荷的辐射功率：
\[
P = \frac{q^2 a^2}{6\pi\eps_0 c^3}
\]
其中 $a$ 是加速度。对非相对论性粒子成立。

[量纲校验] $[q^2/(\eps_0 c^3)] = \mathrm{C}^2 \cdot (\mathrm{N\cdot m^2/C^2}) / (\mathrm{m^3/s^3}) = \mathrm{N\cdot s^2/m} = \mathrm{kg}$，乘 $[a^2] = \mathrm{m^2/s^4}$ 得 $\mathrm{kg\cdot m^2/s^3} = \mathrm{W}$。\checkmark
\end{theorem}

（相对论带电粒子的辐射公式见第10章。）

对时谐电偶极子，平均辐射功率：
\[
\langle P \rangle = \frac{\mu_0 p_0^2 \omega^4}{12\pi c}
\]

辐射功率与频率的\textbf{四次方}成正比——这就是为什么天空是蓝色的（瑞利散射，$\omega^4$ 依赖）。

\toggleButton{挑战题}{ocg挑战6}
\begin{ocg}{挑战}{ocg挑战6}{off}
\begin{example}{LC振荡电路的辐射}{lc-radiation-ex-v2}
\difficulty{挑战}
一个LC电路，电感 $L$，电容 $C$，振荡频率 $\omega = 1/\sqrt{LC}$。电荷振幅 $Q_0$。估算其作为电偶极天线的辐射功率。
\tcblower
\textbf{解答：}

电偶极矩振幅：$p_0 = Q_0 d$，其中 $d$ 是电容板间距。

平均辐射功率：
\[
\langle P \rangle = \frac{\mu_0 (Q_0 d)^2 \omega^4}{12\pi c} = \frac{\mu_0 Q_0^2 d^2}{12\pi c (LC)^2}
\]

电路储存的能量：$U = Q_0^2/(2C)$

辐射导致的品质因数下降：$Q_{\text{rad}} = \omega U / \langle P \rangle$

对于普通电路参数，这个辐射功率极小——但这就是无线电天线的工作原理！
\end{example}
\end{ocg}

\subsection*{拓展：李纳--维谢尔势简介}
\addcontentsline{toc}{subsection}{拓展：李纳--维谢尔势简介}

\begin{tip}{李纳--维谢尔势（进阶选学）}{lienard-wiechert-tip}
拉莫尔公式只适用于\textbf{非相对论}粒子（$v \ll c$）。当源电荷高速运动时，必须使用\textbf{李纳--维谢尔势}（Liénard--Wiechert potentials）——它是推迟势对"运动点电荷"的严格解：
\begin{align}
\phi(\vect{r},t) &= \frac{q}{4\pi\eps_0}\left[\frac{1}{(1 - \vect{n}\cdot\bm{\beta})R}\right]_{t_r} \\
\vecA(\vect{r},t) &= \frac{\bm{\beta}}{c}\,\phi(\vect{r},t)
\end{align}
其中 $\bm{\beta} = \vect{v}/c$ 是源电荷的无量纲速度，$\vect{n}$ 是从源点指向场点的单位矢量，方括号 $[\cdots]_{t_r}$ 表示在推迟时刻 $t_r$ 取值。与静止推迟势的关键区别在于分母中的\textbf{多普勒因子} $(1-\vect{n}\cdot\bm{\beta})$——源朝向观察者运动时辐射被"聚束"增强，这是同步辐射前向峰值强、侧向弱的根源。

由该势可导出\textbf{李纳公式}（拉莫尔公式的相对论推广）：
\[
P = \frac{q^2\gamma^6}{6\pi\eps_0 c^3}\left(a^2 - \frac{|\vect{v}\times\vect{a}|^2}{c^2}\right)
\]
低速极限 $v\to 0$（$\gamma\to 1$，$|\vect{v}\times\vect{a}|/c \to 0$）时退化为拉莫尔公式。

\textbf{学习建议}：本小节为选学内容，第一次阅读只需记住两点——(1) 推迟势对点电荷有严格解，分母多一个多普勒因子；(2) 相对论修正使辐射功率最多放大 $\gamma^6$ 倍（纵向加速度时），同步辐射光源、自由电子激光都源于此。完整推导见 Jackson 第 14 章。
\end{tip}

\begin{review}{第9章要点回顾}{ch9-review}
\begin{center}
\begin{tabular}{@{}ll@{}}
\toprule
\textbf{公式/概念} & \textbf{内容} \\
\midrule
$t_r = t - R/c$ & 推迟时间 \\
$\vecE \propto \sin\theta/r$ & 偶极辐射场 \\
$P = q^2a^2/(6\pi\eps_0c^3)$ & 拉莫尔公式 \\
$P \propto \omega^4$ & 辐射功率频率依赖 \\
\bottomrule
\end{tabular}
\end{center}
\end{review}

\begin{tip}{本章核心框架}{ch9-summary}
\textbf{一句话总结}：加速运动的电荷辐射电磁波，辐射功率与加速度平方成正比。\\
\textbf{核心内容}：推迟势（$t_r=t-R/c$）、拉莫尔公式（$P\propto|\ddot{\vect{p}}|^2$）、电偶极辐射（$dP/d\Omega\propto\sin^2\theta$）、磁偶极辐射、天线方向性。\\
\textbf{关键概念}：远场区（辐射场$\propto 1/r$）、近场区（准静态）、推迟效应。\\
\textbf{衔接}：第10章将电动力学纳入狭义相对论框架，实现电磁理论的四维协变形式。
\end{tip}

% ---------- 第9章随堂自测 ----------
\section*{第9章随堂自测}
\addcontentsline{toc}{section}{第9章随堂自测}

\begin{miniquiz}
\textbf{一、选择题}
\begin{enumerate}[label=\arabic*.]
\item 电偶极辐射的功率与频率的关系是：
  \begin{enumerate}[label=(\Alph*),nosep]
  \item $P \propto \omega$
  \item $P \propto \omega^2$
  \item $P \propto \omega^4$
  \item $P \propto \omega^6$
  \end{enumerate}
\begin{ocg}{答案}{ocg答案9}{off}
  \textit{答案：C。拉莫尔公式 $P \propto |\ddot{\vect{p}}|^2 \propto \omega^4$。}
\end{ocg}

\item 推迟势 $\phi(\vect{r},t) = \frac{1}{4\pi\eps_0}\int \frac{\rho(\vect{r}',t_r)}{|\vect{r}-\vect{r}'|}\diff^3r'$ 中的 $t_r$ 表示：
  \begin{enumerate}[label=(\Alph*),nosep]
  \item 观测时刻
  \item 源点处的推迟时刻 $t - |\vect{r}-\vect{r}'|/c$
  \item 源点处的超前时刻
  \item 恒定时刻
  \end{enumerate}
\begin{ocg}{答案}{ocg答案9}{off}
  \textit{答案：B。推迟效应体现了电磁相互作用的有限传播速度。}
\end{ocg}

\item 天线辐射的方向性图描述了：
  \begin{enumerate}[label=(\Alph*),nosep]
  \item 辐射功率随时间的变化
  \item 辐射功率随方向的变化
  \item 辐射功率随频率的变化
  \item 辐射功率随距离的变化
  \end{enumerate}
\begin{ocg}{答案}{ocg答案9}{off}
  \textit{答案：B。方向性图$F(\theta,\phi)$描述辐射功率的角分布。}
\end{ocg}
\end{enumerate}

\vspace{0.5em}
\textbf{二、判断题}
\begin{enumerate}[label=\arabic*.,resume]
\item 静止的电荷不辐射电磁波。
\begin{ocg}{答案}{ocg答案9}{off}
  \textit{答案：正确。只有加速运动的电荷才辐射（拉莫尔公式要求$\ddot{\vect{p}}\neq 0$）。}
\end{ocg}

\item 磁偶极辐射的强度与电偶极辐射同数量级。
\begin{ocg}{答案}{ocg答案9}{off}
  \textit{答案：错误。磁偶极辐射通常比电偶极辐射弱$(v/c)^2$倍。}
\end{ocg}
\end{enumerate}

\vspace{0.5em}
\textbf{三、计算题}
\begin{enumerate}[label=\arabic*.,resume]
\item 一个半波天线近似为电偶极子，工作频率 $f = \SI{100}{MHz}$，偶极矩振幅 $p_0 = \SI{1e-9}{C.m}$。估算其平均辐射功率。

\begin{ocg}{答案}{ocg答案9}{off}
  \textit{解答：} $\omega = 2\pi f = 6.28\times 10^8 \,\si{rad/s}$。

  $\langle P \rangle = \dfrac{\muz p_0^2\omega^4}{12\pi c} = \dfrac{4\pi\times 10^{-7} \times (10^{-9})^2 \times (6.28\times 10^8)^4}{12\pi \times 3\times 10^8} = \dfrac{1.96\times 10^{11}}{1.13\times 10^{10}} \approx \SI{17}{W}$

  \textit{量纲校验}：$[\mu_0 p_0^2 \omega^4/c] = (\mathrm{H/m})(\mathrm{C\cdot m})^2\mathrm{s}^{-4}/(\mathrm{m/s}) = \mathrm{W}$。$\checkmark$

  \textit{物理解读}：按球面扩散，$1$ km 处能流密度 $\langle S\rangle = \langle P\rangle/(4\pi r^2) \approx 1.4\times10^{-6}\,\si{W/m^2}$，仍远高于典型接收机灵敏度（$10^{-15}\sim10^{-13}\,\si{W/m^2}$）——这正是 100 MHz 波段通信覆盖半径可达数十公里的原因；真正苛刻的是深空通信：距离增到 $10^{6}$ km 量级时 $\langle S\rangle$ 才降到 $10^{-18}\,\si{W/m^2}$ 上下。
\end{ocg}

\item 电偶极子 $\vect{p}(t) = p_0\ee^{-\ii\omega t}\unitvec{z}$ 沿 $z$ 轴放置。求辐射最强的方向和辐射为零的方向，并说明 $\theta = \pi/2$ 处 $\vecE$ 的振动方向。

\begin{ocg}{答案}{ocg答案9}{off}
  \textit{解答：} 角分布 $dP/d\Omega \propto \sin^2\theta$。

  最强：$\theta = \pi/2$（赤道面，垂直于偶极轴）；为零：$\theta = 0, \pi$（沿偶极轴）。

  $\theta = \pi/2$ 处 $\vecE$ 沿 $\unitvec{\theta}$ 方向——在赤道面上 $\unitvec{\theta} = -\unitvec{z}$，即电场\textbf{平行于偶极轴}振动。这就是半波天线在赤道面接收信号最强的偏振特性来源。
\end{ocg}
\end{enumerate}
\end{miniquiz}
\toggleButton{显示答案}{ocg答案9}

\begin{chapterreview}
\begin{itemize}[nosep]
\item \textbf{核心内容}：推迟势、拉莫尔公式、电偶极辐射、磁偶极辐射、天线方向性。
\item \textbf{关键公式}：拉莫尔公式 $P = \mu_0|\ddot{\vect{p}}|^2/(6\pi c)$；推迟时间 $t_r = t - R/c$。
\item \textbf{自检验}：能否从推迟势出发推导电偶极辐射场？
\item \textbf{衔接预告}：第10章将用电动力学为狭义相对论奠基——四维协变形式。
\end{itemize}
\end{chapterreview}

\begin{tip}{本章自检清单}{ch9-checklist}
学完本章后对照下列条目逐条自测，全部通过即可进入下一章；有未通过项时，回到对应小节复习：
\begin{itemize}[label=$\square$, nosep]
\item 写出推迟势并解释推迟时间的物理意义
\item 用拉莫尔公式计算辐射功率
\item 说明电偶极辐射的角分布与偏振
\item 对比电偶极与磁偶极辐射
\end{itemize}
\end{tip}

% ---------- 第9章综合练习题 ----------
\section*{第9章综合练习题}
\addcontentsline{toc}{section}{第9章综合练习题}
\begin{enumerate}[label=\textbf{\arabic*.}]
\item \textbf{基础} 一个振荡电偶极子 $p(t)=p_0\cos\omega t$，求远区辐射场 $\vecE$、$\vecB$ 的平均辐射功率角分布。
\item \textbf{基础} 一个小圆电流环（半径 $a$，电流 $I=I_0\cos\omega t$）的磁偶极矩是多少？写出其远区辐射场的表达式。
\item \textbf{提高} 比较电偶极和磁偶极辐射的角分布：对于相同频率和相同尺度，两者的辐射功率比是多少？（提示：$m\sim Ia^2$，$p\sim qa$）
\item \textbf{提高} 证明推迟势 $G(\vect{r},t)=\frac{\delta(t-r/c)}{4\pi r}$ 是达朗贝尔算子的格林函数，即满足 $\left(\lap-\frac{1}{c^2}\frac{\partial^2}{\partial t^2}\right)G=-\delta^{(3)}(\vect{r})\delta(t)$。
\item \textbf{挑战} 一个非相对论电子在均匀磁场中做圆周运动，用拉莫尔公式计算其同步辐射功率，并讨论辐射频率与回旋频率的关系。
\end{enumerate}
\newpage
\section*{习题解答}
\addcontentsline{toc}{section}{习题解答}

\textbf{题1解析：}

[分析] 本题推导振荡电偶极子的远区辐射场和平均辐射功率，是电偶极辐射的标准结果。

[解答] 时变电偶极矩 $\vect{p}(t) = p_0\cos\omega t\,\unitvec{z}$。

远区辐射场（$r \gg \lambda$）：
\begin{align*}
\vecE &= \frac{k^2}{4\pi\eps_0}\frac{\sin\theta}{r}\,p_0\cos(kr - \omega t)\,\unitvec{\theta} \\
\vecB &= \frac{\muz\omega^2}{4\pi c}\frac{\sin\theta}{r}\,p_0\cos(kr - \omega t)\,\unitvec{\phi} = \frac{\vecE}{c}
\end{align*}

坡印廷矢量：$S = \dfrac{1}{\muz}EB = \dfrac{\muz p_0^2\omega^4\sin^2\theta}{16\pi^2 c r^2}\cos^2(kr - \omega t)$。

平均辐射功率角分布：$\langle dP/d\Omega \rangle = \dfrac{\muz p_0^2\omega^4\sin^2\theta}{32\pi^2 c}$。

总平均功率：$\langle P \rangle = \dfrac{\muz p_0^2\omega^4}{12\pi c}$。\qedsymbol

[量纲校验] $[P] = [W]$，$[\muz p_0^2\omega^4/c] = [H/m][C^2L^2][T^{-4}]/[LT^{-1}] = [ML^2T^{-3}] = [W]$，一致。\qedsymbol

\textbf{题2解析：}

[分析] 本题计算小圆电流环的磁偶极矩和辐射场，是磁偶极辐射的基本公式。

[解答] 小圆电流环（$a \ll \lambda$）的磁偶极矩：
\[
\vect{m}(t) = I(t)\pi a^2\,\unitvec{n} = I_0\pi a^2\cos\omega t\,\unitvec{n}
\]

远区辐射场（磁偶极辐射）：
\begin{align*}
E_\phi &= -\frac{\muz m_0\omega^2}{4\pi c r}\sin\theta\,\sin(kr - \omega t) \\
B_\theta &= -\frac{\muz m_0\omega^2}{4\pi c^2 r}\sin\theta\,\sin(kr - \omega t)
\end{align*}

其中 $m_0 = I_0\pi a^2$。\qedsymbol

[量纲校验] $[m] = [A][L^2]$，$[E_\phi] = [H/m][A L^2][T^{-2}]/[L T^{-1} L] = [MLT^{-3}I^{-1}] = [V/m]$，一致。\qedsymbol

\textbf{题3解析：}

[分析] 本题比较电偶极和磁偶极辐射的相对强度，是理解两种辐射机制的关键。

[解答] 电偶极辐射功率：$P_E = \dfrac{\muz p_0^2\omega^4}{12\pi c}$。

磁偶极辐射功率：$P_M = \dfrac{\muz m_0^2\omega^4}{12\pi c^3}$。

功率比：$\dfrac{P_M}{P_E} = \dfrac{m_0^2}{c^2 p_0^2}$。

对于相同尺度 $a$ 和相同频率：$p_0 \sim qa$，$m_0 \sim Ia^2$。
电流 $I \sim qv/a$（电荷运动），$m_0 \sim qa v$。

$\dfrac{P_M}{P_E} \sim \dfrac{(qav)^2}{c^2(qa)^2} = \dfrac{v^2}{c^2} \ll 1$（非相对论情形）。

物理意义：磁偶极辐射比电偶极辐射弱 $(v/c)^2$ 因子，在低速极限下可忽略。\qedsymbol

[量纲校验] $P_M/P_E$ 无量纲，$v^2/c^2$ 无量纲，一致。\qedsymbol

\textbf{题4解析：}

[分析] 本题验证推迟格林函数满足达朗贝尔方程，是辐射理论的基础。

[解答] 格林函数 $G(\vect{r},t) = \dfrac{\delta(t - r/c)}{4\pi r}$。

计算 $\lap G$：对 $r > 0$，$\lap(1/r) = 0$。但 $r = 0$ 处需要处理 $\delta$ 函数。

$\lap G = \dfrac{1}{4\pi r}\lap\delta(t - r/c) + 2\grad\left(\dfrac{1}{4\pi r}\right)\cdot\grad\delta(t - r/c) + \delta(t - r/c)\lap\left(\dfrac{1}{4\pi r}\right)$。

利用 $\lap(1/r) = -4\pi\delta^{(3)}(\vect{r})$：

$\lap G = \dfrac{1}{4\pi r}\dfrac{\delta''(t - r/c)}{c^2} - \dfrac{2}{4\pi r^2 c}\delta'(t - r/c) + \dfrac{2}{4\pi r^2 c}\delta'(t - r/c) - \delta^{(3)}(\vect{r})\delta(t - r/c)$（两个含 $\delta'$ 的项相互抵消）。

时间导数：$\partial_t^2 G = \delta''(t - r/c)/(4\pi r)$。

达朗贝尔算子：$\lap G - \dfrac{1}{c^2}\partial_t^2 G = -\delta^{(3)}(\vect{r})\delta(t)$。\qedsymbol

[量纲校验] $[G] = [L^{-1}T^{-1}]$，$[\lap G] = [L^{-3}T^{-1}]$，$[\delta^{(3)}(\vect{r})\delta(t)] = [L^{-3}T^{-1}]$，一致。\qedsymbol

\textbf{题5解析：}

[分析] 本题用拉莫尔公式计算非相对论电子在磁场中的同步辐射功率。

[解答] 非相对论电子在均匀磁场 $B$ 中做圆周运动，回旋频率 $\omega_c = eB/m$。

向心加速度：$a = v\omega_c = eBv/m$。

拉莫尔公式：$P = \dfrac{e^2 a^2}{6\pi\eps_0 c^3} = \dfrac{e^4 B^2 v^2}{6\pi\eps_0 m^2 c^3}$。

辐射频率：非相对论电子做匀速圆周运动，其偶极矩按单一频率 $\omega_c$ 振荡（$|\vect{v}|$ 恒定），辐射谱就是\textbf{单频} $\omega_c$，不存在高次谐波——谐波只出现在相对论情形（同步辐射的聚束效应）或磁场非均匀时，本题范围内不涉及。\qedsymbol

[量纲校验] $[P] = [W]$，$[e^2 a^2/(\eps_0 c^3)] = [C^2][L^2T^{-4}]/([F/m][L^3T^{-3}]) = [ML^2T^{-3}] = [W]$，一致。\qedsymbol



% ============================================================

\begin{center}\small \hyperlink{ch8-goal}{上一章：第8章} \quad|\quad \hyperlink{ch10-goal}{下一章：第10章}\end{center}

% 第10章：相对论电动力学（优化版）
% ============================================================
\chapter{相对论电动力学}
\epigraph{时间和空间不是绝对的，但物理定律是。}{--- 阿尔伯特·爱因斯坦}

\hypertarget{ch10-goal}{}\begin{insight}{本章核心目标}{ch10-goal-v2}
将电动力学纳入狭义相对论框架，揭示电场和磁场是同一物理实体——电磁场张量——的不同表现。理解洛伦兹变换下的场变换规律，掌握四维形式的麦克斯韦方程组。
\end{insight}

1905年，爱因斯坦的狭义相对论彻底改变了我们对时空的理解。而电动力学与相对论的关系尤为紧密——事实上，爱因斯坦创立相对论的动机之一就是解决电动力学中的矛盾。

\begin{history}{相对论与电动力学}{relativity-history}
麦克斯韦方程组预言电磁波以光速 $c$ 传播——但相对于什么参考系？19 世纪物理学家假设"以太"作为载体，然而 1887 年\textbf{迈克尔逊-莫雷实验}未能检测到地球相对以太的运动。\textbf{爱因斯坦}在 1905 年给出根本解答：物理定律在所有惯性系中形式相同、光速不变，不需要以太——电磁场的洛伦兹变换自然地包含其中。电动力学\textbf{需要}相对论才能自洽。
\end{history}

\section{四维记号基础}

在学习电磁场的相对论变换之前，我们需要先建立四维时空的数学框架。本节是后续所有协变表述的基石。

\begin{definition}{闵氏度规与四维坐标}{minkowski-metric-v2}
定义\textbf{四维坐标}（逆变形式）：
\[
x^\mu = (ct, x, y, z) = (ct, \vect{r})
\]
其中 $\mu = 0, 1, 2, 3$。时空间隔定义为：
\[
s^2 = \eta_{\mu\nu} x^\mu x^\nu = c^2t^2 - x^2 - y^2 - z^2
\]
\textbf{闵氏度规}采用 $(+,-,-,-)$ 号差：
\[
\eta_{\mu\nu} = \mathrm{diag}(1, -1, -1, -1)
\]
\end{definition}

\begin{tip}{前置小复习：升降指标，先用数字算一遍}{index-mini-review}
协变/逆变是本节最大的拦路虎。先忘掉抽象定义，用\textbf{具体数字}算一遍：

\textbf{第一步}：四维坐标上标（逆变）：$x^\mu = (2, 3, 4, 5)$（即 $ct=2, x=3, y=4, z=5$，单位略）。

\textbf{第二步}：度规 $\eta_{\mu\nu} = \mathrm{diag}(1,-1,-1,-1)$ 是一台"翻号机"——\textbf{时间分量不动，空间分量翻号}。降指标：
\[
x_\mu = \eta_{\mu\nu}x^\nu = (2, -3, -4, -5)
\]
就这一步。上标变下标 = 空间分量乘 $-1$。

\textbf{第三步}：验算不变量（时空间隔的平方）：
\[
x^\mu x_\mu = 2\times 2 + 3\times(-3) + 4\times(-4) + 5\times(-5) = 4 - 9 - 16 - 25 = -46
\]
一个上标配一个下标求和（Einstein 约定），结果与度规正负号约定一致。

\textbf{记住三条}：(1) 上标逆变、下标协变；(2) 同一对象升降指标只差空间分量的正负号；(3) 求和永远一上一下。下面的正式定义只是这套操作的符号化。
\end{tip}

\begin{insight}{协变与逆变指标}{cov-contra-v2}
上标 $x^\mu$ 称为\textbf{逆变}分量，下标 $x_\mu = \eta_{\mu\nu} x^\nu$ 称为\textbf{协变}分量：
\[
x_\mu = (ct, -x, -y, -z)
\]
 Einstein求和约定：重复的上、下指标自动求和。例如 $a_\mu b^\mu = a_0 b^0 + a_1 b^1 + a_2 b^2 + a_3 b^3$。
\end{insight}

\begin{warning}{号差选择}{signature-warning}
本书\textbf{统一采用 $(+,-,-,-)$ 号差}。这是粒子物理和量子场论中的主流约定。但需注意：许多广义相对论教材（如温伯格、卡罗尔）采用 $(-,+,+,+)$ 号差。两者在物理上等价，但公式中的符号会有差异。阅读不同教材时务必先确认其号差约定，否则极易在符号上出错！
\end{warning}

\begin{example}{升降指标运算}{index-raising-lowering-v2}
\difficulty{基础}
已知 $x^\mu=(ct,x,y,z)$，$\eta_{\mu\nu}=\mathrm{diag}(1,-1,-1,-1)$。计算协变坐标 $x_\mu$，验证时空间隔不变量 $x^\mu x_\mu=s^2$，并给出梯度算符的升降指标形式。
\tcblower
\textbf{解答：}

\textbf{步骤1：降指标得到 $x_\mu$}

用度规降指标：
\[
x_\mu = \eta_{\mu\nu}x^\nu
\]
逐分量计算（Einstein求和约定，重复上、下指标求和）：
\begin{align*}
x_0 &= \eta_{00}x^0 + \eta_{01}x^1 + \eta_{02}x^2 + \eta_{03}x^3 = 1\cdot ct + 0 + 0 + 0 = ct \\
x_1 &= \eta_{10}x^0 + \eta_{11}x^1 + \eta_{12}x^2 + \eta_{13}x^3 = 0 + (-1)\cdot x + 0 + 0 = -x \\
x_2 &= \eta_{20}x^0 + \eta_{21}x^1 + \eta_{22}x^2 + \eta_{23}x^3 = -y \\
x_3 &= -z
\end{align*}
因此：
\[
x_\mu = (ct, -x, -y, -z)
\]

\textbf{步骤2：验证不变量 $x^\mu x_\mu = s^2$}

\[
x^\mu x_\mu = x^0 x_0 + x^1 x_1 + x^2 x_2 + x^3 x_3 = ct\cdot ct + x\cdot(-x) + y\cdot(-y) + z\cdot(-z) = c^2t^2 - x^2 - y^2 - z^2 = s^2
\]
这正是闵氏时空间隔，在所有惯性系中保持不变。

\textbf{步骤3：梯度算符的升降指标}

协变梯度定义为 $\partial_\mu = \partial/\partial x^\mu$，因此：
\[
\partial_\mu = \left(\frac{1}{c}\frac{\partial}{\partial t}, \frac{\partial}{\partial x}, \frac{\partial}{\partial y}, \frac{\partial}{\partial z}\right) = \left(\frac{\partial_t}{c}, \nabla\right)
\]
用逆度规 $\eta^{\mu\nu}=\mathrm{diag}(1,-1,-1,-1)$（度规是自逆的）升指标：
\[
\partial^\mu = \eta^{\mu\nu}\partial_\nu = \left(\frac{\partial_t}{c}, -\nabla\right)
\]
于是达朗贝尔算子可写成：
\[
\Box = \lap - \frac{1}{c^2}\frac{\partial^2}{\partial t^2} = -\partial_\mu\partial^\mu
\]
\end{example}

\begin{example}{电磁场张量分量计算}{field-tensor-components-v2}
\difficulty{提高}
已知四维电磁势 $A^\mu=(\phi/c, A_x, A_y, A_z)$，电磁场张量定义为 $F^{\mu\nu}=\partial^\mu A^\nu - \partial^\nu A^\mu$。计算各分量，并验证其与 $\vecE$、$\vecB$ 的关系。
\tcblower
\textbf{解答：}

\textbf{步骤1：计算 $F^{0i}$（时间-空间分量）}

取 $\mu=0, \nu=i$（$i=1,2,3$对应 $x,y,z$）：
\[
F^{0i} = \partial^0 A^i - \partial^i A^0 = \frac{1}{c}\frac{\partial A_i}{\partial t} + \frac{1}{c}\frac{\partial\phi}{\partial x_i} = -\frac{E_i}{c}
\]
其中用到了 $\vecE=-\nabla\phi-\partial\vecA/\partial t$ 的分量形式 $E_i = -\partial_i\phi - \partial_t A_i$，故 $\partial_t A_i + \partial_i\phi = -E_i$。注意 $\partial^i = -\partial_i$（度规 $(+,-,-,-)$ 升指标），得 $F^{0i} = -E_i/c$——与 $F^{\mu\nu}$ 矩阵第一行 $F^{01} = -E_x/c$ 等一致。

\textbf{步骤2：计算 $F^{ij}$（空间-空间分量）}

取 $\mu=i, \nu=j$（$i,j=1,2,3$）：
\[
F^{ij} = \partial^i A^j - \partial^j A^i = -\frac{\partial A_j}{\partial x_i} + \frac{\partial A_i}{\partial x_j} = -\varepsilon_{ijk}B^k
\]
其中 $B^k = (\nabla\times\vecA)_k = \varepsilon_{klm}\partial_l A_m$。

具体验证：
\[
F^{12} = -\partial_x A_y + \partial_y A_x = -(\nabla\times\vecA)_z = -B_z
\]
同理 $F^{23}=-B_x$，$F^{31}=-B_y$。由反对称性 $F^{ij}=-F^{ji}$。

\textbf{步骤3：显式矩阵形式}

综合以上结果，电磁场张量的逆变形式为：
\[
F^{\mu\nu} = \begin{pmatrix}
0 & -E_x/c & -E_y/c & -E_z/c \\
E_x/c & 0 & -B_z & B_y \\
E_y/c & B_z & 0 & -B_x \\
E_z/c & -B_y & B_x & 0
\end{pmatrix}
\]
这正是定义\ref{def:field-tensor-v2}中给出的矩阵形式。可见，电磁场张量把六个独立分量（$E_x,E_y,E_z$ 与 $B_x,B_y,B_z$）统一编码为一个紧凑的反对称四维二阶张量。
\end{example}

\section{洛伦兹变换}

\begin{theorem}{洛伦兹变换}{lorentz-transform-v2}
设参考系 $S'$ 以速度 $\vect{v} = v\unitvec{x}$ 相对于 $S$ 运动。则时空坐标变换为：
\begin{align}
t' &= \gamma\left(t - \frac{vx}{c^2}\right) \\
x' &= \gamma(x - vt) \\
y' &= y \\
z' &= z
\end{align}
其中 $\gamma = 1/\sqrt{1 - v^2/c^2}$ 是\textbf{洛伦兹因子}。
\end{theorem}

\section{四维矢量与张量}

\subsection{四维时空}

引入\textbf{固有时} $\tau$ 和\textbf{四维速度}：
\[
u^\mu = \frac{\diff x^\mu}{\diff\tau} = \gamma(c, v_x, v_y, v_z)
\]

\subsection{电磁场张量}

\begin{definition}{电磁场张量}{field-tensor-v2}
\textbf{电磁场张量} $F^{\mu\nu}$ 是反对称的二阶张量：
\[
F^{\mu\nu} = \begin{pmatrix}
0 & -E_x/c & -E_y/c & -E_z/c \\
E_x/c & 0 & -B_z & B_y \\
E_y/c & B_z & 0 & -B_x \\
E_z/c & -B_y & B_x & 0
\end{pmatrix}
\]
\end{definition}

\begin{theorem}{电磁场张量的洛伦兹不变量}{lorentz-invariants-v2}
电磁场张量 $F^{\mu\nu}$ 的两个独立洛伦兹不变量可构造如下：
\begin{align}
I_1 &= \frac{1}{2}F_{\mu\nu}F^{\mu\nu} = B^2 - \frac{E^2}{c^2} \label{eq:invariant1} \\
I_2 &= \frac{1}{8}\epsilon_{\mu\nu\lambda\sigma}F^{\mu\nu}F^{\lambda\sigma} = \frac{\vecE\cdot\vecB}{c} \label{eq:invariant2}
\end{align}
其中 $\epsilon_{\mu\nu\lambda\sigma}$ 是四维 Levi-Civita 符号（$\epsilon_{0123}=+1$），在洛伦兹变换下按赝张量规律变换，保证 $I_2$ 为真标量。

\textbf{符号约定（全书统一）}：度规 $(+,-,-,-)$、$A^\mu = (\phi/c, \vecA)$ 下 $F^{0i} = -E_i/c$、$F^{ij} = -\varepsilon_{ijk}B_k$；由此 $F_{\mu\nu}F^{\mu\nu} = 2(B^2 - E^2/c^2)$、$\frac12 F_{\mu\nu}F^{\mu\nu} = B^2 - E^2/c^2$。按上述约定逐项展开得 $\epsilon_{\mu\nu\rho\sigma}F^{\mu\nu}F^{\rho\sigma} = +8\vecE\cdot\vecB/c$（注意：Jackson 取 $F^{0i}=+E_i$，同一运算给出 $-8\vecE\cdot\vecB/c$，切勿混用），本书 $I_2$ 按上式定义（系数 $1/8$）使 $I_2 = \vecE\cdot\vecB/c$。注意：其他教材（如 Griffiths）可能定义 $I_1$ 不带 $1/2$ 系数或 $I_2$ 反号——物理结论（$I_1$ 判场型、$I_2$ 判正交性）不受影响，对照阅读时以定义式为准。
\end{theorem}

\begin{insight}{两个不变量的物理意义}{invariant-physics-v2}
$I_1$ 和 $I_2$ 在所有惯性系中取值相同，这是电磁场最深刻的内禀性质之一。它们的符号和大小决定了场在不同参考系中的"可消除性"：
\begin{itemize}
\item \textbf{当 $I_1<0$（即 $E>cB$）}：场以电场为主导。总可以找到一个洛伦兹参考系，使得其中 $B=0$（纯电场）。这对应于 $E>cB$ 的类时场型。
\item \textbf{当 $I_1>0$（即 $B>E/c$）}：场以磁场为主导。总可以找到一个洛伦兹参考系，使得其中 $E=0$（纯磁场）。这对应于 $B>E/c$ 的类空场型（与上一条“类时场型”配对）。
\item \textbf{当 $I_2=0$（即 $\vecE\perp\vecB$）}：电场与磁场互相垂直。由于 $I_2$ 是洛伦兹不变量，若在某一个参考系中 $\vecE\perp\vecB$，则在\textbf{任何}惯性参考系中二者都保持垂直。
\item \textbf{当 $I_2\neq 0$（即 $\vecE$ 与 $\vecB$ 不垂直）}：总可以找到参考系使二者平行（$\vecE\parallel\vecB$）。当 $I_1<0$ 时，平行系中 $E>cB$；当 $I_1>0$ 时，平行系中 $B>E/c$。
\end{itemize}

特别地，若 $I_1=0$ 且 $I_2=0$，则场为\textbf{类光型}（如平面电磁波）：在任何参考系中都有 $E=cB$ 且 $\vecE\perp\vecB$，不可能通过洛伦兹变换消除任何分量。
\end{insight}

\begin{insight}{电场和磁场的统一}{em-unification-v2}
电磁场张量 $F^{\mu\nu}$ 将电场和磁场统一为一个几何对象。在洛伦兹变换下，$F^{\mu\nu}$ 按张量规则变换，导致 $\vecE$ 和 $\vecB$ 的\textbf{混合}：
\begin{align*}
E'_x &= E_x, \quad B'_x = B_x \\
E'_y &= \gamma(E_y - vB_z), \quad B'_y = \gamma(B_y + vE_z/c^2) \\
E'_z &= \gamma(E_z + vB_y), \quad B'_z = \gamma(B_z - vE_y/c^2)
\end{align*}

这意味着：在一个参考系中纯是电场的场，在另一个参考系中会具有磁场分量，反之亦然。电场和磁场不是独立的实体，而是同一电磁场张量的不同"投影"。
\end{insight}

\begin{example}{纯电场在运动参考系中感应出磁场}{em-transform-ex-v2}
\difficulty{提高}
在参考系 $S$ 中，空间某区域只有均匀电场 $\vecE = E_0 \unitvec{y}$，$\vecB = \vect{0}$。求相对于 $S$ 以速度 $v\unitvec{x}$ 运动的参考系 $S'$ 中的电磁场。
\tcblower
\textbf{解答：}

使用电磁场的洛伦兹变换公式，代入 $E_y = E_0$，$E_x = E_z = 0$，$B_x = B_y = B_z = 0$：
\begin{align*}
E'_x &= E_x = 0 \\
E'_y &= \gamma(E_y - vB_z) = \gamma E_0 \\
E'_z &= \gamma(E_z + vB_y) = 0 \\
B'_x &= B_x = 0 \\
B'_y &= \gamma\left(B_y + \frac{vE_z}{c^2}\right) = 0 \\
B'_z &= \gamma\left(B_z - \frac{vE_y}{c^2}\right) = -\gamma \frac{vE_0}{c^2}
\end{align*}
因此：
\[
\vecE' = \gamma E_0\,\unitvec{y}, \qquad \vecB' = -\gamma \frac{v E_0}{c^2}\,\unitvec{z}
\]

\textbf{结论}：在一个参考系中纯粹的静电场，在\textbf{垂直于电场方向}运动的参考系中会表现出\textbf{非零的磁场}（$\vecB' \neq 0$）。注意：若 $S'$ 沿电场方向（$y$）运动，则电场不混合出磁场——只有\textbf{横向}运动才产生感应磁场。这直观地证明了 $\vecE$ 和 $\vecB$ 不是独立的物理实体，而是同一电磁场张量的不同分量。
\end{example}

\begin{example}{低速极限退化}{low-speed-limit-v2}
\difficulty{基础}
证明：当 $v\ll c$ 时，电磁场的洛伦兹变换公式退化为
\[
\vecE'\approx\vecE+\vect{v}\times\vecB, \quad \vecB'\approx\vecB-\frac{1}{c^2}\vect{v}\times\vecE
\]
并说明此结果与第3章运动介质中电磁场的非相对论形式的关系。
\tcblower
\textbf{解答：}

从电磁场的洛伦兹变换（设 $\vect{v}=v\unitvec{x}$，$v\ll c$）：
\begin{align*}
E'_x &= E_x \\
E'_y &= \gamma(E_y - vB_z) \\
E'_z &= \gamma(E_z + vB_y) \\
B'_x &= B_x \\
B'_y &= \gamma\left(B_y + \frac{vE_z}{c^2}\right) \\
B'_z &= \gamma\left(B_z - \frac{vE_y}{c^2}\right)
\end{align*}

当 $v\ll c$ 时，洛伦兹因子 $\gamma\approx 1 + \frac{v^2}{2c^2} + \cdots \approx 1$（保留到一阶），于是：
\begin{align*}
E'_y &\approx E_y - vB_z = E_y + (\vect{v}\times\vecB)_y \\
E'_z &\approx E_z + vB_y = E_z + (\vect{v}\times\vecB)_z
\end{align*}
而 $E'_x = E_x$，且注意到当 $\vect{v}$ 沿 $x$ 方向时 $(\vect{v}\times\vecB)_x = 0$，因此综合三个分量：
\[
\vecE' \approx \vecE + \vect{v}\times\vecB
\]

同理，对磁场分量：
\begin{align*}
B'_y &\approx B_y + \frac{vE_z}{c^2} = B_y - \frac{1}{c^2}(\vect{v}\times\vecE)_y \\
B'_z &\approx B_z - \frac{vE_y}{c^2} = B_z - \frac{1}{c^2}(\vect{v}\times\vecE)_z
\end{align*}
而 $B'_x = B_x$ 且 $(\vect{v}\times\vecE)_x=0$，因此：
\[
\vecB' \approx \vecB - \frac{1}{c^2}\vect{v}\times\vecE
\]

\textbf{物理意义：}这正是第3章运动介质中电磁场的非相对论形式。它表明：相对论的电磁场变换公式在低速极限下自动还原为经典结果——相对论电动力学不是对经典电动力学的否定，而是对其在更宽速度范围内的自然推广。这种"低速退化"是检验相对论理论自洽性的重要判据。
\end{example}

% ---------- 电磁场张量彩色分块图 ----------
\begin{center}
\begin{tikzpicture}[scale=1.1, every node/.style={font=\small}]
  % 格宽 1.3，避免 $-E_x/c$ 标签溢出
  \def\cw{1.3}
  % Matrix frame
  \draw[thick] (0,0) rectangle (4*\cw,4*\cw);

  % Grid lines
  \foreach \i in {1,2,3} {
    \draw[thick] (\i*\cw,0) -- (\i*\cw,4*\cw);
    \draw[thick] (0,\i*\cw) -- (4*\cw,\i*\cw);
  }

  % Red blocks for E components (first row and first col, excluding 00)
  \fill[red!20] (0,3*\cw) rectangle (\cw,4*\cw);
  \fill[red!20] (\cw,3*\cw) rectangle (2*\cw,4*\cw);
  \fill[red!20] (2*\cw,3*\cw) rectangle (3*\cw,4*\cw);
  \fill[red!20] (3*\cw,3*\cw) rectangle (4*\cw,4*\cw);
  \fill[red!20] (0,2*\cw) rectangle (\cw,3*\cw);
  \fill[red!20] (0,\cw) rectangle (\cw,2*\cw);

  % Blue blocks for B components (spatial submatrix)
  \fill[blue!20] (\cw,\cw) rectangle (2*\cw,2*\cw);
  \fill[blue!20] (2*\cw,2*\cw) rectangle (3*\cw,3*\cw);
  \fill[blue!20] (\cw,2*\cw) rectangle (2*\cw,3*\cw);
  \fill[blue!20] (2*\cw,\cw) rectangle (3*\cw,2*\cw);

  % Labels（居中于各格）
  \node at (0.5*\cw,3.5*\cw) {$0$};
  \node[red!70!black] at (1.5*\cw,3.5*\cw) {$-E_x/c$};
  \node[red!70!black] at (2.5*\cw,3.5*\cw) {$-E_y/c$};
  \node[red!70!black] at (3.5*\cw,3.5*\cw) {$-E_z/c$};

  \node[red!70!black] at (0.5*\cw,2.5*\cw) {$E_x/c$};
  \node at (1.5*\cw,2.5*\cw) {$0$};
  \node[blue!70!black] at (2.5*\cw,2.5*\cw) {$-B_z$};
  \node[blue!70!black] at (3.5*\cw,2.5*\cw) {$B_y$};

  \node[red!70!black] at (0.5*\cw,1.5*\cw) {$E_y/c$};
  \node[blue!70!black] at (1.5*\cw,1.5*\cw) {$B_z$};
  \node at (2.5*\cw,1.5*\cw) {$0$};
  \node[blue!70!black] at (3.5*\cw,1.5*\cw) {$-B_x$};

  \node[red!70!black] at (0.5*\cw,0.5*\cw) {$E_z/c$};
  \node[blue!70!black] at (1.5*\cw,0.5*\cw) {$-B_y$};
  \node[blue!70!black] at (2.5*\cw,0.5*\cw) {$B_x$};
  \node at (3.5*\cw,0.5*\cw) {$0$};

  % Axis labels
  \foreach \i/\t in {0/0, 1/1, 2/2, 3/3} {
    \node[font=\small] at ({(\i+0.5)*\cw},{4*\cw+0.3}) {$\t$};
    \node[font=\small,anchor=east] at (-0.15,{((3-\i)+0.5)*\cw}) {$\t$};
  }

  % Legend
  \draw[red!50, fill=red!20, thick] ({4*\cw+0.9},3*\cw) rectangle ({4*\cw+1.4},{3*\cw+0.5});
  \node[right] at ({4*\cw+1.5},{3*\cw+0.25}) {电场分量 $E$};
  \draw[blue!50, fill=blue!20, thick] ({4*\cw+0.9},{3*\cw-0.7}) rectangle ({4*\cw+1.4},{3*\cw-0.2});
  \node[right] at ({4*\cw+1.5},{3*\cw-0.45}) {磁场分量 $B$};
  \node[right] at ({4*\cw+0.9},{3*\cw-1.1}) {$F^{\mu\nu} = -F^{\nu\mu}$};
  \node[right] at ({4*\cw+0.9},{3*\cw-1.55}) {反对称：$F^{00}=F^{ii}=0$};

  % Title
  \node[font=\bfseries] at (2*\cw,{4*\cw+0.8}) {电磁场张量 $F^{\mu\nu}$ 的彩色分块};
\end{tikzpicture}
\end{center}


\section{四维电磁势}

\begin{theorem}{四维电磁势与协变方程}{4d-potential-v2}
定义\textbf{四维电磁势}：
\[
A^\mu = \left(\frac{\phi}{c}, \vecA\right)
\]
定义\textbf{四维电流密度}：
\[
j^\mu = (c\rho, \vecJ)
\]
则洛伦兹规范可写成协变形式：
\[
\partial_\mu A^\mu = 0
\]
达朗贝尔方程可写成：
\[
\Box A^\mu = -\muz j^\mu
\]
其中 $\Box = \lap - \frac{1}{c^2}\frac{\partial^2}{\partial t^2} = -\partial_\mu\partial^\mu$ 是达朗贝尔算子（本书统一约定，见第7章 §7.1.5）。
\end{theorem}

\begin{insight}{为什么四维势的协变形式如此优美}{4d-potential-beauty-v2}
在三维语言中，洛伦兹规范 $\divg\vecA + \frac{1}{c^2}\frac{\partial\phi}{\partial t} = 0$ 看起来像是两个不同物理量之间的"拼凑"。但在四维语言中，它简化为 $\partial_\mu A^\mu = 0$——一个标量方程，在洛伦兹变换下保持不变。同样，两个达朗贝尔方程（一个标量、一个矢量）合并为一个四维矢量方程 $\Box A^\mu = -\muz j^\mu$。这种简洁性正是协变形式的力量。
\end{insight}

\section{协变形式的麦克斯韦方程组}

\begin{theorem}{协变麦克斯韦方程组}{covariant-maxwell-v2}
引入四维电流密度 $j^\mu = (c\rho, j_x, j_y, j_z)$，麦克斯韦方程组可写为：
\begin{align}
\partial_\mu F^{\mu\nu} &= \muz j^\nu \label{eq:cov1} \\
\partial_\lambda F_{\mu\nu} + \partial_\mu F_{\nu\lambda} + \partial_\nu F_{\lambda\mu} &= 0 \label{eq:cov2}
\end{align}
\end{theorem}

\begin{review}{协变方程的三维展开}{covariant-expansion-v2}
协变方程 $\partial_\mu F^{\mu\nu} = \muz j^\nu$ 虽然紧凑，但对自学者来说可能过于抽象。下面列出其四个分量方程与三维麦克斯韦方程的对应关系：
\begin{itemize}
\item $\nu = 0$（时间分量）：
\[
\partial_i F^{i0} = \muz j^0 \quad\Longrightarrow\quad \divg\vecE = \frac{\rho}{\eps_0}
\]
这正是\textbf{高斯定律}。
\item $\nu = i$（空间分量）：
\[
\partial_\mu F^{\mu i} = \muz j^i \quad\Longrightarrow\quad (\curl\vecB)_i = \muz j_i + \muz\eps_0\left(\frac{\partial\vecE}{\partial t}\right)_i
\]
这正是\textbf{安培-麦克斯韦定律}的三个分量。
\end{itemize}
将协变方程展开，看到熟悉的 $\divg$ 和 $\curl$，有助于建立张量语言与前面各章三维表述之间的联系。
\end{review}

\begin{review}{三维与四维麦克斯韦方程逐项对照}{covariant-table-v2}
协变麦克斯韦方程组将前面9章的全部静态与时变电磁学规律，浓缩为两张紧凑的张量方程。下表给出逐行逐项的完整对应，是全书知识框架的核心交汇点：

\begin{center}
\small
\setlength{\tabcolsep}{2pt}
\begin{tabular}{@{}llll@{}}
\toprule
\textbf{四维方程} & \textbf{分量选取} & \textbf{对应三维方程} & \textbf{物理意义} \\
\midrule
$\partial_\mu F^{\mu 0}=\muz j^0$ & $\nu=0$ & $\divg\vecE=\rho/\eps_0$ & 高斯定律 \\
$\partial_\mu F^{\mu i}=\muz j^i$ & $\nu=1,2,3$ & $(\curl\vecB)_i=\muz j_i+\muz\eps_0(\partial\vecE/\partial t)_i$ & 安培-麦克斯韦 \\
$\partial_\lambda F_{\mu\nu}+\partial_\mu F_{\nu\lambda}+\partial_\nu F_{\lambda\mu}=0$ & $\lambda\mu\nu=0ij$ & $(\curl\vecE)_i=-(\partial\vecB/\partial t)_i$ & 法拉第定律 \\
$\partial_\lambda F_{\mu\nu}+\partial_\mu F_{\nu\lambda}+\partial_\nu F_{\lambda\mu}=0$ & $\lambda\mu\nu=ijk$ & $\divg\vecB=0$ & 磁场无源 \\
\bottomrule
\end{tabular}
\end{center}

\begin{insight}{全书知识框架的核心交汇点}{covariant-summary-v2}
这张对照表是整本教材逻辑结构的"中央枢纽"：
\begin{itemize}
\item 第1--2章的静电学（$\divg\vecE=\rho/\eps_0$）和第3--4章的静磁学（$\divg\vecB=0$）是它的\textbf{低速静态特例}；
\item 第5章的法拉第电磁感应定律和第6章的安培-麦克斯韦定律是它的\textbf{时变展开}；
\item 第7--8章的电磁波解是它的\textbf{无源波动解}；
\item 第9章的辐射场是它的\textbf{源项展开}。
\end{itemize}
原本分散在各章的八个标量方程，如今被两个协变张量方程统一概括。这不仅体现了数学的简洁之美，更深刻揭示了电磁场作为统一物理实体的本质——电场与磁场不是独立的对象，而是同一四维张量在不同参考系中的不同投影。
\end{insight}
\end{review}

\begin{insight}{协变形式的美妙之处}{covariant-beauty-v2}
方程\eqref{eq:cov1}（四个分量）对应 $\divg\vecE = \rho/\eps_0$ 和 $\curl\vecB = \muz\vecJ + \muz\eps_0\partial\vecE/\partial t$。

方程\eqref{eq:cov2}（四个独立方程）对应 $\divg\vecB = 0$ 和 $\curl\vecE = -\partial\vecB/\partial t$。

原本的八个标量方程，现在被紧凑地表达为两个张量方程。更重要的是：这两个方程在洛伦兹变换下保持\textbf{形式不变}——这正是相对论要求的协变性。
\end{insight}

\section{带电粒子在电磁场中的运动}

\begin{theorem}{相对论性洛伦兹力}{rel-lorentz-v2}
\textbf{四维力}：
\[
\frac{\diff p^\mu}{\diff\tau} = q F^{\mu\nu} u_\nu
\]
空间分量给出相对论性洛伦兹力：
\[
\frac{\diff(\gamma m\vect{v})}{\diff t} = q(\vecE + \vect{v} \times \vecB)
\]
\end{theorem}

\begin{review}{第10章要点回顾}{ch10-review}
\begin{center}
\begin{tabular}{@{}ll@{}}
\toprule
\textbf{概念/公式} & \textbf{内容} \\
\midrule
$F^{\mu\nu}$ & 电磁场张量（统一E和B） \\
$\vecE$ 和 $\vecB$ 混合变换 & 洛伦兹变换下场变换 \\
$\partial_\mu F^{\mu\nu} = \muz j^\nu$ & 协变麦克斯韦方程（含源） \\
$\partial_{[\lambda}F_{\mu\nu]} = 0$ & 协变麦克斯韦方程（齐次） \\
$\vect{F} = q(\vecE + \vect{v}\times\vecB)$ & 洛伦兹力 \\
\bottomrule
\end{tabular}
\end{center}
\end{review}

\begin{tip}{本章核心框架}{ch10-summary}
\textbf{一句话总结}：狭义相对论将电动力学统一为四维协变理论，证明电磁场是同一物理实体的不同表现。\\
\textbf{核心内容}：洛伦兹变换（时空统一）、四维矢量（$x^\mu, j^\mu, A^\mu$）、电磁场张量$F^{\mu\nu}$（$E$和$B$的协变统一）、麦克斯韦方程组的协变形式。\\
\textbf{关键公式}：$E=\gamma m_0c^2$、$\partial_\mu F^{\mu\nu}=\mu_0j^\nu$、$F_{\mu\nu}F^{\mu\nu}=2(B^2-E^2/c^2)$。\\
\textbf{里程碑}：你已完整掌握经典电动力学。下一步：量子电动力学（QED）！
\end{tip}

\subsection{电磁场的洛伦兹不变量}

\begin{theorem}{两个基本不变量}{lorentz-invariants}
由电磁场张量 $F^{\mu\nu}$ 可以构造两个洛伦兹不变量（按本书约定 $F^{0i}=-E_i/c$、$\epsilon_{0123}=+1$，与本章前述不变量定义式 \eqref{eq:invariant1}、\eqref{eq:invariant2} 完全一致）：
\begin{align*}
\mathcal{I}_1 &= \tfrac{1}{2}F_{\mu\nu}F^{\mu\nu} = B^2 - \frac{E^2}{c^2} \\
\mathcal{I}_2 &= \tfrac{1}{8}\epsilon_{\mu\nu\rho\sigma}F^{\mu\nu}F^{\rho\sigma} = +\frac{1}{c}\vecE\cdot\vecB
\end{align*}
\end{theorem}

\begin{insight}{不变量的物理意义}{invariant-meaning}
$\mathcal{I}_1 = B^2 - E^2/c^2$ 的符号决定了能否通过洛伦兹变换消去某一种场：\\
- 若 $\mathcal{I}_1 > 0$（即 $cB > E$，磁场占优）：存在参考系使 $E=0$（纯磁场）\\
- 若 $\mathcal{I}_1 < 0$（即 $cB < E$，电场占优）：存在参考系使 $B=0$（纯电场）\\
- 若 $\mathcal{I}_1 = 0$（即 $E=cB$）：在所有参考系中都有 $E=cB$（如平面电磁波），两种场都无法被消去\\
$\mathcal{I}_2 = 0$ 时，$\vecE\perp\vecB$，也是平面波的特征。
\end{insight}

\subsection{三维与四维语言对照表}

\begin{center}
\begin{tabular}{lll}
\toprule
\textbf{物理量} & \textbf{三维语言} & \textbf{四维语言} \\
\midrule
时空坐标 & $\vect{r}, t$ & $x^\mu = (ct, \vect{r})$ \\
电流 & $\rho, \vecJ$ & $j^\mu = (c\rho, \vecJ)$ \\
电磁势 & $\phi, \vecA$ & $A^\mu = (\phi/c, \vecA)$ \\
电场 & $\vecE = -\grad\phi - \partial\vecA/\partial t$ & $F^{0i} = -E^i/c$ \\
磁场 & $\vecB = \curl\vecA$ & $F^{ij} = -\epsilon_{ijk}B^k$ \\
高斯定理 & $\divg\vecE = \rho/\eps_0$ & $\partial_\mu F^{\mu 0} = \mu_0 j^0$ \\
法拉第定律 & $\curl\vecE = -\partial\vecB/\partial t$ & $\partial_\mu F^{\nu\rho} + \partial_\nu F^{\rho\mu} + \partial_\rho F^{\mu\nu} = 0$ \\
安培-麦克斯韦 & $\curl\vecB = \mu_0\vecJ + \mu_0\eps_0\partial\vecE/\partial t$ & $\partial_\mu F^{\mu\nu} = \mu_0 j^\nu$ \\
无磁单极 & $\divg\vecB = 0$ & $\partial_\mu {}^*\!F^{\mu\nu} = 0$（对偶张量） \\
连续性方程 & $\partial\rho/\partial t + \divg\vecJ = 0$ & $\partial_\mu j^\mu = 0$ \\
洛伦兹规范 & $\nabla\cdot\vecA + \mu_0\eps_0\partial\phi/\partial t = 0$ & $\partial_\mu A^\mu = 0$ \\
达朗贝尔方程 & $\Box\phi = -\rho/\eps_0$, $\Box\vecA = -\mu_0\vecJ$ & $\Box A^\mu = -\mu_0 j^\mu$ \\
\bottomrule
\end{tabular}
\end{center}

\begin{tip}{记忆诀窍}{4d-memory}
三维→四维的"翻译规则"：\\
1. 标量+矢量 → 拼成四维矢量（$\phi/c$ 是 $A^0$ 的时间分量）\\
2. 梯度+时间导数 → 协变导数 $\partial_\mu = (\partial_t/c, \grad)$\\
3. 散度+旋度 → 统一为 $F^{\mu\nu}$ 的缩并和循环求和\\
4. 所有方程的右侧 → 电流四维矢量 $j^\mu$ 的对应分量
\end{tip}

% ---------- 第10章随堂自测 ----------
\section*{第10章随堂自测}
\addcontentsline{toc}{section}{第10章随堂自测}

\begin{miniquiz}
\textbf{一、选择题}
\begin{enumerate}[label=\arabic*.]
\item 洛伦兹变换保持不变的量是：
  \begin{enumerate}[label=(\Alph*),nosep]
  \item 时间间隔
  \item 空间间隔
  \item 时空间隔 $s^2 = c^2t^2 - x^2$
  \item 同时性
  \end{enumerate}
\begin{ocg}{答案}{ocg答案10}{off}
  \textit{答案：C。时空间隔是洛伦兹不变量。}
\end{ocg}

\item 四维电流密度 $j^\mu = (c\rho, \vecJ)$ 的连续性方程在协变形式下写为：
  \begin{enumerate}[label=(\Alph*),nosep]
  \item $\partial_\mu j^\mu = 0$
  \item $\partial_\mu j^\mu = \rho$
  \item $\Box j^\mu = 0$
  \item $\partial_\mu F^{\mu\nu} = 0$
  \end{enumerate}
\begin{ocg}{答案}{ocg答案10}{off}
  \textit{答案：A。$\partial_\mu j^\mu = \partial\rho/\partial t + \divg\vecJ = 0$。}
\end{ocg}

\item 电磁场张量 $F^{\mu\nu}$ 是：
  \begin{enumerate}[label=(\Alph*),nosep]
  \item 对称张量
  \item 反对称张量
  \item 对角张量
  \item 标量
  \end{enumerate}
\begin{ocg}{答案}{ocg答案10}{off}
  \textit{答案：B。$F^{\mu\nu} = -F^{\nu\mu}$，包含6个独立分量（3个电场+3个磁场）。}
\end{ocg}
\end{enumerate}

\vspace{0.5em}
\textbf{二、填空题}
\begin{enumerate}[label=\arabic*.,resume]
\item 洛伦兹因子 $\gamma = $ \underline{$1/\sqrt{1-v^2/c^2}$}。
\item 达朗贝尔算符 $\Box = $ \underline{$\lap - \frac{1}{c^2}\frac{\partial^2}{\partial t^2} = -\partial^\mu\partial_\mu$}。
\item 洛伦兹规范在协变形式下写为 \underline{$\partial_\mu A^\mu = 0$}。
\end{enumerate}
\end{miniquiz}
\toggleButton{显示答案}{ocg答案10}

\begin{chapterreview}
\begin{itemize}[nosep]
\item \textbf{核心内容}：洛伦兹变换、四维矢量、电磁场张量、麦克斯韦方程组的协变形式、相对论力学。
\item \textbf{关键公式}：$F^{\mu\nu}$的矩阵表示、$\partial_\mu F^{\mu\nu} = \mu_0 j^\nu$、相对论性能量$E = \gamma m_0c^2$。
\item \textbf{自检验}：能否写出电磁场张量的显式矩阵？能否验证$\vecE'$和$\vecB'$的变换公式？
\item \textbf{里程碑}：至此，你已经完整掌握了经典电动力学的核心框架。恭喜你！
\end{itemize}
\end{chapterreview}

\begin{tip}{本章自检清单}{ch10-checklist}
学完本章后对照下列条目逐条自测，全部通过即可进入下一章；有未通过项时，回到对应小节复习：
\begin{itemize}[label=$\square$, nosep]
\item 用洛伦兹变换导出电磁场变换
\item 写出电磁场张量 $F^{\mu\nu}$ 并验证其协变性
\item 把麦克斯韦方程组写成协变形式 $\partial_\mu F^{\mu\nu} = \muz j^\nu$
\item 说明号差约定对公式符号的影响
\end{itemize}
\end{tip}

% ---------- 第10章综合练习题 ----------
\section*{第10章综合练习题}
\addcontentsline{toc}{section}{第10章综合练习题}
\begin{enumerate}[label=\textbf{\arabic*.}]
\item \textbf{基础} 设 $x^\mu=(ct,x,y,z)$，$\eta_{\mu\nu}=\text{diag}(1,-1,-1,-1)$。计算 $x_\mu=\eta_{\mu\nu}x^\nu$ 的各分量，并验证 $x^\mu x_\mu=c^2t^2-x^2-y^2-z^2$。
\item \textbf{基础} 在静止系 $S$ 中只有电场 $\vecE=(0,E_0,0)$，求沿 $x$ 方向以速度 $v$ 运动的参考系 $S'$ 中的 $\vecE'$ 和 $\vecB'$。
\item \textbf{提高} 写出四维电磁势 $A^\mu=(\phi/c,\vecA)$，证明在洛伦兹规范下 $\partial_\mu A^\mu=0$ 等价于 $\nabla\cdot\vecA+\muz\eps_0\partial\phi/\partial t=0$。
\item \textbf{提高} 从协变麦克斯韦方程 $\partial_\mu F^{\mu\nu}=\muz j^\nu$ 出发，取 $\nu=0$ 和 $\nu=i$，分别还原高斯定律和安培-麦克斯韦定律的三维形式。
\item \textbf{挑战} 证明：电磁场张量 $F^{\mu\nu}$ 的两个洛伦兹不变量
\[
\frac{1}{2}F_{\mu\nu}F^{\mu\nu}=B^2-\frac{E^2}{c^2}, \qquad \frac{1}{8}\epsilon_{\mu\nu\lambda\sigma}F^{\mu\nu}F^{\lambda\sigma}=\frac{\vecE\cdot\vecB}{c}
\]
在任何惯性系中取值相同。讨论 $\vecE\perp\vecB$ 与 $\vecE\parallel\vecB$ 的物理意义。
\end{enumerate}
\newpage
\begin{center}\small \hyperlink{ch9-goal}{上一章：第9章}\end{center}

\section*{习题解答}
\addcontentsline{toc}{section}{习题解答}

\textbf{题1解析：}

[分析] 本题考查闵氏空间中协变分量与逆变分量的升降关系，以及时空间隔不变量的验证。

[解答] 由度规定义 $x_\mu = \eta_{\mu\nu} x^\nu$，采用爱因斯坦求和约定：
\begin{align*}
x_0 &= \eta_{00}x^0 + \eta_{01}x^1 + \eta_{02}x^2 + \eta_{03}x^3 = 1 \cdot ct = ct \\
x_1 &= \eta_{10}x^0 + \eta_{11}x^1 = -1 \cdot x = -x \\
x_2 &= -y, \quad x_3 = -z
\end{align*}
即 $x_\mu = (ct, -x, -y, -z) = (ct, -\vect{r})$。

验证不变量：
\begin{align*}
x^\mu x_\mu &= x^0 x_0 + x^1 x_1 + x^2 x_2 + x^3 x_3 \\
&= ct \cdot ct + x \cdot (-x) + y \cdot (-y) + z \cdot (-z) \\
&= c^2t^2 - x^2 - y^2 - z^2 = c^2t^2 - r^2
\end{align*}
这正是时空间隔 $s^2$，在洛伦兹变换下保持不变。

[量纲校验] $[x^\mu x_\mu] = [L]^2$（$ct$ 和 $x,y,z$ 均为长度量纲），结果 $c^2t^2-r^2$ 量纲为 $[L]^2$，一致。\qedsymbol

\textbf{题2解析：}

[分析] 本题考查电磁场在洛伦兹变换下的混合效应——纯电场在运动参考系中会感应出磁场。

[解答] 已知 $S$ 系中只有 $\vecE=(0,E_0,0)$，$\vecB=\vect{0}$。$S'$ 系以速度 $v$ 沿 $x$ 方向运动。

利用电磁场洛伦兹变换（$\gamma = 1/\sqrt{1-v^2/c^2}$）：
\begin{align*}
E'_x &= E_x = 0 \\
E'_y &= \gamma(E_y - vB_z) = \gamma E_0 \\
E'_z &= \gamma(E_z + vB_y) = 0 \\
B'_x &= B_x = 0 \\
B'_y &= \gamma(B_y + vE_z/c^2) = 0 \\
B'_z &= \gamma(B_z - vE_y/c^2) = -\gamma v E_0/c^2
\end{align*}
即 $S'$ 中观察到 $\vecE' = \gamma E_0\,\unitvec{y}$，$\vecB' = -\gamma vE_0/c^2\,\unitvec{z}$。物理上，这是因为在 $S$ 系中静止的电荷在 $S'$ 系中以速度 $-v\unitvec{x}$ 运动，形成等效电流，从而激发磁场。

注意：若 $S'$ 沿电场方向（$y$）运动，则沿运动方向的电场分量不混合出磁场（$B'_z = 0$）——只有\textbf{横向}运动才产生感应磁场。

[量纲校验] $[B] = [MT^{-2}I^{-1}]$，$[vE/c^2] = [LT^{-1}][MLT^{-3}I^{-1}]/[L^2T^{-2}] = [MT^{-2}I^{-1}]$，一致。\qedsymbol

\textbf{题3解析：}

[分析] 本题考查协变语言与三维语言之间的等价性，特别是洛伦兹规范的四维表述。

[解答] 四维电磁势 $A^\mu = (\phi/c, A_x, A_y, A_z) = (\phi/c, \vecA)$。

协变导数 $\partial_\mu = (\partial_0, \partial_1, \partial_2, \partial_3) = (\frac{1}{c}\partial_t, \nabla)$。

计算 $\partial_\mu A^\mu$：
\begin{align*}
\partial_\mu A^\mu &= \partial_0 A^0 + \partial_1 A^1 + \partial_2 A^2 + \partial_3 A^3 \\
&= \frac{1}{c}\partial_t \left(\frac{\phi}{c}\right) + \partial_x A_x + \partial_y A_y + \partial_z A_z \\
&= \frac{1}{c^2}\partial_t\phi + \nabla\cdot\vecA
\end{align*}

在SI单位制中，$\muz\eps_0 = 1/c^2$，因此：
\[
\partial_\mu A^\mu = \muz\eps_0\partial_t\phi + \nabla\cdot\vecA = 0
\]
这正是洛伦兹规范的三维形式。\qedsymbol

[量纲校验] $[\partial_\mu A^\mu] = [L^{-1}][L^{-1}MT^{-2}I^{-1}]$（在SI中 $A^\mu$ 各分量量纲为 $[MLT^{-2}I^{-1}]$），$[\nabla\cdot\vecA] = [L^{-1}][MLT^{-2}I^{-1}]$，一致。\qedsymbol

\textbf{题4解析：}

[分析] 本题将协变麦克斯韦方程展开为三维形式，建立张量语言与矢量分析之间的联系。

[解答] 协变方程：$\partial_\mu F^{\mu\nu} = \muz j^\nu$。

（1）$\nu = 0$（时间分量）：
\begin{align*}
\partial_\mu F^{\mu 0} &= \muz j^0 \\
\partial_i F^{i0} &= \muz c\rho \quad (i=1,2,3)
\end{align*}
由 $F^{i0} = -F^{0i} = +E^i/c$（场张量反对称性），得：
\begin{align*}
\partial_i \left(+\frac{E^i}{c}\right) &= \muz c\rho \\
\frac{1}{c}\divg\vecE &= \muz c\rho \\
\divg\vecE &= \muz c^2\rho = \frac{\rho}{\eps_0}
\end{align*}
这正是高斯定律。

（2）$\nu = i$（空间分量，以 $i=1$ 即 $x$ 分量为例）：
\begin{align*}
\partial_\mu F^{\mu 1} &= \muz j^1 \\
\partial_0 F^{01} + \partial_2 F^{21} + \partial_3 F^{31} &= \muz j_x
\end{align*}
由 $F^{01} = -E_x/c$，$F^{21} = B_z$，$F^{31} = -B_y$：
\begin{align*}
\frac{1}{c}\partial_t\left(-\frac{E_x}{c}\right) + \partial_y B_z - \partial_z B_y &= \muz j_x \\
-\frac{1}{c^2}\partial_t E_x + (\curl\vecB)_x &= \muz j_x \\
(\curl\vecB)_x &= \muz j_x + \muz\eps_0\partial_t E_x
\end{align*}
三个空间分量合起来就是安培--麦克斯韦定律。\qedsymbol

[量纲校验] $[\partial_\mu F^{\mu\nu}] = [L^{-1}][MLT^{-2}I^{-1}] = [MT^{-3}I^{-1}]$，$[\muz j^\nu] = [MLT^{-2}I^{-2}][IL^{-2}] = [MT^{-3}I^{-1}]$，一致。\qedsymbol

\textbf{题5解析：}

[分析] 本题考查电磁场张量的洛伦兹不变量，以及这两个不变量的物理意义。

[解答] 展开 $F_{\mu\nu}F^{\mu\nu}$。由 $F^{0i} = -E_i/c$ 与 $F_{0i} = -F^{0i} = E_i/c$（度规 $(+,-,-,-)$ 降指标），时间-空间分量的贡献：
\begin{align*}
F_{\mu\nu}F^{\mu\nu} &= 2\left(F_{01}F^{01} + F_{02}F^{02} + F_{03}F^{03}\right) + 2\left(F_{12}F^{12} + F_{23}F^{23} + F_{31}F^{31}\right) \\
&= 2\left(\frac{E_x}{c}\cdot\frac{-E_x}{c} + \frac{E_y}{c}\cdot\frac{-E_y}{c} + \frac{E_z}{c}\cdot\frac{-E_z}{c}\right) \\
&\quad + 2\left((-B_z)(-B_z) + (-B_x)(-B_x) + (-B_y)(-B_y)\right) \\
&= -2\frac{E^2}{c^2} + 2B^2 = 2\left(B^2 - \frac{E^2}{c^2}\right)
\end{align*}
因此 $\frac{1}{2}F_{\mu\nu}F^{\mu\nu} = B^2 - E^2/c^2 \equiv I_1$。

第二个不变量利用 Levi-Civita 符号计算：
\[
\epsilon_{\mu\nu\lambda\sigma}F^{\mu\nu}F^{\lambda\sigma} = +\frac{8\vecE\cdot\vecB}{c}
\]
故 $\frac{1}{8}\epsilon_{\mu\nu\lambda\sigma}F^{\mu\nu}F^{\lambda\sigma} = \vecE\cdot\vecB/c \equiv I_2$。

物理意义：
\begin{itemize}
\item 若 $I_2 = 0$（$\vecE\perp\vecB$），则存在一个参考系使场变为纯场：$I_1 < 0$（即 $E > cB$）时纯\textbf{磁场}变为零，只剩纯电场；$I_1 > 0$（即 $E < cB$）时纯\textbf{电场}变为零，只剩纯磁场。验证：纯电场系中 $I_1 = -E^2/c^2 < 0$，纯磁场系中 $I_1 = B^2 > 0$，与不变量取值一致。
\item 若 $I_2 \neq 0$（$\vecE$ 与 $\vecB$ 不垂直），则不存在任何参考系能使其中一个为零。
\item 当 $\vecE\parallel\vecB$ 时，$I_2$ 取最大值，电磁场结构最"稳定"。
\end{itemize}

[量纲校验] $[B^2] = [M^2T^{-4}I^{-2}]$，$[E^2/c^2] = [M^2L^2T^{-6}I^{-2}]/[L^2T^{-2}] = [M^2T^{-4}I^{-2}]$，量纲一致。$[\vecE\cdot\vecB/c] = [MLT^{-3}I^{-1}][MT^{-2}I^{-1}]/[LT^{-1}] = [M^2T^{-4}I^{-2}]$，与前两者同量纲。\qedsymbol



\begin{warning}{第10章常见误区}{ch10-warnings}
\begin{enumerate}
\item 相对论不是"额外"加在电动力学上的——电动力学本身\textbf{要求}相对论才能自洽。
\item 在低速极限（$v \ll c$）下，相对论公式退化为经典结果——电动力学是相对论在低能极限下的表现。
\item 电磁场张量 $F^{\mu\nu}$ 中的 $E/c$ 和 $B$ 具有相同的量纲——这在高斯单位制中更为明显（$E$ 和 $B$ 量纲相同）。
\item 洛伦兹力公式 $\vect{F} = q(\vecE + \vect{v}\times\vecB)$ 在相对论中仍然成立——它是协变方程的空间分量。
\end{enumerate}
\end{warning}

% ============================================================
% 附录
% ============================================================
\appendix

\chapter{常用矢量恒等式}

\begin{align}
\vecA \cdot (\vecB \times \vect{C}) &= \vecB \cdot (\vect{C} \times \vecA) = \vect{C} \cdot (\vecA \times \vecB) \\
\vecA \times (\vecB \times \vect{C}) &= \vecB(\vecA\cdot\vect{C}) - \vect{C}(\vecA\cdot\vecB) \\
(\vecA \times \vecB) \cdot (\vect{C} \times \vect{D}) &= (\vecA\cdot\vect{C})(\vecB\cdot\vect{D}) - (\vecA\cdot\vect{D})(\vecB\cdot\vect{C}) \\
\grad(fg) &= f\grad g + g\grad f \\
\divg(f\vecA) &= f\divg\vecA + \vecA\cdot\grad f \\
\curl(f\vecA) &= f\curl\vecA + \grad f \times \vecA \\
\divg(\vecA \times \vecB) &= \vecB\cdot\curl\vecA - \vecA\cdot\curl\vecB \\
\curl(\vecA \times \vecB) &= \vecA(\divg\vecB) - \vecB(\divg\vecA) + (\vecB\cdot\grad)\vecA - (\vecA\cdot\grad)\vecB \\
\curl(\curl\vecA) &= \grad(\divg\vecA) - \lap\vecA \\
\oint_S \vecA\cdot\diff\vect{S} &= \int_V \divg\vecA\,\diff V \\
\oint_C \vecA\cdot\diff\vect{l} &= \int_S \curl\vecA\cdot\diff\vect{S}
\end{align}

\chapter{坐标系转换公式}

\textbf{直角坐标 $(x, y, z)$}：
\begin{align*}
\lap\psi &= \frac{\partial^2\psi}{\partial x^2} + \frac{\partial^2\psi}{\partial y^2} + \frac{\partial^2\psi}{\partial z^2} \\
\diff^3r &= \diff x\,\diff y\,\diff z
\end{align*}

\textbf{球坐标 $(r, \theta, \phi)$}：
\begin{align*}
\lap\psi &= \frac{1}{r^2}\frac{\partial}{\partial r}\left(r^2\frac{\partial\psi}{\partial r}\right) + \frac{1}{r^2\sin\theta}\frac{\partial}{\partial\theta}\left(\sin\theta\frac{\partial\psi}{\partial\theta}\right) + \frac{1}{r^2\sin^2\theta}\frac{\partial^2\psi}{\partial\phi^2} \\
\diff^3r &= r^2\sin\theta\,\diff r\,\diff\theta\,\diff\phi
\end{align*}

\textbf{柱坐标 $(\rho, \phi, z)$}：
\begin{align*}
\lap\psi &= \frac{1}{\rho}\frac{\partial}{\partial\rho}\left(\rho\frac{\partial\psi}{\partial\rho}\right) + \frac{1}{\rho^2}\frac{\partial^2\psi}{\partial\phi^2} + \frac{\partial^2\psi}{\partial z^2} \\
\diff^3r &= \rho\,\diff\rho\,\diff\phi\,\diff z
\end{align*}

\chapter{基本物理常数}

\begin{center}
\begin{tabular}{@{}lll@{}}
\toprule
\textbf{常数} & \textbf{符号} & \textbf{数值} \\
\midrule
真空光速 & $c$ & $2.998 \times 10^8 \, \si{m/s}$ \\
真空介电常数 & $\eps_0$ & $8.854 \times 10^{-12} \, \si{F/m}$ \\
真空磁导率 & $\muz$ & $4\pi \times 10^{-7} \, \si{H/m}$（近似） \\
基本电荷 & $e$ & $1.602 \times 10^{-19} \, \si{C}$ \\
电子质量 & $m_e$ & $9.109 \times 10^{-31} \, \si{kg}$ \\
精细结构常数 & $\alpha$ & $1/137.036$ \\
\bottomrule
\end{tabular}
\end{center}

\noindent{\small 注：2019 年 SI 修订后，$\eps_0$ 与 $\muz$ 不再是精确定义值（$\muz = 4\pi\times10^{-7}$ H/m 在旧 SI 中为精确定义），改为由精细结构常数 $\alpha$、基本电荷 $e$、普朗克常数 $h$、光速 $c$ 推导的\textbf{测量值}（与 $4\pi\times10^{-7}$ 的偏差约 $2\times10^{-9}$ 量级）。本书作为教学用书统一沿用 $4\pi\times10^{-7}$ 近似值，对全部计算结果的影响可忽略。}


% ============================================================
\chapter{单位制转换手册}
% ============================================================

本书统一采用SI（国际单位制）。但外文经典教材（如Griffiths、Jackson）部分使用高斯CGS单位制，初学者极易混淆。本附录给出两套单位制的核心公式对照，并标注常见陷阱。

\begin{warning}{不要直接照搬外文教材的公式}{unit-copied-tip}
\textbf{Griffiths 与 Jackson 的部分公式采用高斯单位制，直接抄进 SI 计算会出错}。遇到外文教材公式时先做三步检查：
\begin{enumerate}[nosep]
\item 公式中是否出现 $4\pi$ 在\textbf{分子}（高斯制特征，如 $\divg\vecE = 4\pi\rho$）或 $1/4\pi\eps_0$ 在\textbf{分母}（SI 特征）——这是最快的高斯/SI 识别标志
\item 是否显式含 $c$（高斯制中 $c$ 常出现在含磁场的公式中，如 $\vecB = \hat{\vect{k}}\times\vecE$ 无 $1/c$；SI 中 $\vecB = \frac{1}{c}\hat{\vect{k}}\times\vecE$）
\item 用本书"单位制转换手册"的对照表逐项换算后再代入数值
\end{enumerate}
两本教材的具体情况：\textbf{Griffiths} 主文用 SI，个别专题用高斯制并明确标注；\textbf{Jackson}（第 3 版）第 1--10 章用高斯制、第 11 章起才转 SI——读 Jackson 前几章时要特别小心。
\end{warning}

\section{SI与高斯制的基本公式对照}

\begin{center}
\begin{tabular}{@{}lll@{}}
\toprule
\textbf{物理量} & \textbf{SI} & \textbf{高斯CGS} \\
\midrule
库仑力 & $\displaystyle F=\frac{1}{4\pi\eps_0}\frac{q_1 q_2}{r^2}$ & $\displaystyle F=\frac{q_1 q_2}{r^2}$ \\
电场强度 & $\displaystyle \vecE=\frac{1}{4\pi\eps_0}\frac{q\unitvec{r}}{r^2}$ & $\displaystyle \vecE=\frac{q\unitvec{r}}{r^2}$ \\
电势 & $\displaystyle \phi=\frac{1}{4\pi\eps_0}\frac{q}{r}$ & $\displaystyle \phi=\frac{q}{r}$ \\
高斯定律 & $\displaystyle \divg\vecE=\frac{\rho}{\eps_0}$ & $\displaystyle \divg\vecE=4\pi\rho$ \\
电位移 & $\displaystyle \vect{D}=\eps_0\vecE+\vect{P}$ & $\displaystyle \vect{D}=\vecE+4\pi\vect{P}$ \\
毕奥-萨伐尔 & $\displaystyle d\vecB=\frac{\muz}{4\pi}\frac{I d\vect{l}\times\unitvec{r}}{r^2}$ & $\displaystyle d\vecB=\frac{1}{c}\frac{I d\vect{l}\times\unitvec{r}}{r^2}$ \\
安培力 & $\displaystyle d\vect{F}=I d\vect{l}\times\vecB$ & $\displaystyle d\vect{F}=\frac{1}{c}I d\vect{l}\times\vecB$ \\
洛伦兹力 & $\displaystyle \vect{F}=q(\vecE+\vect{v}\times\vecB)$ & $\displaystyle \vect{F}=q\left(\vecE+\frac{\vect{v}}{c}\times\vecB\right)$ \\
麦克斯韦（真空） & $\begin{cases}\divg\vecE=\rho/\eps_0\\\divg\vecB=0\\\curl\vecE=-\partial\vecB/\partial t\\\curl\vecB=\muz\vecJ+\muz\eps_0\partial\vecE/\partial t\end{cases}$ & $\begin{cases}\divg\vecE=4\pi\rho\\\divg\vecB=0\\\curl\vecE=-\frac{1}{c}\partial\vecB/\partial t\\\curl\vecB=\frac{4\pi}{c}\vecJ+\frac{1}{c}\partial\vecE/\partial t\end{cases}$ \\
\bottomrule
\end{tabular}
\end{center}

\begin{warning}{高斯制中的$\muz$与$\eps_0$}{gauss-unit-trap}
在高斯CGS单位制中，$\muz$和$\eps_0$均不出现（或视为$1$）。$1/4\pi\eps_0$在SI中是一个独立常数，在高斯制中因$\eps_0=1/(4\pi)$而被吸收。同时高斯制中光速$c$显式出现在含磁场的公式中。最危险的混淆点：高斯制中$\vecB$的量纲与$\vecE$相同（均为$\sqrt{\text{g}}/\text{cm}^{1/2}/\text{s}$），而SI中$\vecB$的量纲是特斯拉，$\vecE$是伏特/米。
\end{warning}

\section{常数对照}

\begin{center}
\begin{tabular}{@{}lll@{}}
\toprule
\textbf{常数} & \textbf{SI数值} & \textbf{高斯CGS数值} \\
\midrule
$\eps_0$ & $8.854\times 10^{-12}\,\text{F/m}$ & $1/(4\pi)$（无单位） \\
$\muz$ & $4\pi\times 10^{-7}\,\text{H/m}$ & $4\pi/c^2$ \\
$c$ & $2.998\times 10^8\,\text{m/s}$ & $2.998\times 10^{10}\,\text{cm/s}$ \\
$e$ & $1.602\times 10^{-19}\,\text{C}$ & $4.803\times 10^{-10}\,\text{esu}$ \\
\bottomrule
\end{tabular}
\end{center}



% ============================================================

\begin{example}{SI单位与高斯单位的换算}{unit-conversion-example}
\difficulty{基础}
将以下高斯单位制的公式转换为SI单位制：
\begin{enumerate}[label=(\arabic*)]
\item 高斯制中的库仑定律：$F = q_1 q_2 / r^2$
\item 高斯制中的点电荷电势：$\phi = q/r$
\end{enumerate}
\tcblower
\textbf{解答：}

\textbf{(1) 库仑定律的转换}

高斯制中 $F = q_1 q_2 / r^2$，力的单位是达因（dyn）。

SI制中 $F = q_1 q_2 / (4\pi\eps_0 r^2)$，力的单位是牛顿（N）。

转换关系：$1\,\text{N} = 10^5\,\text{dyn}$；电荷单位换算为 $1\,\text{esu} = 3.33564095\times 10^{-10}\,\text{C}$（即 $1\,\text{C} = 2.99792458\times 10^{9}\,\text{esu}$，记作 $1\,\text{C} = c_{\text{cgs}}\,\text{esu}$，其中 $c_{\text{cgs}} = c/(1\,\text{m/s})\times 10^{-2} = 2.99792458\times10^{9}$）。

\textbf{电荷与公式的关系（最易错处）}：单次把\textbf{电荷数值}从高斯制换到 SI 制，用
\[
q_{\text{SI}} = 3.33564095\times 10^{-10}\,q_{\text{G}}\qquad(\text{数值关系})
\]
而若要\textbf{整体替换库仑定律这个公式}，则
\[
F = \frac{q_1 q_2}{r^2}\ \text{(高斯)}\ \longrightarrow\ F = \frac{1}{4\pi\eps_0}\frac{(q_1/\sqrt{4\pi\eps_0}\,)\,(\,q_2/\sqrt{4\pi\eps_0}\,)}{r^2}\ \text{(SI)}
\]
即公式级的替换是「把每个 $q$ 换成 $q/\sqrt{4\pi\eps_0}$」。两种写法并不矛盾：前者的 $3.3356\times10^{-10}$ 已经把 $1/\sqrt{4\pi\eps_0}$ 与 cm→m 的因子合并进去了，\textbf{不可再把 $1/\sqrt{4\pi\eps_0}$ 单独乘一次}（否则量级会差 $10^{5}$ 倍以上）。

因此：$F_{\text{SI}} = \dfrac{1}{4\pi\eps_0}\dfrac{q_{\text{SI},1}q_{\text{SI},2}}{r^2}$，其中 $1/(4\pi\eps_0) \approx 8.99\times 10^9\,\text{N}\cdot\text{m}^2/\text{C}^2$。

\textbf{验证}：量纲 $[q_{\text{SI}}] = \text{C}$，$[\eps_0] = \text{F/m} = \text{C}^2/(\text{N}\cdot\text{m}^2)$，则 $[q^2/(\eps_0 r^2)] = \text{C}^2/[\text{C}^2/(\text{N}\cdot\text{m}^2)\cdot\text{m}^2] = \text{N}$ \checkmark

\textbf{(2) 点电荷电势的转换}

高斯制中 $\phi = q/r$，电势单位是静伏（statV）。

SI制中 $\phi = q/(4\pi\eps_0 r)$。

转换关系：$1\,\text{statV} = 299.79\,\text{V} \approx 300\,\text{V}$。

因此：$\phi_{\text{SI}} = \dfrac{1}{4\pi\eps_0}\dfrac{q_{\text{SI}}}{r}$，单位是伏特（V）。

\textbf{量纲验证}：$[q/(\eps_0 r)] = \text{C}/[\text{C}^2/(\text{N}\cdot\text{m}^2)\cdot\text{m}] = \text{N}\cdot\text{m}/\text{C} = \text{J}/\text{C} = \text{V}$ \checkmark
\end{example}

\chapter{全书高频易错点汇总}
% ============================================================

按章节排序，集中整理自学者最容易混淆的概念、符号和计算陷阱。建议复习时逐项核对。

\section{第1章：数学预备}

\begin{enumerate}
\item \textbf{单位矢量求导陷阱}：球坐标和柱坐标中的单位矢量不是常矢量，$\partial\unitvec{r}/\partial\theta\neq 0$。只有在直角坐标系中单位矢量才是常量。
\item \textbf{旋度的旋度}：$\nabla\times(\nabla\times\vecA)\neq 0$ 一般不为零，它等于 $\grad(\divg\vecA)-\lap\vecA$。
\item \textbf{梯度为零 vs 标量场为零}：$\grad\phi=0$ 只说明 $\phi$ 是常数，不一定是零。电势的零点可以任意选取。
\item \textbf{$\delta$函数的三维性}：$\delta^{(3)}(\vect{r})$ 是体密度，量纲是 $[\text{长度}]^{-3}$，不能直接与一维$\delta(x)$混用。
\end{enumerate}

\section{第2--4章：静电学}

\begin{enumerate}
\item \textbf{电势的零点}：$\phi$ 的绝对值没有物理意义，只有电势差有意义。通常取无穷远或接地为参考点。
\item \textbf{D vs E}：$\vect{D}$ 只与自由电荷有关，但 $D=\eps E$ 仅在各向同性线性介质中成立。各向异性介质中 $D_i=\eps_{ij}E_j$ 是张量关系。
\item \textbf{边界条件方向}：法向条件用 $D_\perp$（因 $\divg\vect{D}=\rho_f$），切向条件用 $E_\parallel$（因 $\curl\vecE=0$ 静场中）。$\vect{D}$ 的切向和 $\vecE$ 的法向一般不连续。
\item \textbf{多极展开的收敛域}：展开式在 $r>a$（$a$为电荷分布最大尺度）外才收敛。在电荷分布内部不能直接使用。
\item \textbf{镜像法的适用范围}：只适用于特定对称边界（无限平面、球面、圆柱面）。对于一般边界，必须使用分离变量或数值方法。
\end{enumerate}

\section{第5--6章：静磁学}

\begin{enumerate}
\item \textbf{磁矢势的规范自由度}：$\vecA$ 本身不可测量，只有 $\vecB=\curl\vecA$ 有物理意义。不同规范给出相同的磁场。
\item \textbf{P标量 vs M矢量}：极化 $\vect{P}$ 产生标量束缚电荷 $\rho_b=-\nabla\cdot\vect{P}$；磁化 $\vect{M}$ 产生矢量束缚电流 $\vecJ_M=\nabla\times\vect{M}$。电磁不对称根源：无磁单极。
\item \textbf{B与H的边界}：法向条件用 $B_\perp$（因 $\divg\vecB=0$），切向条件用 $H_\parallel$（因 $\curl\vect{H}=\vecJ_f$）。不要反过来用。
\item \textbf{磁化电流方向}：$\vecJ_M$ 的方向由 $\nabla\times\vect{M}$ 的旋涡方向决定，与磁矩方向遵循右手定则。
\end{enumerate}

\section{第7章：麦克斯韦方程组}

\begin{enumerate}
\item \textbf{位移电流不是真实电流}：$\muz\eps_0\partial\vecE/\partial t$ 是磁场旋度的修正项，不是电荷的流动。它确保安培定律在时变情况下仍满足电荷守恒。
\item \textbf{感应电场的非保守性}：$\curl\vecE\neq 0$ 时，$\oint\vecE\cdot d\vect{l}\neq 0$，电势 $\phi$ 不再具有全局意义。时变场中应使用 $\vecE=-\nabla\phi-\partial\vecA/\partial t$。
\item \textbf{坡印廷矢量的方向}：$\vect{S}$ 指向能量传播方向。在直流导线中，能量沿径向从导线表面流入，而非沿导线内部传播。
\item \textbf{磁单极子的方程对称性}：如果磁单极存在，麦克斯韦方程可写为完全对称形式。但迄今实验未证实，故 $\divg\vecB=0$ 严格成立。
\end{enumerate}

\section{第8章：电磁波}

\begin{enumerate}
\item \textbf{E、B、k的右手关系}：平面波中 $\vecE\times\vecB$ 指向传播方向。记忆口诀：E-B-k 构成右手系。
\item \textbf{反射相移}：从光疏到光密介质反射时，电场可能有$\pi$相移；磁场则相反。反射率公式中已自动包含此相位信息。
\item \textbf{倏逝波不是传播波}：全反射时的透射波振幅随深度指数衰减，不携带净能量穿过界面。但它是真实的物理场，在光纤和近场光学中至关重要。
\item \textbf{趋肤深度 vs 波长}：良导体中 $\delta\ll\lambda$，电磁波几乎不进入导体内部。这与介质中的传播完全不同。
\end{enumerate}

\section{第9章：辐射}

\begin{enumerate}
\item \textbf{推迟势的因果性}：$t_r=t-R/c$ 保证了电磁相互作用以光速传播，不存在超距作用。这是相对论的基本要求。
\item \textbf{电偶极 vs 磁偶极辐射}：电偶极辐射功率 $\propto p_0^2\omega^4$，磁偶极 $\propto m_0^2\omega^4/c^2$。对于相同尺度的源，磁偶极辐射远弱于电偶极（因子约 $(v/c)^2$）。
\item \textbf{近场 vs 远场}：近场区（$r\ll\lambda$）中感应场主导，电场行为类似静电；远场区（$r\gg\lambda$）中辐射场主导，场 $\propto 1/r$。
\item \textbf{拉莫尔公式的适用范围}：仅适用于非相对论粒子（$v\ll c$）。高速粒子的辐射需使用相对论李纳特-维谢尔势。
\end{enumerate}

\section{第10章：相对论电动力学}

\begin{enumerate}
\item \textbf{号差的统一性}：本书采用 $(+,-,-,-)$。若使用 $(-,+,+,+)$，度规分量的正负号全部反转，所有降标/升标的结果也会改变。\textbf{必须全教材统一}。
\item \textbf{协变与逆变}：上标逆变（如 $x^\mu$），下标协变（如 $x_\mu=\eta_{\mu\nu}x^\nu$）。时空分量的逆变/协变只差一个负号，但混合指标的张量运算规则不变。
\item \textbf{E和B的相对性}：在某一参考系中纯静电场，换到另一参考系会出现磁场；反之亦然。$\vecE$ 和 $\vecB$ 是统一的电磁场张量 $F^{\mu\nu}$ 的不同分量。
\item \textbf{四维势的规范}：洛伦兹规范 $\partial_\mu A^\mu=0$ 在协变形式中极其简洁，是\textbf{相对论协变}的规范选择；库仑规范 $\nabla\cdot\vecA=0$ 反而\textbf{不是}协变的（其标势由电荷瞬时确定），两者是并列的两种规范而非推广关系，详见 §7.1.5。
\item \textbf{光速不可超越}：任何携带信息或能量的传播速度均不能超过 $c$。电磁场本身以 $c$ 传播，但相速度（$v_p=\omega/k$）可以超光速（如波导中），这不携带信息。
\end{enumerate}


\chapter{学习计划建议}

\begin{enumerate}
\item \textbf{第1周}：第1章（数学预备知识）。重点掌握梯度、散度、旋度的几何意义和计算，高斯定理和斯托克斯定理。
\item \textbf{第2周}：第2章（静电学基础）。掌握库仑定律、高斯定理、电势、泊松方程。完成所有基础例题。
\item \textbf{第3周}：第3章（电介质与边界条件）。理解极化机制、D和E的关系、边界条件。这是后续章节的关键。
\item \textbf{第4周}：第4章（边值问题）。镜像法和分离变量法。每种方法至少独立完成2道例题。
\item \textbf{第5周}：第5章（静磁学）。毕奥-萨伐尔定律、安培定理、磁矢势。注意与静电学的对比。
\item \textbf{第6周}：第6章（磁介质）。磁化机制、H场、磁学边界条件。
\item \textbf{第7周}：第7章（麦克斯韦方程组）。法拉第定律、位移电流。理解方程的物理内涵。
\item \textbf{第8周}：第8章（电磁波）。波动方程、平面波、反射折射。理解光的电磁本质。
\item \textbf{第9周}：第9章（辐射）。推迟势、电偶极辐射、拉莫尔公式。
\item \textbf{第10周}：第10章（相对论电动力学）。电磁场张量、协变方程。理解E和B的统一。
\end{enumerate}

\vspace{2em}
\noindent\rule{\textwidth}{0.4pt}
\vspace{1em}



\section*{每章自检清单}
\addcontentsline{toc}{section}{每章自检清单}

以下清单帮助你判断每章是否真正掌握。如果某一项无法独立完成，请回到对应章节重新复习。

\begin{center}
\begin{tabular}{@{}p{0.15\textwidth}p{0.8\textwidth}@{}}
\toprule
\textbf{章节} & \textbf{必须掌握的内容} \\
\midrule
第1章 & 能独立计算球/柱坐标下的梯度、散度、旋度；理解并矢运算；会用$\delta$函数处理点电荷/面电荷/线电荷；能验证至少3个矢量恒等式。 \\
第2章 & 能熟练运用高斯定理求对称电荷分布的电场；会求电偶极子电势和电场；能独立证明静电唯一性定理。 \\
第3章 & 能区分 $\vect{D}$、$\vecE$、$\vect{P}$ 的物理含义；会处理导体-介质混合边界问题；能正确应用法向/切向边界条件。 \\
第4章 & 掌握导体球/平面的镜像法；能在球/柱坐标中写分离变量通解；会计算电荷分布的多极矩（至少到四极）。 \\
第5章 & 能推导规范变换不变性；会选取库仑规范；理解磁矢势的物理意义；能处理磁镜像问题。 \\
第6章 & 能区分 $\vecB$、$\vect{H}$、$\vect{M}$；理解电磁不对称根源；会应用磁学边界条件解决分界面问题。 \\
第7章 & 能从法拉第定律推导时变电磁势；会计算坡印廷矢量；理解位移电流的物理本质；能写达朗贝尔方程。 \\
第8章 & 能推导平面波在介质/导体中的传播特性；会计算波导截止频率；理解全反射和倏逝波；会推导色散关系。 \\
第9章 & 能推导推迟势；会计算电偶极和磁偶极辐射功率；理解近场/远场区别；会用拉莫尔公式计算辐射功率。 \\
第10章 & 理解协变/逆变指标；会写四维电磁势和电磁场张量；能在不同参考系中变换电磁场；能还原协变麦克斯韦方程。 \\
\bottomrule
\end{tabular}
\end{center}

\begin{center}
\textit{本书完}
\end{center}


% ============================================================
\chapter{名词中英文对照表}
% ============================================================

本附录按章节顺序整理全书核心术语的中英文对照，方便读者查阅外文教材（如 Griffiths、Jackson）。

\section{第1章：数学预备}

\begin{center}
\begin{tabular}{@{}lll@{}}
\toprule
\textbf{中文} & \textbf{英文} & \textbf{符号} \\
\midrule
标量 & scalar & $\phi$, $f$ \\
矢量 & vector & $\vecA$, $\vecE$ \\
张量 & tensor & $T_{ij}$, $F^{\mu\nu}$ \\
梯度 & gradient & $\grad\phi$ \\
散度 & divergence & $\divg\vecA$ \\
旋度 & curl / rotation & $\curl\vecA$ \\
拉普拉斯算子 & Laplacian & $\lap$ \\
高斯定理 & Gauss's theorem & $\oint\vecA\cdot d\vect{S}=\int\divg\vecA\,dV$ \\
斯托克斯定理 & Stokes' theorem & $\oint\vecA\cdot d\vect{l}=\int(\curl\vecA)\cdot d\vect{S}$ \\
狄拉克$\delta$函数 & Dirac delta function & $\delta(x)$ \\
并矢 & dyad / dyadic & $\bm{A}\bm{B}$ \\
单位矢量 & unit vector & $\unitvec{x}$, $\unitvec{r}$ \\
正交坐标系 & orthogonal curvilinear coordinates & cylindrical / spherical \\
\bottomrule
\end{tabular}
\end{center}

\section{第2--4章：静电学}

\begin{center}
\begin{tabular}{@{}lll@{}}
\toprule
\textbf{中文} & \textbf{英文} & \textbf{符号} \\
\midrule
电场强度 & electric field & $\vecE$ \\
电势 & electric potential & $\phi$ / $V$ \\
电位移 & electric displacement & $\vect{D}$ \\
极化强度 & polarization & $\vect{P}$ \\
介电常数 & permittivity & $\eps$ \\
相对介电常数 & relative permittivity & $\eps_r$ \\
真空介电常数 & vacuum permittivity & $\eps_0$ \\
自由电荷密度 & free charge density & $\rho_f$ \\
束缚电荷密度 & bound charge density & $\rho_b$ \\
电偶极矩 & electric dipole moment & $\vect{p}$ \\
电四极矩 & electric quadrupole moment & $Q_{ij}$ \\
泊松方程 & Poisson's equation & $\lap\phi=-\rho/\eps_0$ \\
拉普拉斯方程 & Laplace's equation & $\lap\phi=0$ \\
镜像法 & method of images & — \\
分离变量法 & separation of variables & — \\
多极展开 & multipole expansion & — \\
唯一性定理 & uniqueness theorem & — \\
格林函数 & Green's function & $G(\vect{r},\vect{r}')$ \\
\bottomrule
\end{tabular}
\end{center}

\section{第5--6章：静磁学}

\begin{center}
\begin{tabular}{@{}lll@{}}
\toprule
\textbf{中文} & \textbf{英文} & \textbf{符号} \\
\midrule
磁感应强度 & magnetic induction / magnetic flux density & $\vecB$ \\
磁场强度 & magnetic field strength & $\vect{H}$ \\
磁矢势 & magnetic vector potential & $\vecA$ \\
磁化强度 & magnetization & $\vect{M}$ \\
磁导率 & permeability & $\muz$ / $\mu$ \\
相对磁导率 & relative permeability & $\mu_r$ \\
真空磁导率 & vacuum permeability & $\mu_0$ \\
磁偶极矩 & magnetic dipole moment & $\vect{m}$ \\
磁化电流 & magnetization current & $\vecJ_M$ \\
规范变换 & gauge transformation & $\vecA'=\vecA+\nabla\chi$ \\
库仑规范 & Coulomb gauge & $\nabla\cdot\vecA=0$ \\
洛伦兹规范 & Lorenz gauge & $\nabla\cdot\vecA+\muz\eps_0\partial\phi/\partial t=0$ \\
毕奥-萨伐尔定律 & Biot-Savart law & — \\
安培环路定理 & Amp\`ere's circuital law & — \\
迈斯纳效应 & Meissner effect & — \\
\bottomrule
\end{tabular}
\end{center}

\section{第7章：麦克斯韦方程组}

\begin{center}
\begin{tabular}{@{}lll@{}}
\toprule
\textbf{中文} & \textbf{英文} & \textbf{符号} \\
\midrule
法拉第电磁感应定律 & Faraday's law of induction & $\curl\vecE=-\partial\vecB/\partial t$ \\
位移电流 & displacement current & $\vecJ_D=\eps_0\partial\vecE/\partial t=\partial\vect{D}/\partial t$ \\
坡印廷矢量 & Poynting vector & $\vect{S}=\vecE\times\vect{H}$ \\
麦克斯韦应力张量 & Maxwell stress tensor & $T_{ij}$ \\
电磁场动量 & electromagnetic momentum & $\vect{g}=\muz\eps_0\vect{S}$ \\
达朗贝尔方程 & d'Alembert equation & $\Box\phi=-\rho/\eps_0$ \\
推迟势 & retarded potential & — \\
\bottomrule
\end{tabular}
\end{center}

\section{第8章：电磁波}

\begin{center}
\begin{tabular}{@{}lll@{}}
\toprule
\textbf{中文} & \textbf{英文} & \textbf{符号} \\
\midrule
平面电磁波 & plane electromagnetic wave & — \\
波数 & wave number & $k$ \\
角频率 & angular frequency & $\omega$ \\
相速度 & phase velocity & $v_p=\omega/k$ \\
群速度 & group velocity & $v_g=d\omega/dk$ \\
折射率 & refractive index & $n$ \\
反射系数 & reflection coefficient & $r$ / $R$ \\
透射系数 & transmission coefficient & $t$ / $T$ \\
趋肤深度 & skin depth & $\delta$ \\
倏逝波 & evanescent wave & — \\
波导 & waveguide & — \\
截止频率 & cutoff frequency & $\omega_c$ \\
色散 & dispersion & — \\
正常色散 & normal dispersion & $dn/d\omega>0$ \\
反常色散 & anomalous dispersion & $dn/d\omega<0$ \\
\bottomrule
\end{tabular}
\end{center}

\section{第9章：电磁辐射}

\begin{center}
\begin{tabular}{@{}lll@{}}
\toprule
\textbf{中文} & \textbf{英文} & \textbf{符号} \\
\midrule
电偶极辐射 & electric dipole radiation & — \\
磁偶极辐射 & magnetic dipole radiation & — \\
拉莫尔公式 & Larmor formula & $P=q^2a^2/(6\pi\eps_0 c^3)$ \\
同步辐射 & synchrotron radiation & — \\
辐射功率 & radiated power & $P$ \\
近场区 & near-field zone & $r\ll\lambda$ \\
远场区 & far-field zone & $r\gg\lambda$ \\
辐射电阻 & radiation resistance & $R_{\text{rad}}$ \\
天线 & antenna & — \\
半波天线 & half-wave antenna & — \\
\bottomrule
\end{tabular}
\end{center}

\section{第10章：相对论电动力学}

\begin{center}
\begin{tabular}{@{}lll@{}}
\toprule
\textbf{中文} & \textbf{英文} & \textbf{符号} \\
\midrule
洛伦兹变换 & Lorentz transformation & — \\
洛伦兹因子 & Lorentz factor & $\gamma$ \\
固有时 & proper time & $\tau$ \\
闵氏时空 & Minkowski spacetime & — \\
闵氏度规 & Minkowski metric & $\eta_{\mu\nu}$ \\
协变坐标 & covariant coordinates & $x_\mu$ \\
逆变坐标 & contravariant coordinates & $x^\mu$ \\
四维速度 & four-velocity & $u^\mu$ \\
四维电流 & four-current & $j^\mu$ \\
四维势 & four-potential & $A^\mu=(\phi/c,\vecA)$ \\
电磁场张量 & electromagnetic field tensor & $F^{\mu\nu}$ \\
协变麦克斯韦方程 & covariant Maxwell equations & $\partial_\mu F^{\mu\nu}=\muz j^\nu$ \\
洛伦兹不变量 & Lorentz invariant & $I_1, I_2$ \\
光锥 & light cone & — \\
类时 & timelike & $s^2>0$ \\
类空 & spacelike & $s^2<0$ \\
类光 & lightlike / null & $s^2=0$ \\
\bottomrule
\end{tabular}
\end{center}

\section{常用物理常数名称}

\begin{center}
\begin{tabular}{@{}lll@{}}
\toprule
\textbf{中文} & \textbf{英文} & \textbf{数值（SI）} \\
\midrule
真空光速 & speed of light in vacuum & $c=2.998\times 10^8\,\text{m/s}$ \\
真空介电常数 & vacuum permittivity & $\eps_0=8.854\times 10^{-12}\,\text{F/m}$ \\
真空磁导率 & vacuum permeability & $\mu_0=4\pi\times 10^{-7}\,\text{H/m}$ \\
元电荷 & elementary charge & $e=1.602\times 10^{-19}\,\text{C}$ \\
电子质量 & electron mass & $m_e=9.109\times 10^{-31}\,\text{kg}$ \\
\bottomrule
\end{tabular}
\end{center}


% ============================================================
% 附录：随堂自测答案速查
% ============================================================
\chapter{随堂自测答案速查}

本附录汇总全书 10 章随堂自测的答案，供\textbf{打印版读者}和离线阅读时快速自检使用（PDF 电子版可用正文中的折叠按钮展开详解）。

\begin{tip}{使用说明}{quiz-answer-tip}
建议先独立完成各章随堂自测，再对照本表核对。选择题给出选项字母，判断题给出"正确/错误"，填空题给出答案关键词。详细解析见各章正文自测框。
\end{tip}

\begin{center}
\small
\begin{tabular}{@{}cll>{\raggedright\arraybackslash}p{4.2cm}@{}}
\toprule
\textbf{章} & \textbf{选择题} & \textbf{判断题} & \textbf{填空题} \\
\midrule
第1章 数学预备知识 & 1-B \quad 2-C \quad 3-C & 4-正确 \quad 5-错误 \quad 6-错误 & — \\
第2章 静电学基础 & 1-B \quad 2-A \quad 3-B & — & 4-$\curl\vecE=\vect{0}$；5-$\vecE=-\grad\phi$；6-$\lap\phi=-\rho/\eps_0$；7-$\sigma/\eps_0$（法向） \\
第3章 电介质与边界 & 1-B \quad 2-C \quad 3-B & 4-正确 \quad 5-错误 & — \\
第4章 边值问题 & 1-C \quad 2-C \quad 3-B & — & 4-电势/电势梯度 \quad 5-贝塞尔 \\
第5章 静磁学 & 1-A \quad 2-B \quad 3-B & 4-正确 \quad 5-错误 & — \\
第6章 磁介质 & 1-B \quad 2-B \quad 3-A & — & — \\
第7章 麦克斯韦方程组 & 1-B \quad 2-B \quad 3-A & 4-正确 \quad 5-错误 & — \\
第8章 电磁波传播 & 1-B \quad 2-A \quad 3-B & — & 4-坡印廷（$\vecE\times\vecB/\muz$）；5-垂直、垂直 \\
第9章 电磁波辐射 & 1-C \quad 2-B \quad 3-B & 4-正确 \quad 5-错误 & — \\
第10章 相对论电动力学 & 1-C \quad 2-A \quad 3-B & — & 4-$1/\sqrt{1-v^2/c^2}$；5-$\lap-\partial_t^2/c^2$；6-$\partial_\mu A^\mu=0$ \\
\bottomrule
\end{tabular}
\end{center}


% ============================================================
% 附录I：延伸阅读与参考书目
% ============================================================
\chapter{延伸阅读与参考书目}

\begin{tip}{阅读建议}{reading-tip}
以下书目按难度分层，读者可根据自身基础选择：\textbf{入门级}适合第一遍学习，\textbf{进阶级}适合巩固提高，\textbf{专业级}适合深入研究和考研复习。
\end{tip}

\section*{中文教材（入门级）}
\begin{enumerate}[label={[\arabic*]}]
\item \textbf{郭硕鸿，《电动力学（第三版）》}，高等教育出版社。\par
  国内最广泛使用的电动力学教材，篇幅适中，例题丰富，适合本科生第一遍学习。配套习题解答也已出版。

\item \textbf{赵凯华、陈熙谋，《电磁学》}，高等教育出版社。\par
  虽然名为"电磁学"，但其深度已接近电动力学入门水平。赵凯华先生的物理图像清晰，文字优美，强烈推荐作为预习或补充阅读。

\item \textbf{梁灿彬、周彬，《电磁学》}，高等教育出版社。\par
  梁灿彬先生是广义相对论专家，此书对电磁学的数学结构（尤其是微分形式语言）有独到见解。
\end{enumerate}

\section*{英文经典教材（进阶级）}
\begin{enumerate}[label={[\arabic*]},resume]
\item \textbf{David J. Griffiths, \textit{Introduction to Electrodynamics}, 4th ed.}，Pearson。\par
  国际最畅销的电动力学教材，以物理图像清晰、例题详尽著称。第4版已包含相对论的四维协变形式。强烈建议阅读英文原版。

\item \textbf{John David Jackson, \textit{Classical Electrodynamics}, 3rd ed.}，Wiley。\par
  电动力学领域的"圣经"，以数学严谨性和内容深度闻名。适合研究生水平和考研复习。第1-6章是核心，第11-14章涉及辐射和狭义相对论。

\item \textbf{L. D. Landau and E. M. Lifshitz, \textit{The Classical Theory of Fields}, 4th ed.}，Pergamon。\par
  朗道理论物理教程第2卷，从相对论原理出发构建电动力学，风格独特、物理深刻。适合已掌握基础、追求理论深度的读者。

\item \textbf{Andrew Zangwill, \textit{Modern Electrodynamics}}，Cambridge University Press。\par
  较新的教材（2013年），内容涵盖传统电动力学 plus 现代主题（如超材料、拓扑电磁学），数学处理比Jackson更友好。
\end{enumerate}

\section*{专题深化（专业级）}
\begin{enumerate}[label={[\arabic*]},resume]
\item \textbf{J. D. Jackson, \textit{Classical Electrodynamics}, Chap. 14--16}\par
  辐射的相对论性处理、多极辐射的严格理论。

\item \textbf{Misner, Thorne \& Wheeler, \textit{Gravitation}, Chap. 3--4}\par
  用微分几何语言重新表述电动力学，为理解广义相对论做准备。

\item \textbf{M. Schwartz, \textit{Quantum Field Theory and the Standard Model}}\par
  第2--3章从经典场论过渡到量子场论，电动力学是QED的古典极限。
\end{enumerate}

\section*{在线资源}
\begin{itemize}
\item MIT OpenCourseWare: 8.07 Electromagnetism II（MIT OCW官网，含完整讲义、习题和视频）
\item Stanford Physics 210/211: Classical Electrodynamics（含Jackson教材的习题解答讨论）
\item Physics Stack Exchange: 电动力学标签下的问答社区，适合解决具体疑惑
\item 3Blue1Brown的\textit{Essence of linear algebra}系列（巩固矢量和张量的几何直觉）
\end{itemize}

\section*{教材章节对照表}

\begin{tip}{郭硕鸿 / Griffiths 章节对应关系}{textbook-map-tip}
交叉阅读不同教材时，可按下列对应关系定位同一主题（郭硕鸿《电动力学》第三版；Griffiths \textit{Introduction to Electrodynamics} 4th ed.）：
\begin{center}
\small
\begin{tabular}{@{}p{4.2cm}p{4.6cm}p{4.6cm}@{}}
\toprule
\textbf{主题} & \textbf{本书章节} & \textbf{郭硕鸿 / Griffiths} \\
\midrule
矢量分析、 $\delta$ 函数 & 第 1 章 & 第 1 章 / 第 1 章 \\
静电学、泊松方程 & 第 2 章 & 第 1 章 / 第 2 章 \\
介质与边界条件 & 第 3 章 & 第 2 章 / 第 4 章 \\
边值问题（镜像、分离变量、多极） & 第 4 章 & 第 2 章 / 第 3--4 章 \\
静磁学、矢势 & 第 5 章 & 第 3 章 / 第 5 章 \\
磁介质 & 第 6 章 & 第 3 章 / 第 6 章 \\
麦克斯韦方程组、规范 & 第 7 章 & 第 4 章 / 第 7 章 \\
电磁波传播、波导 & 第 8 章 & 第 5 章 / 第 8--9 章 \\
辐射（推迟势、偶极辐射） & 第 9 章 & 第 6 章 / 第 10--11 章 \\
相对论电动力学 & 第 10 章 & 第 7 章 / 第 12 章 \\
\bottomrule
\end{tabular}
\end{center}
\textbf{阅读提示}：郭硕鸿篇幅紧凑、推导精炼，适合对照复习；Griffiths 例题丰富、物理图像详尽，适合第一遍学习卡壳时查阅。两书的符号约定与本书基本一致（SI 单位制），$x^\mu=(ct,\vect{r})$ 与 $A^\mu=(\phi/c,\vecA)$ 的排列也相同；唯一要留心的是 $F^{\mu\nu}$ 中电场的符号约定（不同教材可能取 $F^{0i}=\pm E_i/c$），对照时一律以定义式为准。
\end{tip}


\begin{center}\small \hyperlink{ch9-goal}{上一章：第9章}\end{center}


\end{document}
